%% =====================================================================
%%  Notas de Análisis de Sistemas de Control -- archivo principal
%%  Compilar desde esta carpeta (notas/):
%%      latexmk -pdf NotasControl.tex
%%  o   pdflatex NotasControl; bibtex NotasControl; pdflatex x2
%% =====================================================================
\documentclass{apuntescontrol}   % opciones: borrador, sinqr
\usepackage{notacion}
%% Dirección de la página web y del repositorio (cambiar al publicar):
\definirsitio{https://quijanon-bot.github.io/NQuijano.github.io/curso-control}{https://github.com/quijanon-bot/NQuijano.github.io/tree/main/fuentes/curso-control}[https://github.com/quijanon-bot/NQuijano.github.io/blob/main/fuentes/curso-control/codigo]

\begin{document}
\pagenumbering{arabic}  % (sin \frontmatter: babel-spanish + números romanos rompen \label)
\portada{Notas de Análisis de Sistemas de Control}{Nicanor Quijano\\[2pt]
  {\large Departamento de Ingeniería Eléctrica y Electrónica}}{Versión piloto \textperiodcentered{} noviembre de 2026}
\tableofcontents
%% =====================================================================
%%  Convenciones y notación del curso (capítulo sin número)
%%  Fuente única de la notación: notacion.sty. Esta tabla la documenta.
%% =====================================================================
\chapter*{Cómo leer estas notas}
\addcontentsline{toc}{chapter}{Cómo leer estas notas}
\markboth{Cómo leer estas notas}{}

\epigrafe{Caminante, son tus huellas el camino y nada más; caminante, no hay camino, se hace camino al andar.}{Antonio Machado, \emph{Proverbios y cantares}}

Estas notas acompañan el curso \emph{Análisis de Sistemas de Control} (IELE 2300). Cada capítulo corresponde a una parte del curso y tiene una página web con los códigos de \matlab/\simulink, los videos y las actividades. El código QR al comienzo de cada capítulo lleva a esa página.

\section*{Las cajas}

Cada tipo de contenido tiene su caja y su color. Así se puede leer el capítulo a dos velocidades: la primera vez, completo; para repasar, solo las cajas azules y los procedimientos.

\begin{definicion}[title={\textcolor{ctrlazul}{\faBookOpen}\hspace{0.45em}Definición 10.2: así se ve una definición}]{}
Las definiciones, los ejemplos y las pausas están numerados por capítulo y se citan con ese número.
\end{definicion}

\begin{ideaclave}
Lo que hay que recordar de una sección. Si solo tiene cinco minutos antes de un examen, lea estas cajas.
\end{ideaclave}

\begin{ejemplo}[title={\textcolor{ctrlvioleta}{\faPencil*}\hspace{0.45em}Ejemplo 10.5: título del ejemplo}]{}
El enunciado va primero.
\solucion Después de la línea viene el desarrollo. Intente resolverlo antes de leerlo.
\end{ejemplo}

\begin{cuidado}
Los errores que más se repiten en las evaluaciones y las confusiones de notación o de lenguaje.
\end{cuidado}

\begin{pausa}[title={\textcolor{ctrlverde}{\faPauseCircle}\hspace{0.45em}Pausa activa 10.3}]
Una pregunta corta para comprobar que entendió antes de seguir. Las respuestas están al final de cada capítulo.
\end{pausa}

\begin{historia}
De dónde salió cada idea, con la referencia exacta del trabajo original.
\end{historia}

\begin{profundiza}
Material opcional, marcado con (*): no se evalúa, pero ayuda a entender.
\end{profundiza}

\begin{receta}[Procedimiento]
Los procedimientos de diseño, paso a paso, en una caja oscura para encontrarlos rápido.
\end{receta}

\begin{modelo}[Nombre.slx]
Un modelo de \simulink. En el curso los modelos se \emph{construyen por código} con un script \cmd{crearModelo\_*.m}: así cualquiera puede ver, versionar y mejorar exactamente qué bloques tiene.
\end{modelo}

\paragraph{El código}
Los bloques de código verde azulado salen \emph{directamente} de los archivos del repositorio del curso (la ruta aparece en el encabezado y es un enlace). Si un código se mejora en el repositorio, las notas se actualizan solas al compilar. Los números que aparecen en los comentarios de los códigos son los mismos del texto, y cada capítulo tiene un script \cmd{verificarCapNN.m} que comprueba que siguen coincidiendo. En el margen, \textcolor{ctrlteal}{\faFileCode} indica un script y \textcolor{ctrlcoral}{\faPlayCircle} un video. Los problemas marcados \etiquetaM{} requieren \matlab/\simulink.

\paragraph{Antes de correr los códigos}
Ejecute una vez por sesión \cmd{iniciarCurso} (carpeta \texttt{codigo/} del repositorio). Agrega al \emph{path} las funciones comunes (\cmd{bodeAsintotico}, \cmd{margenesTabla}, \cmd{faseEn}, \cmd{estiloFigura}, \ldots) y permite correr cualquier script desde cualquier carpeta.

\section*{Notación}

La notación se define en un solo lugar, el archivo \texttt{notacion.sty}, y todos los capítulos la usan. Las tablas siguientes la documentan, con su equivalente en \matlab{} y el nombre más común en los libros en inglés.

{\small
\begin{xltabular}{\linewidth}{@{}>{\raggedright\arraybackslash}p{2.9cm} >{\raggedright\arraybackslash}X >{\raggedright\arraybackslash\ttfamily\footnotesize}p{2.5cm}@{}}
\caption{Notación del curso. La última columna es el nombre que usan los códigos. (¡!): \texttt{margin} entrega \texttt{Gm} en veces y llama \texttt{Wcg} a la frecuencia del $\GM$ y \texttt{Wcp} a la del $\PM$.}\label{tab:notacion}\\
\toprule
\textbf{Símbolo} & \textbf{Significado} \textcolor{ctrlgris}{\textit{(en inglés)}} & \textbf{\normalfont\small\bfseries En MATLAB}\\
\midrule
\endhead
\multicolumn{3}{@{}l}{\sffamily\bfseries\color{ctrlazul}Señales y bloques}\\
$r,\,y,\,e,\,u$ & referencia, salida, error $e=r-y$, señal de control \textcolor{ctrlgris}{\textit{(reference, output, error, input)}} & \texttt{r, y, e, u}\\
$d,\,n$ & perturbación, ruido de medición \textcolor{ctrlgris}{\textit{(disturbance, noise)}} & \texttt{d, n}\\
$\Gp(s)$, $\Gc(s)$ & proceso (planta), controlador o compensador \textcolor{ctrlgris}{\textit{(plant, controller)}} & \texttt{Gp, Gc}\\
$\Gv,\Gs,\Gt$ & elemento final de control, sensor, transmisor \textcolor{ctrlgris}{\textit{(actuator, sensor)}} & \texttt{Gv, Gs, Gt}\\
$\Hs(s)$ & retroalimentación \textcolor{ctrlgris}{\textit{(feedback)}} & \texttt{H}\\
$\Gma(s)=G\Hs$ & \textbf{malla abierta}: una vuelta completa al lazo, sin el signo del comparador \textcolor{ctrlgris}{\textit{(loop transfer function $L$)}} & \texttt{Gma}\\
$\Gla(s)$ & \textbf{lazo abierto}: el lazo reescrito con retroalimentación unitaria; tipo y error \textcolor{ctrlgris}{\textit{((open-loop t.f.))}} & \texttt{Gol}\\
$\Tlc(s)=Y/R$ & lazo cerrado \textcolor{ctrlgris}{\textit{(closed loop)}} & \texttt{T = feedback(...)}\\
$\Sens(s)=1/(1+\Gma)$ & sensitividad \textcolor{ctrlgris}{\textit{(sensitivity)}} & \texttt{S}\\
\midrule
\multicolumn{3}{@{}l}{\sffamily\bfseries\color{ctrlazul}Modelado y linealización}\\
$\Lap{f(t)}=F(s)$ & transformada de Laplace (unilateral); $s=\sigma+\jw$, con $j$ \textcolor{ctrlgris}{\textit{(Laplace transform)}} & \texttt{laplace}\\
$\escalon(t)$, $\impulso(t)$ & escalón unitario e impulso (no $u(t)$: $u$ es la entrada) \textcolor{ctrlgris}{\textit{(unit step, impulse)}} & \texttt{heaviside}, \texttt{dirac}\\
$\eq{x},\ \eq{u}$ & punto de equilibrio (de operación) \textcolor{ctrlgris}{\textit{(equilibrium point)}} & \texttt{xe, ue}\\
$\dev{x}=x-\eq{x}$ & \textbf{variable desviación}; en $s$: $\Dev{X}(s)$ \textcolor{ctrlgris}{\textit{(deviation variable)}} & \texttt{xd}\\
$\Delta$ & \textbf{discriminante} del diagrama (ley de Mason) \textcolor{ctrlgris}{\textit{(graph determinant)}} & --\\
\midrule
\multicolumn{3}{@{}l}{\sffamily\bfseries\color{ctrlazul}Espacio de estado}\\
$\Am,\Bm,\Cm,\Dm$ & matrices de la representación de estado \textcolor{ctrlgris}{\textit{(state-space matrices)}} & \texttt{A, B, C, D}\\
$\Phit$ & matriz de transición, $e^{\Am t}$ \textcolor{ctrlgris}{\textit{(transition matrix)}} & \texttt{expm(A*t)}\\
$\Ctrb,\ \Obsv$ & matrices de controlabilidad y observabilidad \textcolor{ctrlgris}{\textit{(controllability, observability)}} & \texttt{ctrb, obsv}\\
$\Km,\ \Lm$ & ganancias de retroalimentación de estado y del observador \textcolor{ctrlgris}{\textit{(state feedback, observer gain)}} & \texttt{K, L}\\
\midrule
\multicolumn{3}{@{}l}{\sffamily\bfseries\color{ctrlazul}Respuesta transitoria y error}\\
$\zeta,\ \wn,\ \wdamp$ & factor de amortiguamiento, frecuencia natural, frecuencia amortiguada \textcolor{ctrlgris}{\textit{(damping ratio, natural freq.)}} & \texttt{zeta, wn, wd}\\
$\OS$ & \textbf{sobrepico} porcentual \textcolor{ctrlgris}{\textit{(percent overshoot)}} & \texttt{S.Overshoot}\\
$\ts,\ \tp,\ \tr$ & tiempo de establecimiento (2\,\%), de pico, de subida \textcolor{ctrlgris}{\textit{(settling, peak, rise time)}} & \texttt{stepinfo}\\
$\ess$ & error en estado estacionario \textcolor{ctrlgris}{\textit{(steady-state error)}} & \texttt{ess}\\
$\kp,\ \kv,\ \ka$ & constantes estáticas de posición, velocidad y aceleración \textcolor{ctrlgris}{\textit{(static error constants)}} & \texttt{kp, kv, ka}\\
\midrule
\multicolumn{3}{@{}l}{\sffamily\bfseries\color{ctrlazul}PID}\\
$\Kc,\ \Ti,\ \Td$ & ganancia del controlador, tiempo integral, tiempo derivativo \textcolor{ctrlgris}{\textit{(controller gain)}} & \texttt{Kc, Ti, Td}\\
$\Kproc$ & ganancia del \emph{proceso} (modelos de primer orden con tiempo muerto) \textcolor{ctrlgris}{\textit{(process gain)}} & \texttt{Kp}\\
$\Kcr,\ \Pcr,\ \wcr$ & ganancia, periodo y frecuencia críticos (últimos) \textcolor{ctrlgris}{\textit{(ultimate gain, period)}} & \texttt{Kcr, Pcr, wcr}\\
\midrule
\multicolumn{3}{@{}l}{\sffamily\bfseries\color{ctrlazul}Respuesta frecuencial}\\
$\Kb$ & ganancia en forma de Bode (cada factor como $s/\omega_i+1$) \textcolor{ctrlgris}{\textit{(Bode gain)}} & \texttt{Kb}\\
$\GM,\ \wgm$ & margen de ganancia (dB) y su frecuencia (fase $=-180\grados$) \textcolor{ctrlgris}{\textit{(gain margin)}} & \texttt{Gm, Wcg} (¡!)\\
$\PM,\ \wpm$ & margen de fase y su frecuencia ($|\Gma|=0\dB$) \textcolor{ctrlgris}{\textit{(phase margin)}} & \texttt{Pm, Wcp} (¡!)\\
$\Mr,\ \wres$ & pico de resonancia del lazo cerrado y su frecuencia \textcolor{ctrlgris}{\textit{(resonant peak)}} & \texttt{Mr, wr}\\
$\wbw$ & ancho de banda ($-3\dB$ del lazo cerrado) \textcolor{ctrlgris}{\textit{(bandwidth)}} & \texttt{bandwidth}\\
$\alpha,\ \phiM,\ \wM,\ \wL,\ \wH$ & parámetros de los compensadores en adelanto y atraso \textcolor{ctrlgris}{\textit{(lead/lag parameters)}} & \texttt{alpha, phiM, wM, wL, wH}\\
\bottomrule
\end{xltabular}}

\section*{Decisiones de lenguaje}

Algunas palabras tienen varias versiones en los libros y en el habla del laboratorio. En el curso usamos siempre la primera columna.

\begin{center}\small
\begin{tabularx}{\linewidth}{@{}>{\raggedright\arraybackslash}p{3.6cm} >{\raggedright\arraybackslash}p{3.0cm} >{\raggedright\arraybackslash}X@{}}
\toprule
\textbf{Usamos} & \textbf{En lugar de} & \textbf{Por qué}\\
\midrule
retroalimentación & realimentación & término del curso y de la industria en Colombia\\
malla abierta $\Gma$ y lazo abierto $\Gla$ & ``open loop'' para ambas & son funciones distintas cuando $\Hs\neq1$\\
variable desviación & variable incremental & ``incremental'' se reserva para el algoritmo digital del PID\\
discriminante (Mason) & determinante & no es el determinante de una matriz\\
sobrepico ($\OS$) & overshoot, sobreimpulso, $M_p$ & $M_p$ se confundía con el pico de resonancia\\
pico de resonancia $\Mr$ & $M_P$ & ídem\\
tiempo de establecimiento & tiempo de asentamiento & una sola palabra para $\ts$ (criterio del 2\,\%)\\
factor de amortiguamiento $\zeta$ & coeficiente o razón de amortiguamiento & una sola palabra\\
$\Kc$ (controlador) & $K_p$ para el controlador & $K_p$ es la ganancia del proceso\\
$\Pcr$ & $T_{cr}$, $P_u$ & $T$ ya nombra al lazo cerrado y a varios tiempos\\
fórmula de Ackermann & Ackerman & el apellido es Ackermann\\
\bottomrule
\end{tabularx}
\end{center}

\begin{profundiza}[Dónde estaban las inconsistencias]
Al unificar las notas se encontraron, entre otras: ``realimentación'' en los capítulos de modelado, PID y Bode; el punto de operación escrito como $\bar{x}$ en modelado y espacio de estado, y como $\eq{x}$ en otros lugares; ``determinante'' para el $\Delta$ de Mason; ``overshoot'' y $MP$ para el sobrepico en el capítulo de segundo orden y $M_P$ para el pico de resonancia en el de Bode; $K_p$ para la ganancia del controlador en la tabla de Ziegler-Nichols; ``Ackerman''; ``tiempo de asentamiento'' y ``de establecimiento'' mezclados. El detalle, capítulo por capítulo, está en \texttt{notas/MIGRACION.md} del repositorio. A medida que cada capítulo se migre al formato nuevo, se reemplazan los símbolos escritos a mano por las macros de \texttt{notacion.sty}.
\end{profundiza}




%% ---------------------------------------------------------------------
%% Capítulos. Los que dicen "por migrar" todavía están en el formato
%% anterior (LibroCtrl2026_V2/Chapters); compilan con esta clase gracias a
%% las cajas de compatibilidad, pero conviene migrarlos uno a uno siguiendo
%% el piloto (capítulo 10) y la guía notas/MIGRACION.md.
%% ---------------------------------------------------------------------
% Primera versión completa (noviembre de 2026): los diez capítulos en el formato nuevo.
%% =====================================================================
%%  Capítulo 1 -- Introducción: qué es el análisis de sistemas de control
%%  y de dónde viene.
%%  Migrado al formato nuevo (noviembre de 2026) desde Introduccion.tex.
%%  Cambios de fondo respecto a la versión anterior (ver MIGRACION.md):
%%    la historia se reorganizó por épocas, se verificaron fechas y nombres
%%    contra fuentes primarias o reseñas especializadas y se corrigieron
%%    varios datos (Watt no patentó el regulador; Sperry voló en 1914;
%%    Bode publicó en 1940--45; IFAC se fundó en 1957; RIAS en 1955...);
%%    el tanque en lazo abierto ahora se compara con el lazo cerrado,
%%    con números verificados; "malla abierta" y "lazo abierto" se
%%    distinguen de "control en lazo abierto".
%% =====================================================================

\capitulo[cap:in]{Introducción: el análisis de sistemas de control}
  {Es que con los ojos del presente ningún pasado es tolerable\ldots{} Nada hay más ingrato e implacable que la eternidad; nada se purga y se corrige con más violencia que la historia.}
  {Juan Esteban Constaín}

\begin{objetivos}[partes/01-introduccion.html]
  \item Explicar qué se entiende por \emph{análisis}, \emph{sistema} y \emph{control}, y qué hace la retroalimentación.
  \item Distinguir el control en lazo abierto del control en lazo cerrado y medir, con un ejemplo, qué gana y qué cuesta retroalimentar.
  \item Plantear un problema de control: objetivos, variables manipuladas y perturbaciones, mediciones, restricciones y estructura del controlador.
  \item Ubicar en el tiempo las ideas del curso, desde los reguladores de flotador de la Antigüedad hasta el control digital, y saber dónde leer más.
  \item Saber cómo están organizados el curso, las notas, los códigos y la página web.
\end{objetivos}

%% ---------------------------------------------------------------------
\section{Análisis, sistema y control}
\label{sec:in-palabras}
%% ---------------------------------------------------------------------

El nombre del curso tiene tres palabras y cada una carga con un significado preciso. Vale la pena detenerse en ellas antes de escribir la primera ecuación.

\paragraph{Análisis}
El diccionario de la Real Academia da varias acepciones de \emph{análisis}: la distinción y separación de las partes de algo para conocer su composición, el estudio detallado de algo y, la que nos interesa, la \emph{parte de las matemáticas basada en los conceptos de límite, convergencia y continuidad}, que da origen al cálculo diferencial e integral y a la teoría de funciones. En el curso las dos ideas van juntas: \emph{separamos} un sistema en sus partes (planta, actuador, sensor, controlador) y \emph{estudiamos} cada una con herramientas del análisis matemático: ecuaciones diferenciales, transformada de Laplace, variable compleja.

\paragraph{Sistema}
Un sistema es un conjunto de cosas (personas, células, moléculas, tuberías, transistores) interconectadas de tal forma que producen un patrón de comportamiento a medida que pasa el tiempo; puede ser activado, restringido o impulsado por fuerzas externas, pero su respuesta depende de sus propias características~\citep{MeWr08}. Más corto: una combinación de componentes que se percibe como un todo~\citep{HaCh08}. Lo que hace interesantes a los sistemas que estudiaremos es que son \emph{dinámicos}: lo que hacen ahora depende de lo que les pasó antes. En el capítulo~\ref{cap:po} precisaremos estas ideas con las nociones de señal y de sistema dinámico.

\paragraph{Control}
La Real Academia también da varias acepciones de \emph{control} (del francés \emph{contrôle}): comprobación, dominio, el examen parcial con el que se comprueba la marcha de los alumnos\ldots{} y la que nos interesa: \emph{regulación, manual o automática, sobre un sistema}.

\begin{definicion}[label=def:in-sistemacontrol]{Sistema de control}
Un \clave{sistema de control} es un conjunto de dispositivos que administra, ordena, dirige o regula el comportamiento de otro sistema para obtener un resultado deseado~\citep{WikiControl}. Está formado por subsistemas (el proceso o planta, los actuadores, los sensores, el controlador) ensamblados para que una salida siga una referencia o rechace perturbaciones, con un desempeño especificado.
\end{definicion}

El \emph{análisis de sistemas de control}, entonces, estudia cómo evolucionan en el tiempo esos elementos interconectados y qué tan bien cumplen su propósito: si el lazo es estable, qué tan rápido y con cuánto sobrepaso responde, cuánto error deja, qué tanto le afectan las perturbaciones y los errores del modelo.

\subsection{La idea central: retroalimentar}

\begin{definicion}[label=def:in-retro]{Retroalimentación}
\clave{Retroalimentar} es medir la variable que se quiere controlar, compararla con el valor deseado y usar la diferencia (el \emph{error}) para decidir la acción. Un sistema así se llama \clave{retroalimentado} o \emph{de lazo cerrado}.
\end{definicion}

La retroalimentación es el concepto fundamental del curso y su gran virtud es la \clave{robustez}: como la acción depende de lo que realmente pasa, el sistema puede cumplir sus requisitos aunque el modelo sea imperfecto, aparezcan perturbaciones o las mediciones lleguen con ruido. Su peligro es la \alerta{sobrecorrección}: si se reacciona demasiado fuerte o demasiado tarde, el lazo oscila y puede volverse inestable. Casi todo el curso trata de cómo aprovechar la virtud sin caer en el peligro.

Richard Murray resumió el área en su conferencia Bode de 2016 con una ecuación de cuatro términos, \emph{control = dynamics, uncertainty, feedforward, feedback}~\citep{Mur16}, cuyas iniciales (DUFF) recuerdan la cerveza de \emph{Los Simpson}. Leída en orden: (i) hay unas \clave{dinámicas} que queremos modificar; (ii) las conocemos con \clave{incertidumbre}; (iii) con la \clave{prealimentación} (\emph{feedforward}) actuamos con lo que sabemos de antemano: la referencia, o una perturbación que podemos medir; y (iv) con la \clave{retroalimentación} (\emph{feedback}) corregimos lo que no sabíamos.

\begin{ideaclave}
\textbf{Señales y sistemas.}
La retroalimentación se entiende con dos conceptos entrelazados~\citep{AlMa10}: las \emph{señales}, porque es midiéndolas como se decide qué hacer, y los \emph{sistemas}, porque la retroalimentación existe para cambiar su comportamiento. Cuando hablamos de un sistema hablamos de algo que produce o modifica una señal; cuando hablamos de una señal, de algo que algún sistema produjo.
\end{ideaclave}

\begin{pausa}
Cuando usted conduce un carro por una carretera con curvas, ¿quién hace de sensor, de controlador y de actuador? ¿Qué es la referencia y qué la perturbación?
\respuesta{Los ojos (y el oído interno) son el sensor; el cerebro compara la posición del carro con el centro del carril (la referencia) y decide: es el controlador; los brazos y el timón son el actuador. El viento lateral, los baches y la inclinación de la vía son perturbaciones. Si usted cierra los ojos, el lazo se abre.}
\end{pausa}

%% ---------------------------------------------------------------------
\section{Lazo abierto y lazo cerrado}
\label{sec:in-lazos}
%% ---------------------------------------------------------------------

Hay dos maneras de mandar sobre un sistema (\cref{fig:in-lazos}). En el \clave{control en lazo abierto} el controlador decide la acción sin mirar la salida: un actuador mueve el proceso según una calibración hecha de antemano. En el \clave{control en lazo cerrado} un sensor mide la salida, el comparador calcula el error $e=r-y_m$ y el controlador actúa sobre ese error.

\begin{figure}[ht]
\centering
\includegraphics[width=\linewidth]{cap01/lazos_ab}
\caption{(a) Control en lazo abierto: nadie mira la salida. (b) Control en lazo cerrado: el sensor cierra el lazo. Notación del curso: $r$ referencia (\emph{set point}), $e$ error, $u$ señal de control, $d$ perturbación, $y$ variable controlada, $y_m$ medición, $n$ ruido de medición.}
\label{fig:in-lazos}
\end{figure}

Una tostadora, un semáforo de tiempos fijos o una lavadora que sigue su programa trabajan en lazo abierto: son baratos y no necesitan sensor, pero dependen de que el mundo se comporte como se supuso al calibrarlos. Un termostato, el control de crucero de un carro o el regulador de una máquina de vapor trabajan en lazo cerrado. El ejemplo siguiente pone números a la diferencia.

\begin{ejemplo}[label=ej:in-tanque]{Llenar un tanque sin mirarlo y mirándolo}
Un tanque de área $A=0.5$\,m$^2$ recibe un caudal $u(t)$ (m$^3$/min) que fijamos con una válvula, puede recibir lluvia $d(t)$ y descarga por una válvula de resistencia $R=2$\,min/m$^2$. Alrededor del punto de operación, el nivel $h$ (m) cumple
\begin{equation}
  A\,\frac{\dif h}{\dif t}=u+d-\frac{h}{R}
  \quad\Longrightarrow\quad
  \tau\frac{\dif h}{\dif t}+h=K\,(u+d),\qquad K=R=2,\quad \tau=AR=1\text{ min}.
  \label{eq:in-tanque}
\end{equation}
Queremos $h=h_r=1$\,m. Compare dos estrategias: (i) lazo abierto calibrado, $u=h_r/K$; (ii) lazo cerrado, $u=h_r/K+K_c\,(h_r-h)$ con $K_c=5$. Evalúe el nivel final en tres casos: sin perturbación, con lluvia $d=0.2$\,m$^3$/min y con la válvula de salida parcialmente obstruida ($R$ sube un $20\%$, a $2.4$).
\solucion En estado estacionario $\dif h/\dif t=0$ y \eqref{eq:in-tanque} da $h=R\,(u+d)$.

\emph{Lazo abierto.} Con $u=h_r/2=0.5$: sin perturbación $h=2(0.5)=1$, perfecto; con lluvia $h=2(0.5+0.2)=1.4$\,m; con la válvula obstruida $h=2.4(0.5)=1.2$\,m. El lazo abierto es exacto \emph{solo} si el modelo es exacto y no hay perturbaciones.

\emph{Lazo cerrado.} Ahora $u$ depende de $h$: $h=R\bigl(0.5+K_c(1-h)+d\bigr)$, de donde
\begin{equation}
  h=\frac{R\,(0.5+K_c+d)}{1+RK_c}.
  \label{eq:in-tanquecerrado}
\end{equation}
Sin perturbación $h=1$; con lluvia $h=11.4/11=1.036$\,m; con la válvula obstruida $h=2.4(5.5)/13=1.015$\,m. La lluvia movía el nivel $0.4$\,m en lazo abierto y ahora lo mueve $0.4/11=0.036$\,m: \emph{once veces menos}. El $11$ es $1+KK_c$, y al producto $KK_c=10$ lo llamaremos \clave{ganancia de lazo}. La \cref{fig:in-tanque} muestra las respuestas completas.
\end{ejemplo}

\begin{figure}[ht]
\centering
\includegraphics[width=\linewidth]{cap01/f_tanque_lazos}
\caption{Ejemplo~\ref{ej:in-tanque}. El lazo abierto (coral) llega exacto mientras el mundo coincide con la calibración; la retroalimentación (azul) reduce el efecto de la lluvia y de la obstrucción por un factor $1+KK_c$ y, de paso, responde más rápido.}
\label{fig:in-tanque}
\end{figure}

\codigorepo{cap01_introduccion/C1_Ej1_TanqueLazos.m}{cerrado}{ejemplo \ref{ej:in-tanque}}

Hay un segundo regalo, menos evidente: el lazo cerrado es más \emph{rápido}. Si se reemplaza la ley de control en \eqref{eq:in-tanque} se obtiene $\tau\dot h+(1+KK_c)\,h=\ldots$, una constante de tiempo de $\tau/(1+KK_c)=1/11$\,min$\approx5.5$\,s en lugar de $1$\,min. Nada es gratis, sin embargo: en el instante del escalón la válvula pasa de $0.5$ a $0.5+K_c=5.5$\,m$^3$/min (\cref{fig:in-esfuerzo}). Si la válvula no da más de, digamos, $2$\,m$^3$/min, se \emph{satura} y la respuesta real será más lenta que la calculada.

\begin{figure}[ht]
\centering
\includegraphics[width=0.7\linewidth]{cap01/f_tanque_esfuerzo}
\caption{El precio de la retroalimentación del ejemplo~\ref{ej:in-tanque} con lluvia: el caudal pedido salta a $5.5$\,m$^3$/min y luego se reduce para compensar la lluvia.}
\label{fig:in-esfuerzo}
\end{figure}

\begin{ideaclave}
\textbf{Lo que compra la retroalimentación.}
Con una ganancia de lazo $KK_c$: las perturbaciones y los errores del modelo se reducen aproximadamente por $1+KK_c$, y la respuesta se acelera en la misma proporción. El precio: más esfuerzo de control, un sensor (con su ruido) y el riesgo de inestabilidad cuando la planta tiene más dinámica que la de este tanque. Con un controlador proporcional queda un error pequeño pero no nulo; anularlo es trabajo de la acción integral (capítulos~6 a~8).
\end{ideaclave}

\begin{modelo}[TanqueLazos.slx]
\texttt{crearModelo\_TanqueLazos.m} construye el modelo de la versión anterior de estas notas, ahora con las dos ramas: el mismo tanque en lazo abierto (bloque \texttt{Calibracion}, que usted ajusta a mano) y en lazo cerrado (bloque \texttt{Kc}). La lluvia se activa con \texttt{Lluvia} $=1$ y la obstrucción cambiando \texttt{R} en el espacio de trabajo. Primero ajuste \texttt{Calibracion} para que el nivel llegue a $1$ sin lluvia; luego active la lluvia e intente, \emph{a mano}, mantener el nivel. ¿Pudo? Compare con la rama retroalimentada.
\end{modelo}

\begin{pausa}
En el ejemplo~\ref{ej:in-tanque}, ¿qué $K_c$ hace falta para que la lluvia mueva el nivel menos de $1$\,cm? ¿Cuánto vale entonces el caudal en el instante del escalón?
\respuesta{La desviación es $Kd/(1+KK_c)=0.4/(1+2K_c)<0.01$, es decir $K_c>19.5$. Con $K_c=20$ la desviación es $0.0098$\,m y el caudal inicial es $0.5+20=20.5$\,m$^3$/min: cuarenta veces el de lazo abierto. Pedir mucha precisión con un controlador proporcional sale caro.}
\end{pausa}

\subsection{Tres nombres parecidos}
\label{sec:in-nombres}

En español usamos ``lazo'' y ``malla'' para cosas distintas, y conviene fijarlas desde el principio porque los libros en inglés usan \emph{open loop} para las tres.

\begin{definicion}[label=def:in-nombres]{Control en lazo abierto, malla abierta y lazo abierto}
\begin{itemize}[leftmargin=1.2em]
  \item \clave{Control en lazo abierto} (\emph{open-loop control}): una \emph{estrategia} sin retroalimentación, como la de la \cref{fig:in-lazos}(a).
  \item \clave{Malla abierta} $\Gma(s)$: en un lazo cerrado, el producto de todo lo que recorre una señal al dar la vuelta completa, sin el signo del comparador. Si el camino directo es $G(s)$ y la retroalimentación $\Hs(s)$, $\Gma(s)=G(s)\Hs(s)$. Es la que se usa para estudiar estabilidad y la que se dibuja en el diagrama de Bode (en inglés, \emph{loop transfer function}, $L$).
  \item \clave{Lazo abierto} $\Gla(s)$: la función que relaciona el error con la salida cuando el lazo se reescribe con retroalimentación unitaria. Se usa para el tipo del sistema y los errores en estado estacionario.
\end{itemize}
\end{definicion}

Con $\Hs(s)=1$ la malla abierta y el lazo abierto coinciden, y por eso muchos textos los confunden; en el capítulo~\ref{cap:rf} (\cref{sec:rf-lenguaje}) se ve por qué no son lo mismo cuando el sensor tiene dinámica. La palabra ``abierto'' en los dos últimos nombres no quiere decir que el sistema funcione sin retroalimentación: describe qué parte del lazo cerrado se está mirando.

%% ---------------------------------------------------------------------
\section{Desempeño, análisis y diseño}
\label{sec:in-desempeno}
%% ---------------------------------------------------------------------

¿Cómo se juzga si un sistema de control es bueno? Tomemos un ascensor~\citep{Nis11}. Oprimir el botón del cuarto piso es la \emph{entrada}; la posición de la cabina es la \emph{salida}. El ascensor arranca, acelera, frena y se detiene. Nos importan tres cosas (\cref{fig:in-ascensor}):

\begin{itemize}
  \item La \clave{respuesta transitoria}: cómo llega. Si es muy rápida, el pasajero sufre aceleraciones incómodas o un sobrepaso (la cabina se pasa del piso y regresa); si es muy lenta, se agota su paciencia.
  \item El \clave{error en estado estacionario}: dónde se queda. Si la cabina se detiene unos centímetros por debajo del piso, alguien se va a tropezar.
  \item La \clave{estabilidad}: que se quede. Un ascensor que oscila alrededor del piso sin detenerse no sirve, por bueno que sea lo demás.
\end{itemize}

\begin{figure}[ht]
\centering
\includegraphics[width=0.85\linewidth]{cap01/f_ascensor}
\caption{Respuesta de un ascensor (modelo de segundo orden con valores exagerados para que se vean): sobrepaso de $9.5\%$, tiempo de establecimiento de $7.4$\,s y un error final de $0.1$ piso.}
\label{fig:in-ascensor}
\end{figure}

\codigorepo{cap01_introduccion/C1_Ej2_Ascensor.m}{ascensor}{respuesta del ascensor de la \cref{fig:in-ascensor}}

Estas tres cosas tienen nombre y número, y aprender a calcularlos es buena parte del curso: la estabilidad en el capítulo~5, el transitorio de primer y segundo orden en los capítulos~2 y~3 y el error en el capítulo~6.\nota{El ejemplo del ascensor tiene un video en el MOOC \emph{Introducción al análisis de sistemas de control con MATLAB}, en Coursera: \url{https://www.coursera.org/learn/introduccion-al-analisis-de-sistemas-de-control-con-matlab}.}

\begin{definicion}[label=def:in-diseno]{Análisis y diseño}
El \clave{análisis} determina el desempeño de un sistema dado: su estabilidad, su respuesta transitoria y su error. El \clave{diseño} es el medio por el cual ese desempeño se logra o se modifica: cambiando parámetros del sistema o agregando un subsistema nuevo (un controlador, un compensador) hasta cumplir las \clave{especificaciones}.
\end{definicion}

Análisis y diseño son dos direcciones del mismo camino: el análisis va del sistema al desempeño; el diseño, del desempeño deseado al sistema. Por eso el curso se llama \emph{análisis} pero termina diseñando controladores: no se puede diseñar lo que no se sabe analizar.

\begin{pausa}
El fabricante quiere eliminar el sobrepaso del ascensor de la \cref{fig:in-ascensor} sin cambiar nada más. ¿Qué cree que le pasará al tiempo de llegada?
\respuesta{Aumentará. Para no pasarse, el ascensor debe frenar antes y llegar ``despacio'' al piso; con la misma fuerza disponible tarda más. Transitorio rápido y sin sobrepaso es un compromiso, que en el capítulo~3 se cuantifica con el factor de amortiguamiento $\zeta$.}
\end{pausa}

%% ---------------------------------------------------------------------
\section{Cómo se plantea un problema de control}
\label{sec:in-problema}
%% ---------------------------------------------------------------------

\subsection{Entradas, salidas y mediciones}

Un sistema se separa de su entorno con una frontera imaginaria; lo que cruza la frontera hacia adentro son \emph{entradas}, lo que sale e influye sobre el entorno son \emph{salidas}~\citep{Bol15}. En los diagramas de bloques la frontera es la caja y las señales son las flechas. Pero para controlar hace falta distinguir más: no todas las entradas se pueden mover ni todas las salidas se pueden medir (\cref{fig:in-entradas}).

\begin{figure}[ht]
\centering
\includegraphics[width=0.8\linewidth]{cap01/proceso_entradas}
\caption{Un proceso con sus entradas y salidas clasificadas (adaptado de~\citep{Beq03}).}
\label{fig:in-entradas}
\end{figure}

\begin{definicion}[label=def:in-variables]{Variable manipulada y perturbación}
Una \clave{variable manipulada} es una entrada que el sistema de control (o un operador) puede ajustar. Una \clave{perturbación} es una entrada que afecta la salida pero que el sistema de control no puede ajustar; a veces se puede medir y a veces no. La salida que se quiere llevar a un valor deseado es la \clave{variable controlada}.
\end{definicion}

Lo que el controlador ``ve'' de la salida llega a través de un sistema de instrumentación con tres partes: el \emph{sensor}, en contacto con la variable física; el \emph{procesamiento de señal}, que convierte la salida del sensor en algo que se pueda transmitir (en la industria, típicamente una corriente de $4$ a $20$\,mA); y la \emph{presentación} o almacenamiento del dato. Al conjunto sensor--acondicionador se le llama a menudo \clave{transductor}: convierte una variable física en otra que se puede usar, casi siempre eléctrica. La medición nunca es exactamente la variable: hay errores de calibración, ruido y dinámica del propio sensor. Por eso en la \cref{fig:in-lazos} distinguimos $y$ de $y_m$.

\subsection{Los elementos de una estrategia de control}

\begin{receta}[Cómo plantear un problema de control~\citep{Beq03}]
\begin{enumerate}
  \item \textbf{Objetivos.} ¿Qué se quiere lograr? Mantener una variable en un valor, seguir una trayectoria, minimizar un costo.
  \item \textbf{Entradas.} Clasificarlas en manipuladas y perturbaciones. Pueden variar de forma continua o cambiar solo en ciertos instantes.
  \item \textbf{Salidas.} Clasificarlas en medidas y no medidas.
  \item \textbf{Restricciones.} \emph{Duras}, que no se pueden violar (un caudal entre un mínimo y un máximo, una válvula entre $0$ y $100\%$), y \emph{blandas}, que conviene respetar pero cuya violación no es peligrosa (la composición de un producto).
  \item \textbf{Operación.} \emph{Continua}, con periodos largos en condiciones parecidas, o \emph{por lotes} (\emph{batch}), con periodos cortos en los que las condiciones cambian mucho, como la fermentación de la cerveza.
  \item \textbf{Seguridad, ambiente y economía.} Qué no puede pasar nunca, qué se puede emitir, cuánto cuesta cada unidad de energía o de producto.
  \item \textbf{Estructura del controlador.} \emph{Retroalimentación}: se mide la variable controlada, se compara con la referencia y se ajusta la manipulada. \emph{Prealimentación}: se mide la perturbación y se ajusta la manipulada \emph{antes} de que la salida se afecte.
\end{enumerate}
\end{receta}

Escoger qué mover y qué medir no es trivial. Una buena \textbf{entrada de control} está \emph{disponible} (influye directa e inmediatamente sobre la salida), \emph{responde rápido}, permite ajustes \emph{finos} y tiene un efecto de \emph{signo conocido} y constante: si en un punto de operación abrir la válvula sube el nivel y en otro lo baja, el controlador no sabe para dónde ir. Una buena \textbf{salida de control} se puede medir de forma precisa, rápida y confiable, es \emph{relevante} (una buena medida de lo que realmente nos importa), responde rápido a lo que pasa en el proceso y es \emph{suave}, es decir, no necesita filtros fuertes.

\begin{definicion}[label=def:in-ganancia]{Ganancia del proceso}
La \clave{ganancia del proceso} mide la sensibilidad de la salida a la entrada: cuánto cambia la salida en estado estacionario por unidad de cambio en la entrada. Si un aumento de $0.5$\,kW en la potencia de un calentador sube la temperatura de un fluido $10$\,°C, la ganancia es $10/0.5=20$\,°C/kW.
\end{definicion}

La ganancia tiene unidades y signo, y ambos importan: en el ejemplo~\ref{ej:in-tanque} era $K=2$\,m por m$^3$/min. Una vez escogida la estructura, falta el \emph{algoritmo} del controlador y sus parámetros. En este curso los escogeremos \emph{a partir de un modelo} del proceso: por eso los capítulos~2 a~4 están dedicados a construir modelos.

\subsection{Ejemplo: darse una ducha}

Darse una ducha es un problema de control multivariable que resolvemos a diario (\cref{fig:in-ducha}).

\begin{figure}[ht]
\centering
\includegraphics[width=\linewidth]{cap01/ducha}
\caption{La ducha como problema de control: dos válvulas, dos variables que importan (temperatura y caudal) y un sensor muy bueno, la piel.}
\label{fig:in-ducha}
\end{figure}

\begin{itemize}
  \item \textbf{Objetivos:} estar limpio, refrescarse y sentirse bien; tomamos el último: una temperatura y un caudal agradables.
  \item \textbf{Manipuladas:} la apertura de las válvulas de agua caliente y fría (y, si se quiere, la posición del cuerpo).
  \item \textbf{Perturbaciones:} caídas de presión, cambios en la temperatura del calentador.
  \item \textbf{Salidas medidas:} temperatura y caudal del agua cuando toca la piel.
  \item \textbf{Restricciones:} las válvulas abren entre un mínimo y un máximo; por lo tanto hay temperaturas mínima y máxima alcanzables.
  \item \textbf{Estructura:} es \emph{multivariable}: cada válvula mueve la temperatura \emph{y} el caudal, así que hay que coordinarlas. Retroalimentamos cuando la piel siente frío y abrimos la caliente; prealimentamos cuando oímos que alguien descarga el sanitario y, antes de quemarnos, cerramos un poco la caliente.
\end{itemize}

\subsection{Ejemplo: un tanque de almacenamiento}
\label{sec:in-surge}

En una planta, los tanques de almacenamiento (\emph{surge tanks}) amortiguan las diferencias entre lo que produce una unidad y lo que consume la siguiente~\citep{Beq03}. Entra un caudal $F_1$ desde el proceso~1 y sale un caudal $F_2$ hacia el proceso~2. El \textbf{objetivo} es mantener el nivel $h$ dentro de unos límites alrededor de un valor deseado $h_{sp}$ (\emph{set point}); las \textbf{entradas} son $F_1$ y $F_2$; la \textbf{salida} es el nivel, que suponemos medible; la \textbf{operación} es continua; y la \textbf{seguridad} pide no rebosar ni vaciar el tanque y que el fluido no dañe tuberías ni paredes. Cuál entrada es manipulada y cuál perturbación \emph{depende de quién manda} (\cref{fig:in-escenarios}).

\begin{figure}[ht]
\begin{anchocompleto}
\centering
\includegraphics[width=\linewidth]{cap01/tanque_escenarios}
\end{anchocompleto}
\caption{Tres estructuras para el mismo tanque, en símbolos de la norma ISA-5.1: LT transmisor de nivel, LC controlador de nivel, FT transmisor de flujo, FY cálculo o relé; en línea discontinua, las señales.}
\label{fig:in-escenarios}
\end{figure}

\paragraph{Escenario 1: el proceso 2 fija $F_2$}
Si el proceso~2 es, por ejemplo, una caldera que toma el agua que necesita según la demanda de vapor, $F_2$ es una perturbación para el tanque y la variable manipulada es $F_1$. El transmisor LT envía el nivel al controlador LC, que lo compara con $h_{sp}$ y manda una señal de presión $P_v$ a la válvula de entrada. La \cref{fig:in-esquemas}(a) muestra el mismo lazo como diagrama de bloques. Note la diferencia: el diagrama de instrumentación dibuja \emph{equipos}; el diagrama de bloques dibuja \emph{flujo de información}, y cada bloque es un sistema dinámico.

\paragraph{Escenario 2: el proceso 1 fija $F_1$}
Si el proceso~1 produce un compuesto que el proceso~2 consume, $F_1$ es la perturbación y se manipula $F_2$ con la válvula de salida. El diagrama de bloques es el mismo intercambiando $F_1$ y $F_2$ (y el signo con el que entran al tanque).

\paragraph{Escenario 3: prealimentación}
Si se puede medir $F_1$, se puede mover $F_2$ en la misma cantidad \emph{antes} de que el nivel cambie (\cref{fig:in-esquemas}(b)). Es la ventaja de la prealimentación: corrige antes del daño. Su debilidad es que no mira el nivel: cualquier error en la medición de $F_1$ o en el modelo se acumula sin que nadie lo corrija. Por eso en la práctica \emph{se combinan} las dos estructuras.

\begin{figure}[ht]
\centering
\includegraphics[width=0.92\linewidth]{cap01/esquema_escenarios}
\caption{Los escenarios 1 y 3 de la \cref{fig:in-escenarios} como diagramas de bloques.}
\label{fig:in-esquemas}
\end{figure}

\begin{ejemplo}[label=ej:in-prealimentacion]{Prealimentación con un sensor imperfecto}
Un tanque de área $A=2$\,m$^2$ cumple $A\,\dot h=F_1-F_2$ (la salida la impone una bomba). En $t=2$\,min, $F_1$ sube $0.5$\,m$^3$/min. Compare, en desviación $\dev{h}$ respecto al nivel inicial: (i) prealimentación con un transmisor de flujo que lee un $10\%$ menos, $\dev{F}_2=0.9\,\dev{F}_1$; (ii) retroalimentación proporcional, $\dev{F}_2=K_c\dev{h}$ con $K_c=1$\,m$^2$/min; (iii) las dos juntas.
\solucion (i) $A\dot{\dev{h}}=0.5-0.45=0.05$: el nivel sube $0.025$\,m/min \emph{para siempre}; en $18$\,min ya subió $0.45$\,m. (ii) $A\dot{\dev{h}}=0.5-K_c\dev{h}$: el nivel se estabiliza en $0.5/K_c=0.5$\,m con constante de tiempo $A/K_c=2$\,min. (iii) $A\dot{\dev{h}}=0.05-K_c\dev{h}$: se estabiliza en $0.05$\,m, diez veces menos que con retroalimentación sola. La \cref{fig:in-prealimentacion} lo muestra: la prealimentación hace casi todo el trabajo y la retroalimentación corrige lo que aquella no sabía.
\end{ejemplo}

\begin{figure}[ht]
\centering
\includegraphics[width=0.85\linewidth]{cap01/f_prealimentacion}
\caption{Ejemplo~\ref{ej:in-prealimentacion}: la prealimentación sola deriva, la retroalimentación sola deja una desviación grande, juntas funcionan.}
\label{fig:in-prealimentacion}
\end{figure}

\codigorepo{cap01_introduccion/C1_Ej3_Prealimentacion.m}{combinada}{ejemplo \ref{ej:in-prealimentacion}}

\begin{pausa}
En el escenario~1, si $F_2$ es una perturbación que \emph{se puede medir}, ¿cómo combinaría prealimentación y retroalimentación?
\respuesta{El controlador LC sigue corrigiendo con el nivel, y a su salida se le suma un término que mueve $F_1$ en la misma cantidad que cambie el $F_2$ medido. Así un cambio de $F_2$ se compensa de inmediato y LC solo corrige lo que la compensación no logró.}
\end{pausa}

%% ---------------------------------------------------------------------
\section{Una breve historia del control}
\label{sec:in-historia}
%% ---------------------------------------------------------------------

La historia del control automático ha sido contada con cuidado por Otto Mayr, para los orígenes~\citep{May71a,May12}; por Stuart Bennett, para el periodo industrial~\citep{Ben86,Ben93}; y por David Mindell, para los años de la guerra~\citep{Min02}. Hay números especiales de revistas dedicados al tema~\citep{Bus96,BeBu02}, una reseña breve y muy completa de Chris Bissell~\citep{Bis09} y una mirada de conjunto de Karl J.\ Åström y P.\,R.\ Kumar, \emph{Control: a perspective}~\citep{AsKu14}. Lo que sigue es un resumen de esas fuentes, con la referencia del trabajo original donde apareció cada idea cuando la hay. La \cref{fig:in-linea} resume la cronología.

\begin{figure}[ht]
\begin{anchocompleto}
\centering
\includegraphics[width=0.92\linewidth]{cap01/f_linea_tiempo}
\end{anchocompleto}
\caption{Algunos hitos de la historia del control. Arriba, de los reguladores de flotador a la teoría de reguladores y los pilotos automáticos; abajo, del amplificador de Black al control digital.}
\label{fig:in-linea}
\end{figure}

\notamargen{\textbf{Un patrón que se repite.} Primero llega el invento; luego aparece un problema de estabilidad que nadie entiende; y solo entonces llega la teoría. Watt $\to$ Maxwell, los barcos $\to$ Minorsky, el amplificador de Black $\to$ Nyquist.}

%% ---------------------------------------------------------------------
\subsection{De los relojes de agua a la máquina de vapor}

Como tecnología, el control tiene más de dos mil años. El primer mecanismo retroalimentado del que hay registro es el regulador de flotador de la \emph{clepsidra} (``ladrona de agua'') de Ktesibios de Alejandría, en el siglo~III a.C.~\citep{May71a}. Las sociedades que crecían necesitaban horarios, y los relojes de sol no sirven de noche ni con nubes. Un reloj de agua mide el tiempo por el volumen que cae a un recipiente, y para que sea preciso el caudal debe ser constante, lo que exige un nivel constante en el recipiente que gotea. Ktesibios lo logró con un flotador cónico que tapa la entrada cuando el nivel sube y la destapa cuando baja (\cref{fig:in-clepsidra}): sensor, comparador y actuador en una sola pieza.\nota{Ktesibios era hijo de un barbero y, según la tradición, empezó inventando un espejo ajustable para la barbería de su padre. Algunas fuentes lo hacen primer director del Museo de Alejandría.} Filón de Bizancio usó la misma idea, hacia el 250~a.C., para mantener constante el nivel de aceite de una lámpara.

\begin{figure}[ht]
\centering
\includegraphics[width=0.6\linewidth]{cap01/regulador_flotador}
\caption{El principio de la clepsidra de Ktesibios: el flotador mantiene constante el nivel del recipiente regulador, el caudal de salida es constante y el nivel del recipiente receptor sube de manera uniforme con el tiempo.}
\label{fig:in-clepsidra}
\end{figure}

En el siglo~I, Herón de Alejandría describió en su \emph{Pneumatica} varios reguladores de nivel y válvulas de flotador. Los ingenieros del mundo islámico los perfeccionaron: los hermanos Banū Mūsā en el siglo~IX y al-Jazarī en su \emph{Libro del conocimiento de los ingeniosos mecanismos} (1206) describen relojes y fuentes con reguladores de flotador~\citep{May71a}. En Europa, en cambio, estas ideas casi desaparecen entre la Edad Media y el Barroco: la \emph{Pneumatica} de Herón no se imprimió en latín sino hasta 1575~\citep{May12}.

El primer sistema retroalimentado de la Europa moderna que no copia los antiguos es el regulador de temperatura de Cornelis Drebbel (1572--1633), hacia 1620, para un horno que incubaba huevos: la dilatación del líquido de un termómetro empujaba una varilla que cerraba la entrada de aire al fuego~\citep{Bis09}. En 1681 Denis Papin agregó a su olla a presión una válvula de seguridad, antecedente de los reguladores de presión~\citep{May71a}.

\paragraph{Los molinos y el regulador centrífugo}
El siglo~XVIII trajo dos problemas de regulación en los molinos de viento: orientar las aspas hacia el viento y mantener la velocidad de las piedras. Para lo primero, Edmund Lee patentó en 1745 la \emph{cola de abanico}, un rotor auxiliar que gira el molino hasta enfrentarlo al viento. Para lo segundo se usaron péndulos centrífugos; Thomas Mead los patentó en 1787 para separar las piedras del molino según la velocidad~\citep{Bis09}. Hacia 1788 James Watt, a sugerencia de su socio Matthew Boulton, adaptó el regulador centrífugo a la máquina de vapor (\cref{fig:in-watt}): dos bolas que giran con el eje se abren cuando la máquina se acelera y, mediante un juego de palancas, cierran la válvula del vapor. El regulador combina sensado, actuación y control en un mismo mecanismo, y se convirtió en el símbolo de la Revolución Industrial. Watt no lo patentó: el principio ya se usaba en los molinos~\citep{May71a,Bis09}.

\begin{figure}[ht]
\centering
\includegraphics[width=\linewidth]{cap01/foto_watt}
\caption{Arriba a la izquierda, un modelo de la máquina de vapor de Watt; a la derecha, las bolas del regulador centrífugo; abajo, el autor junto a la máquina de vapor que conserva el Science Museum de Londres. Fotos del autor.}
\label{fig:in-watt}
\end{figure}

Los reguladores de Watt tenían un defecto: la velocidad de equilibrio cambiaba con la carga, porque la corrección era proporcional al exceso de velocidad. En el siglo~XIX se buscaron reguladores que llevaran la velocidad al mismo valor con cualquier carga. El \emph{regulador cronométrico} de los hermanos Werner y William Siemens lo logró con un engranaje diferencial que integraba la diferencia entre la velocidad de la máquina y la de un reloj: es la \clave{acción integral}~\citep{Bis09}. Además, a medida que las máquinas se volvieron más rápidas y los reguladores más sensibles, apareció un fenómeno inesperado, la \emph{caza} (\emph{hunting}): la máquina oscilaba alrededor de la velocidad deseada sin detenerse. George Biddell Airy, astrónomo real, publicó en 1840 el primer análisis matemático de un sistema retroalimentado al estudiar las oscilaciones del mecanismo que movía un telescopio~\citep{Air40}.

%% ---------------------------------------------------------------------
\subsection{La teoría de los reguladores (1868--1895)}

La teoría empieza con el artículo de James Clerk Maxwell \emph{On governors}, recibido por la Royal Society el 20 de febrero de 1868~\citep{Max68}. Maxwell escribió las ecuaciones diferenciales del regulador y de la máquina, las \emph{linealizó} alrededor del punto de operación y mostró que la estabilidad depende de que las raíces de la ecuación característica tengan parte real negativa. Dedujo las condiciones para sistemas de segundo y tercer orden, mostró las ventajas de la acción integral y distinguió entre los \emph{moderadores}, como el de Watt, que dejan un error que depende de la carga, y los verdaderos \emph{reguladores}, que no lo dejan~\citep{Kan16}. Es, en esencia, la diferencia entre la acción proporcional y la integral que veremos en los capítulos~6 y~7.

\begin{historia}[Maxwell, los anillos de Saturno y la estabilidad]
Antes de los reguladores, Maxwell había pensado en estabilidad por otra razón. En 1857 ganó el premio Adams de la Universidad de Cambridge con un ensayo sobre los anillos de Saturno, en el que descartó los modelos de anillos sólidos o fluidos porque sus movimientos no eran estables: si los anillos existen desde hace tanto tiempo, el modelo correcto tiene que ser estable. Usar la estabilidad como criterio para aceptar o rechazar un modelo es una idea que hoy usamos en identificación de sistemas~\citep{Bit15}. Norbert Wiener lo consideró el precursor de la teoría de la retroalimentación~\citep{May71b}.
\end{historia}

El artículo de Maxwell era difícil y tuvo poco efecto sobre los ingenieros de su tiempo; pero planteó una pregunta matemática: ¿cómo saber si todas las raíces de un polinomio de grado $n$ tienen parte real negativa sin calcularlas? Edward J.\ Routh la respondió y ganó con ello el premio Adams de 1877~\citep{Rou77}.\nota{Routh y Maxwell nacieron el mismo año, 1831, y presentaron juntos el examen de matemáticas de Cambridge (los \emph{tripos}) en 1854. Routh quedó primero (\emph{Senior Wrangler}) y Maxwell segundo. Ambos compartieron ese año el premio Smith, la primera vez que se dividió entre dos candidatos. (La anécdota del sirviente que, enviado a ver quién había quedado segundo, responde ``¡usted, señor!'' se cuenta de William Thomson ---lord Kelvin--- en 1845, no de Maxwell.)} En Rusia, Ivan Vyshnegradskii analizó en 1876 el regulador de la máquina de vapor sin conocer el trabajo de Maxwell, despreciando la fricción, linealizando y resolviendo el caso de tercer orden con un método gráfico~\citep{Vys76}; terminaría su carrera como ministro de finanzas del Imperio ruso. En Suiza, Aurel Stodola necesitaba en 1893 diseñar el regulador de una turbina hidráulica de alta presión para Davos; sus modelos eran de orden alto y le pidió ayuda a su colega del Politécnico de Zúrich, Adolf Hurwitz, quien publicó en 1895 un criterio equivalente al de Routh, escrito con determinantes~\citep{Bis89,Hur95}.\nota{Charles Hermite había resuelto en 1856 un problema matemático equivalente, sin relacionarlo con los reguladores.} En 1892 Aleksandr Lyapunov publicó su teoría de la estabilidad del movimiento, válida también para sistemas no lineales~\citep{Lya92}. Volveremos a todos ellos en el capítulo~5.

%% ---------------------------------------------------------------------
\subsection{Barcos, aviones, fábricas y teléfonos (1900--1940)}

El término \emph{servomoteur} lo acuñó en 1873 el ingeniero francés Joseph Farcot, que construía mecanismos hidráulicos para mover el timón de los barcos~\citep{Bis09}. A comienzos del siglo~XX el giroscopio cambió la navegación: Elmer Sperry fundó en 1910 la Sperry Gyroscope Company, que fabricó girocompases y luego pilotos automáticos para barcos que combinaban el giroscopio con un motor eléctrico conectado al timón. Su hijo Lawrence llevó la idea a los aviones: el 18 de junio de 1914, en el \emph{Concours de la Sécurité en Aéroplane} de París, voló con las manos en alto mientras su mecánico caminaba sobre el ala y el estabilizador giroscópico mantenía el avión nivelado~\citep{Bis09}.

Para gobernar barcos, Nicolas Minorsky observó a los timoneles de la Marina de Estados Unidos y en 1922 propuso una ley de control con tres términos: uno proporcional al error de rumbo, uno a su integral y uno a su derivada; estudió además la estabilidad del lazo con un modelo lineal~\citep{Min22,Ben84}. Es la primera formulación teórica del controlador \clave{PID}.

En la industria de procesos, los controladores neumáticos se volvieron comunes en los años veinte. La Foxboro anunció en 1931 el \emph{Stabilog}, cuya salida era proporcional al error y a su integral, y la Taylor Instrument incorporó en 1935 la acción derivativa (\emph{pre-act}) a su controlador Fulscope. Una encuesta del gobierno de Estados Unidos estimó que entre 1925 y 1935 se vendieron allí más de $75\,000$ controladores automáticos, y a mediados de los treinta había unas $600$ empresas fabricantes de instrumentos~\citep{BenPID93}. Faltaba saber cómo ajustar los parámetros: en 1942 John Ziegler y Nathaniel Nichols, de la Taylor, publicaron las reglas de sintonización que estudiaremos en el capítulo~8~\citep{ZiNi42}.

En los teléfonos estaba pasando algo parecido. Las líneas transcontinentales necesitaban cientos de amplificadores en cascada, cuya distorsión se acumulaba. Harold S.\ Black resolvió el problema en 1927 con el amplificador con retroalimentación negativa~\citep{black34}; la patente, presentada en 1928, tardó nueve años en concederse porque la idea parecía contraria a la intuición~\citep{blackpat,Bri11}. Con mucha ganancia de lazo los amplificadores oscilaban, y Harry Nyquist dio en 1932 el criterio para predecir la estabilidad del lazo cerrado a partir de la respuesta en frecuencia de la malla abierta~\citep{Nyq32,Bri10}. Hendrik Bode convirtió esas ideas, entre 1940 y 1945, en un método de diseño con márgenes de ganancia y de fase~\citep{Bod40,bode45}. El capítulo~\ref{cap:rf} cuenta esta parte en detalle.

En el MIT, mientras tanto, el grupo de Vannevar Bush construyó entre 1928 y 1931 el \emph{analizador diferencial}, un computador analógico mecánico para resolver ecuaciones diferenciales (\cref{fig:in-bush}). Harold Hazen, discípulo y colaborador de Bush durante más de veinte años, publicó en 1934 la \emph{Theory of servo-mechanisms}, la primera presentación sistemática de cómo funcionan y cómo se diseñan los servomecanismos~\citep{Haz34}. Su estudiante Gordon Brown dirigió desde los primeros años de la guerra el Laboratorio de Servomecanismos del MIT. La \cref{fig:in-personajes} reúne a varios de los protagonistas de esta sección.

\begin{figure}[ht]
\centering
\includegraphics[width=0.85\linewidth]{cap01/foto_bush}
\caption{Analizador diferencial del Science Museum de Londres, construido en Mánchester en 1935 siguiendo el diseño de Vannevar Bush en el MIT. Foto del autor.}
\label{fig:in-bush}
\end{figure}

\begin{figure}[ht]
\centering
\includegraphics[width=0.85\linewidth]{cap01/personajes}
\caption{Arriba, de izquierda a derecha: Edward J.\ Routh, Vannevar Bush, James C.\ Maxwell y Elmer A.\ Sperry. Abajo: Aurel Stodola, Nicolas Minorsky, Hendrik W.\ Bode y Adolf Hurwitz. Imágenes de Wikimedia Commons.}
\label{fig:in-personajes}
\end{figure}

%% ---------------------------------------------------------------------
\subsection{La guerra y la posguerra (1940--1960)}

La Segunda Guerra Mundial juntó a las tres comunidades: los ingenieros de procesos, los de servomecanismos y los de comunicaciones. Dirigir un cañón antiaéreo con un radar es un problema de servomecanismos con ruido, y en el Laboratorio de Radiación del MIT y en Bell Labs se unificaron los métodos de frecuencia. El resultado se publicó al terminar la guerra en el volumen \emph{Theory of Servomechanisms}~\citep{JaNiPh47}, el primer libro de texto moderno del área. En 1948 Norbert Wiener publicó \emph{Cybernetics}, que extendió la idea de retroalimentación a los seres vivos y a la sociedad~\citep{Wie48}, y Walter Evans presentó el lugar de las raíces~\citep{Eva50}.

Los computadores digitales traían otra pregunta: qué pasa cuando el controlador solo mide en ciertos instantes. En la Universidad de Columbia, John Ragazzini y Lotfi Zadeh introdujeron en 1952 la transformada $z$ para los sistemas muestreados~\citep{RaZa52}; por ese grupo pasaron también Eliahu Jury, Gene Franklin y Rudolf Kalman. La optimización avanzaba en paralelo: George Dantzig creó en 1947 el método símplex de programación lineal, y en la RAND Corporation Richard Bellman desarrolló en los años cincuenta la programación dinámica~\citep{Bel57}. En 1955 la Glenn L.\ Martin Company fundó en Baltimore el Research Institute for Advanced Studies (RIAS), donde Solomon Lefschetz reunió un grupo de matemáticos dedicados a la teoría de sistemas dinámicos y de control, entre ellos Kalman.\nota{Lefschetz tiene un lazo con la Universidad de los Andes. Mario Laserna, su fundador, organizó a finales de los años cuarenta en Princeton una reunión en honor de la nueva universidad, a la que asistieron Albert Einstein, John von Neumann, Oskar Morgenstern y Lefschetz, entre otros~\citep{IASLaserna}; en el verano de 1950 invitó a von Neumann y a Lefschetz a dictar conferencias en Bogotá~\citep{Les19}.}

El área se organizó institucionalmente. En 1945 se fundó en Estados Unidos la Instrument Society of America (ISA). En 1954 el Institute of Radio Engineers creó su grupo profesional de control automático, origen de la actual IEEE Control Systems Society. Y el 12 de septiembre de 1957, delegados de 18 países fundaron en París la International Federation of Automatic Control (IFAC); su primer presidente fue Harold Chestnut, y su primer congreso mundial se reunió en Moscú en 1960, en plena Guerra Fría. Desde entonces se reúne cada tres años~\citep{IFAChist,Bis09}.

%% ---------------------------------------------------------------------
\subsection{La edad de oro (1960--2000)}

La carrera espacial necesitaba guiar cohetes y satélites con muchas variables, restricciones de combustible y mediciones ruidosas, y los computadores empezaban a permitir cálculos que antes eran impensables. En el congreso de Moscú, Kalman presentó la teoría de sistemas en espacio de estado, con los conceptos de controlabilidad y observabilidad~\citep{Kal60}; el mismo año publicó el filtro que lleva su nombre~\citep{Kal60b}, extendido en 1961 con Richard Bucy~\citep{KaBu61}. El grupo de Lev Pontryagin publicó en 1962 el principio del máximo~\citep{Pon62}, que junto con la programación dinámica de Bellman fundó el \emph{control óptimo}. A esta manera de trabajar, con ecuaciones de estado en el dominio del tiempo, se le llamó \emph{control moderno}, para distinguirla del \emph{clásico} de Nyquist, Bode y Evans. El capítulo~\ref{cap:re} es nuestra puerta de entrada.

Del tronco común salieron muchas ramas: identificación de sistemas, control adaptativo, control no lineal, control robusto, sistemas de eventos discretos, control predictivo. En la industria el cambio fue el paso a lo digital. En 1968--1969 Modicon diseñó para la división Hydramatic de General Motors el primer controlador lógico programable (PLC), que reemplazó gabinetes enteros de relés; en 1975 Honeywell (TDC~2000) y Yokogawa (CENTUM) lanzaron los primeros sistemas de control distribuido (DCS); y la señal de $4$ a $20$\,mA reemplazó a la presión de $3$ a $15$\,psi de la instrumentación neumática. En los automóviles llegaron los frenos antibloqueo (ABS) y luego el control de estabilidad. El proyecto PATH de California, creado en 1986, demostró el 7 de agosto de 1997, en la autopista I-15 de San Diego, un pelotón de ocho automóviles que viajaban automáticamente a velocidad de autopista separados por unos $6$\,m~\citep{PATH97}.

%% ---------------------------------------------------------------------
\subsection{El control hoy}

Hoy el control está en casi todo, la mayoría de las veces escondido: en los discos duros y los teléfonos, en las redes eléctricas con fuentes renovables, en los procesos biológicos, en la logística y en los vehículos autónomos~\citep{AsKu14}. Los problemas nuevos tienen que ver con sistemas en red, con muchos agentes que deciden por su cuenta y con datos abundantes; pero las preguntas de fondo son las de Maxwell: ¿es estable?, ¿qué tan bien sigue lo que se le pide?, ¿qué tan robusto es ante lo que no sabemos? Quijano y Cotrino recogen la historia del control y la automatización en Colombia y el estado actual del área en el país~\citep{QuCo23}.

\begin{pausa}
¿Qué tienen en común la clepsidra de Ktesibios, el regulador de Watt y el amplificador de Black? Identifique en cada uno el sensor, el comparador y el actuador.
\respuesta{Los tres son lazos cerrados con retroalimentación negativa. Clepsidra: el flotador mide el nivel, su posición respecto al orificio hace de comparador y él mismo abre o cierra la entrada (actuador). Watt: las bolas miden la velocidad, el punto de equilibrio de las palancas fija la referencia y la válvula de vapor es el actuador. Black: la red resistiva $\beta$ mide la salida, la resta a la entrada se hace en el propio amplificador, y los tubos de vacío son el actuador.}
\end{pausa}

%% ---------------------------------------------------------------------
\section{El mapa del curso}
\label{sec:in-ruta}
%% ---------------------------------------------------------------------

Estas notas reúnen en un solo lugar lo que está disperso en muchos libros, con una visión propia, en castellano, que pone el énfasis en el análisis y en la simulación. No pretenden reemplazar a los textos clásicos: en cada capítulo se citan los libros y artículos de donde salen las ideas, para que usted pueda ir más lejos.

\subsection{La ruta}

El curso sigue el orden natural de un proyecto de control: \emph{modelar}, \emph{analizar}, \emph{diseñar} (\cref{fig:in-mapa}).

\begin{figure}[ht]
\centering
\includegraphics[width=\linewidth]{cap01/mapa_curso}
\caption{Los capítulos del curso y cómo dependen unos de otros. En azul, \emph{modelar}: la planta como ecuación diferencial y función de transferencia; en verde, \emph{analizar}: estabilidad, transitorio y error; en coral y violeta, \emph{diseñar}: el controlador.}
\label{fig:in-mapa}
\end{figure}

\begin{itemize}
  \item \textbf{Modelar} (capítulos~2 a~4). Para controlar un proceso hay que entenderlo, y lo entendemos con un modelo físico. Trabajaremos con sistemas \emph{lineales e invariantes en el tiempo} (LTI), descritos por ecuaciones diferenciales ordinarias. Los sistemas reales tienen parámetros que cambian, no linealidades y mediciones tomadas en instantes discretos; pero alrededor de un punto de operación, casi siempre un equilibrio, un modelo lineal describe bien su comportamiento. La transformada de Laplace convierte esas ecuaciones en funciones de transferencia, que se manipulan con álgebra, y la representación de estado las organiza para trabajar en el tiempo.
  \item \textbf{Analizar} (capítulos~5, 6 y~10). Con el modelo preguntamos si el lazo es estable (polinomio característico, valores propios, criterio de Routh--Hurwitz), cuánto error deja en estado estacionario y cómo responde a señales de distintas frecuencias (diagramas de Bode, márgenes de estabilidad, robustez).
  \item \textbf{Diseñar} (capítulos~7 a~10). El controlador PID, el más usado en la industria: qué hace cada término, sus configuraciones, sus problemas típicos (como la saturación del integrador) y cómo se sintonizan sus parámetros. La retroalimentación de estado, que ubica todos los polos y abre la puerta al control moderno, con los conceptos de controlabilidad y observabilidad. Y el diseño de compensadores en frecuencia con los diagramas de Bode.
\end{itemize}

\begin{receta}[Lo que usted sabrá hacer al final del curso]
\begin{enumerate}
  \item \textbf{Respuesta transitoria:} modelar, simular y analizar señales y sistemas en el tiempo (espacio de estado) y en la frecuencia (transformada de Laplace, función de transferencia).
  \item \textbf{Error en estado estacionario:} calcularlo según el tipo de entrada y el tipo del sistema, con las constantes de error $\kp$ y $\kv$.
  \item \textbf{Estabilidad:} decidirla con el polinomio característico o los valores propios, y estudiar cómo cambia con los parámetros (robustez).
  \item \textbf{Controladores:} diseñar PID y retroalimentaciones de estado, simularlos y mejorarlos (\emph{anti-windup}, acción integral, control por modelo interno, compensadores en frecuencia).
\end{enumerate}
\end{receta}

\subsection{Cómo usar estas notas}

Cada capítulo empieza con sus objetivos y un código QR que lleva a su página web, donde están el PDF, los códigos, los videos y un explorador interactivo. A lo largo del texto encontrará cajas de colores:

\begin{itemize}
  \item \textbf{Definición} y \textbf{Ejemplo} (numeradas): lo que hay que saber y cómo se usa. Los ejemplos tienen solución.
  \item \textbf{Idea clave}: el resumen de una sección en pocas líneas. \textbf{Cuidado}: un error frecuente.
  \item \textbf{Pausa activa}: una pregunta corta para comprobar que entendió; las respuestas están al final del capítulo.
  \item \textbf{Nota histórica}, \textbf{Para profundizar} y \textbf{Procedimiento}: contexto, material opcional y recetas paso a paso.
  \item \textbf{Código} y \textbf{Simulink}: fragmentos de los programas del repositorio, con un enlace al archivo completo. Antes de usarlos, ejecute \cmd{iniciarCurso} en \matlab. Cada carpeta tiene un \cmd{verificarCapNN} que comprueba los números del capítulo.
\end{itemize}

Los problemas del final de cada capítulo están marcados con \etiquetaM{} si requieren \matlab/\simulink y con \etiquetaEstrella{} si son de profundización. La notación es la misma en todo el curso y está resumida en la sección de convenciones al comienzo de las notas.

\subsection{Unas lecciones para volver a leer}

Dennis Bernstein escribió una guía de lecciones que todo estudiante de control debería ir afianzando~\citep{Ber97}. Muchas de sus palabras no tendrán sentido todavía; la idea es que vuelva a esta lista al terminar cada capítulo y vea cuántas ya entiende. Esta es una selección, adaptada al curso:

\begin{enumerate}
  \item \textbf{La retroalimentación está en todas partes}: en los procesos industriales, en la biología, en la economía.
  \item \textbf{Un diagrama de bloques no es un circuito.} Sus flechas son flujo de información, no conexiones físicas.
  \item \textbf{Encuentre los equilibrios y linealice.} Los sistemas reales son no lineales, pero el modelo lineal describe bien lo que pasa cerca del punto de operación.
  \item \textbf{Verifique primero la estabilidad}, con el polinomio característico o con los valores propios del modelo linealizado.
  \item \textbf{Los sistemas estables tienen respuesta en frecuencia}: una senoide de entrada produce, en régimen permanente, una senoide de la misma frecuencia con otra amplitud y otra fase.
  \item \textbf{No olvide la malla abierta} $\Gma(s)$: todo lo que hay en el lazo (controlador, actuador, planta, sensor, transmisor), su estabilidad y el signo de la retroalimentación.
  \item \textbf{Después de la estabilidad, el desempeño lo es todo.} Si la planta ya es estable, el desempeño es la única razón para retroalimentar.
  \item \textbf{La ganancia alta tiene beneficios}: rechaza perturbaciones y reduce la sensibilidad a los errores del modelo, como vimos en el tanque.
  \item \textbf{Pero tiene precios}: amplifica el ruido del sensor, satura los actuadores y puede desestabilizar el lazo.
  \item \textbf{Cuidado con los polos poco amortiguados}, que producen resonancias y oscilaciones largas.
  \item \textbf{Asegure la estabilidad nominal y luego la robusta}: que el lazo siga estable cuando la planta real no es la del modelo.
  \item \textbf{Respete lo inestable}~\citep{Ste03}: estabilizar una planta inestable impone límites a lo que se puede lograr.
  \item \textbf{Las no linealidades siempre están ahí}, el control multivariable no es trivial y, sobre todo, \textbf{la seguridad de las personas siempre está en juego}.
\end{enumerate}

%% ---------------------------------------------------------------------
\section{Para llevar}
\label{sec:in-cierre}
%% ---------------------------------------------------------------------

\begin{parallevar}
\begin{enumerate}
  \item Un sistema de control administra el comportamiento de otro sistema para que una salida siga una referencia o rechace perturbaciones. Analizar es determinar su desempeño; diseñar es lograrlo.
  \item Retroalimentar es medir, comparar con lo deseado y actuar sobre el error. Con una ganancia de lazo $KK_c$, las perturbaciones y los errores del modelo se reducen aproximadamente por $1+KK_c$. El precio: esfuerzo de control, ruido de medición y riesgo de inestabilidad.
  \item El control en lazo abierto es exacto solo si el modelo es exacto y no hay perturbaciones. La prealimentación corrige antes del daño; la retroalimentación corrige lo que no se sabía. En la práctica se combinan.
  \item ``Control en lazo abierto'' es una estrategia; la \emph{malla abierta} $\Gma=G\Hs$ y el \emph{lazo abierto} $\Gla$ son funciones de transferencia de un lazo cerrado.
  \item El desempeño se mide con la estabilidad, la respuesta transitoria y el error en estado estacionario.
  \item Un problema de control se plantea con objetivos, variables manipuladas y perturbaciones, salidas medidas y no medidas, restricciones, tipo de operación, consideraciones de seguridad y estructura del controlador.
  \item La teoría del control nació casi siempre de un problema de estabilidad: el regulador de Watt llevó a Maxwell, Routh y Hurwitz; los barcos, a Minorsky; los amplificadores, a Nyquist y Bode; la carrera espacial, a Kalman.
\end{enumerate}
\end{parallevar}

\begin{profundiza}[Correcciones respecto a la versión anterior del capítulo]
(i) Watt no patentó el regulador centrífugo; Mead lo había patentado para los molinos en 1787. (ii) El vuelo de Lawrence Sperry en París fue en 1914, no en 1912. (iii) Minorsky propuso el PID y analizó su estabilidad con un modelo lineal; no usó ubicación de polos. (iv) Bode publicó sus resultados entre 1940 y 1945, no en 1934; ese año fue el de la teoría de servomecanismos de Hazen (antes ``Hanzen''). (v) IFAC se fundó en 1957 con 18 países; RIAS en 1955. (vi) ``Stabilorg'' $\to$ Stabilog; ``Danzig'' $\to$ Dantzig; ``Vaneevar'' $\to$ Vannevar. (vii) Se retiró la atribución de la acción integral a los hermanos Périer, que no se encontró en las fuentes; se cita en su lugar el regulador cronométrico de los Siemens. (viii) El ejemplo del tanque ahora compara el lazo abierto con el cerrado y tiene números verificados con \texttt{herramientas/cap01\_numeros.py} y \cmd{verificarCap01.m}. (ix) La función $G(s)H(s)$ se llama \emph{malla abierta}, no ``función de transferencia en lazo abierto''.
\end{profundiza}

%% ---------------------------------------------------------------------
\section{Problemas}
\label{sec:in-problemas}
%% ---------------------------------------------------------------------

Los problemas marcados con \etiquetaM{} requieren \matlab/\simulink; los marcados con \etiquetaEstrella{} son de profundización.

\begin{problema}[lazo abierto o cerrado]\label{prob:in-clasificar}
Clasifique como control en lazo abierto o en lazo cerrado y, en los cerrados, identifique la variable controlada, el sensor, el actuador y una perturbación: (a) una tostadora con temporizador; (b) la nevera de su casa; (c) un semáforo de tiempos fijos; (d) un semáforo que cuenta los carros con una cámara; (e) la regulación de la temperatura del cuerpo humano; (f) una persona que se ducha (\cref{fig:in-ducha}); (g) el control de crucero de un carro. ¿Alguno usa prealimentación?
\end{problema}

\begin{problema}[el tanque, en detalle]\label{prob:in-tanque}
\etiquetaM{} Para el tanque del ejemplo~\ref{ej:in-tanque}: (a) reemplace la ley de lazo cerrado en \eqref{eq:in-tanque} y obtenga la ecuación diferencial del lazo cerrado y su constante de tiempo. (b) Calcule el nivel final con lluvia y con la válvula obstruida para $K_c=1$ y $K_c=10$, y compare con el factor $1+KK_c$. (c) Verifique (b) con \cmd{C1\_Ej1\_TanqueLazos}. (d) Limite el caudal a $0\le u\le2$\,m$^3$/min y repita la simulación con $K_c=10$: ¿qué cambia en el transitorio? ¿Y en el valor final? (e) Quite el término de calibración, $u=K_c(h_r-h)$: ¿cuánto vale el nivel final sin perturbación con $K_c=5$? ¿Por qué?
\end{problema}

\begin{problema}[plantear un problema]\label{prob:in-plantear}
Escoja uno: (a) un invernadero en la sabana de Bogotá; (b) un calentador de agua a gas; (c) la dosificación de insulina en una persona con diabetes. Plantee el problema de control con la receta de la \cref{sec:in-problema}: objetivos, manipuladas, perturbaciones, salidas medidas y no medidas, restricciones duras y blandas, tipo de operación, seguridad y una estructura de control con un diagrama de bloques.
\end{problema}

\begin{problema}[ganancia del proceso]\label{prob:in-ganancia}
(a) Un calentador eléctrico pasa de $2.0$ a $2.5$\,kW y la temperatura de un fluido pasa de $40$ a $50$\,°C: halle la ganancia con unidades. (b) En un enfriador, abrir la válvula de agua fría del $30$ al $40\%$ baja la temperatura de $25$ a $21$\,°C: halle la ganancia. ¿Qué significa su signo para el controlador? (c) Si la ganancia de (a) cambia a la mitad con la temperatura ambiente, ¿por qué eso complica el control en lazo abierto?
\end{problema}

\begin{problema}[prealimentación]\label{prob:in-prealimentacion}
\etiquetaM{} Para el ejemplo~\ref{ej:in-prealimentacion}: (a) resuelva analíticamente las tres ecuaciones diferenciales y verifique los valores de la \cref{fig:in-prealimentacion}. (b) ¿Qué error del transmisor de flujo se puede tolerar con prealimentación sola si la deriva debe ser menor que $1$\,cm/min? (c) Con las dos estructuras combinadas, ¿qué $K_c$ deja una desviación final menor que $1$\,cm? (d) Verifique con \cmd{C1\_Ej3\_Prealimentacion}.
\end{problema}

\begin{problema}[moderadores y reguladores]\label{prob:in-maxwell}
\etiquetaEstrella{} Maxwell distinguía los \emph{moderadores}, que dejan un error que depende de la carga, de los \emph{reguladores}, que no lo dejan. Para el tanque del ejemplo~\ref{ej:in-tanque} considere $u=h_r/K+K_ce+K_i\int_0^te\,\dif\tau$, con $e=h_r-h$. (a) Suponiendo que el lazo llega a un equilibrio, muestre que en él $h=h_r$ para cualquier lluvia constante y cualquier $R$. (b) ¿Cuál de las dos clases de Maxwell es el controlador del ejemplo~\ref{ej:in-tanque}? (c) \etiquetaM{} Simule con $K_c=5$ y $K_i=10$ y compare con la \cref{fig:in-tanque}.
\end{problema}

\begin{problema}[el ascensor]\label{prob:in-ascensor}
\etiquetaM{} Con \cmd{C1\_Ej2\_Ascensor}, cambie el factor de amortiguamiento a $\zeta=1$ y a $\zeta=0.3$ manteniendo $\wn=0.8$\,rad/s. Reporte el sobrepaso y el tiempo de establecimiento ($2\%$) y discuta cuál ascensor preferiría usted como pasajero y cuál como persona apurada.
\end{problema}

\begin{problema}[un invento, una fecha, una referencia]\label{prob:in-historia}
\etiquetaEstrella{} Escoja un dispositivo o una idea de la \cref{sec:in-historia}, busque la publicación o patente original y escriba una página que responda: ¿qué problema resolvía?, ¿dónde estaban el sensor, el comparador y el actuador?, ¿qué problema de estabilidad o de desempeño apareció?, ¿quién lo resolvió y cómo? Cite la fuente original y al menos una reseña histórica de las de esta sección.
\end{problema}

%% ---------------------------------------------------------------------
\section{Soluciones seleccionadas}
%% ---------------------------------------------------------------------

\paragraph{Problema~\ref{prob:in-clasificar}}
Lazo abierto: (a) y (c). Lazo cerrado: (b) temperatura, termostato, compresor, abrir la puerta; (d) longitud de las colas, cámara, tiempos del semáforo, llegada de carros; (e) temperatura corporal, termorreceptores, sudor y escalofríos, el clima; (f) y (g) velocidad, sensor de velocidad, acelerador, pendientes y viento. La ducha usa además prealimentación cuando uno oye el sanitario.

\paragraph{Problema~\ref{prob:in-tanque}}
(a) $\tau\dot h+(1+KK_c)h=K\bigl(h_r/K+K_ch_r+d\bigr)$; constante de tiempo $\tau/(1+KK_c)$. (b) Con $K_c=1$: lluvia $h=1.133$, obstrucción $h=1.059$; con $K_c=10$: $1.019$ y $1.008$. La desviación por lluvia es $0.4/(1+2K_c)$. (d) El caudal se queda en $2$ al comienzo y el nivel sube más despacio; el valor final no cambia, porque en equilibrio el caudal pedido es menor que $2$. (e) $h=RK_c/(1+RK_c)=10/11=0.909$: sin calibración el controlador proporcional necesita un error para producir caudal.

\paragraph{Problema~\ref{prob:in-ganancia}}
(a) $10/0.5=20$\,°C/kW. (b) $-4/10=-0.4$\,°C por punto porcentual: la ganancia es negativa, y el controlador debe abrir la válvula cuando la temperatura está \emph{por encima} de la referencia (acción inversa). (c) La calibración $u=r/K$ queda mal y el error es del $100\%$ del cambio.

\paragraph{Problema~\ref{prob:in-prealimentacion}}
(a) $\dev{h}=0.025(t-2)$; $\dev{h}=0.5\bigl(1-e^{-(t-2)/2}\bigr)$; $\dev{h}=0.05\bigl(1-e^{-(t-2)/2}\bigr)$, para $t\ge2$. (b) $0.5\,\varepsilon/2<0.01$: error menor que el $4\%$. (c) $0.05/K_c<0.01$: $K_c>5$\,m$^2$/min.

\paragraph{Problema~\ref{prob:in-maxwell}}
(a) En equilibrio todas las señales son constantes; si $\int e$ es constante, su derivada $e$ es cero, luego $h=h_r$, sin importar $d$ ni $R$ (siempre que el lazo sea estable). (b) Un moderador: deja un error que depende de la perturbación.

\paragraph{Problema~\ref{prob:in-ascensor}}
$\zeta=1$: sin sobrepaso, $\ts=7.29$\,s. $\zeta=0.6$: $9.48\%$ y $7.43$\,s. $\zeta=0.3$: $37.2\%$ y $14.0$\,s. Menos amortiguamiento no siempre es más rápido: con $\zeta=0.3$ llega antes al piso ($\tp=4.12$\,s, contra $4.91$\,s) pero tarda el doble en quedarse quieto.

\respuestaspausas

                   % migrado (Introduccion.tex)
%% =====================================================================
%%  Capítulo 2 -- Modelado I: Laplace, función de transferencia y primer orden
%%  Migrado al formato nuevo (noviembre de 2026) desde Handout2_CtrlParte1.tex
%%  y las presentaciones Modelado1Prof, Modelado2 y Modelado3.
%%  Cambios de fondo respecto a la versión anterior (ver MIGRACION.md):
%%    equilibrio T^* (antes \bar T); escalón \escalon(t) (antes u(t), que
%%    chocaba con la entrada u); polos complejos -1/3 ± j sqrt(2)/3 (faltaba
%%    la j); error ante rampa e(t) = K1 tau (1 - e^{-t/tau}) (signo); se
%%    eliminó la "inversión del tiempo" (no aplica a la transformada
%%    unilateral); el ejemplo de dsolve usa y'(0)=0 en texto y código.
%% =====================================================================

\capitulo[cap:po]{Modelado I: Laplace, función de transferencia y primer orden}
  {La ilusión es el mejor alimento espiritual, el motor del alma. No engaña, entusiasma.}
  {Jorge Barraza}

\begin{objetivos}[partes/02-modelado-primer-orden.html]
  \item Distinguir señal, sistema estático y sistema dinámico, y describir un sistema con un diagrama de bloques de flujo de información.
  \item Escribir el modelo de un proceso a partir de un balance de conservación, con sus suposiciones, y expresarlo en variables desviación.
  \item Calcular transformadas y antitransformadas de Laplace con la tabla, las propiedades y las fracciones parciales, y usarlas para resolver ecuaciones diferenciales lineales.
  \item Definir función de transferencia, polos, ceros, grado relativo y ecuación característica, y usar (bien) el teorema del valor final.
  \item Reconocer un sistema de primer orden, leer $K_1$ y $\tau$ de su respuesta al escalón y predecir su error ante una rampa, verificando todo en \matlab/\simulink.
\end{objetivos}

%% ---------------------------------------------------------------------
\section{Señales y sistemas}
\label{sec:po-senales}
%% ---------------------------------------------------------------------

En el capítulo anterior vimos qué significa ``análisis de sistemas de control'' y cómo nació el área. Ahora empezamos el trabajo de fondo: \emph{construir modelos} de sistemas dinámicos. Todo lo que hagamos en el curso (estabilidad, error, PID, Bode) se apoya en dos palabras que conviene precisar desde ya: señal y sistema.

\begin{definicion}[label=def:po-senal]{Señal}
Una \clave{señal} es una función del tiempo que describe cómo cambia una variable. Se escribe con el tiempo entre paréntesis: $T(t)$, $V(t)$, $y(t)$. Su valor en un instante particular se obtiene evaluándola: $T(2)$ es la temperatura en $t=2$\,s.
\end{definicion}

La \cref{fig:po-senal} muestra la temperatura $T(t)=20-3e^{-0.15t}$ de un recinto que arranca en $17$\,°C y se acerca a $20$\,°C. Esa curva es una señal; la pregunta que nos interesa es \emph{qué la produce}.

\begin{figure}[ht]
\centering
\includegraphics[width=0.8\linewidth]{cap02/f_senal_temperatura}
\caption{Una señal de temperatura. Es la salida de un sistema dinámico de primer orden: en la \cref{sec:po-primerorden} sabremos leer, de esta misma curva, su ganancia y su constante de tiempo.}
\label{fig:po-senal}
\end{figure}

Considere un bombillo dentro de un contenedor aislado (\cref{fig:po-bombillo}). Dos señales nos importan: el voltaje $V(t)$ aplicado al bombillo y la temperatura $T(t)$ dentro del contenedor. Intuimos que a más voltaje, más temperatura: $T(t)$ \emph{depende} de $V(t)$. Esa relación de causa y efecto es lo que llamamos sistema.

\begin{figure}[ht]
\centering
\begin{minipage}{0.36\linewidth}\centering
\includegraphics[width=\linewidth]{cap02/bombillo}
\end{minipage}\hfill
\begin{minipage}{0.6\linewidth}\centering
\includegraphics[width=\linewidth]{cap02/senal_sistema}
\end{minipage}
\caption{Izquierda: el bombillo en un contenedor aislado. Derecha: (a)--(b) el diagrama de bloques de flujo de información: la flecha que entra es la causa, la que sale es el efecto; (c) un sistema dinámico, descrito por una ecuación diferencial; (d) un sistema estático.}
\label{fig:po-bombillo}
\end{figure}

\begin{definicion}[label=def:po-sistema]{Sistema, sistema dinámico}
Un \clave{sistema} es una transformación que toma una señal de entrada $u(t)$ y devuelve una señal de salida $y(t)$. Es \clave{dinámico} cuando la relación entre $u$ e $y$ es una \emph{ecuación diferencial ordinaria} (ODE): involucra a $u$, a $y$ y al menos una de sus derivadas. A esa ecuación la llamamos las \clave{dinámicas} del sistema. Si no hay derivadas, el sistema es \clave{estático}: la salida en $t$ solo depende de la entrada en ese mismo $t$.
\end{definicion}

El bombillo es dinámico (\cref{fig:po-bombillo}c): su temperatura no salta cuando se cambia el voltaje, sino que ``recuerda'' su historia a través de $\dif T/\dif t$. La conversión de unidades del ejemplo siguiente, en cambio, es estática.

\begin{ejemplo}[label=ej:po-fc]{De Fahrenheit a Celsius: un sistema estático}
Si $u(t)$ es una temperatura en °F, la temperatura en °C es $y(t)=\tfrac59\,(u(t)-32)$ (\cref{fig:po-bombillo}d). ¿Cuánto es $90$\,°F?
\solucion $y=\tfrac59(90-32)=\tfrac{290}{9}\approx32.22$\,°C. No hay derivadas: si la entrada cambia de golpe, la salida cambia de golpe. En \matlab se evalúa sobre un vector de temperaturas, numérica o simbólicamente:
\codigorepo{cap02_modelado_primer_orden/C2_Ej1_Senales.m}{fc}{ejemplo \ref{ej:po-fc}}
\end{ejemplo}

\begin{pausa}
La señal de la \cref{fig:po-senal} cumple $\dif T/\dif t+0.15\,T=V$ con $V$ constante. ¿Cuánto vale $V$? ¿Y $T(0)$?
\respuesta{En estado estacionario $\dif T/\dif t=0$, luego $0.15\cdot20=V$, es decir $V=3$. Y $T(0)=20-3=17$\,°C. Se puede comprobar derivando: $\dot T=0.45e^{-0.15t}$ y $\dot T+0.15T=0.45e^{-0.15t}+3-0.45e^{-0.15t}=3$.}
\end{pausa}

%% ---------------------------------------------------------------------
\section{¿Para qué y cómo se modela?}
\label{sec:po-modelar}
%% ---------------------------------------------------------------------

\subsection{El modelo dentro del proceso de diseño}

Un sistema de control se diseña en cuatro grandes pasos: (i) establecer los objetivos, las variables a controlar y las especificaciones con las que se medirá el desempeño; (ii) definir el sistema y \emph{modelarlo} (la configuración, los modelos del proceso, de los actuadores y de los sensores, y su simulación); (iii) escoger el controlador y los parámetros que se van a ajustar; y (iv) optimizar esos parámetros y analizar el desempeño. Este capítulo y el siguiente son el paso (ii). La \cref{fig:po-intercambiador} muestra un caso típico~\citep{SmCo06}: un intercambiador de calor cuya temperatura de salida se controla manipulando el vapor.

\begin{figure}[ht]
\centering
\includegraphics[width=0.5\linewidth]{cap02/intercambiador_pid}\\[-1mm]{\footnotesize (a)}\\[2mm]
\includegraphics[width=\linewidth]{cap02/lazo_proceso}\\[1mm]{\footnotesize (b)}
\caption{(a) Diagrama de instrumentación de un intercambiador de calor (adaptado de~\citep{SmCo06}). (b) El mismo lazo como diagrama de bloques: cada equipo es un bloque y cada flecha, una señal con sus unidades.}
\label{fig:po-intercambiador}
\end{figure}

Antes de escribir una sola ecuación conviene hacerse tres preguntas~\citep{Mar00}:
\begin{itemize}
  \item \textbf{¿Qué variables influyen y cuáles son influenciadas?} Las \clave{variables manipuladas} son las que podemos ajustar (el vapor); las \clave{variables controladas} son las que se ven afectadas (la temperatura de salida). Las demás entradas son perturbaciones.
  \item \textbf{¿En qué rango se pueden mover?} Una válvula no abre más del $100\%$; un calentador no enfría.
  \item \textbf{¿Qué tan efectivamente puede la retroalimentación mantener el proceso donde queremos?} Depende de tres rasgos de la respuesta del proceso: el \emph{signo y la magnitud} (que no cambien con el punto de operación), la \emph{velocidad} (si manipulamos muy agresivamente, el lazo puede oscilar o inestabilizarse) y la \emph{forma} (dos procesos con la misma velocidad pueden ser muy distintos de controlar).
\end{itemize}
Un proceso ``fácil'' tiene ganancia de signo y magnitud estables, responde rápido y su comportamiento no cambia mucho con el tiempo.

\subsection{Modelos orientados al control}

\begin{receta}[Cómo se construye un modelo~\citep{Mar00}]
\begin{enumerate}
  \item \textbf{Definir objetivos:} qué decisiones se van a tomar con el modelo, qué valores numéricos y con qué precisión.
  \item \textbf{Preparar la información:} bosquejo del proceso, variables de interés, \emph{suposiciones} y datos.
  \item \textbf{Formular el modelo:} balances de conservación (masa, energía, momento), ecuaciones constitutivas, y verificar los grados de libertad (tantas ecuaciones como incógnitas).
  \item \textbf{Resolver:} analítica o numéricamente.
  \item \textbf{Analizar los resultados:} ¿tienen sentido? ¿qué tan sensibles son a los parámetros? ¿qué pasa ``si\ldots''?
  \item \textbf{Validar:} comparar con experimentos o con modelos más completos.
\end{enumerate}
\end{receta}

Buscamos \clave{modelos orientados al control} (\emph{control-oriented models}): los más simples que capturen lo esencial de la estática y la dinámica, porque se usan dentro de optimizaciones, en lazos de tiempo real y para diseñar controladores.\nota{``Essentially, all models are wrong, but some are useful'', escribió el estadístico George E.\,P.\ Box~\citep[p.~424]{BoDr87}.} Los modelos de \emph{principios físicos} tienen ventajas sobre los basados solo en datos: explican la dinámica interna, se pueden extrapolar a otras condiciones y se pueden escribir antes de que el equipo exista. Los modelos basados en datos brillan cuando no entendemos bien el proceso, pero su gran debilidad es la extrapolación: nada garantiza que valgan fuera de los datos con que se entrenaron.

\subsection{Sistemas lineales e invariantes en el tiempo}

En casi todo el curso trabajaremos con sistemas \clave{lineales e invariantes en el tiempo} (LTI). Si $y=h(u)$ es la salida para la entrada $u$:
\begin{itemize}
  \item \textbf{Linealidad} (homogeneidad y superposición): $h(a\,u_1+b\,u_2)=a\,h(u_1)+b\,h(u_2)$. Doble entrada, doble salida; la respuesta a una suma es la suma de las respuestas.
  \item \textbf{Invarianza en el tiempo:} si la entrada se retrasa $T$ segundos, la salida se retrasa lo mismo: $u(t-T)\mapsto y(t-T)$. El sistema responde igual hoy que mañana.
\end{itemize}
Ningún sistema real es exactamente LTI.\nota{Richard Feynman lo resumía así: ``Linear systems are important because we can solve them''.} Pero muchos se comportan como tales \emph{alrededor de un punto de operación}, y en esa región podemos usar todo el arsenal de herramientas lineales. En el \cref{cap:so} veremos cómo construir esa aproximación (linealización).

%% ---------------------------------------------------------------------
\subsection{Ejemplo motivador: un tanque agitado}
\label{sec:po-termico-modelo}

Tomemos el tanque de la \cref{fig:po-tanque}~\citep{SmCo06}. Un líquido entra con caudal $f$ y temperatura $T_i(t)$, se mezcla y sale con el mismo caudal y temperatura $T(t)$. Suposiciones: (1) caudales de entrada y salida constantes; (2) densidades constantes; (3) capacidades caloríficas constantes; (4) líquido bien mezclado (la temperatura de salida es la del tanque); (5) tanque bien aislado (sin pérdidas al ambiente); (6) se desprecia la energía que aporta el agitador.

\begin{figure}[ht]
\centering
\includegraphics[width=0.55\linewidth]{cap02/tanque_termico}
\caption{Tanque agitado: ¿cómo responde la temperatura de salida $T(t)$ a un cambio en la de entrada $T_i(t)$? Adaptado de~\citep{SmCo06}.}
\label{fig:po-tanque}
\end{figure}

\textbf{Objetivo:} obtener un modelo matemático y una función de transferencia que describan cómo responde $T(t)$ a un cambio en $T_i(t)$.

El balance de energía (lo que entra menos lo que sale es lo que se acumula) es
\[
  f\rho_i h_i(t)-f\rho\,h(t)=\deriv{}{t}\bigl[V\rho\,u(t)\bigr],
\]
con $f$ el caudal (m$^3$/s), $\rho_i,\rho$ las densidades (kg/m$^3$), $V$ el volumen (m$^3$), $h_i,h$ las entalpías (J/kg) y $u$ la energía interna del líquido (J/kg).\nota{Aquí $u$ es energía interna, no la entrada del sistema. Es la única vez que la usamos así.} Con una temperatura de referencia, $h=C_pT$ y $u=C_vT$; con las suposiciones (2) y (3) queda
\begin{equation}
  f\rho C_p\,T_i(t)-f\rho C_p\,T(t)=V\rho C_v\,\deriv{T(t)}{t}.
  \label{eq:po-balance}
\end{equation}
$T_i(t)$ la imponemos nosotros, así que hay una ecuación y una incógnita: el modelo está completo.

\paragraph{Equilibrio y variables desviación} En estado estacionario la derivada es cero y $f\rho C_p\,\eq{T_i}-f\rho C_p\,\eq{T}=0$: el punto de equilibrio es $\eq{T}=\eq{T_i}$. Definimos las \clave{variables desviación}
\begin{equation}
  \dev{T}(t)=T(t)-\eq{T},\qquad \dev{T}_i(t)=T_i(t)-\eq{T_i}.
  \label{eq:po-desviacion}
\end{equation}
Restando la ecuación de equilibrio de \eqref{eq:po-balance}, y como $\dif\dev{T}/\dif t=\dif T/\dif t$,
\begin{equation}
  \tau\,\deriv{\dev{T}(t)}{t}+\dev{T}(t)=\dev{T}_i(t),\qquad
  \tau=\frac{V C_v}{f C_p}\quad\left[\frac{\text{m}^3\cdot\text{J/(kg\,°C)}}{(\text{m}^3/\text{s})\cdot\text{J/(kg\,°C)}}=\text{s}\right].
  \label{eq:po-termico-ode}
\end{equation}

\begin{ideaclave}
Las variables desviación son un \emph{cambio de coordenadas}: ponen el origen en el punto de equilibrio. Así el sistema ``arranca en cero'', con condiciones iniciales nulas, que es justamente lo que exige la función de transferencia (\cref{sec:po-ft}). Para volver a la variable física basta sumar el equilibrio: $T(t)=\eq{T}+\dev{T}(t)$.
\end{ideaclave}

Podríamos resolver \eqref{eq:po-termico-ode} con cualquier herramienta numérica si conociéramos la entrada. Pero hay una herramienta que además nos da una \emph{representación} del sistema que no depende de la entrada: la transformada de Laplace. La introducimos ahora y volvemos al tanque en la \cref{sec:po-termico}.

%% ---------------------------------------------------------------------
\section{La transformada de Laplace}
\label{sec:po-laplace}
%% ---------------------------------------------------------------------

\begin{historia}[De los operadores de Heaviside a la transformada de Laplace]
Oliver Heaviside (1850--1925) resolvía circuitos y líneas de transmisión con un \emph{cálculo operacional}: reemplazaba $\dif/\dif t$ por un símbolo $p$ y manipulaba las ecuaciones diferenciales como si fueran algebraicas~\citep{Mah17}. Funcionaba, pero sin demostraciones, y los matemáticos de su tiempo lo despreciaron. Thomas Bromwich, en Cambridge, mostró que muchas de esas manipulaciones se podían justificar tratando $p$ como un número complejo e integrando en el plano complejo~\citep{Bro17}; John R.\ Carson, ingeniero de AT\&T, dio en los años veinte una fórmula integral que producía la función operacional a partir de la función del tiempo~\citep{Car26}; y el matemático alemán Gustav Doetsch reunió todo en una teoría rigurosa con el nombre que usamos hoy~\citep{Doe37}, en honor de Pierre-Simon Laplace (1749--1827). El escalón unitario se sigue llamando \emph{función de Heaviside}, y la expansión en fracciones parciales que usaremos para regresar al tiempo también es suya. Heaviside, además, dio a las ecuaciones de Maxwell la forma vectorial que hoy se enseña e hizo posibles las líneas telefónicas largas sin distorsión.
\end{historia}

\begin{definicion}[label=def:po-laplace]{Transformada unilateral de Laplace}
Para una señal $x(t)$ definida para $t\geq0$,
\begin{equation}
  \Lap{x(t)}=X(s)=\int_0^\infty x(t)\,e^{-st}\,\dif t,\qquad s=\sigma+j\omega\in\mathbb{C}.
  \label{eq:po-laplace}
\end{equation}
\end{definicion}

El factor $e^{-st}=e^{-\sigma t}e^{-j\omega t}$ tiene dos partes: una exponencial $e^{-\sigma t}$ que ayuda a que la integral converja y una oscilación $e^{-j\omega t}$. Por eso $s$ es complejo: la transformada descompone la señal en exponenciales y senoides, que son exactamente las funciones que aparecen en las soluciones de las ODE lineales. Con $\sigma=0$ se recupera la transformada de Fourier, que volverá en el capítulo de respuesta frecuencial.

\paragraph{El plano-$s$} Un número complejo $p_1=\sigma_1+j\omega_1$ ($j^2=-1$) tiene parte real $\Re p_1=\sigma_1$, parte imaginaria $\Im p_1=\omega_1$, magnitud $|p_1|=\sqrt{\sigma_1^2+\omega_1^2}$ y fase $\angle p_1=\phi_1$ (el ángulo medido desde el semieje real positivo; ojo con el cuadrante, $\tan^{-1}(\omega_1/\sigma_1)$ sola no lo distingue). Su conjugado es $\bar p_1=\sigma_1-j\omega_1$ y $|p_1|^2=p_1\bar p_1$. Todos se dibujan en el \clave{plano complejo-$s$} (\cref{fig:po-planos}).

\begin{figure}[ht]
\centering
\includegraphics[width=0.62\linewidth]{cap02/plano_s}
\caption{Un número complejo y su conjugado en el plano-$s$. La fase se mide desde el semieje real positivo: aquí $\phi_1$ está en el segundo cuadrante.}
\label{fig:po-planos}
\end{figure}

\paragraph{Existencia} Basta que $x(t)$ sea \emph{continua a trozos} (en cada intervalo finito tiene un número finito de saltos, con límites a izquierda y derecha) y de \emph{orden exponencial} (existe $a$ tal que $e^{-at}|x(t)|$ es acotada para $t$ grande). La segunda condición impone $\Re s>\sigma_a$ para algún $\sigma_a$: es la \emph{región de convergencia}. En el curso todas las señales cumplen estas condiciones y no insistiremos en la región de convergencia.

\begin{cuidado}[Todas las señales empiezan en $t=0$]
La transformada unilateral solo ``ve'' $t\geq0$. Por eso, cuando escribimos $x(t)$ queremos decir $x(t)\,\escalon(t)$, donde $\escalon(t)$ es el escalón unitario (1 para $t\geq0$ y 0 antes). En las notas anteriores el escalón se escribía $u(t)$; aquí lo escribimos $\escalon(t)$ porque $u(t)$ es la \emph{entrada} de un sistema.
\end{cuidado}

\begin{ejemplo}[label=ej:po-lap-escalon]{Escalón y exponencial, desde la definición}
Calcule $\Lap{\escalon(t)}$ y $\Lap{e^{-at}}$ con $a>0$.
\solucion
\[
  \Lap{\escalon(t)}=\int_0^\infty e^{-st}\,\dif t=\Bigl[-\tfrac1s e^{-st}\Bigr]_0^\infty=\frac1s\quad(\Re s>0),
\]
porque $e^{-st}\to0$ cuando $\sigma>0$. De igual forma,
\[
  \Lap{e^{-at}}=\int_0^\infty e^{-(s+a)t}\,\dif t=\frac{1}{s+a}\quad(\Re s>-a).
\]
\end{ejemplo}

\Cref{tab:po-transformadas} reúne los pares que usaremos todo el semestre. \matlab los calcula con el \emph{Symbolic Math Toolbox}:
\codigorepo{cap02_modelado_primer_orden/C2_Ej2_Laplace.m}{transformadas}{tabla \ref{tab:po-transformadas}}

\begin{table}[ht]
\centering
\caption{Pares de transformadas de Laplace ($t\geq0$, $a,\omega>0$).}
\label{tab:po-transformadas}
\renewcommand{\arraystretch}{2.3}
\begin{tabular}{@{}>{$}l<{$}>{$}l<{$}@{\qquad}>{$}l<{$}>{$}l<{$}@{}}
\toprule
x(t) & X(s) & x(t) & X(s)\\
\midrule
\impulso(t) & 1 & t\,e^{-at} & \dfrac{1}{(s+a)^2}\\
\escalon(t) & \dfrac1s & t^n e^{-at} & \dfrac{n!}{(s+a)^{n+1}}\\
t & \dfrac1{s^2} & \sin\omega t & \dfrac{\omega}{s^2+\omega^2}\\
t^n & \dfrac{n!}{s^{n+1}} & \cos\omega t & \dfrac{s}{s^2+\omega^2}\\
e^{-at} & \dfrac{1}{s+a} & e^{-at}\sin\omega t & \dfrac{\omega}{(s+a)^2+\omega^2}\\
 & & e^{-at}\cos\omega t & \dfrac{s+a}{(s+a)^2+\omega^2}\\
\bottomrule
\end{tabular}
\end{table}

%% ---------------------------------------------------------------------
\subsection{Propiedades}
\label{sec:po-propiedades}

Con unas pocas propiedades la tabla alcanza para casi todo. Sea $x(t)\leftrightarrow X(s)$.

\begin{enumerate}
  \item \textbf{Linealidad:} $a\,x_1(t)+b\,x_2(t)\ \leftrightarrow\ a\,X_1(s)+b\,X_2(s)$. (La integral es lineal.)
  \item \textbf{Corrimiento en el tiempo} ($t_0>0$): $x(t-t_0)\,\escalon(t-t_0)\ \leftrightarrow\ e^{-st_0}X(s)$.
  \item \textbf{Corrimiento en $s$:} $e^{s_0t}x(t)\ \leftrightarrow\ X(s-s_0)$. Por ejemplo, $e^{-2t}\cos3t\leftrightarrow\dfrac{s+2}{(s+2)^2+9}$, y como $\cos\omega_0t=\tfrac12(e^{j\omega_0t}+e^{-j\omega_0t})$, $\cos(\omega_0t)\,x(t)\leftrightarrow\tfrac12X(s-j\omega_0)+\tfrac12X(s+j\omega_0)$.
  \item \textbf{Escalamiento} ($a>0$): $x(at)\ \leftrightarrow\ \dfrac1a X\!\left(\dfrac sa\right)$. (Con el cambio de variable $\tau=at$.)
  \item \textbf{Diferenciación:} $\dot x(t)\ \leftrightarrow\ sX(s)-x(0)$ y, en general,
  \begin{equation}
    \frac{\dif^n x}{\dif t^n}\ \leftrightarrow\ s^nX(s)-s^{n-1}x(0)-s^{n-2}\dot x(0)-\cdots-x^{(n-1)}(0).
    \label{eq:po-derivada}
  \end{equation}
  \item \textbf{Integración:} $\displaystyle\int_0^t x(\lambda)\,\dif\lambda\ \leftrightarrow\ \frac{X(s)}{s}$.
  \item \textbf{Diferenciación en $s$:} $t\,x(t)\ \leftrightarrow\ -\dfrac{\dif X(s)}{\dif s}$. (De aquí sale $t\,e^{-at}\leftrightarrow1/(s+a)^2$.)
  \item \textbf{Convolución:} $\displaystyle (x_1*x_2)(t)=\int_0^t x_1(\lambda)\,x_2(t-\lambda)\,\dif\lambda\ \leftrightarrow\ X_1(s)\,X_2(s)$.
\end{enumerate}

La propiedad de diferenciación es la que convierte una ecuación diferencial en una algebraica. Su demostración es una integración por partes con $u=x(t)$ y $\dif v=e^{-st}\dif t$:
\[
  \Lap{\dot x}=\int_0^\infty \dot x\,e^{-st}\,\dif t
  =\Bigl[x\,e^{-st}\Bigr]_0^\infty+s\int_0^\infty x\,e^{-st}\,\dif t
  =-x(0)+sX(s),
\]
porque $x\,e^{-st}\to0$ (orden exponencial). Aplicándola dos veces se obtiene la de la segunda derivada, y así sucesivamente. Observe que \alerta{las condiciones iniciales aparecen solas}: no hay que agregarlas al final como con los métodos clásicos.

\codigorepo{cap02_modelado_primer_orden/C2_Ej2_Laplace.m}{propiedades}{propiedades de la transformada}

\begin{profundiza}[¿Y la ``inversión del tiempo''?]
La versión anterior de estas notas listaba la propiedad $x(-t)\leftrightarrow X(-s)$. Es válida para la transformada \emph{bilateral} (integral de $-\infty$ a $\infty$), pero no para la unilateral que usamos: si $x(t)$ solo existe para $t\geq0$, entonces $x(-t)$ solo existe para $t\leq0$ y la integral \eqref{eq:po-laplace} la ve como cero. Por la misma razón el escalamiento solo se enuncia con $a>0$.
\end{profundiza}

\begin{pausa}
Use las propiedades y la tabla para obtener $\Lap{3-2e^{-4t}+5t\,e^{-4t}}$.
\respuesta{Por linealidad, $\dfrac3s-\dfrac{2}{s+4}+\dfrac{5}{(s+4)^2}$.}
\end{pausa}

%% ---------------------------------------------------------------------
\section{Resolver ecuaciones diferenciales con Laplace}
\label{sec:po-odes}
%% ---------------------------------------------------------------------

El método tiene tres pasos: \textbf{(1)} transformar la ODE con \eqref{eq:po-derivada}, incluyendo las condiciones iniciales; \textbf{(2)} despejar $X(s)$, que queda como un cociente de polinomios; \textbf{(3)} regresar al tiempo expandiendo $X(s)$ en fracciones parciales y usando la tabla.

\begin{ejemplo}[label=ej:po-motor]{El calibrador óptico}
Un calibrador óptico que verificaba la alineación de las piezas metálicas de varios sedanes de cuatro puertas de Chrysler tenía dinámicas no forzadas (la ODE igualada a cero)
\[
  \ddot x+5\dot x+6x=0,\qquad x(0)=1,\ \dot x(0)=0.
\]
Halle $X(s)$ y $x(t)$.
\solucion Transformando término a término,
\[
  \bigl[s^2X-s\,x(0)-\dot x(0)\bigr]+5\bigl[sX-x(0)\bigr]+6X=0
  \;\Longrightarrow\;
  X(s)(s^2+5s+6)=s+5,
\]
\[
  X(s)=\frac{s+5}{(s+2)(s+3)}=\frac{3}{s+2}-\frac{2}{s+3}
  \quad\Longrightarrow\quad x(t)=3e^{-2t}-2e^{-3t}.
\]
La \cref{fig:po-motor} muestra la respuesta y sus dos componentes.
\end{ejemplo}

\begin{figure}[ht]
\centering
\includegraphics[width=\linewidth]{cap02/f_chrysler}
\caption{Ejemplo~\ref{ej:po-motor}. Cada polo de $X(s)$ aporta un \emph{modo} $e^{p_it}$; el residuo es su peso. El polo en $-3$ muere más rápido que el de $-2$.}
\label{fig:po-motor}
\end{figure}

\codigorepo{cap02_modelado_primer_orden/C2_Ej3_Motor.m}{simbolico}{ejemplo \ref{ej:po-motor} con Laplace simbólico}

\begin{definicion}[label=def:po-polos]{Polos y ceros}
Si $X(s)=N(s)/D(s)$ es un cociente de polinomios, los \clave{ceros} de $X$ son las raíces de $N(s)=0$ y los \clave{polos} son las raíces de $D(s)=0$. En el plano-$s$ se dibujan los polos con $\times$ y los ceros con $\circ$.
\end{definicion}

La expansión en fracciones parciales, que introdujo Heaviside, depende del tipo de polos. Supondremos $X(s)$ estrictamente propia (grado de $N$ menor que el de $D$); si no lo es, se divide primero.

\subsection{Polos simples}
Si $D(s)=(s-p_1)(s-p_2)\cdots(s-p_n)$ con raíces distintas,
\begin{equation}
  X(s)=\sum_{i=1}^n\frac{k_i}{s-p_i},\qquad
  k_i=\Bigl[(s-p_i)\,X(s)\Bigr]_{s=p_i},\qquad
  x(t)=\sum_{i=1}^n k_i\,e^{p_it}.
  \label{eq:po-residuos}
\end{equation}
En el ejemplo~\ref{ej:po-motor}: $k_1=\bigl[(s+5)/(s+3)\bigr]_{s=-2}=3$ y $k_2=\bigl[(s+5)/(s+2)\bigr]_{s=-3}=-2$. A los $k_i$ se les llama \clave{residuos}, y \matlab los calcula con \cmd{residue}:
\codigorepo{cap02_modelado_primer_orden/C2_Ej3_Motor.m}{residue}{residuos, polos y diagrama de polos y ceros}

\subsection{Polos múltiples}
Si $p$ es un polo de multiplicidad $r$, aporta $r$ términos:
\[
  \frac{k_{1}}{s-p}+\frac{k_{2}}{(s-p)^2}+\cdots+\frac{k_{r}}{(s-p)^r},\qquad
  k_{r-m}=\frac{1}{m!}\,\frac{\dif^m}{\dif s^m}\Bigl[(s-p)^rX(s)\Bigr]_{s=p},\ m=0,\dots,r-1,
\]
y en el tiempo, por la tabla, $k_m\dfrac{t^{m-1}}{(m-1)!}e^{pt}$: aparecen términos $t\,e^{pt}$, $t^2e^{pt}$\ldots

\begin{ejemplo}[label=ej:po-multiples]{Dos polos dobles}
Expanda $X(s)=\dfrac{1}{s^2(s+1)^2}$ y obtenga $x(t)$.
\solucion Con $s=0$ (doble): $k_2=\bigl[1/(s+1)^2\bigr]_0=1$ y $k_1=\bigl[\tfrac{\dif}{\dif s}(s+1)^{-2}\bigr]_0=-2$. Con $s=-1$ (doble): $k_2=\bigl[1/s^2\bigr]_{-1}=1$ y $k_1=\bigl[-2s^{-3}\bigr]_{-1}=2$. Luego
\[
  X(s)=-\frac2s+\frac1{s^2}+\frac2{s+1}+\frac1{(s+1)^2},\qquad
  x(t)=-2+t+2e^{-t}+t\,e^{-t}.
\]
\end{ejemplo}
\codigorepo{cap02_modelado_primer_orden/C2_Ej4_FraccionesParciales.m}{multiples}{ejemplo \ref{ej:po-multiples}}

\subsection{Polos complejos conjugados}
Un par $-\alpha\pm j\beta$ tiene residuos complejos conjugados y produce, juntos, un término real $e^{-\alpha t}(A\cos\beta t+B\sin\beta t)$. En la práctica es más cómodo \emph{completar el cuadrado} en el denominador y usar las dos últimas filas de la \cref{tab:po-transformadas}.

\begin{ejemplo}[label=ej:po-complejos]{Una oscilación amortiguada}
Resuelva $3\ddot x+2\dot x+x=0$ con $x(0)=-1$, $\dot x(0)=2$.
\solucion Transformando, $3(s^2X+s-2)+2(sX+1)+X=0$, de donde
\[
  X(s)=\frac{4-3s}{3s^2+2s+1},\qquad\text{polos: } s=-\tfrac13\pm j\,\tfrac{\sqrt2}{3}.
\]
Completando el cuadrado, $3s^2+2s+1=3\bigl[(s+\tfrac13)^2+\tfrac29\bigr]$, y escribiendo el numerador en términos de $s+\tfrac13$:
\[
  X(s)=\frac{-(s+\frac13)+\frac53}{(s+\frac13)^2+(\frac{\sqrt2}{3})^2}
  \;\Longrightarrow\;
  x(t)=e^{-t/3}\left[-\cos\frac{\sqrt2\,t}{3}+\frac{5\sqrt2}{2}\sin\frac{\sqrt2\,t}{3}\right].
\]
Se comprueba que $x(0)=-1$ y $\dot x(0)=\tfrac13+\tfrac53=2$. La parte real de los polos ($-\tfrac13$) da la envolvente; la imaginaria ($\tfrac{\sqrt2}3$), la frecuencia de oscilación (\cref{fig:po-complejos}).
\end{ejemplo}

\begin{figure}[ht]
\centering
\includegraphics[width=0.85\linewidth]{cap02/f_polos_complejos}
\caption{Ejemplo~\ref{ej:po-complejos}: la parte real de los polos fija la envolvente $e^{-t/3}$; la imaginaria, la frecuencia $\sqrt2/3$\,rad/s.}
\label{fig:po-complejos}
\end{figure}
\codigorepo{cap02_modelado_primer_orden/C2_Ej4_FraccionesParciales.m}{complejos}{ejemplo \ref{ej:po-complejos}}

\subsection{De los polos a la forma de la respuesta}
Los tres casos se resumen en un solo mapa (\cref{fig:po-modos}): cada polo $p$ aporta un modo $e^{pt}$. La parte real dice si el modo crece o muere, y qué tan rápido; la parte imaginaria dice si oscila, y a qué frecuencia. Un polo en el semiplano derecho basta para que la respuesta crezca sin límite.

\begin{figure}[ht]
\centering
\includegraphics[width=\linewidth]{cap02/f_polos_modos}
\caption{Cada ubicación de polo en el plano-$s$ produce un modo distinto. Izquierda del eje imaginario: el modo se extingue (más rápido cuanto más a la izquierda). Sobre el eje: oscila sin amortiguarse. A la derecha: crece.}
\label{fig:po-modos}
\end{figure}

\begin{ideaclave}
Los polos determinan \emph{qué} modos aparecen en la respuesta; los residuos (y por lo tanto los ceros y las condiciones iniciales) determinan \emph{cuánto} pesa cada uno. Volveremos sobre esto en el \cref{cap:so} con los polos dominantes.
\end{ideaclave}

\subsection{Sin transformar: \texttt{dsolve} y \texttt{ode45}}
\matlab puede resolver el problema de valor inicial directamente: \cmd{dsolve} lo hace simbólicamente (internamente, con las mismas ideas) y \cmd{ode45} numéricamente, sin necesidad de que la ecuación sea lineal.
\codigorepo{cap02_modelado_primer_orden/C2_Ej5_dsolve.m}{dsolve}{$\dot y=-2y+t$ y $\ddot y+\dot y+y=\sin t$ con $y(0)=1$, $\dot y(0)=0$}

\begin{pausa}
En el ejemplo~\ref{ej:po-motor}, ¿qué condición inicial $\dot x(0)$ haría desaparecer el modo $e^{-3t}$?
\respuesta{Con $\dot x(0)=v$, $X(s)=\dfrac{s+5+v}{(s+2)(s+3)}$. El residuo en $-3$ es $\dfrac{-3+5+v}{-1}=-(2+v)$, que se anula con $v=-2$. Entonces $X=1/(s+2)$ y $x(t)=e^{-2t}$: la condición inicial ``cae'' justo sobre el modo lento.}
\end{pausa}

%% ---------------------------------------------------------------------
\section{Función de transferencia}
\label{sec:po-ft}
%% ---------------------------------------------------------------------

Las herramientas anteriores nos llevan al concepto central de todo el curso.

\begin{definicion}[label=def:po-ft]{Función de transferencia}
La \clave{función de transferencia} de un sistema LTI es el cociente entre la transformada de Laplace de la salida y la de la entrada, \emph{con condiciones iniciales nulas}:
\begin{equation}
  T(s)=\left.\frac{\Lap{\text{salida}}}{\Lap{\text{entrada}}}\right|_{\text{CI}=0}=\frac{Y(s)}{U(s)}.
\end{equation}
\end{definicion}

Si el sistema se describe con la ecuación entrada-salida
\begin{equation}
  \frac{\dif^n y}{\dif t^n}+a_1\frac{\dif^{n-1}y}{\dif t^{n-1}}+\cdots+a_n y
  =b_0\frac{\dif^m u}{\dif t^m}+\cdots+b_{m-1}\frac{\dif u}{\dif t}+b_m u,
  \label{eq:po-inout}
\end{equation}
la linealidad y la diferenciación con condiciones iniciales nulas dan
\begin{equation}
  T(s)=\frac{Y(s)}{U(s)}=\frac{b_0s^m+b_1s^{m-1}+\cdots+b_m}{s^n+a_1s^{n-1}+\cdots+a_n}.
  \label{eq:po-ft}
\end{equation}

Algunos nombres que usaremos todo el tiempo:
\begin{itemize}
  \item $T(s)$ es \clave{propia} si $n\geq m$, y \clave{estrictamente propia} si $n>m$. Los sistemas físicos son propios: una función impropia derivaría la entrada y amplificaría sin límite el ruido de alta frecuencia.
  \item El \clave{grado relativo} es $n-m$. El \clave{orden} del sistema es $n$.
  \item El polinomio del denominador igualado a cero, $s^n+a_1s^{n-1}+\cdots+a_n=0$, es la \clave{ecuación característica}. Sus raíces son los polos de $T(s)$: los modos naturales del sistema.
\end{itemize}

La función de transferencia permite describir el sistema con un bloque (\cref{fig:po-bloqueft}): la señal de salida es el producto $Y(s)=T(s)\,U(s)$. En el tiempo, por la propiedad de convolución, eso es $y(t)=\int_0^t g(\lambda)\,u(t-\lambda)\,\dif\lambda$, donde $g(t)=\iLap{T(s)}$ es la \clave{respuesta al impulso}: el producto en $s$ es mucho más cómodo que la integral en $t$.

\begin{figure}[ht]
\centering
\includegraphics[width=0.62\linewidth]{cap02/bloque_ft}
\caption{Un bloque de flujo de información: la función de transferencia es un mapeo lineal entre la transformada de la entrada y la de la salida. Las reglas para combinar bloques están en el \cref{cap:so}.}
\label{fig:po-bloqueft}
\end{figure}

\begin{cuidado}[La función de transferencia no ve las condiciones iniciales]
$T(s)$ describe la respuesta \emph{forzada} desde el reposo. Si el sistema no arranca en reposo, la respuesta total es la forzada más la libre (la que producen las condiciones iniciales, como en el ejemplo~\ref{ej:po-motor}). Por eso trabajamos en variables desviación alrededor de un equilibrio: allí las condiciones iniciales son cero.
\end{cuidado}

\begin{pausa}
Para $\dddot y+2\ddot y+3\dot y+4y=\dddot u+5\dot u+6u$, escriba $T(s)$. ¿Es propia? ¿Cuál es su grado relativo?
\respuesta{$T(s)=\dfrac{s^3+5s+6}{s^3+2s^2+3s+4}$. Es propia pero no estrictamente propia ($n=m=3$): grado relativo $0$. Una parte de la entrada pasa directamente a la salida.}
\end{pausa}

%% ---------------------------------------------------------------------
\section{El teorema del valor final}
\label{sec:po-tvf}
%% ---------------------------------------------------------------------

Muchas veces no necesitamos toda la respuesta, solo a dónde llega: el \clave{valor en estado estacionario}. Para eso existe una de las propiedades más útiles de la transformada.

\begin{definicion}[label=def:po-tvf]{Teorema del valor final}
Si \alerta{todos los polos de $sY(s)$ tienen parte real negativa}, entonces
\begin{equation}
  \lim_{t\to\infty}y(t)=\lim_{s\to0}sY(s).
  \label{eq:po-tvf}
\end{equation}
\end{definicion}

La condición es esencial: el teorema supone que el límite en el tiempo existe. Si no existe, la fórmula de la derecha da un número igual, y ese número no significa nada.

\begin{ejemplo}[label=ej:po-tvf]{Dos casos en los que el teorema miente}
Determine $\lim_{t\to\infty}h(t)$ para (a) $H(s)=\dfrac1{s-2}$ y (b) $H(s)=\dfrac{6(s^2-2)}{s^3+s^2-4s-4}$.
\solucion (a) Aplicando \eqref{eq:po-tvf} a ciegas, $\lim_{s\to0}\frac{s}{s-2}=0$. Pero $h(t)=e^{2t}$ crece sin cota: el polo de $sH(s)$ en $+2$ invalida el teorema.

(b) De nuevo $\lim_{s\to0}sH(s)=0$. Pero $s^3+s^2-4s-4=(s+1)(s+2)(s-2)$, y
\[
  H(s)=\frac{2}{s+1}+\frac{3}{s+2}+\frac{1}{s-2}\;\Longrightarrow\;h(t)=2e^{-t}+3e^{-2t}+e^{2t}\to\infty.
\]
Moraleja: \emph{primero} los polos, \emph{después} el límite.
\end{ejemplo}

\codigorepo{cap02_modelado_primer_orden/C2_Ej6_ValorFinal.m}{tvf}{ejemplo \ref{ej:po-tvf}: revisar los polos antes del límite}

\begin{profundiza}[El teorema del valor inicial]
Con una condición parecida (que el límite exista), $\displaystyle\lim_{t\to0^+}y(t)=\lim_{s\to\infty}sY(s)$. Sirve, por ejemplo, para saber si una respuesta al escalón arranca con salto (grado relativo $0$) o arranca en cero con pendiente (grado relativo $1$): aplíquelo a $sY(s)$ para obtener $\dot y(0^+)$.
\end{profundiza}

%% ---------------------------------------------------------------------
\section{De vuelta al tanque: la primera función de transferencia}
\label{sec:po-termico}
%% ---------------------------------------------------------------------

\begin{ejemplo}[label=ej:po-termico]{El tanque agitado ante un escalón}
El modelo \eqref{eq:po-termico-ode} es $\tau\,\dif\dev{T}/\dif t+\dev{T}=\dev{T}_i$. (a) Obtenga la función de transferencia. (b) Si la temperatura de entrada, que estaba en el equilibrio $\eq{T_i}=\eq{T}$, sube bruscamente $M$ grados, ¿cómo responde la salida? (c) ¿A qué valor llega?
\solucion (a) Con condiciones iniciales nulas (variables desviación), $\tau s\,\Dev{T}(s)+\Dev{T}(s)=\Dev{T}_i(s)$, luego
\begin{equation}
  \frac{\Dev{T}(s)}{\Dev{T}_i(s)}=\frac{1}{\tau s+1}.
  \label{eq:po-termico-ft}
\end{equation}
(b) La entrada es $\dev{T}_i(t)=M\,\escalon(t)$, es decir, $\Dev{T}_i(s)=M/s$. Por polos simples,
\[
  \Dev{T}(s)=\frac{M}{s(\tau s+1)}=\frac{M}{s}-\frac{M}{s+1/\tau}
  \;\Longrightarrow\;
  \dev{T}(t)=M\left(1-e^{-t/\tau}\right),\qquad
  T(t)=\eq{T}+M\left(1-e^{-t/\tau}\right).
\]
(c) Todos los polos de $s\Dev{T}(s)=M/(\tau s+1)$ están en el semiplano izquierdo, así que el teorema aplica: $\lim_{s\to0}s\Dev{T}(s)=M$. Es lo mismo que se obtiene de $\lim_{t\to\infty}\dev{T}(t)$. La salida sube exactamente lo que subió la entrada: la ganancia del tanque es $1$, como pide la conservación de la energía sin pérdidas.
\end{ejemplo}

\begin{figure}[ht]
\centering
\includegraphics[width=\linewidth]{cap02/f_termico}
\caption{Ejemplo~\ref{ej:po-termico} con $\eq{T}=\eq{T_i}=25$\,°C, $M=2$\,°C y $\tau=5$\,s. Izquierda: las variables desviación, que arrancan en cero. Derecha: las temperaturas físicas, que se obtienen sumando el equilibrio.}
\label{fig:po-termico}
\end{figure}

\codigorepo{cap02_modelado_primer_orden/C2_Ej7_Termico.m}{termico}{ejemplo \ref{ej:po-termico}}

\begin{cuidado}[\texttt{step} es un escalón \emph{unitario}]
\cmd{step(G)} aplica $\escalon(t)$. Para un escalón de tamaño $M$ se multiplica: \cmd{M*step(G,t)} (por linealidad). Y \cmd{step} devuelve la \emph{desviación}: si quiere la temperatura física, sume el equilibrio.
\end{cuidado}

\begin{modelo}[Termico.slx]
El script \cmd{crearModelo\_Termico.m} construye por código el diagrama escalón (o rampa) $\to$ $K_1/(\tau s+1)$ $\to$ $+\eq{T}$, con un \emph{Scope} para las variables desviación y otro para las físicas. Cambie \cmd{tau}, \cmd{M} y \cmd{Teq} en el \emph{workspace} y vuelva a simular; con el interruptor manual pase a la entrada rampa de la \cref{sec:po-primerorden}.
\end{modelo}

\begin{pausa}
Si se duplica el caudal $f$ del tanque, ¿la respuesta es más rápida o más lenta? ¿Cambia el valor final?
\respuesta{Más rápida: $\tau=VC_v/(fC_p)$ se reduce a la mitad (el líquido se renueva en la mitad del tiempo). El valor final no cambia: la ganancia sigue siendo $1$.}
\end{pausa}

%% ---------------------------------------------------------------------
\section{El sistema de primer orden}
\label{sec:po-primerorden}
%% ---------------------------------------------------------------------

El tanque es el ejemplo típico de un \clave{sistema de primer orden}:
\begin{equation}
  \frac{Y(s)}{U(s)}=\frac{K_1}{\tau s+1},
  \label{eq:po-primerorden}
\end{equation}
con dos parámetros: la \clave{ganancia} $K_1$ (cuánto cambia la salida en estado estacionario por unidad de cambio en la entrada) y la \clave{constante de tiempo} $\tau$ (qué tan rápido). El único polo está en $s=-1/\tau$. Veamos qué hacen ante las dos entradas de prueba más usadas.

\subsection{Respuesta al escalón unitario}
Con $U(s)=1/s$,
\begin{equation}
  Y(s)=\frac{K_1}{s}-\frac{K_1}{s+1/\tau}\quad\Longrightarrow\quad y(t)=K_1\left(1-e^{-t/\tau}\right).
  \label{eq:po-escalon}
\end{equation}
La respuesta tiene dos partes: la de estado estacionario $y_{ss}=K_1$ (el teorema del valor final da lo mismo) y la transitoria $-K_1e^{-t/\tau}$, que se extingue. La \cref{fig:po-escalon} muestra los números que vale la pena memorizar.

\begin{figure}[ht]
\centering
\includegraphics[width=0.92\linewidth]{cap02/f_primer_orden_escalon}
\caption{Respuesta al escalón de $K_1/(\tau s+1)$ con $K_1=10$ y $\tau=5$. Cada constante de tiempo, la salida recorre el $63.2\%$ de lo que le falta. La tangente en el origen llega al valor final justo en $t=\tau$.}
\label{fig:po-escalon}
\end{figure}

\begin{itemize}
  \item En $t=\tau$ la salida llega al $1-e^{-1}=63.2\%$ del valor final; en $2\tau$, al $86.5\%$; en $3\tau$, al $95.0\%$; en $4\tau$, al $98.2\%$; en $5\tau$, al $99.3\%$.
  \item \clave{Tiempo de establecimiento} (2\%): $\ts\approx4\tau$ (exactamente $\tau\ln50=3.91\,\tau$).
  \item La pendiente inicial es $\dot y(0)=K_1/\tau$: si la respuesta siguiera con esa pendiente, llegaría al valor final en $t=\tau$. Constante de tiempo pequeña, respuesta rápida.
\end{itemize}

\codigorepo{cap02_modelado_primer_orden/C2_Ej8_PrimerOrden.m}{escalon}{escalón: 63.2 \%, $\ts$ y \cmd{stepinfo}}

\begin{receta}[Identificar un primer orden con una prueba al escalón]
\begin{enumerate}
  \item Llevar el proceso a un equilibrio $(\eq{u},\eq{y})$ y aplicar un escalón $\Delta u$.
  \item Esperar el nuevo estado estacionario y medir $\Delta y$: $K_1=\Delta y/\Delta u$ (con sus unidades).
  \item Buscar el instante en que la salida recorrió el $63.2\%$ de $\Delta y$: ese tiempo es $\tau$. (O, con un solo dato $y(t_1)$: $\tau=-t_1/\ln\bigl(1-\dev{y}(t_1)/\Delta y\bigr)$.)
  \item Escribir $\Dev{Y}/\Dev{U}=K_1/(\tau s+1)$ y verificar superponiendo la simulación a los datos.
\end{enumerate}
\end{receta}

\subsection{Respuesta a la rampa}
Con $u(t)=t$, es decir $U(s)=1/s^2$,
\[
  Y(s)=\frac{K_1}{s^2}-\frac{K_1\tau}{s}+\frac{K_1\tau}{s+1/\tau}
  \;\Longrightarrow\;
  y(t)=K_1t-K_1\tau+K_1\tau\,e^{-t/\tau},\quad t\geq0.
\]
La parte $K_1t-K_1\tau$ es la de estado estacionario y $K_1\tau e^{-t/\tau}$ la transitoria. Definiendo el error de seguimiento como $e(t)=K_1u(t)-y(t)$ (lo que la salida ``debería'' valer si no tuviera dinámica),
\begin{equation}
  e(t)=K_1\tau\left(1-e^{-t/\tau}\right)\;\xrightarrow{\ t\to\infty\ }\;K_1\tau.
  \label{eq:po-error-rampa}
\end{equation}
La salida sigue a la rampa con la misma pendiente, pero \emph{atrasada} $\tau$ segundos: el error final es $K_1\tau$ (\cref{fig:po-rampa}). Cuanto menor $K_1\tau$, mejor el seguimiento.

\begin{figure}[ht]
\centering
\includegraphics[width=\linewidth]{cap02/f_primer_orden_rampa}
\caption{Respuesta a la rampa con $K_1=1$: con $\tau=10$ el error final es $10$; con $\tau=1$, es $1$.}
\label{fig:po-rampa}
\end{figure}

\codigorepo{cap02_modelado_primer_orden/C2_Ej8_PrimerOrden.m}{rampa}{rampa: el error tiende a $K_1\tau$}

\begin{cuidado}[El signo del error ante la rampa]
En la versión anterior de estas notas aparecía $e(t)=K_1\tau+K_1\tau e^{-t/\tau}$, que vale $2K_1\tau$ en $t=0$. Lo correcto es $e(t)=K_1\tau(1-e^{-t/\tau})$: en $t=0$ el error es nulo (la rampa arranca en cero) y crece hasta $K_1\tau$.
\end{cuidado}

\begin{pausa}
Un primer orden con $K_1=2$ responde a un escalón de $3$ unidades y a los $12$\,s ya está en el $86.5\%$ del cambio final. ¿Cuánto valen $\tau$, el cambio final y el $\ts$?
\respuesta{$86.5\%$ corresponde a $2\tau$, luego $\tau=6$\,s. Cambio final $K_1\cdot3=6$ unidades; $\ts\approx4\tau=24$\,s.}
\end{pausa}

%% ---------------------------------------------------------------------
\section{Para llevar}
\label{sec:po-cierre}
%% ---------------------------------------------------------------------

\begin{parallevar}
\begin{enumerate}
  \item Un sistema transforma una señal de entrada en una de salida. Es dinámico si la relación es una ODE: la salida ``recuerda''.
  \item Un modelo orientado al control parte de un balance con suposiciones explícitas. En el equilibrio se definen las variables desviación $\dev{x}=x-\eq{x}$: arrancan en cero.
  \item Laplace convierte derivadas en productos por $s$ \emph{e incluye las condiciones iniciales}. Para regresar al tiempo: fracciones parciales (polos simples, múltiples o complejos) y la tabla.
  \item Cada polo aporta un modo $e^{pt}$: parte real, crecimiento o decaimiento; parte imaginaria, oscilación. Los residuos dicen cuánto pesa cada modo.
  \item Función de transferencia: $Y(s)/U(s)$ con condiciones iniciales nulas. Propia, grado relativo, ecuación característica.
  \item Teorema del valor final: \emph{primero} verificar que los polos de $sY(s)$ estén en el semiplano izquierdo.
  \item Primer orden $K_1/(\tau s+1)$: $63.2\%$ en $\tau$, $\ts\approx4\tau$, pendiente inicial $K_1/\tau$, error ante rampa $K_1\tau$.
\end{enumerate}
\end{parallevar}

\begin{profundiza}[Correcciones respecto a la versión anterior del capítulo]
(i) El punto de equilibrio se escribe $\eq{T}$ (antes $\bar T$) y el escalón $\escalon(t)$ (antes $u(t)$, que chocaba con la entrada). (ii) Polos complejos del ejemplo~\ref{ej:po-complejos}: $-\tfrac13\pm j\tfrac{\sqrt2}{3}$ (faltaba la $j$). (iii) La fase de un número complejo es $\angle p$ (el texto decía ``magnitud''). (iv) Se eliminó la ``inversión del tiempo'', que no vale para la transformada unilateral. (v) El ejemplo de \cmd{dsolve} usa $\dot y(0)=0$ en el texto y en el código. (vi) El error ante la rampa es $K_1\tau(1-e^{-t/\tau})$. (vii) Se quitó una mención a ``la posición final de la masa'' que no correspondía a ningún ejemplo. Todos los números del capítulo se verificaron con \texttt{herramientas/cap02\_numeros.py} y \cmd{verificarCap02.m}.
\end{profundiza}

%% ---------------------------------------------------------------------
\section{Problemas}
\label{sec:po-problemas}
%% ---------------------------------------------------------------------

Los problemas marcados con \etiquetaM{} requieren \matlab/\simulink; los marcados con \etiquetaEstrella{} son de profundización.

\subsection*{Transformadas y fracciones parciales}

\begin{problema}[transformadas]\label{prob:po-1}
Determine $X(s)$ para: (a) $x(t)=3-2e^{-4t}+5te^{-4t}$; (b) $x(t)=e^{-2t}\cos3t+4t^2$; (c) $x(t)=t\,e^{-t}\sin2t$; (d) $x(t)=\int_0^t e^{-3\lambda}\,\dif\lambda$; (e) el pulso $x(t)=2$ para $0\leq t<3$ y $0$ en otro caso. En (a)--(d) use la tabla y las propiedades; en (e), la definición \eqref{eq:po-laplace}, y diga qué propiedad le habría permitido evitar la integral. \etiquetaM{} Verifique con \cmd{laplace}.
\end{problema}

\begin{problema}[fracciones parciales]\label{prob:po-2}
Para $F_1(s)=\dfrac{s+3}{(s+1)(s+4)}$, $F_2(s)=\dfrac{10}{s(s+2)^2}$ y $F_3(s)=\dfrac{2s+5}{s^2+4s+13}$: (a) ubique polos y ceros en el plano-$s$; (b) expanda en fracciones parciales y obtenga $f_i(t)$; (c) \etiquetaM{} verifique con \cmd{residue} y \cmd{pzplot}; (d) diga qué términos desaparecen cuando $t\to\infty$ y relaciónelo con la ubicación de los polos.
\end{problema}

\begin{problema}[ODE con condiciones iniciales]\label{prob:po-3}
Considere $2\ddot y+10\dot y+12y=24u$ con $u(t)=\escalon(t)$, $y(0)=1$, $\dot y(0)=0$. (a) Obtenga $Y(s)$ sin despreciar las condiciones iniciales. (b) Halle $y(t)$ y verifique las dos condiciones iniciales. (c) Calcule $\lim_{t\to\infty}y(t)$ de dos maneras. (d) Repita (c) para $\ddot y-\dot y-2y=3u$ con condiciones iniciales nulas; obtenga $y(t)$ y explique por qué el teorema del valor final falla. (e) \etiquetaM{} Resuelva ambos con \cmd{ode45} y grafique.
\end{problema}

\subsection*{Función de transferencia}

\begin{problema}[función de transferencia]\label{prob:po-4}
Para (i) $\dddot y+2\ddot y+3\dot y+4y=\dddot u+5\dot u+6u$; (ii) $\ddot y+7\dot y+10y=4\dot u+8u$; (iii) $5\dot y+2y=3u$: (a) escriba $T(s)$; (b) diga si es propia o estrictamente propia y su grado relativo; (c) halle polos y ceros; (d) ¿cuál es de primer orden de la forma \eqref{eq:po-primerorden}? Identifique $K_1$ y $\tau$. (e) Mire con cuidado el sistema (ii): ¿qué le pasa a $T(s)$?
\end{problema}

\begin{problema}[del bloque a la ODE]\label{prob:po-5}
Escriba la ecuación diferencial que representa el bloque de la \cref{fig:po-repaso}.
\end{problema}

\begin{figure}[ht]
\centering
\includegraphics[width=0.6\linewidth]{cap02/bloque_repaso}
\caption{Problema~\ref{prob:po-5}.}
\label{fig:po-repaso}
\end{figure}

\subsection*{Primer orden}

\begin{problema}[exploración]\label{prob:po-6}
\etiquetaM{} Con \cmd{C2\_Ej8\_PrimerOrden.m} y \cmd{crearModelo\_Termico.m}, varíe $K_1$ y $\tau$ y compruebe las afirmaciones de la \cref{sec:po-primerorden} para el escalón y la rampa. ¿Qué parámetro cambia el valor final y cuál la velocidad?
\end{problema}

\begin{problema}[identificación de un proceso térmico]\label{prob:po-7}
Un proceso opera en estado estacionario con temperatura de salida $\eq{T}=40$\,°C y entrada $\eq{u}=10$ unidades. En $t=0$ la entrada pasa a $15$ unidades; la salida se estabiliza en $52$\,°C y a los $30$\,s marca $47.6$\,°C.
(a) Halle $K_1$ y $\tau$ y escriba la función de transferencia en variables desviación.
(b) Calcule la pendiente inicial y el $\ts$ (2\%).
(c) ¿Qué entrada lleva la salida a $55$\,°C?
(d) Si desde el equilibrio la entrada sube como una rampa de $0.01$ unidades/s, halle $y(t)$ y el error final $\lim e(t)$ con $e=K_1\dev{u}-\dev{y}$.
(e) Si ese error no puede superar $0.5$\,°C, ¿cuál es el máximo $\tau$ admisible? ¿Qué implicaría para el diseño del equipo?
(f) \etiquetaM{} Verifique (a)--(d) en \simulink.
\end{problema}

\begin{problema}[teorema del valor inicial]\label{prob:po-8}
\etiquetaEstrella{} Use el teorema del valor inicial para hallar $y(0^+)$ y $\dot y(0^+)$ ante un escalón unitario en (a) $\dfrac{K_1}{\tau s+1}$ y (b) $\dfrac{s+3}{s+1}$. Relacione el resultado con el grado relativo y explique por qué en (b) la salida ``salta''.
\end{problema}

%% ---------------------------------------------------------------------
\section{Soluciones seleccionadas}
%% ---------------------------------------------------------------------

\paragraph{Problema~\ref{prob:po-1}}
(a) $\dfrac3s-\dfrac2{s+4}+\dfrac5{(s+4)^2}$. (b) $\dfrac{s+2}{(s+2)^2+9}+\dfrac8{s^3}$. (c) $\Lap{e^{-t}\sin2t}=\dfrac{2}{(s+1)^2+4}$ y, por diferenciación en $s$, $\dfrac{4(s+1)}{\bigl[(s+1)^2+4\bigr]^2}$. (d) Por integración, $\dfrac{1}{s(s+3)}$. (e) $\int_0^3 2e^{-st}\dif t=\dfrac{2(1-e^{-3s})}{s}$; es $2\escalon(t)-2\escalon(t-3)$ y basta el corrimiento en el tiempo.

\paragraph{Problema~\ref{prob:po-2}}
$f_1=\tfrac23e^{-t}+\tfrac13e^{-4t}$ (polos $-1,-4$; cero $-3$). $f_2=2.5-2.5e^{-2t}-5t\,e^{-2t}$ (polos $0$ y $-2$ doble): el término constante no desaparece porque hay un polo en el origen. $f_3=e^{-2t}\bigl(2\cos3t+\tfrac13\sin3t\bigr)$ (polos $-2\pm3j$, cero $-2.5$).

\paragraph{Problema~\ref{prob:po-3}}
(a) $Y(s)=\dfrac{s^2+5s+12}{s(s+2)(s+3)}$. (b) $y(t)=2-3e^{-2t}+2e^{-3t}$; $y(0)=1$, $\dot y(0)=6-6=0$. (c) $2$ por los dos caminos. (d) $Y=\dfrac{3}{s(s-2)(s+1)}$, el teorema daría $-1.5$; pero $y(t)=-1.5+0.5e^{2t}+e^{-t}$ diverge por el polo en $+2$.

\paragraph{Problema~\ref{prob:po-4}}
(i) $\dfrac{s^3+5s+6}{s^3+2s^2+3s+4}$, propia, grado relativo $0$. (ii) $\dfrac{4(s+2)}{(s+2)(s+5)}$: el cero cancela el polo en $-2$ y queda $\dfrac{4}{s+5}=\dfrac{0.8}{0.2s+1}$; el modo $e^{-2t}$ no se ve desde la entrada (volveremos sobre esto en espacio de estado). (iii) $\dfrac{3}{5s+2}=\dfrac{1.5}{2.5s+1}$: primer orden con $K_1=1.5$ y $\tau=2.5$.

\paragraph{Problema~\ref{prob:po-5}}
$y^{(4)}+5\dddot y+6\ddot y+7\dot y+8y=\dddot u+2\ddot u+3\dot u+4u$.

\paragraph{Problema~\ref{prob:po-7}}
(a) $K_1=12/5=2.4$\,°C/unidad; $1-e^{-30/\tau}=7.6/12$ da $\tau=29.9\approx30$\,s; $\Dev{T}/\Dev{U}=2.4/(29.9s+1)$. (b) $0.40$\,°C/s y $\ts\approx4\tau\approx120$\,s. (c) $\dev{u}=15/2.4=6.25$, es decir $u=16.25$. (d) $\dev{y}=0.024\bigl(t-\tau+\tau e^{-t/\tau}\bigr)$ y $e_\infty=K_1a\tau=0.72$\,°C. (e) $\tau\leq0.5/(2.4\cdot0.01)=20.8$\,s: habría que hacer el proceso un $30\%$ más rápido (menos volumen o más caudal), o, mejor, cerrar un lazo de control.

\script[verificación automática de los números del capítulo]{cap02_modelado_primer_orden/verificarCap02.m}

\respuestaspausas

          % migrado (Handout2_CtrlParte1.tex)
%% =====================================================================
%%  Capítulo 3 -- Modelado II: linealización, segundo orden, bloques y Mason
%%  Migrado al formato nuevo (noviembre de 2026) desde Handout2_CtrlParte2.tex
%%  y las presentaciones Modelado3, Modelado4 y Modelado5.
%%  Cambios de fondo respecto a la versión anterior (ver MIGRACION.md):
%%    el ejemplo numérico de los dos tanques (tau1=1, tau2=4, C2=2) no era
%%    físico (C2 = -C1 siempre, y la planta es SIEMPRE sobreamortiguada): se
%%    usan datos físicos y el subamortiguado se ilustra con masa-resorte;
%%    beta = acos(zeta) (decía cos zeta); MP -> %OS; "determinante" ->
%%    discriminante; "asentamiento" -> establecimiento; los números del
%%    ejemplo con cero y polo adicionales se recalcularon (54.5 % / 38.0 %).
%% =====================================================================

\capitulo[cap:so]{Modelado II: linealización, segundo orden y diagramas de bloques}
  {Mi misión es matar el tiempo, la suya, matarme a mí. Se está perfectamente a gusto entre asesinos.}
  {Emil Michel Cioran}

\begin{objetivos}[partes/03-modelado-segundo-orden.html]
  \item Linealizar un modelo no lineal alrededor de un punto de operación y obtener su función de transferencia en variables desviación.
  \item Reconocer un sistema de segundo orden en forma estándar, clasificarlo por su $\zeta$ y ubicar sus polos.
  \item Calcular $\tr$, $\tp$, $\OS$ y $\ts$ de un subamortiguado, y traducir especificaciones de tiempo en una región del plano-$s$.
  \item Reducir diagramas de bloques con álgebra de bloques y con la ley de Mason, y escribir la ecuación característica.
  \item Decidir cuándo un sistema de orden superior se puede aproximar por sus polos dominantes y simplificarlo, verificando en \matlab/\simulink.
\end{objetivos}

%% ---------------------------------------------------------------------
\section{¿Y si el sistema no es lineal?}
\label{sec:so-nolineal}
%% ---------------------------------------------------------------------

En el \cref{cap:po} todo era lineal: el tanque agitado, con sus suposiciones, daba una ODE lineal de primer orden y la transformada de Laplace hacía el resto. Pero la mayoría de los procesos reales no son lineales: el caudal por una válvula depende de la \emph{raíz cuadrada} de la diferencia de presiones, la velocidad de una reacción del \emph{cuadrado} de una concentración, la fuerza de un resorte endurecido del \emph{cubo} de la elongación. Con un sistema no lineal:
\begin{itemize}
  \item No se puede usar la transformada de Laplace (no es lineal), así que no hay función de transferencia.
  \item El análisis cambia: por lo general se busca entender cualitativamente el comportamiento más que calcularlo exactamente.
\end{itemize}
¿Podemos rescatar algo de la teoría que ya tenemos? Sí: si el proceso opera \emph{cerca} de un punto de operación, se comporta casi como un sistema lineal. Esa aproximación, la \clave{linealización}, es lo primero que haremos. Con ella llegaremos a una función de transferencia de \emph{segundo orden}, el otro gran protagonista de los modelos orientados al control, y a la necesidad de describir sistemas con varios bloques interconectados.

\begin{figure}[ht]
\centering
\includegraphics[width=0.85\linewidth]{cap03/dos_tanques}
\caption{Dos tanques que interactúan. El caudal $f_1$ depende de \emph{las dos} alturas: lo que pasa en un tanque afecta al otro. Adaptado de~\citep{SmCo06}.}
\label{fig:so-tanques}
\end{figure}

\textbf{El problema guía.} En el sistema de la \cref{fig:so-tanques}~\citep{SmCo06}: (1) determinar cómo afecta el caudal de entrada $f_i(t)$ al nivel del segundo tanque $h_2(t)$; (2) escribir el modelo; (3) obtener una función de transferencia entre esa entrada y esa salida; y (4) representar el sistema con un diagrama de bloques. Lo resolvemos en la \cref{sec:so-tanques}, después de aprender a linealizar.

%% ---------------------------------------------------------------------
\section{Linealización}
\label{sec:so-linealizacion}
%% ---------------------------------------------------------------------

La idea es vieja y sencilla: \emph{de cerca, toda curva suave parece una recta}. Si el proceso se mueve poco alrededor de un punto, reemplazamos cada función no lineal por su recta tangente (su serie de Taylor truncada en el primer orden). La suposición de fondo es que la respuesta de la aproximación lineal representa la del proceso en la región cercana al punto de operación~\citep{SmCo06,DoBi08}.

\paragraph{Variables desviación} Como en el \cref{cap:po}, medimos cada variable desde el punto de operación:
\begin{equation}
  \dev{x}(t)=x(t)-\eq{x}.
  \label{eq:so-desviacion}
\end{equation}
Es mover el cero del eje hasta $\eq{x}$; las derivadas no cambian: $\dif^n\dev{x}/\dif t^n=\dif^n x/\dif t^n$. En casi todos los casos el punto de operación será un \clave{punto de equilibrio} (donde las derivadas se anulan); cuando no lo sea, lo diremos.

\paragraph{Una variable} Sea $\dot x=f(x)$ con $f$ suave.\nota{Continuamente diferenciable; para la existencia y unicidad de la solución basta, además, que sea Lipschitz.} Su serie de Taylor alrededor de $x_0$ es
\[
  f(x)=\underbrace{f(x_0)+\left.\deriv{f}{x}\right|_{x_0}(x-x_0)}_{\text{lo que se conserva}}
  +\underbrace{\frac{1}{2!}\left.\frac{\dif^2f}{\dif x^2}\right|_{x_0}(x-x_0)^2+\cdots}_{\text{lo que se desprecia}}
\]
y la aproximación lineal es
\begin{equation}
  f(x)\approx f(x_0)+\left.\deriv{f}{x}\right|_{x_0}\,\dev{x}.
  \label{eq:so-taylor1}
\end{equation}
Si $x_0=\eq{x}$ es un equilibrio, $f(\eq{x})=0$ y queda $\dot{\dev{x}}=a\,\dev{x}$ con $a=f'(\eq{x})$: un modelo lineal en variables desviación.

\begin{ejemplo}[label=ej:so-lin1]{Linealizar alrededor de un punto que no es de equilibrio}
Linealice $\dot x=x\sqrt{x^2+5x}$ alrededor de $x_0=5$.
\solucion La derivada es
\[
  \deriv{f}{x}=\sqrt{x^2+5x}+\frac{x(2x+5)}{2\sqrt{x^2+5x}},
\]
y en $x_0=5$: $f(5)=5\sqrt{50}=35.36$ y $f'(5)=\sqrt{50}+\dfrac{75}{2\sqrt{50}}=12.37$. Entonces
\[
  \dot{\dev{x}}\approx35.36+12.37\,\dev{x}.
\]
Como $x_0=5$ \emph{no} es un equilibrio (los equilibrios son $x=0$ y $x=-5$), queda un término constante: la aproximación es \emph{afín}. La \cref{fig:so-lin1} muestra que solo es buena cerca de $x_0$: el error relativo es menor que el $5\%$ entre $x\approx3.9$ y $x\approx6.7$.
\end{ejemplo}

\begin{figure}[ht]
\centering
\includegraphics[width=0.85\linewidth]{cap03/f_linealizacion}
\caption{Ejemplo~\ref{ej:so-lin1}: la función y su recta tangente en $x_0=5$. En la franja verde el error relativo es menor que el $5\%$.}
\label{fig:so-lin1}
\end{figure}

\codigorepo{cap03_modelado_segundo_orden/C3_Ej1_Linealizacion.m}{simbolica}{ejemplo \ref{ej:so-lin1}: derivada simbólica}

\paragraph{Varias variables} Si $f$ depende de $x_1,\dots,x_n$ (o de estados y entradas), se suma una derivada parcial por variable:
\begin{equation}
  f(x_1,\dots,x_n)\approx f(\eq{x}_1,\dots,\eq{x}_n)+\sum_{i=1}^n\left.\frac{\partial f}{\partial x_i}\right|_{\ast}\dev{x}_i,
  \label{eq:so-taylorn}
\end{equation}
donde ``$\ast$'' indica que las derivadas se evalúan en el punto de operación.

\begin{ejemplo}[label=ej:so-reactor]{Un reactor}
La concentración $x$ de un producto evoluciona según $\dot x=f(x,u)=-kx^2+\tfrac1V\,u\,(c-x)$, donde $u$ es el caudal de alimentación, $k=1$\,L/(mol\,min), $V=1$\,L y $c=2$\,mol/L~\citep{Mar00}. Linealice alrededor de $\eq{x}=1$\,mol/L, $\eq{u}=1$\,L/min.
\solucion Primero, es un equilibrio: $f(1,1)=-1+1\cdot(2-1)=0$. Las derivadas parciales son
\[
  \left.\frac{\partial f}{\partial x}\right|_\ast=-2k\eq{x}-\frac{\eq{u}}{V}=-3,\qquad
  \left.\frac{\partial f}{\partial u}\right|_\ast=\frac{c-\eq{x}}{V}=1,
\]
de modo que $\dot{\dev{x}}=-3\,\dev{x}+\dev{u}$ y, con Laplace,
\[
  \frac{\Dev{X}(s)}{\Dev{U}(s)}=\frac{1}{s+3}=\frac{1/3}{\tfrac13 s+1}:
\]
un primer orden con $K_1=1/3$ y $\tau=1/3$\,min.
\end{ejemplo}

\codigorepo{cap03_modelado_segundo_orden/C3_Ej1_Linealizacion.m}{reactor}{ejemplo \ref{ej:so-reactor}: \cmd{jacobian}}

\begin{receta}[Linealizar un modelo]
\begin{enumerate}
  \item Escribir el modelo como $\dot x_i=f_i(x,u)$ y hallar el punto de operación: resolver $f_i(\eq{x},\eq{u})=0$.
  \item Calcular las derivadas parciales de cada $f_i$ respecto a cada estado y cada entrada, y evaluarlas en el punto.
  \item Escribir el modelo en variables desviación: $\dot{\dev{x}}_i=\sum_j\frac{\partial f_i}{\partial x_j}\big|_\ast\dev{x}_j+\sum_k\frac{\partial f_i}{\partial u_k}\big|_\ast\dev{u}_k$.
  \item Transformar con condiciones iniciales nulas y despejar la función de transferencia.
  \item Verificar contra el modelo no lineal con escalones pequeños \emph{y} grandes: el lineal solo vale cerca del punto.
\end{enumerate}
\end{receta}

\begin{pausa}
Linealice el tanque con descarga libre $A\dot h=q-c\sqrt h$ ($A=2$\,m$^2$, $c=0.5$\,m$^{5/2}$/s) alrededor de $\eq{q}=1$\,m$^3$/s. ¿Cuánto valen $K_1$ y $\tau$?
\respuesta{Equilibrio: $c\sqrt{\eq{h}}=1\Rightarrow\eq{h}=4$\,m. Derivada: $\partial(c\sqrt h)/\partial h=c/(2\sqrt{\eq{h}})=0.125$. Entonces $2\,\dot{\dev{h}}=\dev{q}-0.125\,\dev{h}$ y $\dfrac{\Dev{H}}{\Dev{Q}}=\dfrac{8}{16s+1}$: $K_1=8$\,s/m$^2$ y $\tau=16$\,s.}
\end{pausa}

%% ---------------------------------------------------------------------
\section{Dos tanques que interactúan}
\label{sec:so-tanques}
%% ---------------------------------------------------------------------

\paragraph{Modelo} Los dos tanques de la \cref{fig:so-tanques} están abiertos a la atmósfera y las válvulas descargan a presión atmosférica, así que la diferencia de presión en cada válvula es
\[
  \Delta P_1=(P_A+\rho g h_1)-(P_A+\rho g h_2)=\rho g\,(h_1-h_2),\qquad
  \Delta P_2=(P_A+\rho g h_2)-P_A=\rho g\,h_2,
\]
con $\rho$ la densidad (kg/m$^3$), $g=9.8$\,m/s$^2$ y $h_i$ los niveles (m). Con las aperturas fijas, el caudal por una válvula es $f=C_v\sqrt{\Delta P/G_f}$ ($C_v$: coeficiente de la válvula; $G_f$: gravedad específica), así que
\begin{equation}
  f_1=C'_{v1}\sqrt{h_1-h_2},\qquad f_2=C'_{v2}\sqrt{h_2},\qquad C'_{vi}=C_{vi}\sqrt{\rho g/G_f}.
  \label{eq:so-valvulas}
\end{equation}
El caudal $f_1$ depende de los dos niveles: lo que pasa en un tanque se refleja en el otro y viceversa. El sistema está \emph{acoplado}. Los balances de masa (con densidad constante y áreas $A_1$, $A_2$) son
\begin{equation}
  A_1\deriv{h_1}{t}=f_i-C'_{v1}\sqrt{h_1-h_2},\qquad
  A_2\deriv{h_2}{t}=C'_{v1}\sqrt{h_1-h_2}-C'_{v2}\sqrt{h_2}.
  \label{eq:so-tanques-nl}
\end{equation}
Son no lineales: hay que linealizar.

\paragraph{Equilibrio} Operamos con $\eq{f_i}=\eq{f_1}=\eq{f_2}$, que fija los niveles $\eq{h_1}$ y $\eq{h_2}$ (las derivadas se anulan).

\paragraph{Linealización} Aplicando \eqref{eq:so-taylorn} a \eqref{eq:so-valvulas},
\[
  f_1\approx\eq{f_1}+C_1\,\dev{h}_1+C_2\,\dev{h}_2,\qquad f_2\approx\eq{f_2}+C_3\,\dev{h}_2,
\]
\begin{equation}
  C_1=\left.\frac{\partial f_1}{\partial h_1}\right|_\ast=\frac{C'_{v1}}{2\sqrt{\eq{h_1}-\eq{h_2}}},\qquad
  C_2=\left.\frac{\partial f_1}{\partial h_2}\right|_\ast=-C_1,\qquad
  C_3=\left.\frac{\partial f_2}{\partial h_2}\right|_\ast=\frac{C'_{v2}}{2\sqrt{\eq{h_2}}}.
  \label{eq:so-coeficientes}
\end{equation}
Reemplazando en los balances y restando el equilibrio,
\[
  A_1\,\dot{\dev{h}}_1=\dev{f}_i-C_1\dev{h}_1-C_2\dev{h}_2,\qquad
  A_2\,\dot{\dev{h}}_2=C_1\dev{h}_1+(C_2-C_3)\,\dev{h}_2.
\]

\paragraph{Función de transferencia} Transformando y despejando cada nivel,
\[
  \Dev{H}_1(s)=\frac{K_1}{\tau_1s+1}\bigl(\Dev{F}_i(s)-C_2\Dev{H}_2(s)\bigr),\qquad
  \Dev{H}_2(s)=\frac{K_2}{\tau_2s+1}\,\Dev{H}_1(s),
\]
con $K_1=1/C_1$, $\tau_1=A_1/C_1$, $K_2=C_1/(C_3-C_2)$ y $\tau_2=A_2/(C_3-C_2)$. Eliminando $\Dev{H}_1$:
\begin{equation}
  \frac{\Dev{H}_2(s)}{\Dev{F}_i(s)}=\frac{K_1K_2}{\tau_1\tau_2s^2+(\tau_1+\tau_2)s+1+K_1K_2C_2}.
  \label{eq:so-tanques-ft}
\end{equation}
Las dos ecuaciones de primer orden, acopladas, producen un sistema de \clave{segundo orden}. (La estructura de bloques de esta deducción está en el problema~\ref{prob:so-bloques-tanques}.)

\begin{ejemplo}[label=ej:so-tanques]{Los tanques con números}
Con $A_1=2$\,m$^2$, $A_2=3$\,m$^2$, $C'_{v1}=0.1$ y $C'_{v2}=0.05$\,m$^{5/2}$/s, y el tanque 2 operando en $\eq{h_2}=4$\,m: halle el equilibrio, la función de transferencia y sus polos.
\solucion Equilibrio: $\eq{f}=C'_{v2}\sqrt{4}=0.1$\,m$^3$/s y $\eq{h_1}-\eq{h_2}=(0.1/0.1)^2=1$, luego $\eq{h_1}=5$\,m. Coeficientes: $C_1=0.1/(2\cdot1)=0.05$, $C_2=-0.05$, $C_3=0.05/(2\cdot2)=0.0125$\,m$^2$/s. Entonces $K_1=20$, $\tau_1=40$\,s, $K_2=0.8$, $\tau_2=48$\,s y
\[
  \frac{\Dev{H}_2(s)}{\Dev{F}_i(s)}=\frac{16}{1920s^2+88s+0.2}=\frac{80}{9600s^2+440s+1}.
\]
La ganancia es $80$\,s/m$^2$ ($=1/C_3$: en estado estacionario todo lo que entra sale por la segunda válvula). Los polos son $-0.0024$ y $-0.0434$\,s$^{-1}$: \emph{reales}. La \cref{fig:so-tanques-nl} compara el modelo lineal con el no lineal: con un escalón de $0.01$\,m$^3$/s ($10\%$ del caudal) el error final es de $4$\,cm; con $0.05$\,m$^3$/s, de casi un metro.
\end{ejemplo}

\begin{figure}[ht]
\centering
\includegraphics[width=\linewidth]{cap03/f_tanques}
\caption{Ejemplo~\ref{ej:so-tanques}: nivel del tanque 2 con el modelo no lineal \eqref{eq:so-tanques-nl} y con el linealizado. Cuanto más grande el escalón, más nos alejamos del punto de operación y peor la aproximación.}
\label{fig:so-tanques-nl}
\end{figure}

\codigorepo{cap03_modelado_segundo_orden/C3_Ej2_Tanques.m}{tanques}{ejemplo \ref{ej:so-tanques}}

\begin{cuidado}[Esta planta nunca oscila]
Como $C_2=-C_1$ \emph{siempre} (lo que sube $f_1$ al subir $h_1$ es lo que baja al subir $h_2$), la ecuación característica se puede escribir $s^2+\bigl(\tfrac{C_1+C_3}{A_2}+\tfrac{C_1}{A_1}\bigr)s+\tfrac{C_1C_3}{A_1A_2}$, y la desigualdad entre las medias aritmética y geométrica muestra que $\zeta\geq1$ para cualesquiera parámetros físicos (problema~\ref{prob:so-zeta-tanques}): \alerta{dos tanques que interactúan por gravedad nunca dan una respuesta subamortiguada}. La versión anterior de estas notas usaba $\tau_1=1$, $\tau_2=4$, $C_2=2$, $K_1=K_2=1$ para mostrar un sobrepico de $3.78\%$; esos valores no corresponden a ningún par de tanques reales ($C_2>0$ exigiría que el caudal entre tanques \emph{aumentara} al llenarse el segundo). Para ver oscilaciones usaremos un sistema mecánico (\cref{sec:so-segundoorden}).
\end{cuidado}

\begin{modelo}[Tanques.slx]
El script \cmd{crearModelo\_Tanques.m} calcula el equilibrio y la función de transferencia, y construye por código el modelo no lineal (integradores y bloques \emph{Fcn} con las raíces cuadradas) al lado del linealizado, con el mismo escalón en $f_i$. Pruebe \cmd{df = 0.01} y \cmd{df = 0.05} y cuantifique el error final.
\end{modelo}

%% ---------------------------------------------------------------------
\section{El sistema de segundo orden}
\label{sec:so-segundoorden}
%% ---------------------------------------------------------------------

La función \eqref{eq:so-tanques-ft} y muchas otras se pueden escribir en la \clave{forma estándar del segundo orden}:
\begin{equation}
  T(s)=\frac{Y(s)}{U(s)}=K\,\frac{\wn^2}{s^2+2\zeta\wn s+\wn^2},
  \label{eq:so-estandar}
\end{equation}
con tres parámetros: la \clave{ganancia} $K$ (el valor final ante un escalón unitario), la \clave{frecuencia natural} $\wn$ (rad/s) y el \clave{factor de amortiguamiento} $\zeta$ (adimensional). Para los tanques, $\wn^2=\frac{1+K_1K_2C_2}{\tau_1\tau_2}$ y $2\zeta\wn=\frac{\tau_1+\tau_2}{\tau_1\tau_2}$. Los polos son
\begin{equation}
  p_{1,2}=-\zeta\wn\pm\wn\sqrt{\zeta^2-1},
  \label{eq:so-polos}
\end{equation}
y su naturaleza depende solo de $\zeta$:
\begin{center}
\begin{tabular}{@{}lll@{}}
\toprule
$\zeta$ & polos & respuesta al escalón\\
\midrule
$\zeta=0$ & imaginarios puros $\pm j\wn$ & oscilación sostenida (\emph{no amortiguado})\\
$0<\zeta<1$ & complejos: $-\zeta\wn\pm j\wdamp$ & oscilación que se extingue (\clave{subamortiguado})\\
$\zeta=1$ & reales e iguales: $-\wn$ (doble) & la más rápida sin sobrepico (\clave{críticamente amortiguado})\\
$\zeta>1$ & reales y distintos & lenta, sin sobrepico (\clave{sobreamortiguado})\\
\bottomrule
\end{tabular}
\end{center}
donde $\wdamp=\wn\sqrt{1-\zeta^2}$ es la \clave{frecuencia natural amortiguada}. La \cref{fig:so-zetas} lo resume.

\begin{figure}[ht]
\centering
\includegraphics[width=\linewidth]{cap03/f_zetas}
\caption{Segundo orden con $\wn=1$. Izquierda: al aumentar $\zeta$ los polos recorren la circunferencia de radio $\wn$ hasta juntarse en $-1$ ($\zeta=1$) y luego se separan sobre el eje real. Derecha: las respuestas al escalón correspondientes.}
\label{fig:so-zetas}
\end{figure}

\codigorepo{cap03_modelado_segundo_orden/C3_Ej3_SegundoOrden.m}{zetas}{\cref{fig:so-zetas}}

\subsection{Un segundo orden que sí oscila: masa, resorte y amortiguador}

La \cref{fig:so-masa} es el prototipo mecánico del segundo orden (una suspensión, un edificio en un sismo, el brazo de un disco duro). La segunda ley de Newton da
\begin{equation}
  m\ddot x+b\dot x+kx=F
  \quad\Longrightarrow\quad
  \frac{X(s)}{F(s)}=\frac{1}{ms^2+bs+k}=\frac1k\,\frac{k/m}{s^2+\frac bm s+\frac km},
  \label{eq:so-masa}
\end{equation}
así que $K=1/k$, $\wn=\sqrt{k/m}$ y $\zeta=\dfrac{b}{2\sqrt{km}}$. Con poco amortiguamiento ($b$ pequeño) la masa oscila: la energía va y viene entre el resorte y la masa, y el amortiguador la va disipando. Con $m=1$\,kg, $b=1.2$\,N\,s/m y $k=4$\,N/m resultan $\wn=2$\,rad/s y $\zeta=0.3$.

\begin{figure}[ht]
\centering
\begin{minipage}{0.4\linewidth}\centering
\includegraphics[width=\linewidth]{cap03/masa_resorte}
\end{minipage}\hfill
\begin{minipage}{0.57\linewidth}\centering
\includegraphics[width=\linewidth]{cap03/f_masa_resorte}
\end{minipage}
\caption{Masa-resorte-amortiguador y su respuesta a un escalón de $1$\,N. El valor final es $1/k=0.25$\,m y el sobrepico, $37\%$.}
\label{fig:so-masa}
\end{figure}

\begin{pausa}
En la masa-resorte de la \cref{fig:so-masa}, ¿qué valor de $b$ la vuelve críticamente amortiguada?
\respuesta{$\zeta=1\Rightarrow b=2\sqrt{km}=2\sqrt{4}=4$\,N\,s/m.}
\end{pausa}

%% ---------------------------------------------------------------------
\section{La respuesta del subamortiguado, en números}
\label{sec:so-caracteristicas}
%% ---------------------------------------------------------------------

El escalón es la entrada de prueba clásica: es un cambio brusco, fácil de generar y suficientemente drástico para ver la dinámica (y, conocida su respuesta, se obtienen las demás por linealidad). Tomemos $K=1$, condiciones iniciales nulas y $0<\zeta<1$. Con $U(s)=1/s$,
\[
  Y(s)=\frac{\wn^2}{s(s^2+2\zeta\wn s+\wn^2)}
  =\frac1s-\frac{s+\zeta\wn}{(s+\zeta\wn)^2+\wdamp^2}-\frac{\zeta\wn}{(s+\zeta\wn)^2+\wdamp^2},
\]
y con las dos últimas filas de la tabla de transformadas,
\[
  y(t)=1-e^{-\zeta\wn t}\left(\cos\wdamp t+\frac{\zeta}{\sqrt{1-\zeta^2}}\sin\wdamp t\right).
\]
Con la identidad $C\cos\theta+D\sin\theta=\sqrt{C^2+D^2}\,\sin\bigl(\theta+\tan^{-1}(C/D)\bigr)$ queda
\begin{equation}
  y(t)=1-\frac{e^{-\zeta\wn t}}{\sqrt{1-\zeta^2}}\,\sin\left(\wdamp t+\beta\right),\qquad \beta=\cos^{-1}\zeta.
  \label{eq:so-subamortiguado}
\end{equation}
La salida oscila a la frecuencia $\wdamp$ dentro de las \clave{envolventes} $1\pm e^{-\zeta\wn t}/\sqrt{1-\zeta^2}$, y el error $1-y(t)$ tiende a cero. (Con $\zeta=0$ no: oscila para siempre a $\wn$.)

El ángulo $\beta$ tiene una lectura geométrica muy útil (\cref{fig:so-beta}): es el ángulo entre el semieje real negativo y la recta que une el origen con el polo. Como la distancia del polo al origen es $\wn$ y su parte real es $-\zeta\wn$, $\cos\beta=\zeta$.

\begin{figure}[ht]
\centering
\includegraphics[width=0.55\linewidth]{cap03/plano_s_beta}
\caption{Geometría de los polos subamortiguados: distancia al origen $\wn$, parte real $-\zeta\wn$, parte imaginaria $\wdamp$ y $\cos\beta=\zeta$.}
\label{fig:so-beta}
\end{figure}

\subsection{Las cuatro características}

\begin{figure}[ht]
\centering
\includegraphics[width=0.95\linewidth]{cap03/f_caracteristicas}
\caption{Características de la respuesta al escalón con $\zeta=0.4$ y $\wn=1$. La franja verde es la banda del $\pm2\%$; las líneas punteadas, las envolventes.}
\label{fig:so-caracteristicas}
\end{figure}

\paragraph{Tiempo de subida $\tr$} Hay varias definiciones: del $10$ al $90\%$ del valor final (la de \cmd{stepinfo} y la usual para sobreamortiguados), del $5$ al $95\%$, o del $0$ al $100\%$ (común para subamortiguados). Para la de $0$--$100\%$, $y(\tr)=1$ en \eqref{eq:so-subamortiguado} exige $\sin(\wdamp\tr+\beta)=0$, es decir $\wdamp\tr+\beta=\pi$:
\begin{equation}
  \tr=\frac{\pi-\beta}{\wdamp}\quad(0\text{--}100\%).
  \label{eq:so-tr}
\end{equation}
Para la de $10$--$90\%$ no hay fórmula cerrada; un ajuste a la solución exacta es
\begin{equation}
  \tr\approx\frac{2.3\zeta^2-0.078\zeta+1.12}{\wn}\quad(10\text{--}90\%),
  \label{eq:so-tr1090}
\end{equation}
con error menor que el $3.1\%$ para $0.1\leq\zeta\leq0.9$ (lo comprobamos numéricamente; otra aproximación muy citada, $(1-0.4167\zeta+2.917\zeta^2)/\wn$, se equivoca hasta un $14\%$ con $\zeta$ pequeño).

\paragraph{Tiempo de pico $\tp$} Es el instante del primer máximo. Derivando \eqref{eq:so-subamortiguado} e igualando a cero se llega a $\frac{\wn}{\sqrt{1-\zeta^2}}e^{-\zeta\wn t}\sin\wdamp t=0$; el primer cero positivo es $\wdamp\tp=\pi$:
\begin{equation}
  \tp=\frac{\pi}{\wdamp}.
  \label{eq:so-tp}
\end{equation}

\paragraph{Sobrepico porcentual $\OS$} Es cuánto se pasa el máximo del valor final, en porcentaje: $\OS=\frac{y(\tp)-y(\infty)}{y(\infty)}\times100\%$. Evaluando en $\tp$:
\begin{equation}
  \OS=100\,e^{-\zeta\pi/\sqrt{1-\zeta^2}}\ \%,
  \qquad\text{y al revés,}\qquad
  \zeta=\frac{-\ln(\OS/100)}{\sqrt{\pi^2+\ln^2(\OS/100)}}.
  \label{eq:so-os}
\end{equation}
\emph{Solo depende de $\zeta$} (\cref{fig:so-os}). Más $\zeta$, menos sobrepico, pero respuesta más lenta.

\begin{figure}[ht]
\centering
\includegraphics[width=0.8\linewidth]{cap03/f_os_zeta}
\caption{Sobrepico en función de $\zeta$, con algunos valores que aparecen en el curso ($\zeta=0.4559$ da $20\%$; $\zeta=0.5912$, $10\%$; $\zeta=0.707$, $4.3\%$).}
\label{fig:so-os}
\end{figure}

\paragraph{Tiempo de establecimiento $\ts$} Es el tiempo a partir del cual la respuesta se queda dentro de una banda alrededor del valor final ($\pm2\%$ en el curso; a veces $\pm5\%$). Usando las envolventes, que decaen con constante de tiempo $1/(\zeta\wn)$, y el mismo criterio del primer orden ($4\tau$ para el $2\%$, $3\tau$ para el $5\%$):
\begin{equation}
  \ts\approx\frac{4}{\zeta\wn}\ (2\%),\qquad \ts\approx\frac{3}{\zeta\wn}\ (5\%).
  \label{eq:so-ts}
\end{equation}
Es una estimación por la envolvente: el valor exacto puede ser menor (en la \cref{fig:so-caracteristicas}, $8.41$ en lugar de $10$), porque la oscilación entra a la banda antes que la envolvente.

\begin{table}[ht]
\centering
\caption{Resumen del segundo orden subamortiguado $\wn^2/(s^2+2\zeta\wn s+\wn^2)$.}
\label{tab:so-resumen}
\renewcommand{\arraystretch}{1.9}
\begin{tabular}{@{}lll@{}}
\toprule
característica & fórmula & depende de\\
\midrule
tiempo de subida (0--100\%) & $\tr=\dfrac{\pi-\cos^{-1}\zeta}{\wn\sqrt{1-\zeta^2}}$ & $\zeta$ y $\wn$\\
tiempo de pico & $\tp=\dfrac{\pi}{\wn\sqrt{1-\zeta^2}}$ & parte imaginaria del polo\\
sobrepico & $\OS=100\,e^{-\zeta\pi/\sqrt{1-\zeta^2}}$ & solo $\zeta$ (ángulo del polo)\\
establecimiento (2\%) & $\ts\approx\dfrac{4}{\zeta\wn}$ & parte real del polo\\
\bottomrule
\end{tabular}
\end{table}

\codigorepo{cap03_modelado_segundo_orden/C3_Ej3_SegundoOrden.m}{formulas}{fórmulas contra \cmd{stepinfo}}

\begin{cuidado}[¿Qué $\tr$ reporta \matlab?]
\cmd{stepinfo} usa $10$--$90\%$. Si compara con \eqref{eq:so-tr} (de $0$--$100\%$) le van a dar números distintos y ninguno está mal. Diga siempre qué definición usa.
\end{cuidado}

\subsection{Moverse en el plano-$s$}

La \cref{tab:so-resumen} dice algo que conviene ver geométricamente: con los polos en $-\zeta\wn\pm j\wdamp$,
\begin{itemize}
  \item sobre una \emph{recta vertical} (parte real fija) el $\ts$ es constante;
  \item sobre una \emph{recta horizontal} (parte imaginaria fija) el $\tp$ es constante;
  \item sobre una \emph{recta radial} (ángulo $\beta$ fijo) el $\zeta$, y con él el $\OS$, son constantes.
\end{itemize}
La \cref{fig:so-lineas} muestra qué pasa al moverse por cada una. Por eso las especificaciones ``$\OS$ máximo y $\ts$ máximo'' se convierten en una \emph{región} del plano-$s$ donde deben quedar los polos: a la izquierda de una vertical y dentro de un cono alrededor del semieje real negativo. Esa idea es la base del lugar de las raíces.

\begin{figure}[ht]
\centering
\includegraphics[width=\linewidth]{cap03/f_lineas_plano}
\caption{Moviendo los polos por rectas del plano-$s$. Vertical: mismo $\ts$, más oscilación hacia arriba. Horizontal: mismo $\tp$, menos sobrepico hacia la izquierda. Radial: mismo $\OS$, más rápido lejos del origen.}
\label{fig:so-lineas}
\end{figure}

\subsection{Críticamente amortiguado y sobreamortiguado}

\paragraph{$\zeta=1$} Hay un polo doble en $-\wn$ y
\begin{equation}
  y(t)=1-e^{-\wn t}\,(1+\wn t),
  \label{eq:so-critico}
\end{equation}
que también se obtiene de \eqref{eq:so-subamortiguado} con $\lim_{\zeta\to1}\sin(\wdamp t)/\sqrt{1-\zeta^2}=\wn t$. Es la respuesta más rápida \emph{sin} sobrepico.

\paragraph{$\zeta>1$} Hay dos polos reales en $-s_1$ y $-s_2$, con $s_{1,2}=\zeta\wn\pm\wn\sqrt{\zeta^2-1}$, y
\begin{equation}
  y(t)=1+\frac{\wn}{2\sqrt{\zeta^2-1}}\left(\frac{e^{-s_1t}}{s_1}-\frac{e^{-s_2t}}{s_2}\right).
  \label{eq:so-sobre}
\end{equation}
Si $s_1\gg s_2$, el término en $e^{-s_1t}$ muere muy rápido y la respuesta se parece a la de un primer orden con el polo lento: $T(s)\approx\dfrac{s_2}{s+s_2}$. Es la primera aparición de la idea de \emph{polo dominante} (\cref{sec:so-ordensuperior}). Los tanques del ejemplo~\ref{ej:so-tanques} son así: $s_2=0.0024$ y $s_1=0.0434$, y $4/s_2\approx1670$\,s es una buena estimación del $\ts$ exacto ($1655$\,s).

\begin{pausa}
Un sistema tiene polos en $-2\pm j4$. Sin simular: ¿cuánto valen $\wn$, $\zeta$, $\tp$, $\OS$ y $\ts$?
\respuesta{$\wn=\sqrt{4+16}=4.47$\,rad/s, $\zeta=2/4.47=0.447$, $\tp=\pi/4=0.785$\,s, $\OS=100e^{-0.447\pi/0.894}=20.8\%$ y $\ts\approx4/2=2$\,s.}
\end{pausa}

%% ---------------------------------------------------------------------
\section{Diagramas de bloques}
\label{sec:so-bloques}
%% ---------------------------------------------------------------------

Un proceso real tiene válvulas, sensores, transmisores y controladores, cada uno con su propia función de transferencia. Para representar esa cadena de causas y efectos usamos \clave{diagramas de bloques}: elementos \emph{unidireccionales}, cada uno con la función de transferencia entre su entrada y su salida (\cref{fig:so-elementos}).

\begin{figure}[ht]
\centering
\includegraphics[width=\linewidth]{cap03/bloques_elementos}
\caption{Los elementos de un diagrama de bloques: (a) una señal, (b) un comparador, (c) un punto de bifurcación (la señal se copia, no se reparte), (d) un sumador con varios signos y (e) un bloque.}
\label{fig:so-elementos}
\end{figure}

Tres reglas bastan para reducir casi cualquier diagrama (\cref{fig:so-reglas}):
\begin{itemize}
  \item \textbf{Serie:} $Z=G_1R$ y $Y=G_2Z$ dan $Y=G_1G_2\,R$. Las funciones de transferencia en cascada se multiplican.
  \item \textbf{Paralelo:} $Y=G_1R+G_2R=(G_1+G_2)\,R$.
  \item \textbf{Retroalimentación negativa:} con $E=R-\Hs Y$ y $Y=GE$,
  \begin{equation}
    \frac{Y(s)}{R(s)}=\frac{G(s)}{1+G(s)\Hs(s)}.
    \label{eq:so-retro}
  \end{equation}
  Con retroalimentación positiva ($E=R+\Hs Y$) el signo del denominador cambia: $G/(1-G\Hs)$.
\end{itemize}

\begin{figure}[ht]
\centering
\includegraphics[width=0.95\linewidth]{cap03/bloques_reglas}
\caption{Las tres reglas del álgebra de bloques.}
\label{fig:so-reglas}
\end{figure}

\begin{ideaclave}
El denominador de \eqref{eq:so-retro} igualado a cero, $1+G(s)\Hs(s)=0$, es la \clave{ecuación característica} del lazo cerrado. El producto $G\Hs$ es la \clave{malla abierta} $\Gma(s)$ (\cref{def:rf-malla}): lo que recorre una señal al dar una vuelta completa al lazo. Casi todo el análisis del curso (estabilidad, lugar de las raíces, Bode) gira alrededor de esa función.
\end{ideaclave}

\begin{ejemplo}[label=ej:so-dosentradas]{Dos entradas, dos salidas}
Para el diagrama de la \cref{fig:so-dosentradas}, demuestre que $Y_2(s)=G_3(s)\bigl[G_1(s)R_1(s)+G_2(s)R_2(s)\bigr]$.
\solucion La salida del sumador es $G_1R_1+G_2R_2$; el punto de bifurcación la copia en $Y_1$ y en la entrada de $G_3$. Por la regla de la serie, $Y_2=G_3Y_1=G_3(G_1R_1+G_2R_2)$. Por linealidad, cada entrada contribuye por su propio camino (superposición).
\end{ejemplo}

\begin{figure}[ht]
\centering
\includegraphics[width=0.75\linewidth]{cap03/bloques_dos_entradas}
\caption{Ejemplo~\ref{ej:so-dosentradas}.}
\label{fig:so-dosentradas}
\end{figure}

No siempre se quiere \emph{simplificar}: a veces se busca que el diagrama tenga cierta \emph{estructura}. El caso más importante es reescribir un lazo con retroalimentación $\Hs(s)$ como uno con retroalimentación unitaria, que es la forma en la que se estudia el error en estado estacionario.

\begin{ejemplo}[label=ej:so-equivalente]{Lazo abierto equivalente}
Demuestre que los dos lazos de la \cref{fig:so-equivalente} son equivalentes si
\begin{equation}
  \Gla(s)=\frac{G(s)}{1+G(s)\Hs(s)-G(s)}.
  \label{eq:so-gla}
\end{equation}
\solucion El lazo (b) tiene $T=\Gla/(1+\Gla)$. Con \eqref{eq:so-gla},
\[
  1+\Gla=\frac{1+G\Hs-G+G}{1+G\Hs-G}=\frac{1+G\Hs}{1+G\Hs-G}
  \quad\Longrightarrow\quad
  \frac{\Gla}{1+\Gla}=\frac{G}{1+G\Hs},
\]
que es el lazo cerrado de (a). A $\Gla$ la llamamos \clave{lazo abierto} (\cref{def:rf-malla}); no hay que confundirla con la malla abierta $G\Hs$.
\end{ejemplo}

\begin{figure}[ht]
\centering
\includegraphics[width=0.95\linewidth]{cap03/lazo_equivalente}
\caption{Ejemplo~\ref{ej:so-equivalente}: el mismo lazo cerrado con retroalimentación $\Hs$ y con retroalimentación unitaria.}
\label{fig:so-equivalente}
\end{figure}

En \matlab las tres reglas son \cmd{series}, \cmd{parallel} y \cmd{feedback} (o, simplemente, \cmd{*}, \cmd{+} y \cmd{feedback}); \cmd{minreal} cancela los polos y ceros repetidos que deja el álgebra:
\codigorepo{cap03_modelado_segundo_orden/C3_Ej5_Bloques.m}{algebra}{un lazo interno dentro de uno externo (problema \ref{prob:so-multilazo})}
\codigorepo{cap03_modelado_segundo_orden/C3_Ej5_Bloques.m}{equivalente}{ejemplo \ref{ej:so-equivalente}}

\begin{cuidado}[\texttt{feedback} supone signo negativo]
\cmd{feedback(G,H)} calcula $G/(1+GH)$. Para retroalimentación positiva hay que escribir \cmd{feedback(G,H,+1)}. Y \cmd{feedback(G,1)} es retroalimentación \emph{unitaria}, no ``sin retroalimentación''.
\end{cuidado}

\begin{pausa}
Si en la regla de retroalimentación $G\Hs\gg1$ en las frecuencias de interés, ¿a qué se parece $Y/R$? ¿Qué consecuencia práctica tiene?
\respuesta{$Y/R\approx1/\Hs$: el lazo cerrado depende casi solo del sensor y no de la planta. Por eso un buen sensor es más importante que una planta precisa, y por eso Black pudo construir amplificadores precisos con tubos imprecisos (\cref{cap:rf}).}
\end{pausa}

%% ---------------------------------------------------------------------
\section{Diseñar con especificaciones de tiempo}
\label{sec:so-diseno}
%% ---------------------------------------------------------------------

Las fórmulas de la \cref{sec:so-caracteristicas} sirven al revés: dadas unas especificaciones, encontrar los parámetros que las cumplen. Es nuestro primer diseño.

\begin{ejemplo}[label=ej:so-diseno]{Servo de posición con tacómetro}
En el servo de la \cref{fig:so-kkh}, la posición se mide con un sensor y la velocidad con un tacómetro de ganancia $K_h$. Halle $K$ y $K_h$ para que la respuesta al escalón tenga $\OS=20\%$ y $\tp=1$\,s.
\solucion \textbf{1. Lazo cerrado.} El camino directo es $G=\frac{K}{s}\cdot\frac{1}{s+1}=\frac{K}{s(s+1)}$ y la retroalimentación $\Hs=1+K_hs$. Por \eqref{eq:so-retro},
\[
  \frac{Y(s)}{R(s)}=\frac{K}{s^2+(1+KK_h)\,s+K}.
\]
Es la forma estándar \eqref{eq:so-estandar} con ganancia $1$, $\wn^2=K$ y $2\zeta\wn=1+KK_h$.

\textbf{2. De las especificaciones a $\zeta$ y $\wn$.} Con \eqref{eq:so-os}, $\OS=20\%$ da $\zeta=0.4559$; con \eqref{eq:so-tp}, $\wn=\dfrac{\pi}{\tp\sqrt{1-\zeta^2}}=3.5299$\,rad/s.

\textbf{3. Igualar coeficientes.} $K=\wn^2=12.46$ y $K_h=\dfrac{2\zeta\wn-1}{K}=0.178$\,s. El lazo cerrado queda
\[
  T(s)=\frac{12.46}{s^2+3.219\,s+12.46}.
\]
\textbf{4. Verificar} (\cref{fig:so-diseno}): \cmd{stepinfo} reporta $20.0\%$ y $1.00$\,s.
\end{ejemplo}

\begin{figure}[ht]
\centering
\includegraphics[width=0.85\linewidth]{cap03/lazo_kkh}
\caption{Servo de posición con retroalimentación de posición y de velocidad (tacómetro).}
\label{fig:so-kkh}
\end{figure}

\begin{figure}[ht]
\centering
\includegraphics[width=0.85\linewidth]{cap03/f_diseno_kkh}
\caption{Ejemplo~\ref{ej:so-diseno}. Con $K=12.46$ y sin tacómetro el sobrepico sería del $64\%$: la retroalimentación de velocidad agrega amortiguamiento sin cambiar $\wn$.}
\label{fig:so-diseno}
\end{figure}

\codigorepo{cap03_modelado_segundo_orden/C3_Ej4_DisenoKKh.m}{diseno}{ejemplo \ref{ej:so-diseno}}

\begin{modelo}[Servo.slx]
El script \cmd{crearModelo\_Servo.m} arma el servo como $K\to\frac{1}{s+1}\to\frac1s$: la salida del motor es la velocidad, que el tacómetro realimenta sin necesidad de un bloque derivador. Cambie $K$ y $K_h$ por separado y observe qué hace cada uno con el sobrepico y el tiempo de pico.
\end{modelo}

\begin{pausa}
En el ejemplo~\ref{ej:so-diseno}, ¿qué pasa con el $\OS$ y el $\tp$ si se duplica $K$ dejando $K_h=0.178$?
\respuesta{$\wn=\sqrt{24.92}=4.99$ y $2\zeta\wn=1+24.92\cdot0.178=5.44$, luego $\zeta=0.545$. El sobrepico baja (a $13\%$) y el $\tp$ también ($\pi/(4.99\sqrt{1-0.297})=0.75$\,s): con tacómetro, subir $K$ no necesariamente empeora el sobrepico.}
\end{pausa}

%% ---------------------------------------------------------------------
\section{La ley de Mason}
\label{sec:so-mason}
%% ---------------------------------------------------------------------

Con muchos lazos entrelazados, el álgebra de bloques se vuelve un juego de paciencia. Hay una alternativa que da la respuesta de una vez.

\begin{historia}[Samuel J.\ Mason y los grafos de flujo de señal]
En 1953, Samuel J.\ Mason, del MIT, propuso representar un sistema de ecuaciones lineales como un \emph{grafo de flujo de señal} y mostró cómo leer la retroalimentación en él~\citep{Mas53}; en 1956 publicó la fórmula general de la ganancia entre dos nodos, que tiene en cuenta los lazos que no se tocan~\citep{Mas56}. Su fórmula completó lo que Arnold Tustin había dejado sin resolver en un trabajo previo sobre diagramas de flujo. Hoy se la conoce como \emph{Mason's gain formula} o \clave{ley de Mason}.
\end{historia}

En un grafo de flujo, las señales son \emph{nodos} y cada \emph{rama} dirigida lleva una ganancia (una función de transferencia). Un nodo suma las señales de todas las ramas que le llegan y transmite esa suma por todas las que salen. Vocabulario~\citep{ogata10}:

\begin{definicion}[label=def:so-grafo]{Grafo de flujo de señal}
\begin{itemize}
  \item \textbf{Nodo de entrada} (\emph{source}): solo tiene ramas salientes. \textbf{Nodo de salida} (\emph{sink}): solo entrantes. \textbf{Nodo mixto}: ambas.
  \item \textbf{Trayecto}: un recorrido por ramas en el sentido de las flechas. \textbf{Trayecto directo}: va del nodo de entrada al de salida sin pasar dos veces por el mismo nodo.
  \item \textbf{Lazo}: un trayecto cerrado que no repite nodos. Dos lazos son \textbf{disyuntos} si no tienen ningún nodo en común (no se tocan).
  \item \textbf{Ganancia} de un trayecto o de un lazo: el producto de las ganancias de sus ramas.
\end{itemize}
\end{definicion}

\begin{definicion}[label=def:so-mason]{Ley de Mason}
La función de transferencia entre el nodo de entrada y uno de salida es
\begin{equation}
  T(s)=\frac{1}{\Delta}\sum_{k=1}^{n}P_k\,\Delta_k,
  \label{eq:so-mason}
\end{equation}
donde $n$ es el número de trayectos directos, $P_k$ la ganancia del $k$-ésimo, $\Delta$ el \clave{discriminante} del grafo,
\begin{equation}
  \Delta=1-\sum_a L_a+\sum_{b,c}L_bL_c-\sum_{d,e,f}L_dL_eL_f+\cdots
  \label{eq:so-discriminante}
\end{equation}
(uno, menos la suma de las ganancias de todos los lazos, más la suma de los productos de pares de lazos disyuntos, menos la de ternas de lazos disyuntos, y así sucesivamente), y $\Delta_k$ el discriminante de la parte del grafo que \emph{no toca} el trayecto $P_k$ (se calcula como $\Delta$, eliminando los lazos que tocan $P_k$).
\end{definicion}

\begin{cuidado}[Discriminante, no determinante]
En la versión anterior se llamaba ``determinante'' a $\Delta$. Lo llamamos \clave{discriminante} para no confundirlo con el determinante de una matriz (aunque, en el fondo, es el determinante del sistema de ecuaciones de los nodos, y por eso funciona).
\end{cuidado}

Antes de un caso grande, comprobemos con el lazo de siempre: un trayecto $P_1=G$, un lazo $L_1=-G\Hs$ que toca al trayecto, así que $\Delta=1+G\Hs$, $\Delta_1=1$ y $T=G/(1+G\Hs)$, la regla \eqref{eq:so-retro}.

\begin{ejemplo}[label=ej:so-mason]{Dos trayectos y cuatro lazos}
Halle $Y(s)/R(s)$ para el grafo de la \cref{fig:so-mason}.
\solucion \textbf{Trayectos directos:} $P_1=G_1G_2G_3G_4G_5G_6$ (nodos 1-2-3-4-5-6-7) y $P_2=G_1G_2G_7G_8G_9$ (nodos 1-2-3-8-9-7).

\textbf{Lazos:} $L_1=-G_2H_1$ (nodos 2-3), $L_2=-G_2G_3G_4H_4$ (2-3-4-5), $L_3=-G_4H_2$ (4-5) y $L_4=-G_8H_3$ (8-9).

\textbf{Lazos disyuntos:} en pares, $L_1L_3$, $L_1L_4$, $L_2L_4$ y $L_3L_4$ ($L_1$ y $L_2$ comparten los nodos 2 y 3; $L_2$ y $L_3$, el 4 y el 5). En ternas, solo $L_1L_3L_4$.
\[
  \Delta=1-(L_1+L_2+L_3+L_4)+(L_1L_3+L_1L_4+L_2L_4+L_3L_4)-L_1L_3L_4.
\]
\textbf{Cofactores:} $P_1$ toca los nodos 1 a 7, así que solo queda libre $L_4$: $\Delta_1=1-L_4$. $P_2$ toca 1, 2, 3, 8, 9 y 7, así que solo queda libre $L_3$: $\Delta_2=1-L_3$. Finalmente,
\[
  \frac{Y(s)}{R(s)}=\frac{P_1(1-L_4)+P_2(1-L_3)}{\Delta}.
\]
\end{ejemplo}

\begin{figure}[ht]
\centering
\includegraphics[width=0.95\linewidth]{cap03/mason_grafo}
\caption{Ejemplo~\ref{ej:so-mason}: grafo de flujo de señal con nueve nodos, dos trayectos directos y cuatro lazos (en coral, las ramas de retroalimentación).}
\label{fig:so-mason}
\end{figure}

\codigorepo{cap03_modelado_segundo_orden/C3_Ej6_Mason.m}{mason}{ejemplo \ref{ej:so-mason}: Mason contra resolver las ecuaciones de los nodos}

\begin{ideaclave}
El discriminante es siempre el denominador de la función de transferencia: \alerta{$\Delta(s)=0$ es la ecuación característica} y sus raíces son los polos del lazo cerrado. Para el lazo simple, $\Delta=1+G\Hs$: lo mismo que encontramos con el álgebra de bloques.
\end{ideaclave}

\begin{receta}[Aplicar la ley de Mason]
\begin{enumerate}
  \item Numerar los nodos. Listar \emph{todos} los trayectos directos con sus nodos.
  \item Listar todos los lazos con sus nodos y sus ganancias (incluyendo el signo de cada retroalimentación).
  \item Buscar pares, ternas\ldots\ de lazos que no compartan nodos y armar $\Delta$.
  \item Para cada trayecto, tachar los lazos que lo tocan y armar $\Delta_k$ con los que quedan.
  \item Aplicar \eqref{eq:so-mason} y, si se puede, verificar con álgebra de bloques o con \matlab.
\end{enumerate}
\end{receta}

\begin{pausa}
En la \cref{fig:so-mason}, si se elimina la rama $-H_4$, ¿cuánto vale el nuevo $\Delta$?
\respuesta{Quedan $L_1$, $L_3$ y $L_4$, todos disyuntos entre sí: $\Delta=1-(L_1+L_3+L_4)+(L_1L_3+L_1L_4+L_3L_4)-L_1L_3L_4=(1-L_1)(1-L_3)(1-L_4)$. Cuando ningún lazo se toca, el discriminante se factoriza.}
\end{pausa}

%% ---------------------------------------------------------------------
\section{Sistemas de orden superior y polos dominantes}
\label{sec:so-ordensuperior}
%% ---------------------------------------------------------------------

Un lazo cerrado con controlador, válvula y sensor rara vez es de segundo orden. ¿Sirve entonces todo lo anterior? Sí, gracias a una idea: los \clave{polos dominantes}, los que contienen la mayor parte de la información del comportamiento. Si los demás polos están lejos (a la izquierda), sus modos mueren rápido y la respuesta se parece a la de un segundo orden. Así no hace falta antitransformar para estimar $\OS$ y $\ts$.

\paragraph{Cuándo despreciar un polo} Para $T(s)=\dfrac{1}{(s^2+2\zeta\wn s+\wn^2)(\gamma s+1)}$, P.\,R.\ Clement mostró que el $\OS$ y el $\ts$ se parecen a los del segundo orden cuando~\citep{Cle60}
\begin{equation}
  \left|\frac1\gamma\right|\geq10\,|\zeta\wn|,
  \label{eq:so-clement}
\end{equation}
es decir, cuando la parte real del polo adicional es al menos diez veces la de los dominantes. Es una regla conservadora: con cinco veces suele bastar en la práctica.

\begin{ejemplo}[label=ej:so-dominantes]{Un polo lejano}
Para $T(s)=\dfrac{20}{(s+10)(s^2+2s+2)}$ obtenga la respuesta al escalón y compárela con la del segundo orden dominante.
\solucion Los polos dominantes están en $-1\pm j$ ($\zeta\wn=1$) y el tercero en $-10$: se cumple \eqref{eq:so-clement}. Con fracciones parciales,
\[
  Y(s)=\frac1s-\frac{1/41}{s+10}-\frac{\frac{40}{41}(s+1)+\frac{50}{41}}{(s+1)^2+1},
  \qquad
  y(t)=1-\tfrac{40}{41}e^{-t}\cos t-\tfrac{50}{41}e^{-t}\sin t-\tfrac1{41}e^{-10t}.
\]
El término del polo en $-10$ es pequeño ($1/41$) \emph{y} rápido. La aproximación debe conservar la ganancia DC: $T(s)\approx\dfrac{20/10}{s^2+2s+2}$. Los sobrepicos son $4.27\%$ y $4.32\%$ (\cref{fig:so-dominantes}, izquierda).
\end{ejemplo}

\paragraph{Los ceros también cuentan} El análisis anterior supone que no hay ceros cerca de los polos dominantes. Un cero en $-a$ cerca de ellos cambia mucho la respuesta: en $T(s)=\dfrac{(\wn^2/a)(s+a)}{s^2+2\zeta\wn s+\wn^2}$ el sobrepico depende de $a/(\zeta\wn)$, y crece cuanto más cerca está el cero. Con un polo y un cero adicionales, $T(s)=\dfrac{(\wn^2/a)(s+a)}{(s^2+2\zeta\wn s+\wn^2)(\tau s+1)}$, ambos se pueden ignorar si $a\gg\zeta\wn$ y $1/\tau\gg\zeta\wn$; si no, no.

\begin{ejemplo}[label=ej:so-cero]{Un cero y un polo que no se pueden ignorar}
Analice $T(s)=\dfrac{62.5\,(s+2.5)}{(s^2+6s+25)(s+6.25)}$.
\solucion La ganancia DC es $1$. Los dominantes tienen $\zeta=0.6$, $\wn=5$ ($\zeta\wn=3$), el cero está en $-2.5$ y el polo extra en $-6.25$ ($\tau=0.16$). Ninguno está lejos. Comparemos tres modelos (\cref{fig:so-dominantes}, derecha):
\begin{center}\small
\begin{tabular}{@{}lccc@{}}
\toprule
modelo & $\OS$ & $\ts$ (s) & comentario\\
\midrule
completo & $38.0\%$ & $1.59$ & la respuesta real\\
sin el polo: $\dfrac{10(s+2.5)}{s^2+6s+25}$ & $54.5\%$ & $1.47$ & el cero solo dispara el sobrepico\\
solo los dominantes: $\dfrac{25}{s^2+6s+25}$ & $9.5\%$ & $1.19$ & la fórmula del segundo orden\\
\bottomrule
\end{tabular}
\end{center}
El tercer polo amortigua el sobrepico que provoca el cero, pero alarga el $\ts$; ninguna de las dos aproximaciones sirve.
\end{ejemplo}

\begin{figure}[ht]
\centering
\includegraphics[width=\linewidth]{cap03/f_dominantes}
\caption{Izquierda: ejemplo~\ref{ej:so-dominantes}, el polo en $-10$ se puede despreciar. Derecha: ejemplo~\ref{ej:so-cero}, el cero en $-2.5$ y el polo en $-6.25$ no.}
\label{fig:so-dominantes}
\end{figure}

\codigorepo{cap03_modelado_segundo_orden/C3_Ej7_Dominantes.m}{cero}{ejemplo \ref{ej:so-cero}}

\subsection{Ubicación de las raíces y respuesta transitoria}

En general, el lazo cerrado es $T(s)=\dfrac{\sum_kP_k\Delta_k}{\Delta(s)}$ y sus polos son las raíces de $\Delta(s)=0$ (para un lazo simple, $1+G\Hs=0$). Sin raíces repetidas y con ganancia unitaria, la respuesta al escalón tiene la forma
\[
  y(t)=1+\sum_{i}A_i\,e^{-\sigma_it}+\sum_kD_k\,e^{-\alpha_kt}\sin(\omega_kt+\theta_k),
\]
una suma de exponenciales (polos reales $-\sigma_i$) y senoides amortiguadas (polos $-\alpha_k\pm j\omega_k$). Para que la salida sea acotada todas las partes reales deben ser negativas (\cref{fig:po-modos}). Y para leer la forma de la respuesta:
\begin{itemize}
  \item La respuesta está determinada por los polos \emph{y} por los residuos.
  \item Los polos más cercanos al eje imaginario son los modos más lentos: dominan la forma de la respuesta, a menos que su residuo sea pequeño.
  \item Los polos fijan \emph{qué} modos hay; los ceros fijan \emph{cuánto} pesa cada uno. Un cero cerca de un polo reduce la contribución de ese modo (en el límite, lo cancela).
\end{itemize}

\begin{pausa}
En el ejemplo~\ref{ej:so-dominantes}, ¿cuánto pesaría el modo $e^{-10t}$ si el polo extra estuviera en $-3$ en lugar de $-10$ (manteniendo la ganancia DC)? ¿Se podría despreciar?
\respuesta{Con $T=\dfrac{6}{(s+3)(s^2+2s+2)}$, el residuo de $Y=T/s$ en $-3$ es $\dfrac{6}{(-3)(9-6+2)}=-0.4$: casi la mitad del valor final, y el modo $e^{-3t}$ es apenas tres veces más rápido que la envolvente $e^{-t}$. No se puede despreciar.}
\end{pausa}

%% ---------------------------------------------------------------------
\section{Simplificación de sistemas lineales}
\label{sec:so-reduccion}
%% ---------------------------------------------------------------------

A veces conviene \emph{reemplazar} un modelo de orden alto por uno más simple, por ejemplo para diseñar un controlador con las fórmulas del segundo orden. Veremos dos maneras.

\paragraph{1. Despreciar un polo rápido, conservando la ganancia} Si un polo real negativo tiene magnitud mucho mayor que los demás, se elimina, pero dividiendo por su valor para no cambiar la ganancia en baja frecuencia. Por ejemplo,
\[
  G(s)=\frac{K}{s(s+2)(s+30)}\approx\frac{K/30}{s(s+2)},
\]
porque para $|s|\ll30$, $s+30\approx30$. Con $K=60$ y retroalimentación unitaria los sobrepicos son $5.3\%$ y $4.3\%$ y las respuestas casi se superponen (\cref{fig:so-reduccion}, izquierda). La aproximación empeora cuando el lazo es más agresivo: con $K=300$ los polos dominantes del lazo cerrado se alejan del origen, se acercan en magnitud al de $-30$, y los sobrepicos pasan a $42.3\%$ y $35.1\%$.

\codigorepo{cap03_modelado_segundo_orden/C3_Ej8_Reduccion.m}{despreciar}{despreciar el polo en $-30$}

\paragraph{2. Igualación de momentos} Cuando ningún polo es despreciable hay que buscar un modelo de orden menor que se \emph{parezca} en frecuencia al original~\citep{DoBi08}. Escribimos el modelo de orden alto y el de orden bajo con la misma ganancia $K$ (para conservar el estado estacionario) y con término independiente $1$:
\[
  G_H(s)=K\,\frac{a_ms^m+\cdots+a_1s+1}{b_ns^n+\cdots+b_1s+1},\qquad
  G_L(s)=K\,\frac{c_ps^p+\cdots+c_1s+1}{d_gs^g+\cdots+d_1s+1},\quad p\leq g<n.
\]
Sea $\dfrac{G_H(s)}{G_L(s)}=\dfrac{M(s)}{\Delta(s)}$ (aquí $\Delta$ no es el discriminante de Mason, solo un denominador). Se definen
\begin{equation}
  M_{2q}=\sum_{k=0}^{2q}\frac{(-1)^{k+q}\,M^{(k)}(0)\,M^{(2q-k)}(0)}{k!\,(2q-k)!},\qquad q=0,1,2,\dots
  \label{eq:so-momentos}
\end{equation}
($M^{(k)}$ es la $k$-ésima derivada) y lo mismo para $\Delta_{2q}$. Los coeficientes desconocidos $c_i,d_i$ se obtienen igualando
\[
  M_{2q}=\Delta_{2q},\qquad q=1,2,\dots
\]
hasta tener tantas ecuaciones como incógnitas. (Como $M^{(k)}(0)/k!$ es el coeficiente de $s^k$, la fórmula solo usa coeficientes.) La idea: $|G_H(\jw)/G_L(\jw)|^2=\dfrac{\sum_q M_{2q}\omega^{2q}}{\sum_q\Delta_{2q}\omega^{2q}}$, así que igualar los primeros momentos hace que ese cociente se parezca a $1$ en baja frecuencia.

\begin{ejemplo}[label=ej:so-momentos]{De tercer a segundo orden}
Reduzca $G_H(s)=\dfrac{6}{s^3+6s^2+11s+6}=\dfrac{1}{\frac16s^3+s^2+\frac{11}6s+1}$ a $G_L(s)=\dfrac{1}{d_2s^2+d_1s+1}$.
\solucion $\dfrac{G_H}{G_L}=\dfrac{1+d_1s+d_2s^2}{\frac16s^3+s^2+\frac{11}6s+1}$, así que $M(s)=1+d_1s+d_2s^2$ y $\Delta(s)=\frac16s^3+s^2+\frac{11}6s+1$. Las derivadas en cero son $M^{(0)}=1$, $M^{(1)}=d_1$, $M^{(2)}=2d_2$; $\Delta^{(0)}=1$, $\Delta^{(1)}=\frac{11}6$, $\Delta^{(2)}=2$, $\Delta^{(3)}=1$. Con dos incógnitas usamos $q=1,2$:
\[
  M_2=d_1^2-2d_2,\quad \Delta_2=\tfrac{121}{36}-2=\tfrac{49}{36};\qquad
  M_4=d_2^2,\quad \Delta_4=1-2\cdot\tfrac{11}{6}\cdot\tfrac16=\tfrac{7}{18}.
\]
Entonces $d_2=\sqrt{7/18}=0.6236$ y $d_1=\sqrt{49/36+2d_2}=1.615$ (se toman las raíces positivas: con las negativas $G_L$ tendría polos inestables). Así
\[
  G_L(s)=\frac{1}{0.6236s^2+1.615s+1}=\frac{1.604}{s^2+2.590s+1.604},
\]
con polos en $-1.024$ y $-1.565$ (los de $G_H$ estaban en $-1$, $-2$ y $-3$). La \cref{fig:so-reduccion} (derecha) compara los escalones: la diferencia máxima es de $0.10$, al principio, porque $G_L$ tiene grado relativo $2$ y arranca más rápido.
\end{ejemplo}

\begin{figure}[ht]
\centering
\includegraphics[width=\linewidth]{cap03/f_reduccion}
\caption{Izquierda: despreciar el polo en $-30$ (lazo cerrado, $K=60$). Derecha: ejemplo~\ref{ej:so-momentos}, igualación de momentos.}
\label{fig:so-reduccion}
\end{figure}

\codigorepo{cap03_modelado_segundo_orden/C3_Ej8_Reduccion.m}{momentos}{ejemplo \ref{ej:so-momentos} con \cmd{reduccionMomentos}}

%% ---------------------------------------------------------------------
\section{Para llevar}
\label{sec:so-cierre}
%% ---------------------------------------------------------------------

\begin{parallevar}
\begin{enumerate}
  \item Linealizar es reemplazar cada no linealidad por su recta tangente en el punto de operación: $\dot{\dev{x}}=\sum\frac{\partial f}{\partial x_j}\big|_\ast\dev{x}_j+\sum\frac{\partial f}{\partial u_k}\big|_\ast\dev{u}_k$. Solo vale cerca del punto.
  \item Segundo orden: $K\wn^2/(s^2+2\zeta\wn s+\wn^2)$. $\zeta$ decide si oscila; $\wn$, qué tan rápido.
  \item Subamortiguado: $\tp=\pi/\wdamp$ (parte imaginaria), $\OS=100e^{-\zeta\pi/\sqrt{1-\zeta^2}}$ (ángulo), $\ts\approx4/(\zeta\wn)$ (parte real), $\tr=(\pi-\cos^{-1}\zeta)/\wdamp$ ($0$--$100\%$). Especificaciones de tiempo $\leftrightarrow$ región del plano-$s$.
  \item Bloques: serie multiplica, paralelo suma, retroalimentación da $G/(1+G\Hs)$. $1+G\Hs=0$ es la ecuación característica; $G\Hs$ es la malla abierta.
  \item Mason: $T=\frac1\Delta\sum P_k\Delta_k$. El discriminante $\Delta$ es el denominador: $\Delta=0$ es la ecuación característica.
  \item Polos dominantes: un polo extra se desprecia si está unas diez veces más a la izquierda (Clement), \emph{conservando la ganancia DC}; un cero cerca de los dominantes cambia mucho el sobrepico.
  \item Reducir orden: despreciar polos rápidos o igualar momentos. Siempre verificar con una simulación.
\end{enumerate}
\end{parallevar}

\begin{profundiza}[Correcciones respecto a la versión anterior del capítulo]
(i) El ejemplo numérico de los dos tanques ($\tau_1=1$, $\tau_2=4$, $C_2=2$) no era físico: en esa planta $C_2=-C_1$ y $\zeta\geq1$ siempre. Se reemplazó por datos físicos y el subamortiguado se ilustra con la masa-resorte. (ii) $\beta=\cos^{-1}\zeta$ (decía $\beta=\cos\zeta$). (iii) La condición para polos complejos de los tanques estaba mal escrita; con $C_2=-C_1$ nunca se cumple. (iv) El sobrepico se escribe $\OS$ (antes $MP$, que se confundía con el pico de resonancia), ``tiempo de establecimiento'' (no ``de asentamiento'') y $\Delta$ es el \emph{discriminante}. (v) Ejemplo con cero y polo adicionales: $54.5\%$ y $38.0\%$ (antes $55\%$ y $38\%$) y $\ts=1.47$ y $1.59$\,s (antes $1.33$, que era la fórmula $4/(\zeta\wn)$, y $1.6$). (vi) La aproximación de $\tr$ de $10$--$90\%$ se verificó: error menor que el $3.1\%$ para $0.1\leq\zeta\leq0.9$. (vii) El ejercicio de reducción de orden con un cero tenía una respuesta que no satisface las ecuaciones de momentos; se reformuló. Todo se verificó con \texttt{herramientas/cap03\_numeros.py} y \cmd{verificarCap03.m}.
\end{profundiza}

%% ---------------------------------------------------------------------
\section{Problemas}
\label{sec:so-problemas}
%% ---------------------------------------------------------------------

Los problemas marcados con \etiquetaM{} requieren \matlab/\simulink; los marcados con \etiquetaEstrella{} son de profundización.

\subsection*{Linealización}

\begin{problema}[linealización]\label{prob:so-lin}
Linealice alrededor del equilibrio y obtenga la función de transferencia, la ganancia y la constante de tiempo: (a) $\dot x=-x+2\sqrt x+u$ con $\eq{u}=0$ y $\eq{x}>0$; (b) $\dot x=-x^3+2xu$ con $\eq{u}=2$ (equilibrio no trivial). (c) \etiquetaM{} Para el tanque de la pausa de la \cref{sec:so-linealizacion}, grafique $c\sqrt h$ y su aproximación lineal y encuentre el rango de alturas en el que el error relativo es menor que el $5\%$. Compárelo con la \cref{fig:so-lin1}.
\end{problema}

\begin{problema}[estructura de los tanques]\label{prob:so-bloques-tanques}
Dibuje el diagrama de bloques que corresponde a las dos ecuaciones de $\Dev{H}_1$ y $\Dev{H}_2$ de la \cref{sec:so-tanques} (un lazo de retroalimentación con ganancia $C_2$). Redúzcalo con \eqref{eq:so-retro} y con la ley de Mason y verifique \eqref{eq:so-tanques-ft}. ¿De qué signo es la retroalimentación cuando $C_2=-C_1$?
\end{problema}

\begin{problema}[tanques con otros datos]\label{prob:so-tanques}
Con $A_1=1$\,m$^2$, $A_2=2$\,m$^2$, $C'_{v1}=0.2$, $C'_{v2}=0.1$\,m$^{5/2}$/s y $\eq{h_2}=1$\,m: (a) halle $\eq{f_i}$ y $\eq{h_1}$; (b) $C_1$, $C_2$, $C_3$, $K_1$, $\tau_1$, $K_2$, $\tau_2$; (c) la función de transferencia, su ganancia, $\wn$, $\zeta$ y sus polos; (d) la respuesta ante un escalón de $0.01$\,m$^3$/s y su $\ts$; (e) verifique que en el nuevo estado estacionario $h_1>h_2$ (si no, el modelo \eqref{eq:so-valvulas} pierde sentido); (f) \etiquetaM{} adapte \cmd{crearModelo\_Tanques.m} y compare los modelos lineal y no lineal con escalones de $0.01$ y $0.05$\,m$^3$/s.
\end{problema}

\begin{problema}[los tanques nunca oscilan]\label{prob:so-zeta-tanques}
\etiquetaEstrella{} Con $A_1,A_2,C'_{v1},C'_{v2}>0$ arbitrarios: (a) muestre que $C_2=-C_1$ y $C_1,C_3>0$ en cualquier punto de operación admisible; (b) muestre que $\wn^2=\dfrac{C_1C_3}{A_1A_2}$ y $2\zeta\wn=\dfrac{C_1+C_3}{A_2}+\dfrac{C_1}{A_1}$; (c) con la desigualdad de las medias aritmética y geométrica, demuestre que $\zeta\geq1$. ¿Cuándo se alcanzaría $\zeta=1$? ¿Es realizable? (d) Explique por qué el ejemplo con $C_2=2$ de la versión anterior no corresponde a ningún par de tanques. (e) Proponga una modificación física (por ejemplo, una bomba en lugar de la válvula entre tanques) que permita $C_2>0$ y diga cuándo el sistema sería subamortiguado.
\end{problema}

\subsection*{Segundo orden y diseño}

\begin{problema}[cuatro sistemas]\label{prob:so-cuatro}
Para $T_1=\dfrac{50}{s^2+6s+25}$, $T_2=\dfrac{16}{s^2+16}$, $T_3=\dfrac{9}{s^2+6s+9}$ y $T_4=\dfrac{100}{s^2+30s+100}$: (a) identifique $K$, $\wn$, $\zeta$, clasifíquelos y ubique sus polos; (b) para el subamortiguado calcule $\tp$, $\OS$, $\ts$ al $2$ y al $5\%$, y $\tr$ con las dos definiciones; (c) para el crítico, con \eqref{eq:so-critico}, halle el instante en que llega al $95\%$; (d) para el sobreamortiguado, construya la aproximación $s_2/(s+s_2)$ y cuantifique el error máximo ante el escalón; (e) \etiquetaM{} verifique con \cmd{stepinfo} y explique las diferencias en $\tr$.
\end{problema}

\begin{problema}[diseño con tacómetro]\label{prob:so-diseno}
En la estructura de la \cref{fig:so-kkh}, con camino directo $\dfrac{K}{s(s+4)}$: (a) muestre que $\dfrac YR=\dfrac{K}{s^2+(4+KK_h)s+K}$; (b) halle $K$ y $K_h$ para $\OS=10\%$ y $\ts=1.5$\,s; (c) calcule $\tp$, $\tr$ y los polos; (d) si además se pide $\tp\leq0.8$\,s, ¿es compatible? Si no, ¿qué $\ts$ lo permitiría con el mismo $\OS$? Explíquelo con la \cref{fig:so-lineas}; (e) \etiquetaM{} verifique con \cmd{series} y \cmd{feedback}.
\end{problema}

\begin{problema}[región de diseño]\label{prob:so-region}
Dibuje en el plano-$s$ la región donde deben quedar los polos de un segundo orden para que $\OS\leq10\%$ y $\ts\leq2$\,s. ¿Cumplen los polos $-3\pm j3$? ¿Y $-1\pm j1$?
\end{problema}

\subsection*{Bloques y Mason}

\begin{problema}[sistema multilazo]\label{prob:so-multilazo}
Las señales cumplen $E=R-H_2Y$, $V=G_1E$, $W=G_2(V-H_1W)$, $Y=G_3W$, con $G_1=\frac1{s+1}$, $G_2=\frac2s$, $G_3=\frac5{s+3}$, $H_1=1$ y $H_2=0.5$. (a) Dibuje el diagrama de bloques. (b) Obtenga $T=Y/R$ reduciendo primero el lazo interno. (c) Obtenga $T$ con Mason (trayectos, lazos, disyuntos, $\Delta$ y $\Delta_k$). (d) Escriba la ecuación característica y sus raíces. (e) \etiquetaM{} Verifique con \cmd{feedback}.
\end{problema}

\subsection*{Orden superior y reducción}

\begin{problema}[tres aproximaciones]\label{prob:so-tres}
Para $T(s)=\dfrac{10}{s^3+6s^2+11s+11}$ (el del problema anterior): (a) halle los polos y verifique el criterio \eqref{eq:so-clement}; (b) construya la aproximación por polos dominantes (conservando la ganancia DC) y prediga $\OS$ y $\ts$; (c) obtenga otra por igualación de momentos; (d) \etiquetaM{} simule las tres, compare $\OS$ y $\ts$ y decida cuál es preferible; (e) explique por qué el criterio de (a), aunque no se cumpla, no impidió una aproximación aceptable.
\end{problema}

\begin{problema}[momentos con un cero]\label{prob:so-momentos}
\etiquetaM{} Reduzca $G_H(s)=\dfrac{s^3+7s^2+24s+24}{(s+1)(s+2)(s+3)(s+4)}$ a $G_L(s)=\dfrac{1}{d_2s^2+d_1s+1}$ por igualación de momentos (note que ahora $M(s)$ incluye el numerador de $G_H$). Compare los escalones. \etiquetaEstrella{} Intente luego $G_L=\dfrac{1+c_1s}{d_2s^2+d_1s+1}$ con $q=1,2,3$: ¿tiene solución real?
\end{problema}

\begin{problema}[un resorte no lineal]\label{prob:so-resorte}
\etiquetaEstrella{} Una masa con resorte endurecido y amortiguamiento cuadrático cumple $m\ddot x+b\,\dot x|\dot x|+k_1x+k_3x^3=F$, con $m=1$\,kg, $b=0.8$\,N\,s$^2$/m$^2$, $k_1=4$\,N/m y $k_3=1$\,N/m$^3$. (a) Halle el equilibrio para $\eq{F}=5$\,N. (b) Linealice y obtenga $\Dev{X}/\Dev{F}$, $\wn$ y $\zeta$. (c) El $\zeta$ debe llamarle la atención: explique por qué el amortiguamiento desaparece al linealizar y qué predice el modelo lineal. (d) \etiquetaM{} Simule el no lineal con un escalón adicional de $0.5$\,N y compárelo con el lineal; defina un criterio y diga durante cuántos ciclos es aceptable la predicción lineal. (e) Estime un $\zeta$ equivalente con el decremento logarítmico. (f) Repita (a)--(b) con $\eq{F}=16$\,N y explique por qué un solo modelo LTI no describe todo el rango.
\end{problema}

%% ---------------------------------------------------------------------
\section{Soluciones seleccionadas}
%% ---------------------------------------------------------------------

\paragraph{Problema~\ref{prob:so-lin}}
(a) $\eq{x}=4$; $\partial f/\partial x=-1+1/\sqrt{\eq{x}}=-0.5$, $\partial f/\partial u=1$: $\dfrac{\Dev{X}}{\Dev{U}}=\dfrac{1}{s+0.5}=\dfrac{2}{2s+1}$ ($K_1=2$, $\tau=2$). (b) $\eq{x}=2$; $-3(\eq{x})^2+2\eq{u}=-8$ y $2\eq{x}=4$: $\dfrac{4}{s+8}=\dfrac{0.5}{0.125s+1}$. (c) $2.13\leq h\leq7.51$\,m: un rango asimétrico alrededor de $\eq{h}=4$, porque la raíz cuadrada se curva más para alturas bajas.

\paragraph{Problema~\ref{prob:so-bloques-tanques}}
Es un lazo con camino directo $\frac{K_1}{\tau_1s+1}\frac{K_2}{\tau_2s+1}$ y retroalimentación $C_2$ (signo $-$ en el comparador): $T=\dfrac{K_1K_2}{(\tau_1s+1)(\tau_2s+1)+K_1K_2C_2}$. Con $C_2=-C_1<0$ la retroalimentación es \emph{positiva}: al subir $h_2$ disminuye $f_1$, el tanque 1 se llena más y empuja más agua hacia el 2. Pero su ganancia de malla, $K_1K_2C_1=C_1/(C_1+C_3)<1$, no alcanza a inestabilizar el sistema.

\paragraph{Problema~\ref{prob:so-tanques}}
(a) $\eq{f}=0.1$\,m$^3$/s, $\eq{h_1}=1.25$\,m. (b) $C_1=0.2$, $C_2=-0.2$, $C_3=0.05$\,m$^2$/s; $K_1=5$, $\tau_1=5$\,s, $K_2=0.8$, $\tau_2=8$\,s. (c) $\dfrac{20}{200s^2+65s+1}$: ganancia $20$\,s/m$^2$, $\wn=0.0707$\,rad/s, $\zeta=2.30$, polos $-0.0162$ y $-0.309$. (d) $\Delta h_2\to0.2$\,m con $\ts=245$\,s (la estimación con el polo lento, $4/0.0162=247$\,s). (e) $h_1=1.25+0.05+0.2=1.5>h_2=1.2$\,m.

\paragraph{Problema~\ref{prob:so-zeta-tanques}}
(c) $2\zeta\wn=\dfrac{C_1}{A_2}+\dfrac{C_3}{A_2}+\dfrac{C_1}{A_1}\geq\dfrac{C_3}{A_2}+\dfrac{C_1}{A_1}\geq2\sqrt{\dfrac{C_1C_3}{A_1A_2}}=2\wn$, luego $\zeta\geq1$. La primera desigualdad es estricta porque $C_1/A_2>0$: $\zeta=1$ exigiría $C_1\to0$ (válvula entre tanques cerrada), que desacopla los tanques.

\paragraph{Problema~\ref{prob:so-cuatro}}
(a) $T_1$: $K=2$, $\wn=5$, $\zeta=0.6$; $T_2$: $K=1$, $\wn=4$, $\zeta=0$; $T_3$: $K=1$, $\wn=3$, $\zeta=1$; $T_4$: $K=1$, $\wn=10$, $\zeta=1.5$ (polos $-3.82$ y $-26.2$). (b) $\tp=0.785$\,s, $\OS=9.48\%$, $\ts\approx1.33$\,s ($2\%$) y $1.0$\,s ($5\%$), $\tr=0.554$\,s ($0$--$100\%$) y $\approx0.38$\,s ($10$--$90\%$; exacto $0.371$). (c) $1.58$\,s. (d) $\dfrac{3.82}{s+3.82}$; error máximo $0.105$, al principio.

\paragraph{Problema~\ref{prob:so-diseno}}
(b) $\zeta=0.591$, $\zeta\wn=4/1.5=2.67$, $\wn=4.51$\,rad/s; $K=20.35$ y $K_h=(2\zeta\wn-4)/K=0.0655$. (c) $\tp=0.864$\,s. (d) No: $\tp\leq0.8$ con $\zeta=0.591$ exige $\wn\geq4.87$, y entonces $\ts\leq1.39$\,s: hay que mover los polos \emph{por la recta de $\zeta$ constante}, alejándolos del origen.

\paragraph{Problema~\ref{prob:so-region}}
$\zeta\geq0.591$ (un cono de semiángulo $\cos^{-1}0.591=53.8\grados$ alrededor del semieje real negativo) y $\zeta\wn\geq2$ (a la izquierda de $\Re s=-2$). $-3\pm j3$: $\zeta=0.707$ y parte real $-3$, cumple. $-1\pm j1$: $\zeta=0.707$ pero parte real $-1$, no cumple el $\ts$.

\paragraph{Problema~\ref{prob:so-multilazo}}
Lazo interno: $\frac{2/s}{1+2/s}=\frac{2}{s+2}$. $T=\dfrac{10}{(s+1)(s+2)(s+3)+5}=\dfrac{10}{s^3+6s^2+11s+11}$. Con Mason: un trayecto $P_1=G_1G_2G_3$; lazos $L_1=-G_2H_1$ y $L_2=-G_1G_2G_3H_2$, que se tocan; $\Delta=1-L_1-L_2$, $\Delta_1=1$. Polos: $-3.90$ y $-1.05\pm j1.31$.

\paragraph{Problema~\ref{prob:so-tres}}
(a) La relación entre partes reales es $3.90/1.05=3.7<10$. (b) $\dfrac{(10/11)\,2.817}{s^2+2.096s+2.817}$: $\zeta=0.624$, $\wn=1.68$; $\OS=8.1\%$ y $\ts\approx3.6$\,s. (c) $d_1=0.768$, $d_2=0.340$. (d) Original: $\OS=7.2\%$, $\ts=3.84$\,s; dominantes: $8.1\%$ y $3.56$\,s; momentos: $6.4\%$ y $3.51$\,s. (e) El residuo del polo real es pequeño.

\paragraph{Problema~\ref{prob:so-momentos}}
$d_1=1.211$ y $d_2=0.2295$ (polos $-1.03$ y $-4.25$); diferencia máxima entre escalones $\approx0.10$. Con un cero, las tres ecuaciones no tienen solución real exacta.

\paragraph{Problema~\ref{prob:so-resorte}}
(a) $4x+x^3=5\Rightarrow\eq{x}=1$\,m. (b) $k_{\text{ef}}=k_1+3k_3(\eq{x})^2=7$ y la derivada de $b\dot x|\dot x|$ en $\dot x=0$ es cero: $\dfrac{1}{s^2+7}$, $\wn=2.65$\,rad/s, $\zeta=0$. (c) El amortiguamiento cuadrático es de segundo orden en $\dot x$: la linealización lo borra y predice una oscilación eterna. (f) $\eq{x}=2$\,m, $k_{\text{ef}}=16$, $\wn=4$\,rad/s.

\respuestaspausas

         % migrado (Handout2_CtrlParte2.tex)
%% =====================================================================
%%  Capítulo 4 -- Representación en espacio de estado
%%  Migrado al formato nuevo (noviembre de 2026) desde Handout3_Ctrl.tex.
%%  Correcciones de fondo: ver MIGRACION.md y la caja "Correcciones" de la
%%  sección 4.8 (levitador sin la constante C, historia del RIAS, ejemplo del
%%  determinante, funciones de transferencia del bioreactor, truco de canon).
%% =====================================================================

\capitulo[cap:ee]{Representación en espacio de estado}
  {Otra consecuencia es que se está impulsando que la educación sea vista como una capacitación que solo conduce al trabajo, y no como un medio que conduce al pensamiento crítico y a la transformación.}
  {Junot Díaz}

\begin{objetivos}[partes/04-espacio-estado.html]
  \item Definir el estado de un sistema y escribir su representación en espacio de estado, en particular en variables de fase.
  \item Pasar de una función de transferencia a una representación de estado, y de vuelta, y relacionar los valores propios de $\Am$ con los polos.
  \item Resolver la ecuación de estado con la matriz de transición $\Phi(t)=e^{\Am t}$.
  \item Linealizar un modelo no lineal de varias variables alrededor de un equilibrio y escribirlo en espacio de estado.
\end{objetivos}

%% ---------------------------------------------------------------------
\section{Dos maneras de describir un sistema}
\label{sec:ee-intro}
%% ---------------------------------------------------------------------

Hasta aquí hemos descrito los sistemas con \textbf{funciones de transferencia}: una relación algebraica entre la entrada y la salida, en el dominio de Laplace. Es el método \emph{clásico} o \emph{frecuencial}, y tiene ventajas grandes: interconectar subsistemas es multiplicar y sumar bloques (\cref{cap:so}), y la estabilidad y el transitorio se leen casi de inmediato en los polos. Su límite es igual de claro: solo sirve para sistemas lineales e invariantes en el tiempo, o para lo que se pueda aproximar así, y solo describe la relación entre una entrada y una salida, no lo que pasa por dentro.

La \textbf{representación en espacio de estado} describe el sistema con un conjunto de ecuaciones diferenciales de \emph{primer orden}, una por cada variable que hace falta para conocer su evolución. Es más general (admite sistemas no lineales, variantes en el tiempo y con varias entradas y salidas) y muestra la dinámica interna. En este curso la usamos para sistemas lineales, en tres momentos: aquí, para escribir y resolver modelos; en el \cref{cap:es}, para la estabilidad; y en el \cref{cap:re}, para diseñar controladores que usan todo el estado.

\begin{historia}[Del Sputnik al Apolo]
El lanzamiento del Sputnik, en 1957, despertó en Occidente el interés por lo que hacía la ingeniería soviética. El primer congreso mundial de la IFAC se reunió en Moscú en 1960, y allí muchos ingenieros de control constataron que la escuela soviética trabajaba de otra manera: describía los procesos con ecuaciones diferenciales en lugar de funciones de transferencia, analizaba la estabilidad con los métodos de Lyapunov en lugar de los de Bode y Nyquist, y estudiaba el desempeño óptimo con el principio del máximo de Pontryagin en lugar del método de Wiener-Hopf. Así empezó lo que se llamó \textbf{control moderno}.

En Estados Unidos las ideas se difundieron desde la Universidad de Columbia, en el grupo de John R. Ragazzini, donde se formó Rudolf E. Kalman, y desde el RIAS (\emph{Research Institute for Advanced Studies}) de Baltimore, cuyo centro de matemáticas dirigía Solomon Lefschetz y donde Kalman trabajó desde 1958. En el propio congreso de Moscú, Kalman presentó una teoría general de los sistemas de control en espacio de estado~\citep{Kal60}; con J.~E. Bertram aplicó el método de Lyapunov, y con Richard Bucy extendió la teoría de filtrado de Wiener a sistemas variantes en el tiempo~\citep{KaBu61}. En 1963 mostró con precisión cómo se relacionan la función de transferencia y la representación de estado: qué partes del sistema ve la primera y cuáles le quedan ocultas~\citep{Kal63}. El filtro de Kalman voló en la navegación del programa Apolo.
\end{historia}

%% ---------------------------------------------------------------------
\section{Un ejemplo: el ciclo global del carbono}
\label{sec:ee-carbono}
%% ---------------------------------------------------------------------

Antes de definir nada, un modelo. Li, Jarvis y Leedal describen el ciclo global del carbono con cinco reservorios que intercambian carbono entre sí~\citep{LJL09} (\cref{fig:ee-carbono}). Las variables, en gigatoneladas de carbono (GtC), son las masas en la atmósfera $m_A$, la vegetación $m_V$, el suelo $m_S$, la superficie del océano $m_{SO}$ y el océano profundo $m_{DO}$; la entrada es la emisión humana $u_E$, en GtC/año. Cada flujo es proporcional a la masa del reservorio de donde sale, con un coeficiente de intercambio $k$ en 1/año.

\begin{figure}[ht]
\centering
\includegraphics[width=0.85\linewidth]{cap04/carbono}
\caption{Ciclo global del carbono con cinco reservorios~\citep{LJL09}. Cada flecha es un flujo proporcional a la masa del reservorio de origen.}
\label{fig:ee-carbono}
\end{figure}

Un balance de masa por reservorio (lo que entra menos lo que sale) da
\begin{align*}
  \dot m_A    &= u_E-(k_{O1}+k_{L1})m_A+k_{L2}m_V+k_{L4}m_S+k_{O2}m_{SO},\\
  \dot m_V    &= k_{L1}m_A-(k_{L2}+k_{L3})m_V,\\
  \dot m_S    &= k_{L3}m_V-k_{L4}m_S,\\
  \dot m_{SO} &= k_{O1}m_A-(k_{O2}+k_{O3})m_{SO}+k_{O4}m_{DO},\\
  \dot m_{DO} &= k_{O3}m_{SO}-k_{O4}m_{DO}.
\end{align*}
Son cinco ecuaciones de primer orden, y se pueden escribir de una vez en forma matricial, $\dot x=\Am x+\Bm u_E$, con $x=[m_A\;\,m_V\;\,m_S\;\,m_{SO}\;\,m_{DO}]^\top$,
\begin{equation}
  \Am=\begin{bmatrix}
  -(k_{O1}+k_{L1}) & k_{L2} & k_{L4} & k_{O2} & 0\\
  k_{L1} & -(k_{L2}+k_{L3}) & 0 & 0 & 0\\
  0 & k_{L3} & -k_{L4} & 0 & 0\\
  k_{O1} & 0 & 0 & -(k_{O2}+k_{O3}) & k_{O4}\\
  0 & 0 & 0 & k_{O3} & -k_{O4}
  \end{bmatrix},\qquad
  \Bm=\begin{bmatrix}1\\0\\0\\0\\0\end{bmatrix}.
  \label{eq:ee-carbono}
\end{equation}

\begin{ejemplo}[label=ej:ee-carbono]{Lo que dice la estructura de $\Am$}
Con $k_{L1}=0.0842$, $k_{L2}=0.0684$, $k_{L3}=0.0159$, $k_{L4}=0.0332$, $k_{O1}=0.0846$, $k_{O2}=0.0433$, $k_{O3}=0.0011$ y $k_{O4}=0.0001$ (todos en 1/año), ¿qué pasa con $u_E=8$ GtC/año constante?
\solucion Cada columna de $\Am$ suma cero: lo que sale de un reservorio entra a otro. Por eso $\mathbf 1^\top\Am=0$ y, como $\mathbf 1^\top\Bm=1$,
\[
  \deriv{}{t}\bigl(m_A+m_V+m_S+m_{SO}+m_{DO}\bigr)=\mathbf 1^\top\dot x=u_E :
\]
el carbono total crece exactamente al ritmo de las emisiones; los intercambios solo lo redistribuyen. La misma propiedad obliga a que $\Am$ tenga un valor propio en $\lambda=0$ (sus filas son linealmente dependientes, porque suman cero). Los otros cuatro son $-0.227$, $-0.067$, $-0.036$ y $-0.00058$ por año: constantes de tiempo de $4.4$, $15$, $28$ y $1729$ años. Con $u_E=8$ GtC/año durante 500 años, la atmósfera pasa de $590$ a $1833$ GtC (\cref{fig:ee-carbono-sim}): el valor propio nulo actúa como un integrador, y con emisión constante no hay estado estacionario.
\end{ejemplo}

\begin{figure}[ht]
\centering
\includegraphics[width=\linewidth]{cap04/f_carbono}
\caption{Ejemplo~\ref{ej:ee-carbono}: 500 años de emisiones constantes de 8 GtC/año, desde $x(0)=[590\;\,550\;\,1500\;\,900\;\,37000]^\top$ GtC.}
\label{fig:ee-carbono-sim}
\end{figure}

\codigorepo{cap04_espacio_estado/C4_Ej5_Carbono.m}{matrices}{las matrices, la conservación de la masa y los valores propios}

Con este ejemplo en mente veamos qué es el estado, cómo se escoge y qué se gana con él.

%% ---------------------------------------------------------------------
\section{El estado de un sistema}
\label{sec:ee-estado}
%% ---------------------------------------------------------------------

La premisa de la mecánica de Newton es que la evolución futura de un sistema queda determinada por su condición presente y por las fuerzas que actúen sobre él. Esa condición presente es el estado.

\begin{definicion}[label=def:ee-estado]{Estado}
El \textbf{estado} de un sistema dinámico es el conjunto más pequeño de variables cuyo valor en un instante $t_0$, junto con las entradas para $t\ge t_0$, determina por completo el comportamiento del sistema para todo $t\ge t_0$~\citep{Fri86}. Esas variables son las \textbf{variables de estado}; agrupadas forman el \textbf{vector de estado} $x(t)$, y el espacio $n$-dimensional cuyos ejes son las variables de estado es el \textbf{espacio de estado}.
\end{definicion}

El número de variables de estado (el \textbf{orden} del sistema) es el número de condiciones iniciales que hacen falta para resolver sus ecuaciones; es único. Las variables, en cambio, no lo son: el mismo sistema admite muchas elecciones distintas. Una regla práctica: en un sistema físico, cada elemento que almacena energía de manera independiente aporta una variable de estado (la posición de un resorte, la velocidad de una masa, la carga de un condensador, la corriente de una bobina, el nivel de un tanque).

\paragraph{La forma general.} El modelo se escribe como $n$ ecuaciones de primer orden, las \textbf{ecuaciones de estado}, y una \textbf{ecuación de salida} algebraica:
\begin{equation}
  \dot x=f(x,u,t),\qquad y=h(x,u,t),
  \label{eq:ee-no-lineal}
\end{equation}
con $x\in\mathbb R^n$, $u\in\mathbb R^m$ y $y\in\mathbb R^p$. Si el sistema es lineal e invariante en el tiempo,
\begin{equation}
  \boxed{\;\dot x=\Am x+\Bm u,\qquad y=\Cm x+\Dm u\;}
  \label{eq:ee-lineal}
\end{equation}
donde $\Am$ ($n\times n$) es la matriz del sistema, $\Bm$ ($n\times m$) la de entrada, $\Cm$ ($p\times n$) la de salida y $\Dm$ ($p\times m$) la de transmisión directa. Este capítulo trabaja con \eqref{eq:ee-lineal}, salvo en la \cref{sec:ee-linealizacion}, donde se llega a ella desde \eqref{eq:ee-no-lineal}.

\subsection{Variables de fase}

Si la ecuación diferencial no tiene derivadas de la entrada,
\begin{equation}
  y^{(n)}+a_{n-1}y^{(n-1)}+\dots+a_1\dot y+a_0y=b_0u,
  \label{eq:ee-edo}
\end{equation}
la elección más directa es tomar la salida y sus $n-1$ derivadas: las \textbf{variables de fase},
\[
  x_1=y,\quad x_2=\dot y,\quad\dots,\quad x_n=y^{(n-1)} .
\]
Cada variable es la derivada de la anterior, y la última ecuación sale de \eqref{eq:ee-edo}:
\begin{equation}
  \dot x=\underbrace{\begin{bmatrix}0&1&0&\cdots&0\\0&0&1&\cdots&0\\\vdots&&&\ddots&\vdots\\0&0&0&\cdots&1\\-a_0&-a_1&-a_2&\cdots&-a_{n-1}\end{bmatrix}}_{\Am}x
  +\underbrace{\begin{bmatrix}0\\0\\\vdots\\0\\b_0\end{bmatrix}}_{\Bm}u,
  \qquad y=\underbrace{\begin{bmatrix}1&0&\cdots&0\end{bmatrix}}_{\Cm}x .
  \label{eq:ee-fase}
\end{equation}
Los coeficientes del polinomio característico aparecen, cambiados de signo, en la última fila de $\Am$.

\begin{ejemplo}[label=ej:ee-masa]{Masa, resorte y amortiguador}
Para el sistema de la \cref{fig:ee-masa}, con fricción viscosa $b\dot y$ y resorte lineal $ky$, la segunda ley de Newton da $M\ddot y+b\dot y+ky=u$. Escriba su modelo de estado.
\solucion Es de segundo orden: dos variables de estado. Con $x_1=y$ (posición) y $x_2=\dot y$ (velocidad),
\[
  \dot x_1=x_2,\qquad \dot x_2=-\frac kMx_1-\frac bMx_2+\frac1Mu,\qquad y=x_1,
\]
es decir,
\[
  \Am=\begin{bmatrix}0&1\\-\dfrac kM&-\dfrac bM\end{bmatrix},\quad
  \Bm=\begin{bmatrix}0\\\dfrac1M\end{bmatrix},\quad
  \Cm=\begin{bmatrix}1&0\end{bmatrix},\quad \Dm=0 .
\]
Con $M=2$, $k=16$ y un pulso de $5$\,N durante $20$\,s, la \cref{fig:ee-masa-sim} muestra el efecto de $b$: con $b=2$ los valores propios de $\Am$ son $-0.5\pm2.78j$ y la masa oscila; con $b=16$ son reales ($-1.17$ y $-6.83$) y no hay oscilación. En los cuatro casos, la posición se asienta en $5/16=0.3125$\,m mientras dura el pulso.
\end{ejemplo}

\begin{figure}[ht]
\centering
\includegraphics[width=0.55\linewidth]{cap04/masa_resorte_u}
\caption{Masa, resorte y amortiguador: entrada $u$ (fuerza), salida $y$ (posición).}
\label{fig:ee-masa}
\end{figure}

\begin{figure}[ht]
\centering
\includegraphics[width=0.95\linewidth]{cap04/f_masa_resorte}
\caption{Ejemplo~\ref{ej:ee-masa}: posición y velocidad para cuatro valores de $b$, con $M=2$ y $k=16$.}
\label{fig:ee-masa-sim}
\end{figure}

\codigorepo{cap04_espacio_estado/C4_Ej1_MasaResorte.m}{fase}{el modelo de estado de la masa-resorte}
\codigorepo{cap04_espacio_estado/C4_Ej1_MasaResorte.m}{simulacion}{la simulación con \cmd{lsim} para cuatro valores de $b$}

\paragraph{El computador analógico.} Una ventaja práctica del modelo de estado es que se arma con tres tipos de bloques: \textbf{integradores}, \textbf{sumadores} y \textbf{ganancias}. Así se simulaban los sistemas en los computadores analógicos, antes de los digitales: cada variable de estado es la salida de un integrador, y su condición inicial se fijaba con un potenciómetro. La \cref{fig:ee-analogico} muestra la masa-resorte armada de esa manera. Simulink conserva la misma lógica.

\begin{figure}[ht]
\centering
\includegraphics[width=0.85\linewidth]{cap04/computador_analogico}
\caption{La masa-resorte como en un computador analógico: dos integradores en cascada, cuyas salidas son las variables de estado.}
\label{fig:ee-analogico}
\end{figure}

\begin{modelo}[Modelo de Simulink: \texttt{crearModelo\_MasaResorte}]
\texttt{codigo/cap04\_espacio\_estado/crearModelo\_MasaResorte.m} arma la \cref{fig:ee-analogico} con bloques \emph{Integrator}, \emph{Sum} y \emph{Gain}, y un pulso de $5$\,N durante $20$\,s. Cambie \texttt{b} en el espacio de trabajo y vuelva a simular para reproducir la \cref{fig:ee-masa-sim}.
\end{modelo}

\begin{ejemplo}[label=ej:ee-dos-masas]{Dos masas}
En la \cref{fig:ee-dos-masas}, $M_1$ está unida a la pared por un amortiguador $b$ y a $M_2$ por un resorte $k$; la fuerza $u$ actúa sobre $M_2$. Las ecuaciones de Newton son
\[
  M_1\ddot y_1+b\dot y_1+k(y_1-y_2)=0,\qquad M_2\ddot y_2+k(y_2-y_1)=u .
\]
¿Cuántas variables de estado hacen falta? Escriba el modelo si se miden las dos posiciones.
\solucion Las posiciones $y_1$ y $y_2$ no bastan: hay dos ecuaciones de segundo orden y hacen falta cuatro condiciones iniciales. Con $x_1=y_1$, $x_2=y_2$, $x_3=\dot y_1$ y $x_4=\dot y_2$,
\[
  \Am=\begin{bmatrix}0&0&1&0\\0&0&0&1\\-\dfrac k{M_1}&\dfrac k{M_1}&-\dfrac b{M_1}&0\\\dfrac k{M_2}&-\dfrac k{M_2}&0&0\end{bmatrix},\quad
  \Bm=\begin{bmatrix}0\\0\\0\\\dfrac1{M_2}\end{bmatrix},\quad
  \Cm=\begin{bmatrix}1&0&0&0\\0&1&0&0\end{bmatrix},\quad\Dm=0 .
\]
Aquí $p=2$: dos salidas.
\end{ejemplo}

\begin{figure}[ht]
\centering
\includegraphics[width=0.7\linewidth]{cap04/dos_masas}
\caption{Ejemplo~\ref{ej:ee-dos-masas}: dos masas unidas por un resorte.}
\label{fig:ee-dos-masas}
\end{figure}

\begin{pausa}
Un circuito RLC en serie, alimentado por una fuente de voltaje, ¿cuántas variables de estado tiene? Proponga dos elecciones distintas.
\respuesta{Dos, porque hay dos elementos que almacenan energía de forma independiente. Por ejemplo, el voltaje del condensador y la corriente de la bobina; o la carga del condensador y su derivada (que es la corriente), que son variables de fase.}
\end{pausa}

%% ---------------------------------------------------------------------
\section{De la función de transferencia al estado}
\label{sec:ee-ft-estado}
%% ---------------------------------------------------------------------

Una función de transferencia sin ceros,
\[
  \frac{Y(s)}{U(s)}=\frac{b_0}{s^n+a_{n-1}s^{n-1}+\dots+a_1s+a_0},
\]
es la transformada de \eqref{eq:ee-edo} con condiciones iniciales nulas, y las variables de fase \eqref{eq:ee-fase} salen de inmediato.

\begin{ejemplo}[label=ej:ee-fase]{Una función de transferencia de tercer orden}
Escriba en variables de fase $\dfrac{Y(s)}{U(s)}=\dfrac{24}{s^3+9s^2+26s+24}$.
\solucion La ecuación diferencial es $\dddot y+9\ddot y+26\dot y+24y=24u$. Con $x_1=y$, $x_2=\dot y$ y $x_3=\ddot y$,
\[
  \Am=\begin{bmatrix}0&1&0\\0&0&1\\-24&-26&-9\end{bmatrix},\quad
  \Bm=\begin{bmatrix}0\\0\\24\end{bmatrix},\quad
  \Cm=\begin{bmatrix}1&0&0\end{bmatrix},\quad\Dm=0 .
\]
Los valores propios de $\Am$ son $-2$, $-3$ y $-4$: los polos de la función de transferencia.
\end{ejemplo}

\subsection{Cuando hay ceros}

Si el numerador tiene grado $m\ge1$, la ecuación diferencial tiene derivadas de la entrada y la salida y sus derivadas ya no sirven como estado. El truco es partir la función de transferencia en dos (\cref{fig:ee-intermedia}): primero el denominador, que produce una variable intermedia $Z(s)$ sin derivadas de la entrada,
\[
  Z(s)=\frac{1}{s^n+a_{n-1}s^{n-1}+\dots+a_0}\,U(s),
\]
a la que sí se le aplican las variables de fase ($x_1=z$, $x_2=\dot z$, \dots); y después el numerador, que solo combina esas variables:
\[
  Y(s)=(b_ms^m+\dots+b_1s+b_0)Z(s)\quad\Longrightarrow\quad y=b_0x_1+b_1x_2+\dots+b_mx_{m+1}.
\]
Resultado: la misma $\Am$ de \eqref{eq:ee-fase}, $\Bm=[0\;\cdots\;0\;\,1]^\top$ y $\Cm=[b_0\;\,b_1\;\cdots\;b_m\;\,0\;\cdots\;0]$. Los ceros quedan todos en $\Cm$. Esta es la forma que entrega la función del curso \cmd{variablesFase}, y la que se usa en el \cref{cap:re}. (Para $m=0$ coincide con \eqref{eq:ee-fase} salvo en dónde se pone $b_0$: en $\Bm$ o en $\Cm$.)

\begin{figure}[ht]
\centering
\includegraphics[width=0.9\linewidth]{cap04/variable_intermedia}
\caption{(a) Función de transferencia con ceros. (b) Partición con la variable intermedia $Z(s)$: las variables de fase se definen sobre $z$, y el numerador solo forma la salida.}
\label{fig:ee-intermedia}
\end{figure}

\begin{ejemplo}[label=ej:ee-ceros]{Con ceros}
Escriba en variables de fase $\dfrac{Y(s)}{U(s)}=\dfrac{s^2+7s+2}{s^3+9s^2+26s+24}$.
\solucion La misma $\Am$ del ejemplo~\ref{ej:ee-fase}, $\Bm=[0\;\,0\;\,1]^\top$ y $\Cm=[2\;\,7\;\,1]$: $y=2z+7\dot z+\ddot z$.
\end{ejemplo}

\codigorepo{cap04_espacio_estado/C4_Ej2_FTaEstado.m}{fase}{ejemplo~\ref{ej:ee-fase}}
\codigorepo{cap04_espacio_estado/C4_Ej2_FTaEstado.m}{numerador}{ejemplo~\ref{ej:ee-ceros} con la función \cmd{variablesFase}}

\paragraph{Lo que hace \matlab.} El comando \cmd{tf2ss} no devuelve las variables de fase. Para el ejemplo~\ref{ej:ee-ceros} entrega
\[
  \Am_1=\begin{bmatrix}-9&-26&-24\\1&0&0\\0&1&0\end{bmatrix},\quad
  \Bm_1=\begin{bmatrix}1\\0\\0\end{bmatrix},\quad
  \Cm_1=\begin{bmatrix}1&7&2\end{bmatrix} :
\]
las mismas variables, pero numeradas al revés ($x_1$ de \matlab{} es nuestro $x_3$). Invertir el orden de los estados con $\mathbb P=$\cmd{fliplr(eye(3))} recupera nuestra forma: $\mathbb P\Am_1\mathbb P$, $\mathbb P\Bm_1$, $\Cm_1\mathbb P$. Las dos representaciones dan la misma función de transferencia: \alerta{la representación de estado no es única, la función de transferencia sí}. La \cref{sec:ee-estado-ft} explica por qué.

\codigorepo{cap04_espacio_estado/C4_Ej2_FTaEstado.m}{matlab}{\cmd{tf2ss} y el cambio de orden de los estados}

\begin{pausa}
¿Qué pasa con $\Dm$ si el numerador tiene el mismo grado que el denominador, por ejemplo $\dfrac{2s+3}{s+1}$?
\respuesta{Primero se divide: $\dfrac{2s+3}{s+1}=2+\dfrac{1}{s+1}$. La parte constante va directo de la entrada a la salida: $\Dm=2$; el resto, $\Am=-1$, $\Bm=1$, $\Cm=1$.}
\end{pausa}

%% ---------------------------------------------------------------------
\section{Solución de la ecuación de estado}
\label{sec:ee-solucion}
%% ---------------------------------------------------------------------

\paragraph{Respuesta libre.} Empecemos sin entrada: $\dot x=\Am x$, con $x(0)$ conocido. Integrando a ambos lados, $x(t)=x(0)+\int_0^t\Am x(\tau)\,\dif\tau$; reemplazando $x(\tau)$ por esa misma expresión una y otra vez,
\[
  x(t)=\Bigl(\Id+\Am t+\frac{\Am^2t^2}{2!}+\frac{\Am^3t^3}{3!}+\cdots\Bigr)x(0).
\]
La serie es la de la exponencial, ahora con una matriz:
\begin{equation}
  x(t)=e^{\Am t}x(0)=\Phi(t)\,x(0) .
  \label{eq:ee-libre}
\end{equation}

\begin{definicion}[label=def:ee-transicion]{Matriz de transición de estado}
$\Phi(t)=e^{\Am t}=\displaystyle\sum_{k=0}^{\infty}\frac{\Am^kt^k}{k!}$ es la \textbf{matriz de transición de estado}: lleva el estado del instante $0$ al instante $t$. Cumple
\[
  \Phi(0)=\Id,\qquad \Phi^{-1}(t)=\Phi(-t),\qquad \Phi(t_1+t_2)=\Phi(t_1)\Phi(t_2),\qquad \Phi^k(t)=\Phi(kt).
\]
\end{definicion}

Calcular $e^{\Am t}$ no es trivial: Moler y Van Loan reunieron diecinueve maneras de hacerlo, todas con algún defecto numérico~\citep{MoLo78}. A mano, la más útil es la de Laplace (\cref{sec:ee-estado-ft}); en \matlab, \cmd{expm}.

\paragraph{Respuesta forzada.} Con entrada, $\dot x-\Am x=\Bm u$. Multiplicando por $e^{-\Am t}$, el lado izquierdo es $\deriv{}{t}\bigl(e^{-\Am t}x\bigr)$, y al integrar de $0$ a $t$:
\begin{equation}
  x(t)=\underbrace{e^{\Am t}x(0)}_{\text{respuesta libre}}+\underbrace{\int_0^te^{\Am(t-\tau)}\Bm\,u(\tau)\,\dif\tau}_{\text{respuesta forzada}},
  \qquad y(t)=\Cm x(t)+\Dm u(t).
  \label{eq:ee-forzada}
\end{equation}
Es la versión matricial de lo visto en el \cref{cap:po}: la respuesta total es la suma de la respuesta a las condiciones iniciales y la convolución de la entrada con la respuesta al impulso.

\begin{ejemplo}[label=ej:ee-transicion]{Matriz de transición de un segundo orden}
Con $\Am=\begin{bmatrix}0&1\\-2&-3\end{bmatrix}$ y $\Bm=[0\;\,1]^\top$: halle $\Phi(t)$, la respuesta al escalón unitario con $x(0)=0$ y la respuesta libre con $x(0)=[1\;\,{-1}]^\top$.
\solucion Con el procedimiento de la \cref{sec:ee-estado-ft},
\[
  \Phi(t)=\begin{bmatrix}2e^{-t}-e^{-2t} & e^{-t}-e^{-2t}\\-2e^{-t}+2e^{-2t} & -e^{-t}+2e^{-2t}\end{bmatrix},
  \qquad \Phi(0.5)=\begin{bmatrix}0.8452&0.2387\\-0.4773&0.1292\end{bmatrix}.
\]
Con \eqref{eq:ee-forzada} y $u=1$: $x_1(t)=\int_0^t\bigl(e^{-(t-\tau)}-e^{-2(t-\tau)}\bigr)\dif\tau=\frac12-e^{-t}+\frac12e^{-2t}$, y $x_2=\dot x_1=e^{-t}-e^{-2t}$. Con $x(0)=[1\;\,{-1}]^\top$ y sin entrada, $x(t)=e^{-t}[1\;\,{-1}]^\top$: un solo modo, porque esa condición inicial es un vector propio de $\Am$ (con valor propio $-1$).
\end{ejemplo}

\codigorepo{cap04_espacio_estado/C4_Ej3_Transicion.m}{simbolico}{$\Phi(t)$ con la resolvente y \cmd{ilaplace}}
\codigorepo{cap04_espacio_estado/C4_Ej3_Transicion.m}{numerico}{\cmd{expm} y las propiedades de $\Phi$}
\codigorepo{cap04_espacio_estado/C4_Ej3_Transicion.m}{forzada}{la respuesta forzada contra \cmd{lsim}}

\begin{pausa}
Si $x(0)=[1\;\,{-2}]^\top$ en el ejemplo~\ref{ej:ee-transicion}, ¿qué modo aparece?
\respuesta{Solo $e^{-2t}$: $[1\;\,{-2}]^\top$ es vector propio de $\Am$ con valor propio $-2$ ($\Am[1\;\,{-2}]^\top=[-2\;\,4]^\top$). Entonces $x(t)=e^{-2t}[1\;\,{-2}]^\top$.}
\end{pausa}

%% ---------------------------------------------------------------------
\section{Del estado a la función de transferencia}
\label{sec:ee-estado-ft}
%% ---------------------------------------------------------------------

Transformando $\dot x=\Am x+\Bm u$ con condiciones iniciales nulas, $sX(s)=\Am X(s)+\Bm U(s)$, de donde $X(s)=(s\Id-\Am)^{-1}\Bm U(s)$, y con la ecuación de salida,
\begin{equation}
  \boxed{\;G(s)=\frac{Y(s)}{U(s)}=\Cm(s\Id-\Am)^{-1}\Bm+\Dm\;}
  \label{eq:ee-ss2tf}
\end{equation}
siempre que $s\Id-\Am$ sea invertible. Como $(s\Id-\Am)^{-1}=\operatorname{adj}(s\Id-\Am)/\det(s\Id-\Am)$, cada entrada de $G(s)$ es un cociente de polinomios con denominador $\det(s\Id-\Am)$, el \textbf{polinomio característico} de $\Am$. Por eso los polos de $G(s)$ están entre los valores propios de $\Am$.

La matriz $(s\Id-\Am)^{-1}$ se llama \textbf{resolvente}, y es la transformada de la matriz de transición: transformando la serie de $e^{\Am t}$ término a término, $\Lap{e^{\Am t}}=\frac{\Id}{s}+\frac{\Am}{s^2}+\frac{\Am^2}{s^3}+\cdots$, y al multiplicar por $(s\Id-\Am)$ todo se cancela salvo $\Id$. Entonces
\begin{equation}
  \Phi(t)=e^{\Am t}=\iLap{(s\Id-\Am)^{-1}} .
  \label{eq:ee-resolvente}
\end{equation}

\begin{receta}[Matriz de transición por Laplace]
\begin{enumerate}
  \item Forme $s\Id-\Am$ y calcule su determinante, el polinomio característico.
  \item Inviértala: $(s\Id-\Am)^{-1}=\dfrac{\operatorname{adj}(s\Id-\Am)}{\det(s\Id-\Am)}$. Para $2\times2$, $\begin{bmatrix}a&b\\c&d\end{bmatrix}^{-1}=\dfrac{1}{ad-bc}\begin{bmatrix}d&-b\\-c&a\end{bmatrix}$.
  \item Antitransforme cada elemento, con fracciones parciales.
\end{enumerate}
\end{receta}

\paragraph{Adjunta y determinante.} La adjunta es la transpuesta de la matriz de cofactores: $\operatorname{adj}(\Am)=[C_{ij}]^\top$, con $C_{ij}=(-1)^{i+j}M_{ij}$ y $M_{ij}$ el determinante de la matriz que queda al quitar la fila $i$ y la columna $j$ (el \emph{menor}). El determinante se puede desarrollar por cualquier fila $i$ o columna $k$: $\det\Am=\sum_ja_{ij}C_{ij}=\sum_ia_{ik}C_{ik}$. Conviene escoger la fila o la columna con más ceros.

\begin{ejemplo}[label=ej:ee-det]{Desarrollo por cofactores}
Calcule el determinante de $\begin{bmatrix}2&4&1\\3&0&2\\2&0&3\end{bmatrix}$.
\solucion La segunda columna tiene dos ceros. Los menores son $M_{12}=\begin{vmatrix}3&2\\2&3\end{vmatrix}=5$, $M_{22}=\begin{vmatrix}2&1\\2&3\end{vmatrix}=4$ y $M_{32}=\begin{vmatrix}2&1\\3&2\end{vmatrix}=1$, y los cofactores, $C_{12}=-5$, $C_{22}=4$ y $C_{32}=-1$. Entonces $\det=4\cdot C_{12}+0\cdot C_{22}+0\cdot C_{32}=-20$.
\end{ejemplo}

\begin{ejemplo}[label=ej:ee-motor]{Motor de corriente continua}
Un motor con inercia se modela como $\dot x_1=x_2$, $\dot x_2=-\alpha x_2+\beta u$, $y=x_1$, con $\alpha,\beta>0$~\citep{Fri86}. Halle $\Phi(t)$ y $G(s)$.
\solucion $s\Id-\Am=\begin{bmatrix}s&-1\\0&s+\alpha\end{bmatrix}$, $\det=s(s+\alpha)$ y
\[
  (s\Id-\Am)^{-1}=\begin{bmatrix}\dfrac1s&\dfrac{1}{s(s+\alpha)}\\[2mm]0&\dfrac{1}{s+\alpha}\end{bmatrix},
  \qquad
  \Phi(t)=\begin{bmatrix}1&\dfrac{1-e^{-\alpha t}}{\alpha}\\0&e^{-\alpha t}\end{bmatrix}.
\]
Con \eqref{eq:ee-ss2tf}, $G(s)=\dfrac{\beta}{s(s+\alpha)}$: el valor propio en cero es el integrador que convierte velocidad en posición.
\end{ejemplo}

\subsection{¿Valores propios o polos?}

Los polos de $G(s)$ son valores propios de $\Am$, pero \alerta{no siempre todos los valores propios son polos}: puede haber cancelaciones entre el numerador y $\det(s\Id-\Am)$.

\begin{ejemplo}[label=ej:ee-cancelacion]{Una función de transferencia que esconde modos}
Con $\Am=\begin{bmatrix}0&1&0\\0&0&1\\-6&-11&-6\end{bmatrix}$ y $\Bm=[0\;\,0\;\,1]^\top$, compare $\Cm_a=[1\;\,0\;\,0]$ y $\Cm_b=[2\;\,3\;\,1]$.
\solucion Los valores propios son $-1$, $-2$ y $-3$. Con $\Cm_a$, $G_a=\dfrac{1}{(s+1)(s+2)(s+3)}$. Con $\Cm_b$, el numerador es $s^2+3s+2=(s+1)(s+2)$ y
\[
  G_b(s)=\frac{(s+1)(s+2)}{(s+1)(s+2)(s+3)}=\frac{1}{s+3}:
\]
un solo polo para tres valores propios. Los modos $e^{-t}$ y $e^{-2t}$ siguen en el estado (aparecen si $x(0)\neq0$), pero no llegan a la salida. Aquí son estables y la cancelación es inofensiva; si uno de ellos fuera inestable, la función de transferencia lo escondería. Es la pregunta que responde la \textbf{observabilidad}, en el \cref{cap:re}.
\end{ejemplo}

\codigorepo{cap04_espacio_estado/C4_Ej4_EstadoaFT.m}{ss2tf}{\cmd{ss2tf}, \cmd{eig} y \cmd{dcgain}}
\codigorepo{cap04_espacio_estado/C4_Ej4_EstadoaFT.m}{cancelacion}{ejemplo~\ref{ej:ee-cancelacion}}

\paragraph{Cambio de variables.} Si $x=\mathbb T\bar x$ con $\mathbb T$ invertible, el sistema en las nuevas variables es $\bar\Am=\mathbb T^{-1}\Am\mathbb T$, $\bar\Bm=\mathbb T^{-1}\Bm$, $\bar\Cm=\Cm\mathbb T$, $\bar\Dm=\Dm$, y
\[
  \bar\Cm(s\Id-\bar\Am)^{-1}\bar\Bm=\Cm\mathbb T\bigl(\mathbb T^{-1}(s\Id-\Am)\mathbb T\bigr)^{-1}\mathbb T^{-1}\Bm=\Cm(s\Id-\Am)^{-1}\Bm .
\]
Por eso hay infinitas representaciones de estado con la misma función de transferencia, y los valores propios no cambian. Si las columnas de $\mathbb T$ son vectores propios de $\Am$ (con valores propios distintos), $\bar\Am$ es diagonal, cada estado nuevo es un \textbf{modo} desacoplado, y $\Phi(t)=\mathbb T\,\operatorname{diag}(e^{\lambda_it})\,\mathbb T^{-1}$.

\codigorepo{cap04_espacio_estado/C4_Ej4_EstadoaFT.m}{similitud}{la forma diagonal y $\Phi(t)$}

\begin{ideaclave}
\textbf{Lo que es único y lo que no.}\enspace La función de transferencia y los valores propios de $\Am$ no dependen de cómo se escojan las variables de estado; las matrices sí. Y la función de transferencia puede tener menos polos que valores propios: describe solo la parte del sistema que conecta la entrada con la salida.
\end{ideaclave}

%% ---------------------------------------------------------------------
\section{Linealización en espacio de estado}
\label{sec:ee-linealizacion}
%% ---------------------------------------------------------------------

En la \cref{sec:so-linealizacion} linealizamos ecuaciones de una variable. La representación de estado permite hacerlo con todas a la vez. Sea $\dot x=f(x,u)$, $y=h(x,u)$, y sea $(\eq x,\eq u)$ un equilibrio: $f(\eq x,\eq u)=0$. Con las variables desviación $\dev x=x-\eq x$, $\dev u=u-\eq u$ y $\dev y=y-h(\eq x,\eq u)$, la serie de Taylor de primer orden de cada componente da
\begin{equation}
  \dot{\dev x}=\Am\dev x+\Bm\dev u,\qquad \dev y=\Cm\dev x+\Dm\dev u,
  \label{eq:ee-lineal-desv}
\end{equation}
donde cada matriz es un \textbf{jacobiano} evaluado en el equilibrio:
\begin{equation}
  a_{ij}=\left.\frac{\partial f_i}{\partial x_j}\right|_{\eq x,\eq u},\quad
  b_{ij}=\left.\frac{\partial f_i}{\partial u_j}\right|_{\eq x,\eq u},\quad
  c_{ij}=\left.\frac{\partial h_i}{\partial x_j}\right|_{\eq x,\eq u},\quad
  d_{ij}=\left.\frac{\partial h_i}{\partial u_j}\right|_{\eq x,\eq u}.
  \label{eq:ee-jacobianos}
\end{equation}
Cada equilibrio tiene su propio modelo lineal, y en general dinámicas distintas.

\begin{ejemplo}[label=ej:ee-duffing]{La ecuación de Duffing}
Para $\ddot x+\delta\dot x+x^3-x=0$ con $\delta>0$: escriba el modelo de estado, halle los equilibrios y linealice alrededor de cada uno.
\solucion Con $x_1=x$, $x_2=\dot x$: $\dot x_1=x_2$, $\dot x_2=x_1-x_1^3-\delta x_2$. Los equilibrios cumplen $x_2=0$ y $x_1(1-x_1^2)=0$: son $(0,0)$ y $(\pm1,0)$. El jacobiano es
\[
  J(x_1,x_2)=\begin{bmatrix}0&1\\1-3x_1^2&-\delta\end{bmatrix} .
\]
En el origen, $\lambda^2+\delta\lambda-1=0$ y $\lambda=\frac{-\delta\pm\sqrt{\delta^2+4}}{2}$: un valor propio positivo y uno negativo para \emph{todo} $\delta>0$ (su producto es $-1$). Es un \textbf{punto de silla}, inestable. En $(\pm1,0)$, $\lambda^2+\delta\lambda+2=0$: $\omega_n=\sqrt2$ y $\zeta=\delta/(2\sqrt2)$, estables, subamortiguados si $\delta<\sqrt8=2.83$. Con $\delta=0.3$: $0.861$ y $-1.161$ en el origen, y $-0.15\pm1.41j$ en $(\pm1,0)$. La \cref{fig:ee-duffing} muestra el retrato de fase: cada condición inicial termina en uno de los dos pozos, y la frontera entre ellos (en negro) pasa por el origen en la dirección del vector propio estable de la silla. Ningún modelo lineal único reproduce eso.
\end{ejemplo}

\begin{figure}[ht]
\centering
\includegraphics[width=0.6\linewidth]{cap04/f_duffing}
\caption{Ejemplo~\ref{ej:ee-duffing} con $\delta=0.3$: retrato de fase, la frontera entre los dos dominios de atracción (negro) y tres trayectorias.}
\label{fig:ee-duffing}
\end{figure}

\begin{ejemplo}[label=ej:ee-levitador]{Levitador magnético}
Un electroimán sostiene una esfera de acero de masa $M$ (\cref{fig:ee-levitador}). La fuerza magnética es $Ci^2/y$, con $y$ la distancia de la esfera al imán~\citep{barie1996maglev}, y la bobina tiene resistencia $R$ e inductancia $L$:
\[
  M\ddot y=Mg-C\frac{i^2}{y},\qquad e=Ri+L\frac{\dif i}{\dif t}.
\]
Linealice alrededor de la posición $y=y_d$ con $M=0.05$\,kg, $R=10\,\Omega$, $L=0.5$\,H, $C=10^{-3}$\,N\,m/A$^2$, $g=9.81$\,m/s$^2$ y $y_d=0.01$\,m.
\solucion Con $x_1=y$, $x_2=\dot y$, $x_3=i$ y $u=e$:
\[
  \dot x_1=x_2,\qquad \dot x_2=g-\frac CM\frac{x_3^2}{x_1},\qquad \dot x_3=-\frac RLx_3+\frac1Lu .
\]
En el equilibrio, $\eq x_1=y_d$, $\eq x_2=0$, $\eq x_3=\sqrt{Mgy_d/C}=2.215$\,A y $\eq u=R\eq x_3=22.15$\,V. Los jacobianos son
\[
  \Am=\begin{bmatrix}0&1&0\\\dfrac CM\dfrac{(\eq x_3)^2}{(\eq x_1)^2}&0&-\dfrac{2C\eq x_3}{M\eq x_1}\\0&0&-\dfrac RL\end{bmatrix}
  =\begin{bmatrix}0&1&0\\981&0&-8.86\\0&0&-20\end{bmatrix},\qquad
  \Bm=\begin{bmatrix}0\\0\\2\end{bmatrix}.
\]
Note que $a_{21}=g/y_d$: en el equilibrio la fuerza magnética iguala al peso, y si la esfera se acerca, la fuerza crece más rápido que el peso. Los valores propios son $\pm31.3$ y $-20$: \alerta{el equilibrio es inestable}, y la función de transferencia $\dfrac{\dev Y(s)}{\dev E(s)}=\dfrac{-17.72}{(s-31.3)(s+31.3)(s+20)}$ lo confirma. La \cref{fig:ee-levitador-sim} compara el modelo lineal con el no lineal tras un desplazamiento inicial de $0.1$\,mm: coinciden durante unos $0.1$\,s (error del $4$\,\%) y luego se separan, cuando la esfera ya se alejó varios milímetros.
\end{ejemplo}

\begin{figure}[ht]
\centering
\includegraphics[width=0.6\linewidth]{cap04/levitador}
\caption{Levitador magnético: circuito de la bobina y fuerzas sobre la esfera.}
\label{fig:ee-levitador}
\end{figure}

\begin{figure}[ht]
\centering
\includegraphics[width=0.75\linewidth]{cap04/f_levitador}
\caption{Ejemplo~\ref{ej:ee-levitador}: la esfera parte $0.1$\,mm por debajo del equilibrio, con $u=\eq u$.}
\label{fig:ee-levitador-sim}
\end{figure}

\codigorepo{cap04_espacio_estado/C4_Ej6_Linealizacion.m}{duffing}{ejemplo~\ref{ej:ee-duffing}}
\codigorepo{cap04_espacio_estado/C4_Ej6_Linealizacion.m}{levitador}{ejemplo~\ref{ej:ee-levitador}}

\begin{ejemplo}[label=ej:ee-bioreactor]{Un bioreactor}
En un bioreactor, la biomasa $x_1$ (células) consume el sustrato $x_2$ (azúcar, desechos) y se multiplica. Un modelo simplificado es~\citep{Beq03}
\[
  \dot x_1=(\mu-u_1)x_1,\qquad \dot x_2=u_1(u_2-x_2)-\frac{\mu x_1}{Y},\qquad
  \mu(x_2)=\frac{\mu_{\max}x_2}{k_m+x_2+k_1x_2^2},
\]
donde $u_1$ es la tasa de dilución (caudal/volumen, la variable manipulada), $u_2$ la concentración de sustrato en la alimentación (una perturbación), $Y$ el rendimiento y $\mu$ la tasa de crecimiento con \emph{inhibición por sustrato} (con $k_1=0$ sería el modelo de Monod). Con $\mu_{\max}=0.53$\,h$^{-1}$, $k_m=0.12$\,g/L, $k_1=0.4545$\,L/g, $Y=0.4$, $\eq u_1=0.3$\,h$^{-1}$ y $\eq u_2=4$\,g/L, halle los equilibrios y linealice.
\solucion De $\dot x_1=0$: o $x_1=0$, o $\mu(x_2)=u_1$. El primer caso da el \emph{lavado}, $(0,\,4)$: las células se fueron del reactor. El segundo da $0.1362x_2^2-0.23x_2+0.036=0$, con $x_2=0.1746$ y $x_2=1.5122$, y de $\dot x_2=0$, $x_1=Y(u_2-x_2)$: los equilibrios $(1.5302,\,0.1746)$ y $(0.9951,\,1.5122)$. Con $\mu=u_1$ en el equilibrio, los jacobianos son
\[
  \Am=\begin{bmatrix}0&\eq x_1\mu'\\-\dfrac{u_1}{Y}&-u_1-\dfrac{\eq x_1\mu'}{Y}\end{bmatrix},\qquad
  \Bm=\begin{bmatrix}-\eq x_1&0\\u_2-\eq x_2&u_1\end{bmatrix},
\]
con $\mu'=\dif\mu/\dif x_2$. En $(0.9951,\,1.5122)$, $\mu'=-0.068$ y
\[
  \Am=\begin{bmatrix}0&-0.0679\\-0.75&-0.1302\end{bmatrix},\qquad
  \Bm=\begin{bmatrix}-0.9951&0\\2.4878&0.3\end{bmatrix},
\]
con valores propios $0.170$ y $-0.3$: \alerta{inestable}. En $(1.5302,\,0.1746)$, $\mu'=+0.591$ y los valores propios son $-0.3$ y $-2.26$: estable. La diferencia está en el signo de $\mu'$, que se anula en $x_2=\sqrt{k_m/k_1}=0.514$: a la derecha de ese máximo, más sustrato \emph{frena} el crecimiento (\cref{fig:ee-bioreactor}). Con las dos salidas $y=x$, las funciones de transferencia en el punto de operación son
\begin{align*}
  \frac{\dev X_1}{\dev U_1}&=\frac{-0.9951s-0.2985}{s^2+0.1302s-0.0510}, &
  \frac{\dev X_1}{\dev U_2}&=\frac{-0.0204}{s^2+0.1302s-0.0510},\\
  \frac{\dev X_2}{\dev U_1}&=\frac{2.488s+0.7463}{s^2+0.1302s-0.0510}, &
  \frac{\dev X_2}{\dev U_2}&=\frac{0.3\,s}{s^2+0.1302s-0.0510}.
\end{align*}
\end{ejemplo}

\begin{figure}[ht]
\centering
\includegraphics[width=\linewidth]{cap04/f_bioreactor}
\caption{Ejemplo~\ref{ej:ee-bioreactor}. Izquierda: la tasa de crecimiento corta a $u_1=0.3$ en dos puntos. Derecha: desde $(1,1)$ y desde cerca del equilibrio inestable, las trayectorias van al equilibrio estable o al lavado $(0,4)$.}
\label{fig:ee-bioreactor}
\end{figure}

\codigorepo{cap04_espacio_estado/C4_Ej7_Bioreactor.m}{equilibrios}{los tres equilibrios}
\codigorepo{cap04_espacio_estado/C4_Ej7_Bioreactor.m}{jacobiano}{jacobianos, valores propios y funciones de transferencia}

\begin{modelo}[Modelo de Simulink: \texttt{crearModelo\_Bioreactor}]
\texttt{codigo/cap04\_espacio\_estado/crearModelo\_Bioreactor.m} arma el modelo no lineal con dos entradas y los dos estados como salidas. Con él, \cmd{trim} busca el equilibrio para las entradas fijas y \cmd{linmod} entrega $\Am$, $\Bm$, $\Cm$ y $\Dm$ del modelo linealizado (las instrucciones están en el script \texttt{C4\_Ej7\_Bioreactor}).
\end{modelo}

\begin{cuidado}[El modelo lineal vale cerca del equilibrio]
Un modelo linealizado describe las desviaciones \emph{pequeñas}. Qué tan pequeñas, depende del sistema: en el levitador, unas décimas de milímetro; en el bioreactor, la linealización alrededor de un equilibrio no dice nada sobre el lavado. Y un equilibrio inestable en lazo abierto, como el del levitador o el del bioreactor, es justamente el que pide control.
\end{cuidado}

%% ---------------------------------------------------------------------
\section{Para llevar}
\label{sec:ee-cierre}
%% ---------------------------------------------------------------------

\begin{parallevar}
\begin{enumerate}
  \item El estado es el conjunto mínimo de variables que, con las entradas futuras, determina todo el comportamiento futuro. Su número (el orden) es único; las variables no.
  \item Modelo lineal: $\dot x=\Am x+\Bm u$, $y=\Cm x+\Dm u$. En variables de fase, la última fila de $\Am$ lleva los coeficientes del polinomio característico con signo cambiado; si hay ceros, van en $\Cm$ (variable intermedia $Z$). \cmd{tf2ss} numera los estados al revés.
  \item Solución: respuesta libre más respuesta forzada, $x(t)=e^{\Am t}x(0)+\int_0^te^{\Am(t-\tau)}\Bm u(\tau)\,\dif\tau$. La matriz de transición es $\Phi(t)=e^{\Am t}=\iLap{(s\Id-\Am)^{-1}}$.
  \item Función de transferencia: $G(s)=\Cm(s\Id-\Am)^{-1}\Bm+\Dm$. Sus polos son valores propios de $\Am$, pero una cancelación puede esconder valores propios.
  \item Un cambio de variables $x=\mathbb T\bar x$ cambia las matrices pero no la función de transferencia ni los valores propios.
  \item Linealizar es evaluar jacobianos en un equilibrio. Cada equilibrio tiene su modelo, y el modelo solo vale cerca de él.
\end{enumerate}
\end{parallevar}

\begin{profundiza}[Correcciones respecto a la versión anterior del capítulo]
(i) En el levitador, la ecuación de estado había perdido la constante $C$: es $\dot x_2=g-\frac CM\frac{x_3^2}{x_1}$, el equilibrio es $\eq x_3=\sqrt{Mgy_d/C}$ y $a_{21}=\frac CM\frac{(\eq x_3)^2}{(\eq x_1)^2}=g/y_d$ (en la matriz faltaba además el cuadrado de $\eq x_3$). (ii) En el ejemplo del determinante, el segundo cofactor es $C_{22}=4$ (estaba rotulado $C_{12}$). (iii) En el bioreactor, la cuarta función de transferencia repetía la tercera (\texttt{numU2(1,:)} en lugar de \texttt{numU2(2,:)}); los números del \texttt{trim} correspondían a $k_1=0.454$ y no a $0.4545$, y se recalcularon todos con el jacobiano analítico. (iv) El truco con \cmd{canon(\ldots,'companion')} y trasposiciones para recuperar las variables de fase dependía de la versión de \matlab; se reemplazó por la inversión del orden de los estados. (v) Historia: Lefschetz no creó el RIAS, sino que dirigió su centro de matemáticas; Pontryagin aportó el principio del máximo (no el ``cálculo de variaciones''); la conexión entre función de transferencia y estado que mostró Kalman (1963) es la de qué partes del sistema ve cada descripción, no que el espacio de estado cubra todo sistema no lineal. (vi) Erratas: $k_{04}\to k_{O4}$, $x_1=y\to x_1=y_1$ en el ejemplo de las dos masas, y una nota tomada de Wikipedia se reemplazó por la definición. (vii) Se agregaron los valores numéricos de todos los ejemplos, las figuras y el retrato de fase de Duffing; el resumen del capítulo anterior que cerraba el handout se retiró (está en el \cref{cap:so}). (viii) Notación: $\bar x\to\eq x$ para el equilibrio y $\dev x$ para la desviación, como en el \cref{cap:so}.
\end{profundiza}

%% ---------------------------------------------------------------------
\section{Problemas}
\label{sec:ee-problemas}
%% ---------------------------------------------------------------------

Los problemas marcados con \etiquetaM{} requieren \matlab/\simulink; los marcados con \etiquetaEstrella{} son de profundización.

\begin{problema}[variables de fase]\label{prob:ee-fase}
Escriba en variables de fase, con $\Am$, $\Bm$, $\Cm$ y $\Dm$: (a) $\dfrac{10}{s^3+6s^2+11s+6}$; (b) $2\ddot y+6\dot y+8y=4u$; (c) $\dfrac{2s^2+7s+3}{s^3+5s^2+8s+4}$, con la variable intermedia $Z(s)$. (d) \etiquetaM{} Compare con \cmd{tf2ss} y recupere su forma invirtiendo el orden de los estados. (e) \etiquetaM{} Aplique \cmd{ss2tf} a las dos representaciones de (c). ¿Qué propiedad ilustra el resultado?
\end{problema}

\begin{problema}[resolvente y polos]\label{prob:ee-resolvente}
Para (i) $\Am=\begin{bmatrix}0&1\\-6&-5\end{bmatrix}$ y (ii) $\Am=\begin{bmatrix}-2&1&0\\0&-3&1\\0&0&-4\end{bmatrix}$, con $\Bm=[0\;\cdots\;0\;\,1]^\top$, $\Cm=[1\;\,0\;\cdots]$ y $\Dm=0$: (a) calcule $(s\Id-\Am)^{-1}$ con la adjunta; (b) halle $G(s)$, sus polos y ceros; (c) compare con los valores propios de $\Am$; (d) halle la ganancia DC de (i) y el valor final ante escalón unitario; (e) \etiquetaM{} verifique con \cmd{ss2tf}, \cmd{eig} y \cmd{dcgain}.
\end{problema}

\begin{problema}[matriz de transición]\label{prob:ee-transicion}
(a) Calcule $\Phi(t)$ para $\Am=\begin{bmatrix}0&1\\-2&-3\end{bmatrix}$ con los tres pasos de la receta. (b) \etiquetaM{} Verifique numéricamente, con $t_1=0.5$\,s y $t_2=1.2$\,s, que $\Phi(0)=\Id$, $\Phi^{-1}(t_1)=\Phi(-t_1)$ y $\Phi(t_1+t_2)=\Phi(t_1)\Phi(t_2)$. (c) Deduzca la $\Phi(t)$ del motor del ejemplo~\ref{ej:ee-motor}. (d) Con $\Bm=[0\;\,1]^\top$ y $x(0)=0$, halle $x(t)$ ante escalón unitario con \eqref{eq:ee-forzada} y verifique en Laplace. (e) \etiquetaM{} Con $x(0)=[1\;\,{-1}]^\top$ y entrada nula, compare con \cmd{expm} en varios instantes. ¿Por qué aparece un solo modo?
\end{problema}

\begin{problema}[levitador]\label{prob:ee-levitador}
Con el modelo y los datos del ejemplo~\ref{ej:ee-levitador}: (a) escriba el modelo de estado no lineal conservando $C$; (b) halle $(\eq x,\eq u)$; (c) obtenga $\Am$ y $\Bm$ y muestre que $a_{21}=g/y_d$; explique por qué; (d) halle los valores propios; ¿qué concluye? (e) obtenga $\dev Y(s)/\dev E(s)$ y compare sus polos con (d); (f) \etiquetaM{} simule el modelo no lineal desde el equilibrio desplazado $0.1$\,mm y compare con el lineal; ¿hasta cuándo es aceptable la predicción?
\end{problema}

\begin{problema}[ciclo del carbono]\label{prob:ee-carbono}
Con el modelo y los coeficientes del ejemplo~\ref{ej:ee-carbono}: (a) construya $\Am$ y $\Bm$; (b) demuestre que $\mathbf 1^\top\Am=0$ y que $\frac{\dif}{\dif t}(\mathbf 1^\top x)=u_E$; interprete; (c) deduzca, sin calcular, que $\lambda=0$ es valor propio, y \etiquetaM{} halle los cinco y las constantes de tiempo; (d) \etiquetaM{} simule 500 años con $u_E=8$ GtC/año desde $x(0)=[590\;\,550\;\,1500\;\,900\;\,37000]^\top$; ¿hay estado estacionario? (e) halle $M_A(s)/U_E(s)$ con $\Cm=[1\;\,0\;\,0\;\,0\;\,0]$ e identifique el polo responsable del crecimiento.
\end{problema}

\begin{problema}[bioreactor]\label{prob:ee-bioreactor}
Con el modelo del ejemplo~\ref{ej:ee-bioreactor}: (a) derive $\mu'(x_2)$, muestre que se anula en $x_2=\sqrt{k_m/k_1}$ y ubique ese valor respecto a los equilibrios; (b) calcule a mano las ocho derivadas de $\Am$ y $\Bm$ en $(0.9951,\,1.5122)$; (c) halle los valores propios; ¿cuál modo es inestable? (d) \etiquetaM{} obtenga las cuatro funciones de transferencia; (e) repita (b) y (c) en $(1.5302,\,0.1746)$ y explique la diferencia con el signo de $\mu'$; (f) \etiquetaM{} simule desde $x(0)=(1,1)$; ¿a qué equilibrio llega?
\end{problema}

\begin{problema}[valores propios y polos]\label{prob:ee-cancelacion}
Con $\Am$, $\Bm$, $\Cm_a$ y $\Cm_b$ del ejemplo~\ref{ej:ee-cancelacion}: (a) halle los valores propios y las dos funciones de transferencia; (b) explique el resultado con $\Cm_b$ y reformule con precisión la frase ``los valores propios de $\Am$ son los polos''; (c) \etiquetaM{} simule los dos sistemas con $x(0)=[1\;\,0\;\,0]^\top$ y entrada nula, y muestre que hay modos en el estado que no llegan a la salida; (d) con $\mathbb T$ formada por los vectores propios, obtenga la forma diagonal y verifique que $G$ no cambia; (e) demuestre que todo cambio de variables preserva $G$; (f) \etiquetaM{} verifique $\Phi(t)=\mathbb T\operatorname{diag}(e^{\lambda_it})\mathbb T^{-1}$ contra \cmd{expm}.
\end{problema}

\begin{problema}[Duffing]\label{prob:ee-duffing}
Con $\ddot x+\delta\dot x+x^3-x=0$, $\delta>0$: (a) muestre que el origen es una silla para todo $\delta>0$; (b) halle el $\delta$ a partir del cual $(\pm1,0)$ deja de ser subamortiguado, y clasifique esos equilibrios para $\delta=0.3$ y $\delta=3$; (c) para $\delta=0.3$, halle $\Phi(t)$ del modelo lineal en $(1,0)$ y su tiempo de establecimiento; (d) \etiquetaM{} simule desde $(0.9,0)$, $(0.5,0)$ y $(-0.1,\,0.05)$ y compare con el modelo lineal del equilibrio correspondiente; (e) \etiquetaEstrella{} construya el retrato de fase en $[-2,2]^2$ y trace la frontera entre los dominios de atracción; relaciónela con los vectores propios del origen; (f) ¿por qué ningún modelo lineal reproduce lo de (e)?
\end{problema}

%% ---------------------------------------------------------------------
\section{Soluciones seleccionadas}
%% ---------------------------------------------------------------------

\paragraph{Problema~\ref{prob:ee-fase}}
(a) $\Am=\begin{bmatrix}0&1&0\\0&0&1\\-6&-11&-6\end{bmatrix}$, $\Bm=[0\;\,0\;\,10]^\top$, $\Cm=[1\;\,0\;\,0]$. (b) Dividiendo por $2$: $\ddot y+3\dot y+4y=2u$, $\Am=\begin{bmatrix}0&1\\-4&-3\end{bmatrix}$, $\Bm=[0\;\,2]^\top$. (c) $\Am=\begin{bmatrix}0&1&0\\0&0&1\\-4&-8&-5\end{bmatrix}$, $\Bm=[0\;\,0\;\,1]^\top$, $\Cm=[3\;\,7\;\,2]$. (e) La misma $G(s)$: la representación no es única, la función de transferencia sí.

\paragraph{Problema~\ref{prob:ee-resolvente}}
(i) $(s\Id-\Am)^{-1}=\dfrac{1}{s^2+5s+6}\begin{bmatrix}s+5&1\\-6&s\end{bmatrix}$, $G=\dfrac{1}{(s+2)(s+3)}$, sin ceros, ganancia DC $1/6$ y $y_\infty=1/6$. (ii) $G=\dfrac{1}{(s+2)(s+3)(s+4)}$ (los valores propios de una triangular son su diagonal), ganancia DC $1/24$.

\paragraph{Problema~\ref{prob:ee-transicion}}
(a) La del ejemplo~\ref{ej:ee-transicion}. (c) La del ejemplo~\ref{ej:ee-motor}. (d) $x_1=\frac12-e^{-t}+\frac12e^{-2t}$, $x_2=e^{-t}-e^{-2t}$; en Laplace, $X_1=\frac{1}{s(s+1)(s+2)}$. (e) $x(t)=e^{-t}[1\;\,{-1}]^\top$: la condición inicial es un vector propio.

\paragraph{Problema~\ref{prob:ee-levitador}}
(b) $\eq x=[0.01\;\,0\;\,2.215]^\top$, $\eq u=22.15$\,V. (c) $a_{21}=981$\,s$^{-2}$, $a_{23}=-8.86$, $a_{33}=-20$, $\Bm=[0\;\,0\;\,2]^\top$. (d) $\pm31.3$ y $-20$: inestable en lazo abierto. (e) $\dfrac{-17.72}{s^3+20s^2-981s-19620}$. (f) A $0.1$\,s el error del modelo lineal es de $4$\,\%; a $0.15$\,s, de $16$\,\%.

\paragraph{Problema~\ref{prob:ee-carbono}}
(c) $\lambda=0$, $-0.0006$, $-0.0355$, $-0.0673$ y $-0.2274$ por año; constantes de tiempo de $1729$, $28$, $15$ y $4.4$ años. (d) El total crece $8$ GtC/año; no hay estado estacionario. La atmósfera llega a $1833$ GtC en 500 años. (e) El polo en $s=0$.

\paragraph{Problema~\ref{prob:ee-bioreactor}}
(a) $\mu'=\mu_{\max}\dfrac{k_m-k_1x_2^2}{(k_m+x_2+k_1x_2^2)^2}$, nulo en $x_2=0.514$, entre los dos equilibrios. (c) $0.170$ y $-0.3$. (e) $\Am=\begin{bmatrix}0&0.905\\-0.75&-2.562\end{bmatrix}$, valores propios $-0.3$ y $-2.26$. (f) Llega a $(1.5302,\,0.1746)$.

\paragraph{Problema~\ref{prob:ee-cancelacion}}
(a) $-1$, $-2$, $-3$; $G_a=\frac{1}{(s+1)(s+2)(s+3)}$ y $G_b=\frac{1}{s+3}$. (b) Los polos de $G$ son valores propios de $\Am$; los valores propios que se cancelan no aparecen como polos.

\paragraph{Problema~\ref{prob:ee-duffing}}
(a) El producto de los valores propios es $\det J(0,0)=-1<0$. (b) $\delta=\sqrt8$; con $\delta=0.3$: $\zeta=0.106$, $\omega_n=1.41$, focos estables; con $\delta=3$: $\zeta=1.06$, valores propios $-1$ y $-2$, nodos estables. (c) $\ts\approx4/(\zeta\omega_n)=26.7$\,s. (d) $(0.9,0)$ y $(0.5,0)$ van a $(1,0)$; $(-0.1,\,0.05)$ va a $(-1,0)$: su componente sobre el vector propio inestable del origen es negativa.

\respuestaspausas

                 % migrado (Handout3_Ctrl.tex)
%% =====================================================================
%%  Capítulo 5 -- Estabilidad
%%  Migrado al formato nuevo (noviembre de 2026) desde Handout4_Ctrl.tex,
%%  con material de las Actividades 1 (Routh y Kharitonov) y de la clase 2
%%  de Bode (rango de Kc por Routh y por Bode).
%%  Cambios de fondo respecto a la versión anterior (ver MIGRACION.md):
%%    la sección de tiempo muerto (Padé) se conserva, con la fórmula de c_k
%%    corregida y la comparación con el límite exacto; "estable si todos los polos Y CEROS están
%%    en el semiplano izquierdo" -> solo los polos; el ejemplo
%%    s^3+2s^2+s+2 es un caso de FILA de ceros (no de epsilon); se
%%    resolvieron los ejemplos que no tenían solución; en el primer ejemplo
%%    de Kharitonov faltaba el intervalo de a4 y la familia NO es
%%    robustamente estable; el ejemplo de diseño robusto se resolvió y se
%%    muestra el conservadurismo; C(s)/R(s) -> Y(s)/R(s); Kharitanov ->
%%    Kharitonov; anécdota del sirviente (es de Kelvin, no de Maxwell).
%% =====================================================================

\capitulo[cap:es]{Estabilidad}
  {En la cocina, la educación y el amor la rapidez va en contra del placer y la calidad. Necesitamos cocinar, amar y leer a fuego lento.}
  {Julián de Zubiría Samper}

\begin{objetivos}[partes/05-estabilidad.html]
  \item Distinguir estabilidad asintótica, marginal (en el sentido de Lyapunov) y BIBO, y decidirla a partir de los polos o de los valores propios.
  \item Aplicar el criterio de Routh-Hurwitz, incluidos sus dos casos especiales, para contar las raíces con parte real positiva sin calcularlas.
  \item Usar la tabla de Routh para diseñar: rango de ganancia, ganancia y frecuencia críticas, regiones de estabilidad de un PI y estabilidad relativa.
  \item Analizar la estabilidad de un lazo con tiempo muerto mediante la aproximación de Padé, y reconocer cuánto se equivoca.
  \item Decidir la estabilidad de un equilibrio en el espacio de estado, también cuando el modelo viene de linealizar, y reconocer los modos que la función de transferencia esconde.
  \item Verificar la estabilidad robusta de una familia de polinomios de intervalo con el teorema de Kharitonov, y reconocer cuándo el resultado es conservador.
\end{objetivos}

%% ---------------------------------------------------------------------
\section{¿Qué es estabilidad?}
\label{sec:es-concepto}
%% ---------------------------------------------------------------------

Hasta aquí hemos analizado la respuesta transitoria: qué tan rápido responde un sistema y cuánto se pasa. Antes de preguntar \emph{cómo} llega, hay que preguntar \emph{si} llega. Esa es la pregunta de la estabilidad, y es la primera exigencia de cualquier diseño: un controlador que desestabiliza el lazo no sirve, por buenos que sean los demás números.

El diccionario de la Real Academia da varias acepciones de \emph{estable}; la que nos sirve es ``que mantiene o recupera el equilibrio''. La palabra clave es \clave{equilibrio}: la estabilidad no es una propiedad del sistema sino de un \alerta{punto de equilibrio} del sistema. Un péndulo tiene dos: colgado hacia abajo, donde una perturbación pequeña se amortigua, y parado hacia arriba, donde la menor perturbación lo tumba. Usted mismo, de pie, está en un equilibrio inestable que su sistema nervioso estabiliza todo el tiempo.

En un sistema lineal e invariante no forzado hay un solo equilibrio, el origen ($x=0$); con entrada constante, el equilibrio se desplaza, pero en variables desviación sigue siendo el origen. Como hay uno solo, se abusa del lenguaje y se dice ``el sistema es estable''. Lo haremos también, pero recuerde que siempre se habla del equilibrio.\nota{En sistemas no lineales el abuso no se puede cometer: un mismo sistema tiene equilibrios estables e inestables, como el péndulo (\cref{sec:es-espacio}).}

\begin{definicion}[label=def:es-estabilidad]{Tres nociones de estabilidad}
Para un sistema lineal e invariante, con su equilibrio en el origen:
\begin{itemize}
  \item El equilibrio es \clave{asintóticamente estable} si la respuesta natural (sin entrada, desde cualquier condición inicial) tiende a cero.
  \item Es \clave{marginalmente estable} (estable en el sentido de Lyapunov) si la respuesta natural se mantiene acotada, pero no necesariamente tiende a cero.
  \item Es \clave{inestable} si alguna condición inicial produce una respuesta natural que crece sin límite.
\end{itemize}
El sistema es \clave{BIBO estable} (\emph{bounded input, bounded output}) si toda entrada acotada produce una salida acotada.
\end{definicion}

\subsection*{Los polos lo deciden todo}

La respuesta natural es una suma de \clave{modos} $e^{p_it}$, uno por cada polo $p_i$ (capítulo~2). Un polo real negativo da un modo que decae; un par $\sigma\pm j\omega$ da $e^{\sigma t}\sin(\omega t+\phi)$, que decae si $\sigma<0$, oscila con amplitud constante si $\sigma=0$ y crece si $\sigma>0$. Los polos repetidos agregan factores $t,t^2,\dots$: sobre el eje imaginario, un polo doble da $t\sin(\omega t)$, que crece aunque no haya polos en el semiplano derecho (\cref{fig:es-modos}).

\begin{figure}[ht]
\centering
\includegraphics[width=\linewidth]{cap05/f_modos}
\caption{Ubicación de los polos y respuesta natural (al impulso). Solo la primera tiende a cero. La última muestra que los polos repetidos sobre el eje imaginario también desestabilizan.}
\label{fig:es-modos}
\end{figure}

\begin{ideaclave}
\textbf{Estabilidad y ubicación de los polos.}\enspace Sea $\Tlc(s)=N(s)/D(s)$ sin cancelaciones. El equilibrio es
\begin{itemize}
  \item \textbf{asintóticamente estable} si y solo si todos los polos tienen parte real negativa (semiplano izquierdo abierto). Solo en este caso el sistema es además BIBO estable;
  \item \textbf{marginalmente estable} si no hay polos en el semiplano derecho y los que están sobre el eje imaginario son simples;
  \item \textbf{inestable} si hay algún polo con parte real positiva o algún polo repetido sobre el eje imaginario.
\end{itemize}
Los \alerta{ceros no intervienen}: modifican la forma de la respuesta, no su estabilidad.
\end{ideaclave}

\begin{cuidado}[Marginal no es BIBO]
Un sistema marginalmente estable \emph{no} es BIBO estable. El integrador $1/s$ ante un escalón da una rampa, y $1/(s^2+1)$ ante $\sin t$ da $\tfrac12(\sin t-t\cos t)$, que crece como $t/2$: la entrada excita justo el modo que no se amortigua. En diseño, ``estable'' quiere decir siempre \emph{asintóticamente} estable.
\end{cuidado}

\codigorepo{cap05_estabilidad/C5_Ej1_Modos.m}{modos}{\cref{fig:es-modos}: la respuesta natural según la ubicación de los polos}

\subsection*{¿Cómo se decide?}

Para saber si un equilibrio es estable hay varios caminos, y el curso los recorre casi todos:
\begin{itemize}
  \item \textbf{Calcular los polos} o los valores propios (\cmd{roots}, \cmd{pole}, \cmd{eig}). Con computador es lo más directo, pero no dice \emph{para qué valores de un parámetro} el lazo es estable.
  \item \textbf{El criterio de Routh-Hurwitz} (\cref{sec:es-routh}): cuenta las raíces inestables a partir de los coeficientes, sin calcularlas. Como trabaja con los coeficientes, admite parámetros libres: es una herramienta de diseño.
  \item \textbf{El lugar de las raíces} de Evans~\citep{Eva50}, que dibuja cómo se mueven los polos del lazo cerrado cuando cambia una ganancia.
  \item \textbf{Los márgenes de fase y de ganancia} sobre el diagrama de Bode y el criterio de Nyquist~\citep{Nyq32,Bod40}, en el capítulo~10.
  \item \textbf{El método de Lyapunov}~\citep{Lya92}, el más general, que sirve también para sistemas no lineales y variantes en el tiempo; en este curso usaremos su método indirecto, la linealización (\cref{sec:es-espacio}).
\end{itemize}

\begin{pausa}
Clasifique: (a) $\dfrac{s-3}{(s+1)(s+4)}$; (b) $\dfrac{1}{s(s+2)}$; (c) $\dfrac{s+1}{(s^2+4)^2}$; (d) $\dfrac{1}{s^2-s+1}$.
\respuesta{(a) Asintóticamente estable: el cero en $+3$ no importa. (b) Marginalmente estable: polo simple en el origen (un escalón da una rampa: no es BIBO). (c) Inestable: polos dobles en $\pm 2j$. (d) Inestable: polos en $0.5\pm j0.866$.}
\end{pausa}

%% ---------------------------------------------------------------------
\section{Un poco de historia: de los anillos de Saturno a las turbinas de Davos}
\label{sec:es-historia}
%% ---------------------------------------------------------------------

En el capítulo~1 vimos el recorrido general; aquí nos detenemos en la historia de la pregunta que resuelve este capítulo.

\begin{figure}[ht]
\centering
\includegraphics[width=0.9\linewidth]{cap05/personajes}
\caption{De izquierda a derecha: James Clerk Maxwell (1831--1879), Edward John Routh (1831--1907), Aurel Stodola (1859--1942) y Adolf Hurwitz (1859--1919).}
\label{fig:es-personajes}
\end{figure}

\paragraph{Maxwell deja la pregunta abierta.} El primer análisis matemático de la estabilidad de un regulador fue el de George Biddell Airy, astrónomo real, que en 1840 estudió por qué el regulador de fricción de un telescopio oscilaba~\citep{Air40}. Pero la pregunta general la planteó James Clerk Maxwell en \emph{On governors} (1868)~\citep{Max68,Kan16}. Maxwell ya había usado la estabilidad como criterio de verdad: en su ensayo sobre los anillos de Saturno, premiado en 1857, descartó los modelos que no eran estables, porque unos anillos que existen desde hace tanto tiempo tienen que serlo~\citep{Bit15}. En \emph{On governors} linealizó las ecuaciones del regulador y de la máquina y mostró que el equilibrio es estable si todas las raíces de la ecuación característica tienen parte real negativa. Obtuvo las condiciones para segundo y tercer orden ---la de tercer orden es la que veremos en el ejemplo~\ref{ej:es-tercer}--- y confesó que no había podido hacerlo para grados mayores, con la esperanza de que el problema ``atrajera la atención de los matemáticos''.

\paragraph{Routh y el premio Adams.} La Universidad de Cambridge entregaba cada dos años el premio Adams, creado en honor de John Couch Adams, quien había predicho la existencia de Neptuno a partir de las irregularidades de la órbita de Urano. Para el periodo 1875--1877 los examinadores ---entre ellos Maxwell y George Stokes--- escogieron como tema \emph{El criterio de estabilidad dinámica}. Lo ganó en 1877 Edward John Routh, con el tratado que contiene el arreglo que hoy lleva su nombre~\citep{Rou77}. Routh y Maxwell habían nacido el mismo año, entraron juntos a Cambridge y presentaron juntos los exámenes de matemáticas de 1854, en los que Routh quedó primero y Maxwell segundo. Routh fue además el tutor más famoso de Cambridge: preparó a más de seiscientos estudiantes, entre ellos Rayleigh y J.~J.\ Thomson. Estaba casado con una hija de Airy.

\paragraph{Stodola, Hurwitz y las turbinas.} El trabajo de Routh tardó en salir de Inglaterra. En el continente el texto de referencia sobre reguladores era el de Ivan Vyshnegradskii (1876), que había resuelto el caso de tercer orden con un método gráfico~\citep{Vys76}. Sobre esa base trabajaba Aurel Stodola, ingeniero eslovaco que llegó en 1892 como profesor al Politécnico de Zúrich (hoy ETH). En 1893 le encargaron el regulador de las turbinas hidráulicas de una planta en Davos; su modelo era de orden mayor que tres, el criterio de Vyshnegradskii no alcanzaba, y Stodola le preguntó a su colega matemático Adolf Hurwitz cuándo un polinomio tiene todas sus raíces con parte real negativa~\citep{Bis89}. Hurwitz respondió con un criterio en términos de determinantes, publicado en 1895~\citep{Hur95}, sin conocer el de Routh. Stodola publicó su análisis de las turbinas en la revista suiza de ingeniería (1893--94) y presentó los resultados en variables relativas, sin dimensiones, lo que los hizo útiles para otros ingenieros.

\paragraph{Dos criterios, uno solo.} En 1911 el italiano Enrico Bompiani demostró que los criterios de Routh y de Hurwitz son equivalentes~\citep{Bom11}; ese mismo año Luciano Orlando publicó otra demostración y una fórmula que relaciona los determinantes de Hurwitz~\citep{Orl11}. Por eso hoy se habla del criterio de \clave{Routh-Hurwitz}: el arreglo de Routh es la forma cómoda de calcular, y los determinantes de Hurwitz son la forma cómoda de demostrar.

\begin{historia}[Lyapunov y Kharitonov]
En 1892 Aleksandr Lyapunov, entonces profesor en Járkov, presentó su tesis \emph{El problema general de la estabilidad del movimiento}~\citep{Lya92}. Definió con precisión qué es un equilibrio estable y dio dos métodos: el \emph{indirecto}, que decide la estabilidad local linealizando (\cref{sec:es-espacio}), y el \emph{directo}, que usa funciones parecidas a la energía y sirve para sistemas no lineales. Su obra se conoció en Occidente a mediados del siglo~XX. Casi un siglo después, en 1978, Vladimir Kharitonov publicó en una revista soviética de ecuaciones diferenciales un teorema sorprendente: para saber si \emph{toda} una familia de polinomios con coeficientes inciertos es estable, basta revisar cuatro polinomios~\citep{Kha78}. El resultado pasó inadvertido fuera de la Unión Soviética hasta 1983--84, cuando Ross Barmish y otros lo dieron a conocer en las revistas de control occidentales~\citep{Bar84} (\cref{sec:es-kharitonov}).
\end{historia}

%% ---------------------------------------------------------------------
\section{El criterio de Routh-Hurwitz}
\label{sec:es-routh}
%% ---------------------------------------------------------------------

Para el lazo de la \cref{fig:es-lazo}, la función de transferencia en lazo cerrado es
\[
  \frac{Y(s)}{R(s)}=\frac{\Gc(s)\Gp(s)}{1+\Gc(s)\Gp(s)\Hs(s)}
  =\frac{b_0s^m+b_1s^{m-1}+\cdots+b_m}{a_0s^n+a_1s^{n-1}+\cdots+a_{n-1}s+a_n},\qquad m\le n .
\]
El denominador igualado a cero es la \clave{ecuación característica},
\begin{equation}
  1+\Gma(s)=0\quad\Longleftrightarrow\quad a_0s^n+a_1s^{n-1}+\cdots+a_{n-1}s+a_n=0,
  \label{eq:es-caracteristica}
\end{equation}
con $\Gma=\Gc\Gp\Hs$ la malla abierta. Sus raíces son los polos del lazo cerrado: \alerta{la estabilidad del lazo se decide en el denominador de $1+\Gma$}, que es numerador más denominador de $\Gma$ cuando esta es una fracción sin cancelaciones.

\begin{figure}[ht]
\centering
\includegraphics[width=0.62\linewidth]{cap10/lazo_basico}
\caption{Lazo básico del curso.}
\label{fig:es-lazo}
\end{figure}

\subsection*{Una condición necesaria: los signos}

Si todas las raíces tienen parte real negativa, el polinomio es un producto de factores $(s+p)$ y $(s^2+2\zeta\wn s+\wn^2)$ con coeficientes positivos, y al multiplicarlos no aparece ningún signo menos ni ningún término faltante. Entonces:

\begin{ideaclave}
\textbf{Condición necesaria.}\enspace Si a la ecuación característica le \alerta{falta un término} o tiene \alerta{coeficientes de signos distintos}, hay raíces con parte real positiva o nula: el equilibrio no es asintóticamente estable. Para $n=1$ y $n=2$ la condición también es suficiente; para $n\ge3$ \emph{no}: $s^3+s^2+2s+8$ tiene todos sus coeficientes positivos y raíces en $0.5\pm j1.94$.
\end{ideaclave}

\subsection*{La tabla de Routh}

Con todos los coeficientes del mismo signo (tomamos $a_0>0$), se arma la tabla. Las dos primeras filas se llenan alternando los coeficientes; cada fila siguiente se calcula con las dos de arriba:
\[
\begin{array}{c|cccc}
  s^n     & a_0 & a_2 & a_4 & \cdots\\
  s^{n-1} & a_1 & a_3 & a_5 & \cdots\\
  s^{n-2} & b_1 & b_2 & b_3 & \cdots\\
  s^{n-3} & c_1 & c_2 & \cdots\\
  \vdots  & \vdots\\
  s^0     & \ast
\end{array}
\qquad
\begin{aligned}
  b_1&=\frac{a_1a_2-a_0a_3}{a_1}, & b_2&=\frac{a_1a_4-a_0a_5}{a_1}, & \dots\\[3pt]
  c_1&=\frac{b_1a_3-a_1b_2}{b_1}, & c_2&=\frac{b_1a_5-a_1b_3}{b_1}, & \dots
\end{aligned}
\]
Todas las entradas siguen el mismo patrón: ``el de abajo a la izquierda por el de arriba a la derecha, menos el de arriba a la izquierda por el de abajo a la derecha, dividido por el de abajo a la izquierda''. Donde falte un elemento se pone cero. La tabla tiene $n+1$ filas, de $s^n$ a $s^0$, y la última fila es siempre $a_n$.

\begin{definicion}[label=def:es-routh]{Criterio de Routh-Hurwitz}
El número de raíces de la ecuación característica con \alerta{parte real positiva} es igual al número de \clave{cambios de signo} en la primera columna de la tabla de Routh. El equilibrio es asintóticamente estable si y solo si todos los elementos de la primera columna tienen el mismo signo~\citep{Rou77,Hur95}.
\end{definicion}

\begin{cuidado}[Multiplicar una fila]
Una fila completa se puede multiplicar o dividir por un número \emph{positivo} sin cambiar los signos de la primera columna; sirve para evitar fracciones. Nunca por un número negativo, ni por una expresión cuyo signo no conoce (como $K$).
\end{cuidado}

\begin{ejemplo}[label=ej:es-tercer]{Tercer orden: la condición de Maxwell}
Encuentre la condición de estabilidad de $a_0s^3+a_1s^2+a_2s+a_3=0$, con $a_0>0$.
\solucion
\[
\begin{array}{c|cc}
  s^3 & a_0 & a_2\\
  s^2 & a_1 & a_3\\
  s^1 & \dfrac{a_1a_2-a_0a_3}{a_1} & 0\\
  s^0 & a_3
\end{array}
\]
Sin cambios de signo: $a_1>0$, $a_3>0$ y
\begin{equation}
  a_1a_2>a_0a_3 .
  \label{eq:es-tercer}
\end{equation}
Es la condición que encontró Maxwell en 1868. Vale la pena memorizarla: \alerta{``el producto de los del medio mayor que el producto de los extremos''}.
\end{ejemplo}

\begin{ejemplo}[label=ej:es-cuarto]{Cuarto orden}
¿Cuántas raíces inestables tiene $s^4+2s^3+3s^2+4s+5=0$?
\solucion Todos los coeficientes son positivos, así que hay que hacer la tabla:
\[
\begin{array}{c|ccc}
  s^4 & 1 & 3 & 5\\
  s^3 & 2 & 4 & 0\\
  s^2 & \frac{2\cdot3-1\cdot4}{2}=1 & \frac{2\cdot5-1\cdot0}{2}=5\\
  s^1 & \frac{1\cdot4-2\cdot5}{1}=-6 & 0\\
  s^0 & 5
\end{array}
\]
La primera columna es $1,\,2,\,1,\,-6,\,5$: dos cambios de signo ($1\to-6$ y $-6\to5$), luego \alerta{dos raíces con parte real positiva}. En efecto, \cmd{roots([1 2 3 4 5])} da $-1.288\pm j0.858$ y $0.288\pm j1.416$.
\end{ejemplo}

\codigorepo{cap05_estabilidad/C5_Ej2_Routh.m}{routh}{ejemplo~\ref{ej:es-cuarto} y el contraejemplo de los signos con la función \cmd{tablaRouth}}

\begin{receta}[Criterio de Routh-Hurwitz]
\begin{enumerate}
  \item Escribir la ecuación característica $1+\Gma(s)=0$ como polinomio, ordenado en potencias decrecientes.
  \item Revisar los signos: si falta un término o hay signos distintos, ya se sabe que no es asintóticamente estable (la tabla dirá cuántas raíces son inestables).
  \item Armar la tabla. Si aparece un cero en la primera columna o una fila de ceros, ir a la \cref{sec:es-casos}.
  \item Contar los cambios de signo de la primera columna.
  \item Verificar con \cmd{roots}, \cmd{pole} o \cmd{tablaRouth}.
\end{enumerate}
\end{receta}

\begin{pausa}
Use \eqref{eq:es-tercer} para decidir, sin tabla, si $s^3+2s^2+3s+4$ y $2s^3+s^2+3s+4$ son estables.
\respuesta{El primero sí: $2\cdot3=6>1\cdot4$. El segundo no: $1\cdot3=3<2\cdot4=8$; la tabla tiene dos cambios de signo y dos raíces inestables.}
\end{pausa}

%% ---------------------------------------------------------------------
\section{Los dos casos especiales}
\label{sec:es-casos}
%% ---------------------------------------------------------------------

La regla para construir la tabla divide por el primer elemento de la fila anterior. Si ese elemento es cero, hay que hacer algo. Hay dos situaciones distintas, y conviene no confundirlas.

\subsection*{Caso 1: un cero en la primera columna, con el resto de la fila distinto de cero}

Se reemplaza el cero por un número pequeño y positivo $\epsilon$, se sigue construyendo la tabla y se cuentan los cambios de signo haciendo $\epsilon\to0^+$.

\begin{ejemplo}[label=ej:es-caso1]{Un cero que se vuelve $\epsilon$}
Analice $s^5+2s^4+4s^3+8s^2+10s+6=0$.
\solucion
\[
\begin{array}{c|ccc}
  s^5 & 1 & 4 & 10\\
  s^4 & 2 & 8 & 6\\
  s^3 & 0\;\to\;\epsilon & 7 & 0\\
  s^2 & 8-\dfrac{14}{\epsilon} & 6\\[6pt]
  s^1 & \approx 7 & \\
  s^0 & 6
\end{array}
\]
La fila $s^3$ vale $\bigl[\frac{2\cdot4-1\cdot8}{2},\,\frac{2\cdot10-1\cdot6}{2}\bigr]=[0,\,7]$: solo el primer elemento es cero. Con $\epsilon\to0^+$, el elemento de $s^2$ es muy negativo y el de $s^1$, $\dfrac{(8-14/\epsilon)\,7-6\epsilon}{8-14/\epsilon}$, tiende a $7$. La primera columna queda $+,+,+,-,+,+$: \alerta{dos cambios de signo}, dos raíces inestables ($0.484\pm j1.738$).
\end{ejemplo}

\subsection*{Caso 2: una fila completa de ceros}

Una fila entera se anula cuando el polinomio tiene un factor \clave{par} (solo potencias pares de $s$), cuyas raíces son simétricas respecto al origen (\cref{fig:es-simetricas}): reales y opuestas (A), imaginarias puras (B) o en cuadrante, $\pm\sigma\pm j\omega$ (C). Ese factor es el \clave{polinomio auxiliar} $P(s)$, que se lee en la fila \emph{anterior} a la de ceros. Se reemplaza la fila de ceros por los coeficientes de $\dif P/\dif s$ y se sigue.

\begin{figure}[ht]
\centering
\includegraphics[width=0.36\linewidth]{cap05/raices_simetricas}
\caption{Raíces simétricas respecto al origen que producen una fila de ceros: reales (A), imaginarias (B) o en cuadrante (C).}
\label{fig:es-simetricas}
\end{figure}

\begin{ideaclave}
\textbf{Cómo leer el caso 2.}\enspace \begin{itemize}
  \item Las raíces del polinomio auxiliar son raíces del polinomio original, y se pueden calcular: $P(s)$ tiene grado par y solo potencias pares.
  \item Los cambios de signo \emph{desde la fila del auxiliar hacia abajo} cuentan las raíces del auxiliar con parte real positiva; como son simétricas, hay el mismo número con parte real negativa, y las que faltan están sobre el eje imaginario.
  \item Si no hay cambios de signo y el auxiliar tiene raíces imaginarias simples, el equilibrio es \alerta{marginalmente estable}.
\end{itemize}
\end{ideaclave}

\begin{ejemplo}[label=ej:es-caso2]{Raíces sobre el eje imaginario}
Analice $s^5+7s^4+6s^3+42s^2+8s+56=0$.
\solucion
\[
\begin{array}{c|ccc}
  s^5 & 1 & 6 & 8\\
  s^4 & 7 & 42 & 56\\
  s^3 & 0\;\to\;28 & 0\;\to\;84 & \\
  s^2 & 21 & 56\\
  s^1 & 28/3\\
  s^0 & 56
\end{array}
\]
La fila $s^3$ es toda cero: $\frac{7\cdot6-1\cdot42}{7}=0$ y $\frac{7\cdot8-1\cdot56}{7}=0$. El auxiliar es la fila $s^4$: $P(s)=7s^4+42s^2+56=7(s^2+2)(s^2+4)$, y su derivada $28s^3+84s$ da la nueva fila. No hay cambios de signo, así que no hay raíces en el semiplano derecho; las cuatro raíces del auxiliar están sobre el eje: $\pm j\sqrt2$ y $\pm 2j$. La quinta es $-7$. El equilibrio es \alerta{marginalmente estable}: la respuesta natural oscila con dos frecuencias sin amortiguarse.
\end{ejemplo}

\begin{ejemplo}[label=ej:es-caso2b]{Un caso 2 que parece caso 1}
Analice $s^3+2s^2+s+2=0$.
\solucion La fila $s^1$ vale $\frac{2\cdot1-1\cdot2}{2}=0$, y como es su único elemento, es una \emph{fila completa} de ceros: caso 2, no caso 1. El auxiliar es $2s^2+2$, con derivada $4s$:
\[
\begin{array}{c|cc}
  s^3 & 1 & 1\\
  s^2 & 2 & 2\\
  s^1 & 0\;\to\;4\\
  s^0 & 2
\end{array}
\]
Sin cambios de signo y auxiliar $2(s^2+1)$: raíces $\pm j$ y $-2$. Marginalmente estable. (Reemplazar el cero por $\epsilon$ da los mismos signos, pero esconde la información importante: dónde están las raíces del eje y a qué frecuencia oscilará el sistema.)
\end{ejemplo}

\codigorepo{cap05_estabilidad/C5_Ej3_CasosEspeciales.m}{caso2}{ejemplo~\ref{ej:es-caso2}: el polinomio auxiliar y sus raíces}

\begin{pausa}
Analice $s^4+4=0$. Antes de hacer la tabla, ¿qué puede decir solo con mirar los coeficientes?
\respuesta{Faltan los términos en $s^3$, $s^2$ y $s$: no es asintóticamente estable. Toda la fila $s^3$ es cero; el auxiliar es $s^4+4$, con derivada $4s^3$. La fila $s^2$ queda $[0,\,4]$ (caso 1: $\epsilon$) y la fila $s^1$, $-16/\epsilon$. Dos cambios de signo: dos raíces en el semiplano derecho. Son $\pm1\pm j$, en cuadrante (tipo C).}
\end{pausa}

%% ---------------------------------------------------------------------
\section{Routh-Hurwitz para diseñar}
\label{sec:es-diseno}
%% ---------------------------------------------------------------------

La gran ventaja de Routh-Hurwitz sobre \cmd{roots} es que trabaja con los coeficientes, y los coeficientes pueden contener parámetros del controlador. La tabla responde entonces la pregunta de diseño más básica: \emph{¿para qué valores de la ganancia el lazo es estable?}

\subsection{Rango de ganancia, ganancia crítica y frecuencia de oscilación}

\begin{ejemplo}[label=ej:es-rangokc]{Rango de $\Kc$ de un lazo de proceso}
En el lazo de la \cref{fig:es-tres}, un controlador proporcional $\Kc$ actúa sobre una válvula, un proceso y un sensor de primer orden. Halle el rango de $\Kc$ para el cual el lazo es estable, la ganancia crítica y la frecuencia a la que oscila.
\solucion \textbf{1. Ecuación característica.}
\[
  1+\Kc\,\frac{0.016}{3s+1}\cdot\frac{50}{30s+1}\cdot\frac{1}{10s+1}=0
  \;\Longrightarrow\;
  900s^3+420s^2+43s+(1+0.8\Kc)=0 .
\]
\textbf{2. Tabla.}
\[
\begin{array}{c|cc}
  s^3 & 900 & 43\\
  s^2 & 420 & 1+0.8\Kc\\
  s^1 & \dfrac{17\,160-720\Kc}{420} & \\[6pt]
  s^0 & 1+0.8\Kc
\end{array}
\]
\textbf{3. Condiciones.} Fila $s^0$: $1+0.8\Kc>0$, o sea $\Kc>-1.25$. Fila $s^1$: $17\,160-720\Kc>0$, o sea $\Kc<23.83$. Entonces
\[
  \alerta{-1.25<\Kc<23.83}.
\]
\textbf{4. Ganancia y frecuencia críticas.} En $\Kc=\Kcr=23.83$ la fila $s^1$ se anula: es una fila de ceros, y el auxiliar es la fila $s^2$,
\[
  420s^2+(1+0.8\cdot23.83)=420s^2+20.07=0\;\Longrightarrow\;s=\pm j\,0.2186 .
\]
El lazo oscila sin amortiguarse con $\wcr=0.2186$\,rad/s, es decir, con periodo $\Pcr=2\pi/\wcr=28.7$\,s (\cref{fig:es-rangokc}).
\end{ejemplo}

\begin{figure}[ht]
\centering
\includegraphics[width=0.92\linewidth]{cap05/lazo_tres}
\caption{Lazo de proceso del ejemplo~\ref{ej:es-rangokc}.}
\label{fig:es-tres}
\end{figure}

\begin{figure}[ht]
\centering
\includegraphics[width=0.85\linewidth]{cap05/f_rango_kc}
\caption{Ejemplo~\ref{ej:es-rangokc}: por debajo de $\Kcr$ la oscilación se amortigua, en $\Kcr$ se sostiene con periodo $\Pcr=28.7$\,s y por encima crece.}
\label{fig:es-rangokc}
\end{figure}

\begin{ideaclave}
\textbf{Ganancia crítica: tres caras del mismo número.}\enspace $\Kcr$ y $\wcr$ salen de la fila que se anula y del polinomio auxiliar. Son los mismos números que (i) el método de oscilación sostenida de Ziegler y Nichols busca experimentalmente en la planta (capítulo~8)~\citep{ZiNi42}, y (ii) el margen de ganancia y la frecuencia de cruce de fase que se leen en el diagrama de Bode (capítulo~10): para este lazo, \cmd{margin} da $\GM=23.83$ y $\wgm=0.2186$\,rad/s.
\end{ideaclave}

\codigorepo{cap05_estabilidad/C5_Ej4_RangoK.m}{rangokc}{ejemplo~\ref{ej:es-rangokc}: rango de $\Kc$, $\Kcr$, $\wcr$ y la comparación con \cmd{margin}}

\begin{cuidado}[Las ganancias negativas]
El rango incluye $-1.25<\Kc<0$. Con $\Kc$ negativo el controlador empuja en la dirección equivocada; el lazo resulta estable mientras el término independiente siga positivo, pero el error en estado estacionario es peor que en lazo abierto. Que un valor sea estable no lo hace útil: en la práctica se usa $0<\Kc<\Kcr$, y bastante por debajo de $\Kcr$.
\end{cuidado}

\subsection{Estabilizar una planta inestable}

Cuando la planta es inestable, la ganancia ya no tiene una cota superior sola: también aparece una cota \emph{inferior}, porque con poca ganancia el controlador no alcanza a contener el polo inestable.

\begin{ejemplo}[label=ej:es-inestable]{Una ventana estrecha}
La planta $\Gp(s)=\dfrac{1}{(s-1)(s^2+2s+2)}$ se controla con el compensador $\Gc(s)=\dfrac{K(s+2)}{s+3}$, con retroalimentación unitaria (\cref{fig:es-inestable-lazo}). Halle el rango de $K$ que estabiliza el lazo.
\solucion La ecuación característica es $(s+3)(s-1)(s^2+2s+2)+K(s+2)=0$:
\[
  s^4+4s^3+3s^2+(K-2)\,s+(2K-6)=0 .
\]
\[
\begin{array}{c|ccc}
  s^4 & 1 & 3 & 2K-6\\
  s^3 & 4 & K-2 & \\
  s^2 & \dfrac{14-K}{4} & 2K-6\\[6pt]
  s^1 & \dfrac{68-16K-K^2}{14-K} & \\[6pt]
  s^0 & 2K-6
\end{array}
\]
Las condiciones son $K<14$ (fila $s^2$), $2K-6>0$ (fila $s^0$) y $68-16K-K^2>0$ (fila $s^1$, con $14-K>0$). Las raíces de $K^2+16K-68$ son $-8\pm2\sqrt{33}$, es decir $3.489$ y $-19.49$. Juntando:
\[
  \alerta{3<K<3.489}.
\]
Es un rango muy estrecho (\cref{fig:es-inestable}): un error de pocos puntos porcentuales en la ganancia, o en el modelo, desestabiliza el lazo. Este compensador estabiliza, pero no es robusto; dentro del rango, la mejor ganancia es $K\approx3.07$, con los polos más lentos en $-0.28$.
\end{ejemplo}

\begin{figure}[ht]
\centering
\includegraphics[width=0.78\linewidth]{cap05/lazo_inestable}
\caption{Ejemplo~\ref{ej:es-inestable}: compensador para una planta con un polo en $+1$.}
\label{fig:es-inestable-lazo}
\end{figure}

\begin{figure}[ht]
\centering
\includegraphics[width=0.85\linewidth]{cap05/f_inestable}
\caption{Ejemplo~\ref{ej:es-inestable}: parte real del polo más a la derecha en función de $K$. Por debajo de $3$ queda un polo real inestable; por encima de $3.489$ cruzan el eje dos polos complejos.}
\label{fig:es-inestable}
\end{figure}

\subsection{Dos parámetros: la región de estabilidad de un PI}

Con dos parámetros, las condiciones de la tabla dibujan una \emph{región} en el plano de los parámetros.

\begin{ejemplo}[label=ej:es-pi]{Región de estabilidad de un PI}
En el lazo de la \cref{fig:es-pi-lazo}, $\Gc(s)=\Kc+K_i/s$ y $\Gp(s)=\dfrac{2}{s^3+4s^2+5s+2}=\dfrac{2}{(s+1)^2(s+2)}$. Halle la región del plano $(\Kc,K_i)$ en la que el lazo es estable.
\solucion La ecuación característica es $s(s^3+4s^2+5s+2)+2(\Kc s+K_i)=0$:
\[
  s^4+4s^3+5s^2+(2+2\Kc)\,s+2K_i=0 .
\]
\[
\begin{array}{c|ccc}
  s^4 & 1 & 5 & 2K_i\\
  s^3 & 4 & 2+2\Kc & \\
  s^2 & \dfrac{9-\Kc}{2} & 2K_i\\[6pt]
  s^1 & c_1 & \\
  s^0 & 2K_i
\end{array}
\qquad
c_1=(2+2\Kc)-\frac{16K_i}{9-\Kc}.
\]
Las condiciones son $K_i>0$, $\Kc<9$ y $c_1>0$:
\[
  \alerta{0<K_i<\frac{(1+\Kc)(9-\Kc)}{8},\qquad -1<\Kc<9}.
\]
La frontera es una parábola con máximo $K_i=25/8$ en $\Kc=4$ (\cref{fig:es-region}). Sobre ella, la fila $s^1$ se anula y el auxiliar $\frac{9-\Kc}{2}s^2+2K_i$ da la frecuencia de la oscilación: con $\Kc=5$ y $K_i=3$, $\omega=\sqrt{2\cdot3/2}=1.732$\,rad/s (periodo $3.63$\,s).
\end{ejemplo}

\begin{figure}[ht]
\centering
\includegraphics[width=0.85\linewidth]{cap05/lazo_pi_estab}
\caption{Lazo del ejemplo~\ref{ej:es-pi}, con el PI en forma paralela.}
\label{fig:es-pi-lazo}
\end{figure}

\begin{figure}[ht]
\centering
\includegraphics[width=\linewidth]{cap05/f_region_pi}
\caption{Ejemplo~\ref{ej:es-pi}. Izquierda: región de estabilidad. Derecha: con $\Kc=5$, $K_i=1$ (adentro) es estable, aunque con $\OS=51\%$; $K_i=3$ (frontera) oscila sin amortiguarse; $K_i=4$ (afuera) es inestable.}
\label{fig:es-region}
\end{figure}

\codigorepo{cap05_estabilidad/C5_Ej5_RegionPI.m}{region}{ejemplo~\ref{ej:es-pi}: la región de estabilidad en el plano $(\Kc,K_i)$}

\begin{modelo}[LazoPI.slx]
El script \cmd{crearModelo\_LazoPI.m} arma el lazo de la \cref{fig:es-pi-lazo} con $\Kc$ y $K_i$ como variables del espacio de trabajo. Recorra la región: un punto adentro, uno sobre la frontera y uno afuera, y mida el periodo de la oscilación sostenida.
\end{modelo}

\begin{cuidado}[Estable no quiere decir aceptable]
El punto $(\Kc,K_i)=(5,1)$ está dentro de la región y, sin embargo, la respuesta tiene un sobrepico del $51\%$ y tarda $22$\,s en establecerse. La tabla de Routh dice \emph{si} es estable; para saber \emph{qué tan} estable hace falta algo más.
\end{cuidado}

\subsection{Estabilidad relativa}

La tabla se puede usar también para exigir un margen: que todos los polos tengan parte real menor que $-\sigma$, no solo menor que cero. Basta con desplazar el eje: se reemplaza $s=z-\sigma$ en la ecuación característica y se aplica el criterio al polinomio en $z$. Las raíces con $\Re z>0$ son las que tienen $\Re s>-\sigma$.

\begin{ejemplo}[label=ej:es-relativa]{Polos a la izquierda de $-0.05$}
En el ejemplo~\ref{ej:es-rangokc}, ¿para qué valores de $\Kc$ todos los polos tienen parte real menor que $-0.05$ (constante de tiempo de la envolvente menor que $20$\,s)?
\solucion Con $s=z-0.05$, $900s^3+420s^2+43s+(1+0.8\Kc)$ se convierte en
\[
  900z^3+285z^2+7.75z+(0.8\Kc-0.2125)=0 .
\]
\[
\begin{array}{c|cc}
  z^3 & 900 & 7.75\\
  z^2 & 285 & 0.8\Kc-0.2125\\
  z^1 & \frac{160-48\Kc}{19} & \\
  z^0 & 0.8\Kc-0.2125
\end{array}
\qquad\Longrightarrow\qquad
\alerta{0.266<\Kc<3.33}.
\]
El margen exigido reduce mucho el rango: de $(-1.25,\,23.83)$ a $(0.266,\,3.33)$. En los dos extremos hay un polo exactamente en $-0.05$ (\cref{fig:es-relativa}).
\end{ejemplo}

\begin{figure}[ht]
\centering
\includegraphics[width=0.5\linewidth]{cap05/f_relativa}
\caption{Ejemplo~\ref{ej:es-relativa}: polos del lazo para $\Kc=0.266$ y $\Kc=3.33$ (extremos del rango con margen) y para $\Kc=\Kcr$, con polos sobre el eje imaginario.}
\label{fig:es-relativa}
\end{figure}

\begin{pausa}
Para $\Gma(s)=\dfrac{K}{s(s+1)(s+2)}$ con retroalimentación unitaria, halle el rango de $K$, $\Kcr$ y $\wcr$.
\respuesta{$s^3+3s^2+2s+K=0$. Por \eqref{eq:es-tercer}: $3\cdot2>K$, y $K>0$: $0<K<6$. En $\Kcr=6$ el auxiliar es $3s^2+6=0$: $\wcr=\sqrt2=1.414$\,rad/s.}
\end{pausa}

%% ---------------------------------------------------------------------
\section{Tiempo muerto: la aproximación de Padé}
\label{sec:es-retardo}
%% ---------------------------------------------------------------------

En muchos procesos la salida no responde de inmediato: el agua tarda en recorrer la tubería entre la válvula y el sensor, la banda transportadora tarda en llevar el material a la báscula, la señal tarda en viajar por la red. La salida es la entrada \emph{desplazada} un tiempo $T$, el \clave{tiempo muerto} o \clave{retardo de transporte}. Por la propiedad de desplazamiento en el tiempo de la transformada de Laplace (capítulo~2),
\[
  y(t)=u(t-T)\,\escalon(t-T)\quad\Longleftrightarrow\quad Y(s)=e^{-Ts}\,U(s).
\]
Este es el único lugar del curso en que se analiza la estabilidad con retardo, y lo volveremos a encontrar en los modelos de primer orden con tiempo muerto que se usan para sintonizar controladores PID (capítulo~8).

\begin{cuidado}[Routh no se puede aplicar directamente]
Con un retardo, la ecuación característica deja de ser un polinomio: para el lazo de la \cref{fig:es-retardo} es
\[
  s+1+Ke^{-Ts}=0 ,
\]
una ecuación \emph{trascendente}, con infinitas raíces. La tabla de Routh necesita coeficientes de un polinomio, así que primero hay que \emph{aproximar} $e^{-Ts}$ por un cociente de polinomios.
\end{cuidado}

\begin{figure}[ht]
\centering
\includegraphics[width=0.8\linewidth]{cap05/lazo_retardo}
\caption{Lazo con retardo de transporte del ejemplo~\ref{ej:es-retardo}.}
\label{fig:es-retardo}
\end{figure}

\subsection*{La aproximación de Padé}

La \clave{aproximación de Padé} de orden $n$ es el cociente de dos polinomios de grado $n$ cuya serie de Taylor coincide con la de $e^{-Ts}$ en tantos términos como es posible~\citep{Pad92}:
\begin{equation}
  e^{-Ts}\approx P_n(s)=\frac{\sum_{k=0}^{n}(-1)^k c_k\,(Ts)^k}{\sum_{k=0}^{n}c_k\,(Ts)^k},
  \qquad
  c_k=\frac{(2n-k)!\;n!}{(2n)!\;k!\;(n-k)!},\quad k=0,\dots,n.
  \label{eq:es-pade}
\end{equation}
Con $n=1$, $c_0=1$ y $c_1=\tfrac12$; con $n=2$, $c_0=1$, $c_1=\tfrac12$ y $c_2=\tfrac1{12}$:
\[
  e^{-Ts}\approx\frac{1-\frac{T}{2}s}{1+\frac{T}{2}s},
  \qquad
  e^{-Ts}\approx\frac{1-\frac{T}{2}s+\frac{(Ts)^2}{12}}{1+\frac{T}{2}s+\frac{(Ts)^2}{12}} .
\]
Observe la estructura: el numerador tiene los mismos coeficientes que el denominador con los signos alternados, así que tiene \emph{ceros en el semiplano derecho} (simétricos a los polos). Por eso la aproximación tiene magnitud $1$ en todas las frecuencias, como el retardo, y por eso su respuesta al escalón arranca en sentido contrario (\cref{fig:es-pade}, izquierda): es la forma que tiene una función racional de ``esperar'' $T$ segundos.

\begin{historia}[Henri Padé]
El matemático francés Henri Padé (1863--1953) estudió en su tesis doctoral de 1892 cómo aproximar una función por cocientes de polinomios, y organizó las aproximaciones en la tabla que lleva su nombre~\citep{Pad92}. La de $e^{-x}$ fue uno de sus ejemplos. Setenta años después, cuando los ingenieros de control quisieron llevar los retardos a las herramientas de los sistemas racionales (Routh, el lugar de las raíces, el espacio de estado), la tomaron de ahí; en \matlab\ es la función \cmd{pade(T,n)}.
\end{historia}

\begin{ejemplo}[label=ej:es-retardo]{Rango de $K$ con tiempo muerto}
Halle, en función de $T$, el rango de $K$ para el cual el lazo de la \cref{fig:es-retardo} tiene un equilibrio asintóticamente estable.
\solucion \textbf{1. Aproximar.} Con Padé de primer orden, $1+K\dfrac{1}{s+1}\cdot\dfrac{1-\frac{T}{2}s}{1+\frac{T}{2}s}=0$, es decir,
\[
  (s+1)\Bigl(1+\tfrac{T}{2}s\Bigr)+K\Bigl(1-\tfrac{T}{2}s\Bigr)=0
  \;\Longrightarrow\;
  \tfrac{T}{2}\,s^2+\Bigl(1+\tfrac{T}{2}-K\tfrac{T}{2}\Bigr)s+(1+K)=0 .
\]
\textbf{2. Tabla.}
\[
\begin{array}{c|cc}
  s^2 & \frac{T}{2} & 1+K\\[3pt]
  s^1 & 1+\frac{T}{2}-K\frac{T}{2} & \\[3pt]
  s^0 & 1+K
\end{array}
\]
\textbf{3. Condiciones.} $1+K>0$ y $1+\frac{T}{2}(1-K)>0$:
\[
  \alerta{-1<K<1+\frac{2}{T}} .
\]
El rango se \emph{estrecha} a medida que crece el retardo: con $T=1$\,s, $K<3$; con $T=5$\,s, $K<1.4$. Y si $T\to0$, la cota superior desaparece, como debe ser: sin retardo el lazo es de primer orden, estable para todo $K>-1$.
\end{ejemplo}

\begin{ideaclave}
\textbf{El retardo consume estabilidad.}\enspace El retardo no cambia la magnitud de la malla abierta, pero le resta una fase $\omega T$ que crece sin límite con la frecuencia. Por eso un lazo que sin retardo era estable para cualquier ganancia, con retardo tiene una ganancia crítica, y esta baja cuando $T$ crece. Es la misma idea del margen de fase del capítulo~10: el retardo se ``come'' el margen.
\end{ideaclave}

\subsection*{¿Qué tan buena es la aproximación?}

Para este lazo el límite exacto se puede calcular. En el borde de la estabilidad hay raíces en $s=\pm j\omega$; reemplazando en $1+Ke^{-j\omega T}/(j\omega+1)=0$ y separando magnitud y fase,
\[
  K=\sqrt{1+\omega^2},\qquad \omega T+\arctan\omega=\pi .
\]
Con $T=1$: $\omega=2.029$\,rad/s y $K_{\max}=2.262$. Padé de primer orden predice $3$: un error del $33\%$, y \alerta{hacia el lado peligroso}. Con Padé de segundo orden la tabla da $2.292$, apenas $1.3\%$ por encima (\cref{fig:es-pade}, centro).

\begin{figure}[ht]
\centering
\includegraphics[width=\linewidth]{cap05/f_pade}
\caption{Izquierda: escalón retrasado $1$\,s y sus aproximaciones de Padé (arrancan en sentido contrario). Centro: ganancia máxima en función del retardo; Padé~1 sobrestima entre $15$ y $33\%$, Padé~2 alrededor del $1\%$. Derecha: con $K=2.5$ y $T=1$, el modelo de Padé~1 es estable, pero el lazo con el retardo exacto oscila y crece.}
\label{fig:es-pade}
\end{figure}

\begin{cuidado}[Padé de primer orden es optimista]
La aproximación de primer orden sobrestima la ganancia máxima: un diseño que la tabla declara estable puede no serlo (\cref{fig:es-pade}, derecha). Úsela para entender la tendencia ($K_{\max}$ baja con $T$) y para un primer cálculo; para decidir, use Padé de segundo orden o más, o verifique con el retardo exacto: \cmd{margin} y \cmd{allmargin} aceptan modelos con retardo (\texttt{tf(1,[1 1],'InputDelay',T)}), y en \simulink\ está el bloque \emph{Transport Delay}.
\end{cuidado}

\codigorepo{cap05_estabilidad/C5_Ej8_Pade.m}{pade}{ejemplo~\ref{ej:es-retardo}: Padé de orden 1 y 2 contra el retardo exacto}

\begin{modelo}[Retardo.slx]
El script \cmd{crearModelo\_Retardo.m} arma el lazo de la \cref{fig:es-retardo} dos veces: una con el bloque \emph{Transport Delay} y otra con la aproximación de Padé de primer orden. Con $T=1$, suba $K$ desde $2$ hasta $3$ y encuentre la ganancia a la que cada uno empieza a oscilar.
\end{modelo}

\begin{pausa}
Para $\Gma(s)=\dfrac{Ke^{-Ts}}{s+2}$, use Padé de primer orden para hallar el rango de $K$ en función de $T$. ¿Cuánto vale la cota con $T=1$?
\respuesta{$(s+2)(1+\tfrac{T}{2}s)+K(1-\tfrac{T}{2}s)=\tfrac{T}{2}s^2+(1+T-K\tfrac{T}{2})s+(2+K)$: $-2<K<2+\tfrac{2}{T}$. Con $T=1$, $K<4$ (el valor exacto es $3.04$).}
\end{pausa}

%% ---------------------------------------------------------------------
\section{Estabilidad en el espacio de estado}
\label{sec:es-espacio}
%% ---------------------------------------------------------------------

Con el modelo de estado $\dot x=\Am x+\Bm u$, $y=\Cm x+\Dm u$ (capítulo~4), la respuesta natural es $x(t)=e^{\Am t}x(0)$, y sus modos son $e^{\lambda_it}$, con $\lambda_i$ los \clave{valores propios} de $\Am$: las raíces de
\[
  \det(s\Id-\Am)=0 .
\]
Como $\dfrac{Y(s)}{U(s)}=\Cm(s\Id-\Am)^{-1}\Bm+\Dm=\dfrac{\Cm\operatorname{adj}(s\Id-\Am)\Bm}{\det(s\Id-\Am)}+\Dm$, los polos de la función de transferencia están entre los valores propios de $\Am$.

\begin{definicion}[label=def:es-hurwitz]{Matriz Hurwitz}
Una matriz $\Am$ es \clave{Hurwitz} si todos sus valores propios tienen parte real negativa. El equilibrio $x=0$ de $\dot x=\Am x$ es asintóticamente estable si y solo si $\Am$ es Hurwitz. Para decidirlo sin calcular los valores propios, se aplica Routh-Hurwitz al polinomio característico $\det(s\Id-\Am)$.
\end{definicion}

\begin{ejemplo}[label=ej:es-matriz]{Routh sobre el polinomio característico}
¿Es Hurwitz $\Am=\begin{bmatrix}0&1&0\\0&0&1\\-6&-11&-6\end{bmatrix}$?
\solucion Está en variables de fase, así que el polinomio característico se lee en la última fila: $s^3+6s^2+11s+6$. Por \eqref{eq:es-tercer}, $6\cdot11=66>1\cdot6$: es Hurwitz. (Sus valores propios son $-1$, $-2$ y $-3$.)
\end{ejemplo}

\subsection*{Lo que la función de transferencia esconde}

``Los polos están entre los valores propios'' quiere decir que pueden ser \emph{menos}: si hay cancelaciones, la función de transferencia no muestra todos los modos (\cref{sec:re-kalman}). Y un modo que no se ve puede ser inestable.

\begin{ejemplo}[label=ej:es-escondido]{Un modo inestable escondido}
Sea $\Am=\begin{bmatrix}0&1&0\\0&0&1\\2&1&-2\end{bmatrix}$, $\Bm=\begin{bmatrix}0\\0\\1\end{bmatrix}$, $\Cm=[-1\;\;0\;\;1]$.
\solucion El polinomio característico es $s^3+2s^2-s-2=(s-1)(s+1)(s+2)$: hay un signo menos, y el valor propio $+1$ hace que el equilibrio sea \alerta{inestable}. Sin embargo, la función de transferencia es
\[
  \frac{Y(s)}{U(s)}=\frac{s^2-1}{(s-1)(s+1)(s+2)}=\frac{1}{s+2},
\]
que parece perfectamente estable. El modo $e^{t}$ está ahí: cualquier condición inicial que lo excite hace crecer los estados, aunque no se vea en la salida.
\end{ejemplo}

\begin{cuidado}[Routh sobre $\det(s\Id-\Am)$, no sobre la función de transferencia]
Para decidir la estabilidad del \emph{equilibrio} se analiza el polinomio característico completo de $\Am$. Analizar el denominador de una función de transferencia ya simplificada solo dice si el sistema es BIBO estable \emph{entre esa entrada y esa salida}. Por la misma razón, nunca se cancela un polo inestable con un cero del controlador.
\end{cuidado}

\codigorepo{cap05_estabilidad/C5_Ej6_EspacioEstado.m}{escondido}{ejemplo~\ref{ej:es-escondido}: el modo inestable que la función de transferencia no muestra}

\subsection*{Equilibrios de sistemas no lineales: el método indirecto de Lyapunov}

Un sistema no lineal $\dot x=f(x,u)$ puede tener varios equilibrios. Para cada uno, la linealización de la \cref{sec:so-linealizacion} da $\dot{\dev x}=\Am\,\dev x+\Bm\,\dev u$, con $\Am=\partial f/\partial x$ evaluada en el equilibrio. El \clave{método indirecto de Lyapunov} dice:
\begin{itemize}
  \item si $\Am$ es Hurwitz, el equilibrio del sistema no lineal es \emph{localmente} asintóticamente estable;
  \item si $\Am$ tiene algún valor propio con parte real positiva, el equilibrio es inestable;
  \item si los valores propios con parte real máxima están sobre el eje imaginario, la linealización \alerta{no decide}: hay que mirar los términos no lineales.
\end{itemize}

\begin{ejemplo}[label=ej:es-pendulo]{Los dos equilibrios del péndulo}
El péndulo con fricción $\ddot\theta=-\frac{g}{l}\sin\theta-b\,\dot\theta$, con $g/l=9.81$\,s$^{-2}$ y $b=0.5$\,s$^{-1}$, tiene equilibrios en $\theta^*=0$ (abajo) y $\theta^*=\pi$ (arriba). Decida su estabilidad.
\solucion Con $x_1=\theta$, $x_2=\dot\theta$, el jacobiano es $\Am=\begin{bmatrix}0&1\\-\frac{g}{l}\cos\theta^*&-b\end{bmatrix}$.
\begin{itemize}
  \item Abajo, $\cos0=1$: $s^2+0.5s+9.81$, coeficientes positivos y segundo orden: estable. Valores propios $-0.25\pm j3.12$.
  \item Arriba, $\cos\pi=-1$: $s^2+0.5s-9.81$, un signo menos: inestable. Valores propios $2.89$ y $-3.39$.
\end{itemize}
Sin fricción ($b=0$), abajo quedaría $s^2+9.81$, con valores propios sobre el eje: la linealización no decidiría (aunque, por conservación de la energía, el péndulo sin fricción es marginalmente estable).
\end{ejemplo}

\begin{pausa}
¿Para qué valores de $a$ es Hurwitz $\Am(a)=\begin{bmatrix}0&1&0\\0&0&1\\-a&-3&-2\end{bmatrix}$? ¿A qué frecuencia oscila en el límite?
\respuesta{$\det(s\Id-\Am)=s^3+2s^2+3s+a$. Por \eqref{eq:es-tercer}: $0<a<6$. En $a=6$ el auxiliar es $2s^2+6$: $\omega=\sqrt3=1.732$\,rad/s.}
\end{pausa}

%% ---------------------------------------------------------------------
\section{Robustez: el teorema de Kharitonov}
\label{sec:es-kharitonov}
%% ---------------------------------------------------------------------

Routh-Hurwitz decide la estabilidad de \emph{un} polinomio. Pero los parámetros de una planta real nunca se conocen exactamente: una ganancia está entre $0.8$ y $1.3$, un amortiguamiento entre $0.2$ y $0.3$, una frecuencia natural entre $10$ y $12$\,rad/s. El modelo es en realidad una \emph{familia} de modelos, y lo que queremos es que el lazo sea estable para \emph{todos} ellos: que sea \clave{robustamente estable}.

Con la planta de segundo orden de ese ejemplo, $\Gp(s)=\dfrac{\Kproc}{s^2+2\zeta\wn s+\wn^2}=\dfrac{q_0}{r_0s^2+r_1s+r_2}$, los intervalos son $q_0\in[0.8,\,1.3]$, $r_0=1$, $r_1=2\zeta\wn\in[4,\,7.2]$ y $r_2=\wn^2\in[100,\,144]$. Una planta cuyos coeficientes varían en intervalos se llama \clave{planta de intervalo}. Con un controlador $\Gc=N_c/D_c$, la ecuación característica $D_cD_p+N_cN_p=0$ tiene coeficientes que también varían en intervalos, y hay infinitos polinomios que revisar.

\begin{definicion}[label=def:es-intervalo]{Familia de polinomios de intervalo}
Es el conjunto de polinomios
\[
  \chi(s)=a_0s^n+a_1s^{n-1}+\cdots+a_{n-1}s+a_n,\qquad a_k\in[a_k^-,\,a_k^+],
\]
en el que cada coeficiente varía \emph{independientemente} de los demás en su intervalo (con $0\notin[a_0^-,a_0^+]$). La familia es \clave{robustamente estable} si todos sus polinomios son Hurwitz.
\end{definicion}

\begin{definicion}[label=def:es-kharitonov]{Teorema de Kharitonov}
Una familia de polinomios de intervalo es robustamente estable \alerta{si y solo si} son Hurwitz los cuatro polinomios
\[
\begin{aligned}
  K_1(s)&=a_n^-+a_{n-1}^-\,s+a_{n-2}^+\,s^2+a_{n-3}^+\,s^3+a_{n-4}^-\,s^4+a_{n-5}^-\,s^5+\cdots\\
  K_2(s)&=a_n^++a_{n-1}^+\,s+a_{n-2}^-\,s^2+a_{n-3}^-\,s^3+a_{n-4}^+\,s^4+a_{n-5}^+\,s^5+\cdots\\
  K_3(s)&=a_n^-+a_{n-1}^+\,s+a_{n-2}^+\,s^2+a_{n-3}^-\,s^3+a_{n-4}^-\,s^4+a_{n-5}^+\,s^5+\cdots\\
  K_4(s)&=a_n^++a_{n-1}^-\,s+a_{n-2}^-\,s^2+a_{n-3}^+\,s^3+a_{n-4}^+\,s^4+a_{n-5}^-\,s^5+\cdots
\end{aligned}
\]
escritos desde el término independiente hacia arriba, con los patrones de extremos $(-,-,+,+)$, $(+,+,-,-)$, $(-,+,+,-)$ y $(+,-,-,+)$ que se repiten cada cuatro términos~\citep{Kha78}.
\end{definicion}

Infinitas verificaciones se reducen a cuatro, y a cada una se le aplica Routh-Hurwitz. ¿Por qué esos cuatro y no, simplemente, ``todos los coeficientes en el mínimo'' y ``todos en el máximo''? Porque en $s=j\omega$ el polinomio vale
\[
  \chi(j\omega)=\underbrace{(a_n-a_{n-2}\omega^2+a_{n-4}\omega^4-\cdots)}_{\text{parte real: coeficientes pares}}
  +j\,\omega\underbrace{(a_{n-1}-a_{n-3}\omega^2+a_{n-5}\omega^4-\cdots)}_{\text{parte imaginaria: coeficientes impares}} ,
\]
con signos que se alternan de dos en dos. La parte real es mínima con $(a_n^-,a_{n-2}^+,a_{n-4}^-,\dots)$ y máxima con los extremos opuestos, y lo mismo pasa con la parte imaginaria. Para cada $\omega$, los valores de toda la familia llenan un \emph{rectángulo} del plano complejo cuyos vértices son exactamente $K_1,\dots,K_4$ (\cref{fig:es-rectangulos}). La familia pierde la estabilidad cuando algún polinomio tiene una raíz en el eje imaginario, es decir, cuando algún rectángulo toca el origen; y eso no puede pasar sin que uno de los vértices lo cruce antes.

\begin{figure}[ht]
\centering
\includegraphics[width=0.85\linewidth]{cap05/f_kharitonov}
\caption{Valores de $\chi(j\omega)$ para la familia $[1,2]s^3+[4,5]s^2+[3,4]s+[1,2]$: para cada $\omega$ forman un rectángulo cuyos vértices (puntos de colores) son los cuatro polinomios de Kharitonov. Ningún rectángulo toca el origen: la familia es robustamente estable.}
\label{fig:es-rectangulos}
\end{figure}

\begin{ejemplo}[label=ej:es-khar1]{Un compensador sobre una planta de intervalo}
Sean $\Gc(s)=\dfrac{s+2}{s^2+2s+2}$ y $\Gp(s)=\dfrac{q_0}{r_0s^3+r_1s^2+r_2s+r_3}$, con $r_0\in[1,\,1.1]$, $r_1\in[4,\,4.2]$, $r_2\in[6,\,8]$, $r_3\in[10,\,20]$ y $q_0\in[3,\,5]$. ¿Es robustamente estable el lazo?
\solucion La ecuación característica $(s^2+2s+2)(r_0s^3+r_1s^2+r_2s+r_3)+q_0(s+2)=0$ es
\[
  a_0s^5+a_1s^4+a_2s^3+a_3s^2+a_4s+a_5=0,
\]
con
\[
\begin{aligned}
  a_0&=r_0\in[1,\,1.1], & a_1&=2r_0+r_1\in[6,\,6.4], & a_2&=2r_0+2r_1+r_2\in[16,\,18.6],\\
  a_3&=2r_1+2r_2+r_3\in[30,\,44.4], & a_4&=2r_2+2r_3+q_0\in[35,\,61], & a_5&=2r_3+2q_0\in[26,\,50].
\end{aligned}
\]
Los cuatro polinomios de Kharitonov son (de $s^0$ a $s^5$)
\[
\begin{aligned}
  K_1&: (26,\,35,\,44.4,\,18.6,\,6,\,1) &&\text{Hurwitz}\\
  K_2&: (50,\,61,\,30,\,16,\,6.4,\,1.1) &&\text{dos raíces en }0.364\pm j2.01\\
  K_3&: (26,\,61,\,44.4,\,16,\,6,\,1.1) &&\text{dos raíces en }0.041\pm j2.60\\
  K_4&: (50,\,35,\,30,\,18.6,\,6.4,\,1) &&\text{dos raíces en }0.194\pm j1.50
\end{aligned}
\]
Tres de los cuatro son inestables: la familia de intervalo \alerta{no es robustamente estable}. Pero hay que tener cuidado con la conclusión: los seis coeficientes dependen de solo cinco parámetros, así que no varían independientemente, y la ``caja'' de coeficientes contiene polinomios que la planta real no puede producir. ¿Hay algún caso \emph{real} inestable? Sí: con $r_0=1$, $r_1=4$, $r_2=6$, $r_3=20$ y $q_0=5$, el lazo tiene polos en $0.033\pm j2.26$. El lazo con la planta nominal (valores centrales) sí es estable, pero el diseño no es robusto.
\end{ejemplo}

\begin{cuidado}[Cuando los coeficientes no son independientes]
El teorema supone que cada coeficiente varía por su cuenta. Si varios coeficientes dependen de los mismos parámetros físicos, Kharitonov se aplica a una familia \emph{más grande} que la real: si dice ``robustamente estable'', la conclusión vale (la familia real está adentro); si dice ``no'', el resultado es \alerta{conservador} y hay que buscar un caso real inestable para confirmarlo.
\end{cuidado}

\begin{ejemplo}[label=ej:es-khar-diseno]{Rango robusto de ganancia}
En el lazo de la \cref{fig:es-khar}, $r_1\in[6,\,7]$, $r_2\in[100,\,120]$ y $r_3\in[1,\,2]$. Halle el rango de $K$ para el cual el lazo es robustamente estable.
\solucion La ecuación característica es $(s^2+r_1s+r_2)(r_3s+1)+K=0$:
\[
  r_3s^3+(r_1r_3+1)\,s^2+(r_2r_3+r_1)\,s+(r_2+K)=0,
\]
con $a_0\in[1,\,2]$, $a_1\in[7,\,15]$, $a_2\in[106,\,247]$ y $a_3\in[100+K,\,120+K]$. Para tercer orden, la condición \eqref{eq:es-tercer} es $a_1a_2>a_0a_3$, y el caso más desfavorable es el que hace el lado izquierdo mínimo y el derecho máximo: $a_1^-a_2^-$ contra $a_0^+a_3^+$. Ese es precisamente $K_4$ (los otros tres polinomios son más holgados), y da
\[
  7\cdot106=742>2\,(120+K)\;\Longrightarrow\;K<251 .
\]
Con $a_3>0$ para todos los casos, $K>-100$: Kharitonov garantiza $\alerta{-100<K<251}$.

¿Qué tan conservador es? Aquí $r_3$ aparece en tres coeficientes y $r_1$, $r_2$ en dos. Haciendo la cuenta con los parámetros, la condición es $r_1r_2r_3+r_1^2+r_1/r_3>K$, cuyo mínimo (en $r_1=6$, $r_2=100$, $r_3=1$) es $642$. La familia real tolera $K<642$: Kharitonov da menos de la mitad del rango (\cref{fig:es-khar-diseno}). Con la planta nominal ($r_1=6.5$, $r_2=110$, $r_3=1.5$), Routh daría $K<1119$: diseñar con el modelo nominal sería peligroso.
\end{ejemplo}

\begin{figure}[ht]
\centering
\includegraphics[width=0.72\linewidth]{cap05/lazo_khar}
\caption{Lazo con planta y sensor inciertos del ejemplo~\ref{ej:es-khar-diseno}.}
\label{fig:es-khar}
\end{figure}

\begin{figure}[ht]
\centering
\includegraphics[width=0.85\linewidth]{cap05/f_khar_diseno}
\caption{Ejemplo~\ref{ej:es-khar-diseno}: parte real del polo más a la derecha en el peor caso de la familia real. Kharitonov garantiza $K<251$; el límite real es $642$; con el modelo nominal se creería que es $1119$.}
\label{fig:es-khar-diseno}
\end{figure}

\codigorepo{cap05_estabilidad/C5_Ej7_Kharitonov.m}{kharitonov}{los cuatro polinomios con la función \cmd{kharitonov} y la comparación con un muestreo aleatorio}

\begin{profundiza}[Menos de cuatro polinomios]
Para órdenes bajos no hacen falta los cuatro: con $n=3$ basta revisar uno ($K_4$, el que usamos arriba), con $n=4$ dos y con $n=5$ tres; desde $n=6$ hacen falta los cuatro~\citep{AJM87}. El teorema también tiene límites: exige coeficientes independientes, y no dice nada sobre el desempeño, solo sobre la estabilidad.
\end{profundiza}

\begin{pausa}
La familia $a_3s^3+a_2s^2+a_1s+a_0$ tiene $a_3\in[1,2]$, $a_2\in[4,5]$, $a_1\in[3,4]$ y $a_0\in[1,2]$. ¿Cuántos polinomios hay que verificar? ¿Es robustamente estable?
\respuesta{Cuatro (y en tercer orden basta el peor, $K_4$). Peor caso de \eqref{eq:es-tercer}: $a_2^-a_1^-=4\cdot3=12>a_3^+a_0^+=2\cdot2=4$. Sí es robustamente estable.}
\end{pausa}

%% ---------------------------------------------------------------------
\section{Para llevar}
\label{sec:es-cierre}
%% ---------------------------------------------------------------------

\begin{parallevar}
\begin{enumerate}
  \item La estabilidad es una propiedad del \emph{equilibrio}. Asintóticamente estable: la respuesta natural tiende a cero (todos los polos con parte real negativa). Marginal: se mantiene acotada (polos simples sobre el eje imaginario). Solo la asintótica garantiza BIBO. Los ceros no intervienen.
  \item La estabilidad del lazo se decide en la ecuación característica $1+\Gma(s)=0$.
  \item Condición necesaria: todos los coeficientes presentes y del mismo signo. Suficiente solo para $n\le2$; para $n=3$, además, $a_1a_2>a_0a_3$.
  \item Routh-Hurwitz: el número de raíces con parte real positiva es el número de cambios de signo en la primera columna. Cero aislado en la primera columna: $\epsilon$. Fila de ceros: polinomio auxiliar (raíces simétricas, que se pueden calcular).
  \item Con parámetros en los coeficientes, la tabla da rangos y regiones de diseño. La fila que se anula da $\Kcr$ y el auxiliar, $\wcr$: los mismos números de Ziegler-Nichols y del margen de ganancia. Con $s=z-\sigma$ se exige un margen de estabilidad.
  \item En el espacio de estado se analiza $\det(s\Id-\Am)$, no una función de transferencia simplificada: puede esconder modos inestables. Linealizando, $\Am$ Hurwitz implica estabilidad local del equilibrio no lineal.
  \item Kharitonov: una familia de intervalo es robustamente estable si y solo si cuatro polinomios lo son. Si los coeficientes dependen de los mismos parámetros, el resultado es conservador.
\end{enumerate}
\end{parallevar}

\begin{profundiza}[Correcciones respecto a la versión anterior del capítulo]
(i) Se retiró la sección de tiempo muerto (aproximación de Padé) y sus dos ejercicios: el curso no cubre retardos. (ii) ``Estable si todos los polos y ceros están en el semiplano izquierdo'' $\to$ solo los polos. (iii) $s^3+2s^2+s+2$ es un caso de fila de ceros, no de $\epsilon$; el auxiliar da las raíces $\pm j$. (iv) Se resolvieron los ejemplos que no tenían solución ($s^5+7s^4+6s^3+42s^2+8s+56$, el primer ejemplo de Kharitonov y el diseño robusto). (v) En el primer ejemplo de Kharitonov faltaba el intervalo de $a_4\in[35,61]$, y la familia \emph{no} es robustamente estable. (vi) El rango robusto del diseño es $-100<K<251$ por Kharitonov y $K<642$ con la dependencia entre coeficientes. (vii) $C(s)/R(s)\to Y(s)/R(s)$; Kharitanov $\to$ Kharitonov. (viii) Historia: Hurwitz publicó en 1895 en \emph{Mathematische Annalen}; el premio Adams de Routh fue en 1877; la anécdota del sirviente es de lord Kelvin. Todos los números se verificaron con \texttt{herramientas/cap05\_numeros.py} y \cmd{verificarCap05.m}.
\end{profundiza}

%% ---------------------------------------------------------------------
\section{Problemas}
\label{sec:es-problemas}
%% ---------------------------------------------------------------------

Los problemas marcados con \etiquetaM{} requieren \matlab/\simulink; los marcados con \etiquetaEstrella{} son de profundización.

\subsection*{El criterio}

\begin{problema}[tablas]\label{prob:es-tablas}
Construya la tabla de Routh, cuente las raíces con parte real positiva y clasifique el equilibrio (asintóticamente estable, marginal o inestable): (a) $s^4+3s^3+3s^2+2s+1$; (b) $s^5+2s^4+3s^3+6s^2+2s+1$; (c) $s^5+3s^4+5s^3+15s^2+4s+12$; (d) $s^4+2s^3-3s^2+4s+5$. En (d) justifique \emph{sin} construir la tabla completa. En (c) identifique el polinomio auxiliar y calcule las raíces que están sobre el eje imaginario. \etiquetaM{} Verifique con \cmd{tablaRouth} y \cmd{roots}.
\end{problema}

\begin{problema}[ganancia crítica]\label{prob:es-critica}
Con retroalimentación unitaria negativa: (i) $\Gma(s)=\dfrac{K}{s(s+2)(s+5)}$; (ii) $\Gma(s)=\dfrac{K(s+1)}{s(s-1)(s+6)}$. (a) Escriba las ecuaciones características. (b) Halle el rango de $K$ para estabilidad asintótica. (c) Para (i), halle $\Kcr$ y la frecuencia de la oscilación sostenida. (d) \etiquetaM{} Simule (i) ante un escalón con $K=0.5\Kcr$, $\Kcr$ y $1.5\Kcr$ y mida el periodo en el caso crítico. (e) Explique por qué en (ii) aparece una cota \emph{inferior} para $K$ y verifíquelo ubicando los polos para tres valores de $K$.
\end{problema}

\begin{problema}[ganancias negativas y Bode]\label{prob:es-bode}
Para $\Gma(s)=\dfrac{K}{(s+1)(s+3)(s+5)}$, con $K\in\mathbb{R}$: (a) halle el rango de $K$ (incluidos los negativos); (b) $\Kcr>0$ y la frecuencia de oscilación; (c) ¿cuántos polos en el semiplano derecho hay con $K=-20$? (d) \etiquetaM{} Con $K=1$, compare $\Kcr$ y $\wcr$ con lo que da \cmd{margin}. ¿Por qué el margen de ganancia no le dice nada sobre la cota negativa?
\end{problema}

\subsection*{Diseño}

\begin{problema}[región de un PI]\label{prob:es-pi}
La planta $\Gp(s)=\dfrac{1}{(s+1)(s+2)}$ se controla con $\Gc(s)=\Kc+K_i/s$ y retroalimentación unitaria. (a) Halle las condiciones sobre $\Kc$ y $K_i$ para estabilidad asintótica. (b) Dibuje la región en el plano $(\Kc,K_i)$. (c) \etiquetaM{} Escoja un punto adentro, uno sobre la frontera y uno afuera, y simule. (d) Sobre la frontera, calcule la frecuencia de la oscilación con el auxiliar y compárela con la simulación. (e) Exija que todos los polos tengan parte real menor o igual a $-0.5$ con $s=z-0.5$, y superponga la nueva región. (f) \etiquetaM{} Tome un punto de la región de (b) que no esté en la de (e), simule y explique qué hace inaceptable el diseño.
\end{problema}

\begin{problema}[estabilidad local]\label{prob:es-local}
El sistema $\dot x_1=x_2$, $\dot x_2=-2x_1+x_1^2-c\,x_2$ tiene equilibrios en $x^*=(0,0)$ y $x^*=(2,0)$. (a) Linealice alrededor de cada uno. (b) Con $c=1$, decida la estabilidad de cada equilibrio con Routh-Hurwitz y con los valores propios. (c) ¿Qué dice la linealización si $c=0$? (d) \etiquetaM{} Simule el sistema no lineal con $c=1$ desde $(0.5,0)$ y desde $(2.1,0)$ y relacione lo que ve con (b).
\end{problema}

\subsection*{Espacio de estado}

\begin{problema}[matrices Hurwitz]\label{prob:es-matrices}
(a) Decida si $\Am_1=\begin{bmatrix}0&1&0\\0&0&1\\-6&-11&-6\end{bmatrix}$ y $\Am_2=\begin{bmatrix}-1&2\\3&-4\end{bmatrix}$ son Hurwitz, con Routh sobre $\det(s\Id-\Am)$ y con los valores propios. (b) Halle el rango de $a$ para el que $\Am_3(a)=\begin{bmatrix}0&1&0\\0&0&1\\-a&-3&-2\end{bmatrix}$ es Hurwitz. (c) \etiquetaM{} En el límite superior, calcule la frecuencia de oscilación con el auxiliar y verifíquela simulando desde $x(0)=[1\;\;0\;\;0]^\top$. (d) Con $\Am_1$, $\Bm=[0\;\;0\;\;1]^\top$ y $\Dm=0$, encuentre un $\Cm$ para el que la función de transferencia tenga \emph{menos} polos que valores propios tiene $\Am_1$. ¿Por qué Routh sobre esa función de transferencia no basta para hablar del equilibrio?
\end{problema}

\subsection*{Robustez}

\begin{problema}[familia de cuarto orden]\label{prob:es-familia}
Sea $\chi(s)=a_0s^4+a_1s^3+a_2s^2+a_3s+a_4$ con $a_0\in[1,1.2]$, $a_1\in[4,5]$, $a_2\in[6,8]$, $a_3\in[3,4]$, $a_4\in[1,2]$. (a) Escriba los cuatro polinomios de Kharitonov. (b) Aplique Routh a cada uno y concluya. (c) \etiquetaM{} Genere 2000 polinomios aleatorios de la familia y verifique que ninguno es inestable; compare el costo con (b). (d) Repita (a)--(b) con $a_3\in[0.2,4]$. (e) Si la respuesta a (d) es negativa, muestre un polinomio inestable de la familia y sus raíces.
\end{problema}

\begin{problema}[planta incierta]\label{prob:es-incierta}
\etiquetaEstrella{} Con $\Gc=K$ y $\Gp(s)=\dfrac{1}{s^3+r_1s^2+r_2s+r_3}$, $r_1\in[6,7]$, $r_2\in[100,120]$, $r_3\in[1,2]$: (a) escriba el polinomio característico y los intervalos de sus coeficientes; (b) con Kharitonov, halle el rango robusto de $K$ y diga qué polinomio impone la restricción; (c) compare con el rango de la planta nominal ($6.5$, $110$, $1.5$); (d) escoja un $K$ admisible para la nominal pero no para la familia y muestre un caso inestable; (e) \etiquetaM{} recorra los ocho vértices $(r_1,r_2,r_3)$ y verifique el peor caso. (f) ¿Es conservador aquí el resultado de Kharitonov? Compare con el ejemplo~\ref{ej:es-khar-diseno}.
\end{problema}

%% ---------------------------------------------------------------------
\section{Soluciones seleccionadas}
%% ---------------------------------------------------------------------

\paragraph{Problema~\ref{prob:es-tablas}}
(a) Sin cambios: asintóticamente estable. (b) Un cero en la fila $s^3$ ($\epsilon$) y dos cambios: inestable, raíces $0.096\pm j1.62$. (c) Fila $s^3$ de ceros; auxiliar $3s^4+15s^2+12=3(s^2+1)(s^2+4)$: raíces $\pm j$ y $\pm2j$, sin cambios: marginalmente estable ($s^5+\dots=(s+3)(s^2+1)(s^2+4)$). (d) Hay un coeficiente negativo: no es estable (la tabla da dos cambios: $0.956\pm j1.115$).

\paragraph{Problema~\ref{prob:es-critica}}
(i) $s^3+7s^2+10s+K$: $0<K<70$; $\Kcr=70$, auxiliar $7s^2+70$, $\wcr=\sqrt{10}=3.16$\,rad/s ($\Pcr=1.99$\,s). (ii) $s^3+5s^2+(K-6)s+K$: $K>0$ y $\frac{5(K-6)-K}{5}>0$, o sea $K>7.5$. Con poca ganancia el polo en $+1$ no alcanza a moverse al semiplano izquierdo.

\paragraph{Problema~\ref{prob:es-bode}}
(a) $s^3+9s^2+23s+15+K$: $-15<K<192$. (b) $\Kcr=192$, $\wcr=\sqrt{23}=4.80$\,rad/s. (c) Con $K=-20$ el término independiente es $-5$: un cambio de signo, un polo en $+0.20$. (d) \cmd{margin} da $\GM=192$ en $4.80$\,rad/s; la cota negativa corresponde a un polo que cruza por el origen ($\omega=0$), que el margen de ganancia (definido en el cruce de $-180^\circ$) no mira.

\paragraph{Problema~\ref{prob:es-pi}}
(a) $s^3+3s^2+(2+\Kc)s+K_i$: $K_i>0$ y $K_i<3(2+\Kc)$. (d) Sobre la frontera, el auxiliar $3s^2+K_i$ da $\omega=\sqrt{2+\Kc}$. (e) $z^3+1.5z^2+(\Kc-0.25)z+(K_i-\Kc/2-0.375)$: $\frac{\Kc}{2}+0.375<K_i<2\Kc$. (f) Por ejemplo $(\Kc,K_i)=(2,5)$: estable, pero con polos en $-0.39\pm j1.45$, poco amortiguados.

\paragraph{Problema~\ref{prob:es-local}}
(a) $\Am=\begin{bmatrix}0&1\\-2+2x_1^*&-c\end{bmatrix}$. (b) En $(0,0)$: $s^2+s+2$, estable ($-0.5\pm j1.32$). En $(2,0)$: $s^2+s-2$, inestable ($1$ y $-2$). (c) Con $c=0$, en el origen los valores propios son $\pm j\sqrt2$: la linealización no decide.

\paragraph{Problema~\ref{prob:es-matrices}}
(a) $\Am_1$: $s^3+6s^2+11s+6$, Hurwitz. $\Am_2$: $s^2+5s-2$, no ($0.372$ y $-5.372$). (b) $0<a<6$. (c) $\omega=\sqrt3$. (d) Por ejemplo $\Cm=[2\;\;3\;\;1]$: el numerador $s^2+3s+2=(s+1)(s+2)$ cancela dos polos y queda $1/(s+3)$.

\paragraph{Problema~\ref{prob:es-familia}}
(b) Los cuatro son Hurwitz: robustamente estable. (d) Con $a_3\in[0.2,4]$, $K_1$ y $K_4$ son inestables. (e) $K_1$: $s^4+5s^3+8s^2+0.2s+1$, con raíces $0.026\pm j0.350$.

\paragraph{Problema~\ref{prob:es-incierta}}
(a) $s^3+r_1s^2+r_2s+(r_3+K)$, con coeficientes independientes. (b) Peor caso de \eqref{eq:es-tercer}: $6\cdot100>2+K$: $-1<K<598$. (c) Nominal: $-1.5<K<713.5$. (f) No es conservador: cada coeficiente depende de un solo parámetro, y el peor vértice es una planta real.

\respuestaspausas

                    % migrado (Handout4_Ctrl.tex)
%% =====================================================================
%%  Capítulo 6 -- Error en estado estacionario
%%  Migrado al formato nuevo (noviembre de 2026) desde Handout5_Ctrl.tex.
%%  Correcciones de fondo: ver MIGRACION.md y la caja "Correcciones" de la
%%  sección 6.8 (fórmulas de kv y ka, retroalimentación no unitaria con el
%%  lazo abierto G_OL, sensitividad complementaria T en lugar de C).
%% =====================================================================

\capitulo[cap:er]{Error en estado estacionario}
  {Los lectores del tercer mundo quieren que alguien les arme una buena ficción con la pésima realidad que les ha tocado en mala suerte.}
  {Rodrigo Fresán}

\begin{objetivos}[partes/06-error.html]
  \item Relacionar el error de seguimiento con la referencia, las perturbaciones y el ruido a través de las funciones de sensitividad.
  \item Calcular el error en estado estacionario ante escalón, rampa y parábola con el tipo del sistema y las constantes $\kp$, $\kv$ y $\ka$, después de verificar la estabilidad.
  \item Tratar la retroalimentación no unitaria con el lazo abierto $\Gla$ y distinguir el error de seguimiento del error de actuación.
  \item Calcular el error que deja una perturbación, el error en espacio de estado, y reconocer cuándo una especificación de error choca con una de respuesta transitoria.
\end{objetivos}

%% ---------------------------------------------------------------------
\section{Por qué retroalimentar: sensitividad y error}
\label{sec:er-sensitividad}
%% ---------------------------------------------------------------------

Analizar un sistema de control es evaluar tres cosas: la estabilidad (sin ella las otras dos no tienen sentido), la respuesta transitoria y el \textbf{error en estado estacionario}. Las dos primeras ocuparon los capítulos~\ref{cap:po}, \ref{cap:so} y~\ref{cap:es}; este capítulo trata la tercera.

En malla abierta nada corrige una perturbación ni un cambio en los parámetros del proceso. La retroalimentación sí: reduce la sensibilidad a variaciones de la planta, mejora el rechazo de perturbaciones, reduce el error en estado estacionario y facilita ajustar el transitorio. A cambio, hace llegar al lazo el ruido del sensor. La \cref{fig:er-rdn} muestra las tres señales exógenas: la referencia $R$, la perturbación $D$ y el ruido de medición $N$. Definimos el \clave{error de seguimiento} como
\begin{equation}
  E(s)=R(s)-Y(s),
  \label{eq:er-error}
\end{equation}
lo que nos interesa del desempeño. Con $\Hs=1$, por superposición,
\begin{equation}
  E=\underbrace{\frac{1}{1+\Gc\Gp}}_{\Sens(s)}R-\frac{\Gp}{1+\Gc\Gp}D+\underbrace{\frac{\Gc\Gp}{1+\Gc\Gp}}_{\Tlc(s)}N
  =\Sens R-\Sens\Gp D+\Tlc N .
  \label{eq:er-sensitividad}
\end{equation}

\begin{figure}[ht]
\centering
\includegraphics[width=0.8\linewidth]{cap06/lazo_rdn}
\caption{Lazo con referencia $R$, perturbación $D$ a la entrada del proceso y ruido de medición $N$.}
\label{fig:er-rdn}
\end{figure}

\begin{definicion}[label=def:er-sensitividad]{Sensitividad y sensitividad complementaria}
Con retroalimentación unitaria, la \textbf{sensitividad} es $\Sens(s)=\dfrac{1}{1+\Gc\Gp}$ y la \textbf{sensitividad complementaria} es $\Tlc(s)=\dfrac{\Gc\Gp}{1+\Gc\Gp}$, que es la propia función de lazo cerrado $Y/R$. Para todo $s$,
\[
  \Sens(s)+\Tlc(s)=1 .
\]
\end{definicion}

Todos los términos de \eqref{eq:er-sensitividad} tienen el mismo denominador: la ecuación característica $1+\Gc\Gp=0$ es la misma para cualquier entrada. Y la identidad $\Sens+\Tlc=1$ dice que no se pueden hacer pequeñas las dos a la misma frecuencia. El compromiso se resuelve separando frecuencias: $\Gc\Gp$ grande donde viven las perturbaciones y las referencias (bajas frecuencias), para que $\Sens$ sea pequeña; $\Gc\Gp$ pequeña donde vive el ruido (altas frecuencias), para que $\Tlc$ lo sea.

\begin{ejemplo}[label=ej:er-sens]{Dos ganancias}
Con $\Gp=1/((s+1)(s+10))$ y $\Gc=K$, compare $K=10$ y $K=200$ ante una perturbación $d=0.5\sin(0.2t)$ y un ruido $n=0.01\sin(50t)$, con referencia constante.
\solucion La amplitud de cada aporte es la amplitud de la señal por la magnitud de su función a esa frecuencia. Con $K=200$: perturbación $0.5\,|\Sens\Gp|(0.2)=0.0024$ y ruido $0.01\,|\Tlc|(50)=0.00085$. Con $K=10$: $0.025$ y $0.00004$. La ganancia alta rechaza diez veces mejor la perturbación y deja pasar veinte veces más ruido (\cref{fig:er-sens}); aquí gana, porque el ruido es pequeño. Si su amplitud pasara de $0.28$ (sumando los aportes como cota), convendría $K=10$.
\end{ejemplo}

\begin{figure}[ht]
\centering
\includegraphics[width=0.85\linewidth]{cap06/f_sensibilidad}
\caption{Ejemplo~\ref{ej:er-sens}: $|\Sens|$ y $|\Tlc|$ para $K=10$ y $K=200$. Donde una es pequeña, la otra vale casi $1$.}
\label{fig:er-sens}
\end{figure}

\codigorepo{cap06_error/C6_Ej6_Sensibilidad.m}{aportes}{ejemplo~\ref{ej:er-sens}}

La sensitividad y la robustez se estudian a fondo en cursos posteriores. Aquí nos quedamos con su consecuencia más directa: el error que queda cuando todo se asienta. En lo que sigue (salvo en la \cref{sec:er-perturbacion}) el lazo no tiene ni perturbación ni ruido.

%% ---------------------------------------------------------------------
\section{Señales de prueba}
\label{sec:er-senales}
%% ---------------------------------------------------------------------

El \clave{error en estado estacionario} es el límite del error de seguimiento cuando $t\to\infty$, ante una señal de prueba. Las tres tradicionales son (\cref{fig:er-senales}):
\begin{itemize}
  \item \textbf{Escalón}, $r(t)=\escalon(t)$, $R(s)=1/s$: mantener una posición o un valor constante (un satélite geoestacionario, una temperatura).
  \item \textbf{Rampa}, $r(t)=t\,\escalon(t)$, $R(s)=1/s^2$: seguir un blanco que se mueve a velocidad constante (una antena que sigue un satélite).
  \item \textbf{Parábola}, $r(t)=\frac{t^2}{2}\escalon(t)$, $R(s)=1/s^3$: seguir un blanco con aceleración constante.
\end{itemize}

\begin{figure}[ht]
\centering
\includegraphics[width=\linewidth]{cap06/f_senales}
\caption{Las tres señales de prueba.}
\label{fig:er-senales}
\end{figure}

\begin{cuidado}[Primero la estabilidad]
El teorema del valor final solo vale si $sE(s)$ tiene todos sus polos en el semiplano izquierdo, es decir, si el lazo cerrado es estable. Calcular un ``error en estado estacionario'' de un lazo inestable da un número que no significa nada. \alerta{Verifique la estabilidad primero} (\cref{cap:es}).
\end{cuidado}

%% ---------------------------------------------------------------------
\section{Tipo del sistema y constantes de error}
\label{sec:er-tipo}
%% ---------------------------------------------------------------------

Con retroalimentación unitaria (\cref{fig:er-unitario}), $\Tlc=G/(1+G)$ y $E(s)=\dfrac{1}{1+G(s)}R(s)$. Si el lazo cerrado es estable, el teorema del valor final da
\begin{equation}
  \boxed{\;\ess=\lim_{t\to\infty}e(t)=\lim_{s\to0}\,s\,\frac{R(s)}{1+G(s)}\;}
  \label{eq:er-ess}
\end{equation}
Aquí $G(s)$ es el \textbf{lazo abierto}: la función que, con retroalimentación unitaria, va del error a la salida (\cref{def:rf-malla}; si la retroalimentación no es unitaria, ver la \cref{sec:er-no-unitaria}).

\begin{figure}[ht]
\centering
\includegraphics[width=0.55\linewidth]{cap06/lazo_unitario}
\caption{Lazo con retroalimentación unitaria.}
\label{fig:er-unitario}
\end{figure}

Escribamos el lazo abierto en forma de Bode, separando los polos en el origen:
\begin{equation}
  G(s)=\frac{K\,(T_1s+1)(T_2s+1)\cdots(T_ms+1)}{s^N(\tau_1s+1)(\tau_2s+1)\cdots(\tau_{n-N}s+1)} .
  \label{eq:er-bode}
\end{equation}

\begin{definicion}[label=def:er-tipo]{Tipo del sistema}
El \textbf{tipo} es el número $N$ de polos en el origen del lazo abierto, es decir, el número de integradores del lazo. $K$ es su ganancia de Bode: el valor de $s^NG(s)$ en $s=0$.
\end{definicion}

Cada señal de prueba define una \textbf{constante estática de error}; cuanto mayor, menor el error:
\begin{itemize}
  \item \textbf{Escalón}, $R=1/s$: $\ess=\dfrac{1}{1+\lim_{s\to0}G(s)}=\dfrac{1}{1+\kp}$, con la \clave{constante de posición} $\kp=\lim_{s\to0}G(s)$.
  \item \textbf{Rampa}, $R=1/s^2$: $\ess=\lim_{s\to0}\dfrac{1}{s+sG(s)}=\dfrac{1}{\kv}$, con la \clave{constante de velocidad} $\kv=\lim_{s\to0}sG(s)$.
  \item \textbf{Parábola}, $R=1/s^3$: $\ess=\dfrac{1}{\ka}$, con la \clave{constante de aceleración} $\ka=\lim_{s\to0}s^2G(s)$.
\end{itemize}
Con \eqref{eq:er-bode}, $s^iG(s)\to K/s^{N-i}$: la constante vale $K$ si $N=i$, infinito si $N>i$ y cero si $N<i$. Eso da la tabla que resume el capítulo.

\begin{table}[H]
\centering
\caption{Error en estado estacionario según el tipo del lazo abierto (lazo cerrado estable, retroalimentación unitaria). Para una entrada de amplitud $A$ (escalón $A$, rampa $At$, parábola $At^2/2$), el error se multiplica por $A$.}
\label{tab:er-tabla}
\renewcommand{\arraystretch}{1.5}
\begin{tabular}{@{}lcccccc@{}}
\toprule
 & $\kp$ & $\kv$ & $\ka$ & escalón & rampa & parábola\\
\midrule
Tipo 0 & $K$ & $0$ & $0$ & $\dfrac{1}{1+\kp}$ & $\infty$ & $\infty$\\
Tipo 1 & $\infty$ & $K$ & $0$ & $0$ & $\dfrac{1}{\kv}$ & $\infty$\\
Tipo 2 & $\infty$ & $\infty$ & $K$ & $0$ & $0$ & $\dfrac{1}{\ka}$\\
\bottomrule
\end{tabular}
\end{table}

\begin{ideaclave}
\textbf{Cada integrador vale una derivada de la referencia.}\enspace Un lazo de tipo $N$ sigue sin error una referencia polinomial de grado menor que $N$, con error finito una de grado $N$, y pierde una de grado mayor. El error se lee en la tabla con solo dos datos: el tipo y una ganancia.
\end{ideaclave}

\begin{figure}[ht]
\centering
\includegraphics[width=\linewidth]{cap06/f_tipos}
\caption{Lazos abiertos de tipo 0, 1 y 2 ($\frac{4}{(s+1)(s+2)}$, $\frac{4}{s(s+2)}$ y $\frac{4(s+0.5)}{s^2(s+2)}$) ante escalón y rampa unitarios. Errores: $1/3$ y $\infty$; $0$ y $0.5$; $0$ y $0$.}
\label{fig:er-tipos}
\end{figure}

\begin{ejemplo}[label=ej:er-motor]{Un motor con retroalimentación de velocidad}
En el servo de la \cref{fig:er-motor}~\citep{DoBi08}, halle $\ess$ ante una rampa unitaria en función de $K$, $K_b$ y $K_m$, y escoja $K$ para $\ess=1$ con $K_m=10$ y $K_b=0.05$.
\solucion El lazo interno da $\Omega/(KE)=K_m/(s+0.01+K_bK_m)$, y con el integrador,
\[
  G(s)=\frac{KK_m}{s(s+0.01+K_bK_m)} :
\]
tipo 1, con $\kv=\dfrac{KK_m}{0.01+K_bK_m}$ y $\ess=\dfrac{0.01+K_bK_m}{KK_m}$. Con los datos, $\ess=0.51/(10K)=1$ da $K=0.051$. (El lazo cerrado es de segundo orden con coeficientes positivos: estable para todo $K>0$.)
\end{ejemplo}

\begin{figure}[ht]
\centering
\includegraphics[width=0.85\linewidth]{cap06/lazo_motor_kb}
\caption{Ejemplo~\ref{ej:er-motor}: servo de posición con retroalimentación interna de velocidad.}
\label{fig:er-motor}
\end{figure}

\codigorepo{cap06_error/C6_Ej1_TipoError.m}{tipo}{tipo, constantes y estabilidad de cuatro lazos con la función \cmd{constantesError}}
\codigorepo{cap06_error/C6_Ej1_TipoError.m}{simulacion}{el error de rampa simulado}

\begin{pausa}
Un lazo abierto $G=\dfrac{5(s+1)}{s^3(s+2)}$ tiene $\kp=\kv=\ka=\infty$. ¿Sigue sin error una parábola?
\respuesta{No se puede saber con la tabla: el polinomio característico es $s^4+2s^3+5s+5$, al que le falta el término en $s^2$, así que el lazo cerrado es inestable. La salida no sigue nada.}
\end{pausa}

%% ---------------------------------------------------------------------
\section{Retroalimentación no unitaria}
\label{sec:er-no-unitaria}
%% ---------------------------------------------------------------------

Si el sensor tiene dinámica o ganancia, $\Hs(s)\neq1$, y al comparador ya no llega el error de seguimiento, sino el \clave{error de actuación}
\[
  E_a(s)=R(s)-\Hs(s)Y(s),
\]
la señal que entra a $G$. Aplicar la \cref{tab:er-tabla} a la malla abierta $\Gma=G\Hs$ da el error de \emph{actuación} en estado estacionario, no el de seguimiento. Para el de seguimiento, el lazo se reescribe con retroalimentación unitaria (ejemplo~\ref{ej:so-equivalente}) con el \textbf{lazo abierto}
\begin{equation}
  \Gla(s)=\frac{G(s)}{1+G(s)\Hs(s)-G(s)},
  \label{eq:er-gla}
\end{equation}
y con él se usan el tipo, las constantes y la tabla. Si $\Hs=1$, $\Gla=G$.

\begin{ejemplo}[label=ej:er-gh]{Dos errores que no coinciden}
En la \cref{fig:er-gh}, $G=\dfrac{100}{s+1}$ y $\Hs=\dfrac{1}{s+5}$. Halle el error de actuación y el de seguimiento ante un escalón unitario.
\solucion El lazo cerrado es $\Tlc=\dfrac{G}{1+G\Hs}=\dfrac{100(s+5)}{s^2+6s+105}$, estable (polos en $-3\pm9.8j$), con $y_\infty=\Tlc(0)=500/105=4.76$.

\emph{Con la malla abierta}: $G\Hs=\dfrac{100}{(s+1)(s+5)}$, tipo 0, $\kp=20$, y $e_{a,\infty}=1/21=0.048$. Es correcto, pero es el error de actuación: en estado estacionario $\Hs(0)y_\infty=4.76/5=0.952\approx r$.

\emph{Con el lazo abierto}: \eqref{eq:er-gla} da
\[
  \Gla(s)=\frac{100(s+5)}{s^2-94s-395},\qquad \kp=\Gla(0)=-\frac{500}{395}=-1.266,
\]
y $\ess=\dfrac{1}{1+\kp}=\dfrac{1}{1-1.266}=-3.76=1-y_\infty$. La salida termina en $4.76$ veces la referencia (\cref{fig:er-no-unitaria}). El signo negativo dice que la salida \emph{sobrepasa} la referencia.

$\Gla$ tiene un polo en $+98$: no es un sistema físico que haya que estabilizar, sino una función auxiliar que reproduce el mismo lazo cerrado; el lazo cerrado sigue siendo estable. La causa del error es la ganancia DC del sensor, $\Hs(0)=1/5$: el lazo lleva $\Hs y$ a la referencia, no $y$. Si el sensor se recalibra como $\Hs=k/(s+5)$, el error de seguimiento se anula con $\Tlc(0)=500/(5+100k)=1$, es decir, $k=4.95$.
\end{ejemplo}

\begin{figure}[ht]
\centering
\includegraphics[width=0.55\linewidth]{cap06/lazo_gh}
\caption{Ejemplo~\ref{ej:er-gh}: lazo con un sensor dinámico.}
\label{fig:er-gh}
\end{figure}

\begin{figure}[ht]
\centering
\includegraphics[width=0.85\linewidth]{cap06/f_no_unitaria}
\caption{Ejemplo~\ref{ej:er-gh}: el error de actuación tiende a $1/21$ y el de seguimiento a $-3.76$.}
\label{fig:er-no-unitaria}
\end{figure}

\codigorepo{cap06_error/C6_Ej2_NoUnitaria.m}{lazoabierto}{malla abierta, lazo abierto y los dos errores}
\codigorepo{cap06_error/C6_Ej2_NoUnitaria.m}{simulacion}{la simulación y el sensor recalibrado}

\begin{cuidado}[Malla abierta no es lazo abierto]
La estabilidad y el diagrama de Bode se estudian con la malla abierta $G\Hs$; el error de seguimiento, con el lazo abierto $\Gla$. Coinciden solo si $\Hs=1$. Si $\Hs(0)=1$ (un sensor bien calibrado), el error ante escalón también coincide, aunque las funciones sean distintas.
\end{cuidado}

%% ---------------------------------------------------------------------
\section{Error debido a perturbaciones}
\label{sec:er-perturbacion}
%% ---------------------------------------------------------------------

Con $R=0$ y $N=0$, \eqref{eq:er-sensitividad} da el error que deja una perturbación a la entrada del proceso:
\begin{equation}
  E(s)=-\frac{\Gp(s)}{1+\Gc(s)\Gp(s)}\,D(s).
  \label{eq:er-perturbacion}
\end{equation}
Si la perturbación entra a la salida del proceso, $E=-\dfrac{1}{1+\Gc\Gp}D$. En los dos casos, el error ante una perturbación escalón se anula si $\Gc$ tiene un integrador: \alerta{importa dónde está el integrador}. Uno en el controlador, \emph{antes} del punto donde entra la perturbación, la rechaza; uno en la planta, después de ese punto, no.

\begin{ejemplo}[label=ej:er-perturbacion]{Proporcional contra PI}
Con $\Gp=\dfrac{1}{(s+1)(s+2)}$ y una perturbación escalón de $0.5$, halle el error con $\Gc=K$ y con $\Gc=K+K_i/s$.
\solucion \emph{Proporcional, perturbación a la entrada}: $\ess=-0.5\,\dfrac{\Gp(0)}{1+K\Gp(0)}=-\dfrac{0.5}{2+K}$: $-0.167$, $-0.042$ y $-0.0049$ para $K=1$, $10$ y $100$. Baja con $K$ pero no se anula. \emph{A la salida}: $\ess=-\dfrac{0.5}{1+K/2}=-\dfrac{1}{2+K}$, el doble, porque la perturbación no pasa por $\Gp(0)=1/2$.

\emph{PI}: $\ess=-0.5\lim_{s\to0}\dfrac{\Gp}{1+(K+K_i/s)\Gp}=-0.5\lim_{s\to0}\dfrac{s\Gp}{s+(Ks+K_i)\Gp}=0$ para cualquier $K_i\neq0$ que deje el lazo estable. Con $K=10$, la ecuación característica es $s^3+3s^2+12s+K_i=0$ y Routh exige $0<K_i<36$ (\cref{fig:er-perturbacion}).
\end{ejemplo}

\begin{figure}[ht]
\centering
\includegraphics[width=\linewidth]{cap06/f_perturbacion}
\caption{Ejemplo~\ref{ej:er-perturbacion}: error por una perturbación escalón de $0.5$.}
\label{fig:er-perturbacion}
\end{figure}

\codigorepo{cap06_error/C6_Ej3_Perturbacion.m}{proporcional}{el controlador proporcional, perturbación a la entrada y a la salida}
\codigorepo{cap06_error/C6_Ej3_Perturbacion.m}{pi}{el controlador PI}

\begin{modelo}[Modelo de Simulink: \texttt{crearModelo\_Perturbacion}]
\texttt{codigo/cap06\_error/crearModelo\_Perturbacion.m} arma el lazo con un PI y una perturbación escalón de $0.5$ en $t=5$\,s. La variable \texttt{alEntrada} decide si entra a la entrada o a la salida del proceso; con \texttt{Ki = 0} el controlador es proporcional.
\end{modelo}

%% ---------------------------------------------------------------------
\section{Error en espacio de estado}
\label{sec:er-estado}
%% ---------------------------------------------------------------------

Si el proceso está en espacio de estado (\cref{cap:ee}), el error en estado estacionario se obtiene sin pasar por la función de transferencia: con el lazo cerrado y una referencia constante, el estado estacionario cumple $\dot x=0$.

\begin{ejemplo}[label=ej:er-estado]{Estado estacionario directo}
Un proceso con $\Am=\begin{bmatrix}0&1\\-2&-3\end{bmatrix}$, $\Bm=[0\;\,1]^\top$, $\Cm=[2\;\,0]$ se controla con $u=r-y$. Halle $\ess$ ante un escalón unitario. Después agregue un estado integrador $\dot x_3=r-y$ con $u=r-y+K_ix_3$.
\solucion Con la función de transferencia: $G=\dfrac{2}{s^2+3s+2}$, tipo 0, $\kp=1$, $\ess=0.5$. En espacio de estado: $\dot x=(\Am-\Bm\Cm)x+\Bm r$ y, con $\dot x=0$,
\[
  x_{ss}=-(\Am-\Bm\Cm)^{-1}\Bm\,r=\begin{bmatrix}0.25\\0\end{bmatrix},\qquad y_{ss}=\Cm x_{ss}=0.5,\qquad \ess=0.5 .
\]
Con el estado integrador, el sistema aumentado es
\[
  \begin{bmatrix}\dot x\\\dot x_3\end{bmatrix}=
  \begin{bmatrix}\Am-\Bm\Cm&K_i\Bm\\-\Cm&0\end{bmatrix}\begin{bmatrix}x\\x_3\end{bmatrix}+\begin{bmatrix}\Bm\\1\end{bmatrix}r,
\]
con polinomio característico $s^3+3s^2+4s+2K_i$: estable para $0<K_i<6$. En estado estacionario $\dot x_3=r-y=0$: el error se anula para todo $K_i$ que deje el lazo estable. Con $K_i=3$, los polos son $-2.38$ y $-0.31\pm1.56j$.
\end{ejemplo}

\codigorepo{cap06_error/C6_Ej4_Estado.m}{estacionario}{el estado estacionario con $\dot x=0$}
\codigorepo{cap06_error/C6_Ej4_Estado.m}{integral}{el estado integrador}

El estado integrador convierte el lazo en uno de tipo 1. Es la misma idea de la acción integral de un PI, y la que se usa con la retroalimentación de estado en la \cref{sec:re-integral}.

%% ---------------------------------------------------------------------
\section{Especificaciones de error y de transitorio}
\label{sec:er-diseno}
%% ---------------------------------------------------------------------

Una especificación de error se traduce en una cota sobre una constante: ``$\ess\le0.05$ ante rampa unitaria'' es $\kv\ge20$. Con un controlador proporcional, la única perilla es $K$, y la misma perilla mueve los polos del lazo cerrado. Las dos especificaciones pueden chocar.

\begin{ejemplo}[label=ej:er-conflicto]{Un servo que no puede con las dos}
Con $G_p=\dfrac{1}{s(s+4)(s+8)}$ y $\Gc=K$, se pide $\ess\le0.05$ ante rampa unitaria y $\OS\le20$\,\% ante escalón.
\solucion Tipo 1 con $\kv=K/32$: el error pide $K\ge640$. Pero la ecuación característica $s^3+12s^2+32s+K=0$ exige $K<12\cdot32=384$: \alerta{con $K$ solo, ni siquiera es estable}. Y el sobrepico pide un $K$ mucho menor: para $\zeta=0.49$ en los polos dominantes ($17$\,\% estimado, con margen; la simulación da $16$\,\%), $K=67.8$, que da $\kv=2.1$ y $\ess=0.47$ (\cref{fig:er-atraso}, azul).

Un \textbf{compensador de atraso} $\Gc=K\dfrac{s+a}{s+b}$ con $a>b>0$ multiplica $\kv$ por $a/b$. Si el par polo-cero está muy cerca del origen, en la zona de los polos dominantes $(s+a)/(s+b)\approx1$, y el transitorio casi no cambia. Se escoge $K=67.8$ por el transitorio, y $a/b=1.1\cdot640/67.8=10.4$ (con un $10$\,\% de margen); con $b=0.005$, $a=0.0519$. El lazo compensado tiene $\kv=22$ y $\ess=0.045$, con $\OS=19$\,\%; el precio es una cola lenta (un polo en $-0.053$ casi cancelado por el cero en $-0.052$): el tiempo de establecimiento sube de $3.1$ a $4.8$\,s, y el error de rampa tarda unos $116$\,s en quedar a menos del $2$\,\% de su valor final.
\end{ejemplo}

\begin{figure}[ht]
\centering
\includegraphics[width=\linewidth]{cap06/f_atraso}
\caption{Ejemplo~\ref{ej:er-conflicto}: con solo $K=67.8$ (azul) y con el atraso (coral). El atraso deja casi igual la respuesta al escalón y baja el error de rampa de $0.47$ a $0.045$, con una cola lenta.}
\label{fig:er-atraso}
\end{figure}

\codigorepo{cap06_error/C6_Ej5_Atraso.m}{conflicto}{el conflicto entre las dos especificaciones}
\codigorepo{cap06_error/C6_Ej5_Atraso.m}{atraso}{el compensador de atraso}

El diseño de atrasos en el dominio de la frecuencia es tema de la \cref{sec:rf-atraso}; el PI, que lleva $b$ hasta cero, del \cref{cap:pn}.

%% ---------------------------------------------------------------------
\section{Para llevar}
\label{sec:er-cierre}
%% ---------------------------------------------------------------------

\begin{parallevar}
\begin{enumerate}
  \item $E=\Sens R-\Sens\Gp D+\Tlc N$ con $\Sens+\Tlc=1$: lazo abierto grande donde están la referencia y las perturbaciones, pequeño donde está el ruido.
  \item Primero la estabilidad. Después, $\ess=\lim_{s\to0}sR(s)/(1+G(s))$.
  \item El tipo es el número de integradores del lazo abierto. $\kp=\lim G$, $\kv=\lim sG$, $\ka=\lim s^2G$; escalón $1/(1+\kp)$, rampa $1/\kv$, parábola $1/\ka$.
  \item Con retroalimentación no unitaria, el error de seguimiento se calcula con el lazo abierto $\Gla=G/(1+G\Hs-G)$; la malla abierta $G\Hs$ da el error de actuación.
  \item Una perturbación escalón se rechaza por completo si hay un integrador \emph{antes} del punto donde entra.
  \item En espacio de estado, el estado estacionario sale de $\dot x=0$; un estado integrador del error vuelve el lazo de tipo 1.
  \item Error y transitorio pueden chocar con un solo parámetro; un atraso sube $\kv$ por el factor $a/b$ a costa de una cola lenta.
\end{enumerate}
\end{parallevar}

\begin{profundiza}[Correcciones respecto a la versión anterior del capítulo]
(i) En las secciones de $\kv$ y $\ka$ se escribía $\kv=\lim_{s\to0}G(s)$ y $\ka=\lim_{s\to0}G(s)$ antes del desarrollo; son $\lim sG$ y $\lim s^2G$. (ii) La sensitividad complementaria se llamaba $C(s)$; se usa $\Tlc(s)$, que es además la función de lazo cerrado, para no confundirla con $\Cm$. (iii) El error con retroalimentación no unitaria se trata ahora con el lazo abierto $\Gla$ del \cref{cap:so} y la distinción entre error de actuación y de seguimiento, que la versión anterior dejaba para un ejercicio. (iv) La tabla resumen tenía una columna vacía de más y no incluía las constantes; se completó, y se dejó explícito que solo vale con el lazo cerrado estable. (v) Los ejemplos físicos de las constantes (``la temperatura de salida'') se retiraron por imprecisos, y el código de Mathematica de las señales de prueba se reemplazó por una figura generada con el resto. (vi) Se resolvió el ejemplo del lazo con sensor dinámico, que estaba solo enunciado, y se agregaron el error por perturbaciones, el error en espacio de estado y el conflicto entre especificaciones.
\end{profundiza}

%% ---------------------------------------------------------------------
\section{Problemas}
\label{sec:er-problemas}
%% ---------------------------------------------------------------------

Los problemas marcados con \etiquetaM{} requieren \matlab/\simulink; los marcados con \etiquetaEstrella{} son de profundización.

\begin{problema}[tipo y constantes]\label{prob:er-tipos}
Con retroalimentación unitaria y lazo abierto (i) $\dfrac{20(s+2)}{(s+1)(s+5)}$, (ii) $\dfrac{50}{s(s+3)(s+10)}$, (iii) $\dfrac{100(s+4)(s+6)}{s^2(s+5)(s+8)}$ y (iv) $\dfrac{5(s+1)}{s^3(s+2)}$: (a) verifique la estabilidad con Routh; (b) halle el tipo y $\kp$, $\kv$, $\ka$; (c) calcule $\ess$ ante $r=5$, $r=2t$ y $r=\frac32t^2$ y tabule; (d) uno de los cuatro no admite el análisis: ¿cuál y por qué? (e) \etiquetaM{} verifique por simulación los casos admisibles.
\end{problema}

\begin{problema}[error contra estabilidad]\label{prob:er-posicion}
Un servo tiene lazo abierto $\dfrac{K}{s(s+4)(s+8)}$ con retroalimentación unitaria. (a) Halle el tipo y $\kv(K)$. (b) Halle el menor $K$ con $\ess\le0.1$ ante rampa unitaria. (c) Halle el rango de $K$ estable. (d) ¿Hay valores que cumplan las dos? (e) Para el $K$ de (b), halle los polos, $\zeta$ y $\omega_n$ del par dominante, y estime $\OS$ y $\ts$; \etiquetaM{} verifique con \cmd{stepinfo}. (f) ¿Por qué no es aceptable reducir el error solo con la ganancia?
\end{problema}

\begin{problema}[perturbaciones]\label{prob:er-perturbacion}
En la \cref{fig:er-rdn}, con $\Hs=1$, $N=R=0$, $\Gp=\dfrac{1}{(s+1)(s+2)}$ y una perturbación escalón de $0.5$ a la entrada del proceso: (a) halle $\ess$ con $\Gc=K$ para $K=1$, $10$ y $100$; (b) muestre que con $\Gc=K+K_i/s$ el error es nulo para todo $K_i$ estable, y diga qué propiedad de $\Gc$ lo produce; (c) halle el rango de $K_i$ estable con $K=10$; (d) repita (a) con la perturbación a la salida y explique la diferencia; (e) \etiquetaM{} verifique con \texttt{crearModelo\_Perturbacion}.
\end{problema}

\begin{problema}[sensor dinámico]\label{prob:er-gh}
Con $G=\dfrac{100}{s+1}$ y $\Hs=\dfrac{1}{s+5}$: (a) verifique la estabilidad; (b) halle $\kp$ de la malla abierta y el error de actuación ante escalón; (c) halle $y_\infty$ y el error de seguimiento; (d) explique por qué difieren y cuándo coinciden; (e) \etiquetaM{} simule $r$, $y$, $e_a$ y $e$; ¿cuál error da la \cref{tab:er-tabla} con la malla abierta?
\end{problema}

\begin{problema}[lazo abierto equivalente]\label{prob:er-gla}
Con los datos del problema~\ref{prob:er-gh}: (a) muestre que el lazo con $G$ y $\Hs$ tiene el mismo $\Tlc$ que uno unitario con $\Gla$ de \eqref{eq:er-gla}; (b) calcule $\Gla$ y sus polos; (c) explique cómo puede tener un polo en el semiplano derecho si el lazo cerrado es estable; (d) calcule $\kp$ de $\Gla$ y $\ess=1/(1+\kp)$, y compárelo con (c) del problema anterior; (e) interprete el signo; (f) con $\Hs=k/(s+5)$, halle el $k$ que anula el error y discuta la importancia de calibrar el sensor.
\end{problema}

\begin{problema}[espacio de estado]\label{prob:er-estado}
Con $\Am$, $\Bm$, $\Cm$ del ejemplo~\ref{ej:er-estado} y $u=r-y$: (a) halle $G$, el tipo, $\kp$ y $\ess$; (b) obtenga lo mismo con $\dot x=0$; (c) escriba el sistema aumentado con $\dot x_3=r-y$ y $u=r-y+K_ix_3$; (d) halle el rango de $K_i$ estable; (e) muestre que el error es nulo y relaciónelo con el tipo; (f) \etiquetaM{} simule con $K_i$ en la mitad del rango.
\end{problema}

\begin{problema}[atraso]\label{prob:er-atraso}
Con $G_p=\dfrac{1}{s(s+4)(s+8)}$ se pide $\ess\le0.05$ ante rampa unitaria y $\OS\le20$\,\%. (a) Muestre que ningún $\Gc=K$ cumple las dos y cuantifique el conflicto. (b) Muestre que $\Gc=K\frac{s+a}{s+b}$ multiplica $\kv$ por $a/b$ y casi no mueve los polos dominantes si $a$ y $b$ están cerca del origen. (c) Escoja $K$ por el transitorio y $(a,b)$ por el error. (d) Verifique la estabilidad con Routh sobre el polinomio de cuarto orden. (e) \etiquetaM{} Simule escalón y rampa, verifique las especificaciones y cuantifique la cola lenta. (f) ¿Qué precio se paga y cuándo sería inaceptable?
\end{problema}

\begin{problema}[sensitividad]\label{prob:er-sens}
Con $\Gp=\dfrac{1}{(s+1)(s+10)}$ y $\Gc=K$: (a) muestre que $\Sens+\Tlc=1$; (b) \etiquetaM{} tabule $|\Sens|$ y $|\Tlc|$ en $\omega=0.1$, $1$, $10$ y $100$\,rad/s para $K=10$ y $K=200$; (c) halle la frecuencia donde $|\Sens|=|\Tlc|$ y relaciónela con el cruce de ganancia; (d) con $d=0.5\sin(0.2t)$ a la entrada y $n=0.01\sin(50t)$, halle la amplitud de cada aporte al error; (e) ¿cuál $K$ da menor error y por qué? (f) \etiquetaEstrella{} ¿Desde qué amplitud de ruido conviene $K=10$?
\end{problema}

%% ---------------------------------------------------------------------
\section{Soluciones seleccionadas}
%% ---------------------------------------------------------------------

\paragraph{Problema~\ref{prob:er-tipos}}
(a) Polinomios $s^2+26s+45$, $s^3+13s^2+30s+50$, $s^4+13s^3+140s^2+1000s+2400$ (estables) y $s^4+2s^3+5s+5$ (inestable: falta $s^2$). (b) (i) tipo 0, $\kp=8$; (ii) tipo 1, $\kv=5/3$; (iii) tipo 2, $\ka=60$. (c) (i) $5/9$, $\infty$, $\infty$; (ii) $0$, $1.2$, $\infty$; (iii) $0$, $0$, $0.05$. (d) El (iv): la tabla daría error nulo para todo, pero la salida diverge.

\paragraph{Problema~\ref{prob:er-posicion}}
(a) Tipo 1, $\kv=K/32$. (b) $K=320$. (c) $0<K<384$. (d) $320\le K<384$. (e) Polos $-11.6$ y $-0.19\pm5.25j$: $\zeta=0.037$, $\omega_n=5.25$, $\OS\approx89$\,\% estimado ($81$\,\% simulado) y $\ts\approx21$\,s. (f) El lazo queda casi en el límite de estabilidad.

\paragraph{Problema~\ref{prob:er-perturbacion}}
(a) $-0.167$, $-0.042$, $-0.0049$. (b) El integrador de $\Gc$, que está antes de la perturbación. (c) $0<K_i<36$. (d) $-0.333$, $-0.083$, $-0.0098$: a la salida, la perturbación no se atenúa por $\Gp(0)=1/2$.

\paragraph{Problema~\ref{prob:er-gh}}
(a) Polos $-3\pm9.8j$. (b) $\kp=20$, $e_a=1/21$. (c) $y_\infty=100/21=4.76$, $\ess=-3.76$. (d) Coinciden si $\Hs(0)=1$.

\paragraph{Problema~\ref{prob:er-gla}}
(b) $\Gla=\dfrac{100(s+5)}{s^2-94s-395}$, polos $98.0$ y $-4.03$. (d) $\kp=-1.266$ y $\ess=-3.76$. (f) $k=4.95$.

\paragraph{Problema~\ref{prob:er-estado}}
(a) $G=\frac{2}{s^2+3s+2}$, tipo 0, $\kp=1$, $\ess=0.5$. (d) $s^3+3s^2+4s+2K_i$: $0<K_i<6$. (e) $\dot x_3=r-y=0$ en estado estacionario; el lazo es de tipo 1.

\paragraph{Problema~\ref{prob:er-atraso}}
(a) $\ess\le0.05$ pide $K\ge640$; la estabilidad, $K<384$; $\OS\le20$\,\%, $K\lesssim74$. (c) Con $K=67.8$, $b=0.005$ y $a=0.0519$: $\kv=22$. (e) $\OS=19$\,\%, $\ess=0.045$; $\ts$ pasa de $3.1$ a $4.8$\,s y el error de rampa tarda unos $116$\,s en asentarse.

\paragraph{Problema~\ref{prob:er-sens}}
(b) $K=10$: $|\Sens|=0.50$, $0.65$, $1.04$, $1.00$ y $|\Tlc|=0.50$, $0.46$, $0.07$, $0.001$; $K=200$: $|\Sens|=0.048$, $0.068$, $0.91$, $1.02$ y $|\Tlc|=0.95$, $0.96$, $1.29$, $0.020$. (c) Con $K=200$, en $12.5$\,rad/s, el cruce de $|\Gc\Gp|=1$; con $K=10$, $|\Gc\Gp|\le1$ en todas las frecuencias y las dos solo se igualan en $\omega\to0$. (d) $K=10$: $0.025$ y $0.00004$; $K=200$: $0.0024$ y $0.00085$. (f) Ruido de amplitud mayor que $0.28$.

\respuestaspausas

                          % migrado (Handout5_Ctrl.tex)
%% =====================================================================
%%  Capítulo 7 -- Controladores PID: nociones básicas
%%  Migrado al formato nuevo (noviembre de 2026) desde Handout6a_Ctrl.tex.
%%  Correcciones de fondo (ver MIGRACION.md): conversión ideal -> serie
%%  (tau_d^S usa tau_i^I, no tau_d^I); expansión de la forma serie; la forma
%%  paralelo no es "la única que permite quitar la integral"; historia
%%  (Taylor: pre-act fijo en 1936, ajustable en el Fulscope de 1939); se
%%  completaron las figuras y simulaciones marcadas "AÑADIR FIGURA";
%%  notación tau_i, tau_d -> T_i, T_d (convenciones del curso); MP -> %OS.
%% =====================================================================

\capitulo[cap:pn]{Controladores PID: nociones básicas}
  {Nunca, ni en las situaciones más desesperadas, he perdido mi sentido del humor, y soy de los que piensan que esta vida es de risa.}
  {Enrique Vila-Matas}

\begin{objetivos}[partes/07-pid-nociones.html]
  \item Explicar qué hace cada una de las tres acciones (proporcional, integral y derivativa) y cómo afecta la respuesta transitoria, el error y la estabilidad.
  \item Pasar los parámetros de un PID entre las formas ideal, paralelo y serie, y saber cuándo la conversión es posible.
  \item Reconocer los problemas prácticos del PID ---\emph{kicks}, ruido, \emph{wind-up}--- y las modificaciones que los resuelven: filtro derivativo, estructuras I-P e I-PD, anti-windup y ponderación de la referencia.
  \item Escribir el algoritmo digital del PID en forma posicional e incremental.
\end{objetivos}

%% ---------------------------------------------------------------------
\section{El controlador de tres términos}
\label{sec:pn-intro}
%% ---------------------------------------------------------------------

Hasta aquí hemos modelado, analizado la respuesta transitoria, la estabilidad y el error. Ahora empezamos a \emph{diseñar}, y lo hacemos con el controlador más usado en la industria: el de tres términos, \clave{proporcional, integral y derivativo} (PID). Más del noventa por ciento de los lazos de control de procesos son PI o PID~\citep{ACL05}. Los primeros fueron neumáticos, luego electrónicos y, desde los años setenta, digitales; la idea, sin embargo, es la misma desde hace un siglo.

\begin{historia}[De los timoneles a la Taylor]
En 1910 Elmer Sperry fundó la Sperry Gyroscope Company; su piloto automático para barcos, que acoplaba un girocompás al timón, hacía en la práctica lo que hacían los timoneles expertos, y resultó ser un PID mecánico~\citep{AsKu14}. En 1922 Nicolas Minorsky, que había observado a esos timoneles, propuso por escrito la ley de tres términos y analizó la estabilidad del lazo con un modelo lineal (capítulo~1)~\citep{Min22,Ben01}. En la industria de procesos el PID llegó con los controladores neumáticos. La Foxboro vendía desde 1931 el \emph{Stabilog}, con acción proporcional e integral (\emph{reset}), y la Taylor agregó hacia 1936 una acción derivativa fija (\emph{pre-act}) a su controlador. En 1939 aparecieron los dos primeros PID comerciales con las tres acciones ajustables: el \emph{Fulscope} rediseñado de la Taylor y el \emph{Stabilog} con \emph{hyper-reset} de la Foxboro~\citep{Ben84,Ben93}. Entre 1925 y 1935 se vendieron en Estados Unidos unos $75\,000$ controladores, y había unas $600$ empresas de instrumentos. Quedaba un problema: cómo escoger los parámetros. Lo resolvieron en 1942 John Ziegler y Nathaniel Nichols, de la Taylor~\citep{ZiNi42}, y en 1953 Cohen y Coon, también de la Taylor~\citep{CoCo53}; Geraldine Coon comprobó sus reglas en el computador analógico de la empresa. En 1959 se cerró el primer lazo industrial completamente digital, en la unidad de polimerización de la refinería de Texaco en Port Arthur (Texas).
\end{historia}

\begin{figure}[ht]
\centering
\includegraphics[width=0.9\linewidth]{cap07/lazo_pid}
\caption{Lazo de control con filtro de la referencia $G_f$, perturbación de carga $D$ y ruido de medida $N$.}
\label{fig:pn-lazo}
\end{figure}

En el lazo de la \cref{fig:pn-lazo}, $R$ es la referencia (\emph{set point}, SP), $Y$ la variable del proceso (PV), $U_c$ la salida del controlador (OC), $D$ una perturbación de carga y $N$ el ruido de medida. Mientras no digamos otra cosa, $G_f=1$ y $D=N=0$, y el error es $E=R-Y$. El \clave{PID ideal} es
\begin{equation}
  \frac{U_c(s)}{E(s)}=\Kc\left(1+\frac{1}{\Ti s}+\Td s\right),
  \qquad
  u_c(t)=\Kc\left(e(t)+\frac{1}{\Ti}\int_0^t e(\tau)\,\dif\tau+\Td\,\deriv{e}{t}\right),
  \label{eq:pn-ideal}
\end{equation}
con $\Kc$ la ganancia proporcional, $\Ti$ el \clave{tiempo integral} y $\Td$ el \clave{tiempo derivativo}.

\begin{ideaclave}
\textbf{Pasado, presente y futuro.}\enspace La acción proporcional mira el error \emph{ahora}; la integral, el área acumulada, es decir, el \emph{pasado}; la derivativa, la tendencia, es decir, hacia dónde va el error en el \emph{futuro} cercano (\cref{fig:pn-tiempos}).
\end{ideaclave}

\begin{figure}[ht]
\centering
\includegraphics[width=0.8\linewidth]{cap07/f_tiempos}
\caption{Las tres acciones del PID sobre la misma señal de error.}
\label{fig:pn-tiempos}
\end{figure}

%% ---------------------------------------------------------------------
\section{El controlador más sencillo: on/off}
\label{sec:pn-onoff}
%% ---------------------------------------------------------------------

Para entender el PID conviene empezar por el controlador más simple, el \clave{on/off} (o de dos posiciones):
\[
  u_c(t)=\begin{cases} u_{\max}, & e(t)>0,\\ u_{\min}, & e(t)<0. \end{cases}
\]
Solo mira el signo del error. Supone que el proceso tiene ganancia positiva: más salida del controlador, más PV.

\begin{ejemplo}[label=ej:pn-onoff]{Un horno on/off}
Un proceso $e^{-2s}/(10s+1)$ se controla con $r=50\%$, $u_{\max}=90\%$ y $u_{\min}=10\%$ (\cref{fig:pn-onoff}).
\solucion La PV oscila sin asentarse alrededor de la referencia, entre unos $42$ y $58\%$, y la salida del controlador salta entre sus dos valores cada pocos segundos. La oscilación la sostiene el tiempo muerto: cuando el controlador cambia, el proceso tarda $2$\,s en enterarse.
\end{ejemplo}

\begin{figure}[ht]
\centering
\includegraphics[width=0.85\linewidth]{cap07/f_onoff}
\caption{Ejemplo~\ref{ej:pn-onoff}. Arriba, referencia y PV; abajo, salida del controlador y error.}
\label{fig:pn-onoff}
\end{figure}

El on/off es barato, no hay nada que sintonizar y basta un sensor que distinga ``por encima'' de ``por debajo''. Por eso está en los termostatos de hornos, neveras y calentadores. Sus defectos son la oscilación permanente de la PV y los cambios bruscos de la salida, que desgastan el elemento final de control (una válvula, una bomba, un relé). En la práctica se le agrega una histéresis para que no conmute con el ruido.

\codigorepo{cap07_pid_nociones/C7_Ej1_OnOff.m}{onoff}{ejemplo~\ref{ej:pn-onoff}: el control on/off y su oscilación}

%% ---------------------------------------------------------------------
\section{Las tres acciones, una por una}
\label{sec:pn-acciones}
%% ---------------------------------------------------------------------

\subsection{Acción proporcional}

El on/off funciona razonablemente cuando el error es grande: si la PV está lejos, tiene sentido abrir la válvula del todo. El problema aparece cerca de la referencia. La solución natural es que la acción sea \emph{proporcional} al error dentro de una banda:
\[
  u_c(t)=\begin{cases} u_{\max}, & e(t)>e_0,\\ u_0+\Kc\,e(t), & -e_0\le e(t)\le e_0,\\ u_{\min}, & e(t)<-e_0. \end{cases}
\]
$u_0$ es la salida cuando el error es cero (el \emph{bias} o \emph{reset} manual) y $\Kc$ la ganancia proporcional (\cref{fig:pn-banda}). Dentro de la banda, $\Gc(s)=\Kc$.

\begin{figure}[ht]
\centering
\includegraphics[width=0.85\linewidth]{cap07/f_banda}
\caption{Característica estática del on/off y del proporcional con banda.}
\label{fig:pn-banda}
\end{figure}

\paragraph{El error en estado estacionario.} Despejando, $e(t)=\bigl(u_c(t)-u_0\bigr)/\Kc$: el error solo es cero si, en estado estacionario, la salida del controlador coincide con el \emph{bias}. Para un proceso de tipo~0 eso ocurre para una sola referencia. Con $\Gp=\Kproc/(\tau s+1)$ y $u_0=0$, el error ante escalón es $\ess=1/(1+\Kc\Kproc)$ (capítulo~6): baja al subir $\Kc$, pero no se anula, y una ganancia muy alta puede desestabilizar el lazo (capítulo~5).

\begin{cuidado}[Ajustar el \emph{bias} a mano no es robusto]
Con el proceso conocido, el \emph{bias} que anula el error es $u_0=r/\Gp(0)$. Pero basta que la ganancia del proceso cambie un poco (ensuciamiento, otro punto de operación) para que el error reaparezca. Lo que hace falta es un mecanismo que ajuste $u_0$ \emph{solo}: esa es la acción integral.
\end{cuidado}

\paragraph{Banda proporcional.} Muchos controladores industriales no muestran $\Kc$ sino la \clave{banda proporcional}: el porcentaje del rango del error que lleva la salida de un extremo al otro. Con $e_R$ y $u_R$ los rangos del error y de la salida,
\[
  PB\%=\frac{1}{\Kc}\,\frac{u_R}{e_R}\times100\%,
  \qquad\text{y con señales normalizadas al }100\%:\qquad PB\%=\frac{100}{\Kc}.
\]
Una banda estrecha es una ganancia grande. (Algunos fabricantes, como ABB, usan la primera forma; otros, como Foxboro, la segunda~\citep{ACL05}.)

\begin{pausa}
Un controlador tiene $PB=25\%$ con señales normalizadas. ¿Cuánto vale $\Kc$? ¿Y si el error se mide sobre un rango de $200\,\grados$C y la salida sobre $100\%$?
\respuesta{$\Kc=100/25=4$ (adimensional). Con los rangos: $\Kc=\frac{1}{0.25}\cdot\frac{100}{200}=2$\,\%/$\grados$C.}
\end{pausa}

\subsection{Acción integral}

La acción integral es proporcional al área acumulada del error:
\[
  u_c(t)=K_i\int_0^t e(\tau)\,\dif\tau,\qquad K_i=\frac{\Kc}{\Ti},\qquad \Gc(s)=\frac{K_i}{s}.
\]
Mientras haya error, la salida sigue cambiando; solo se detiene cuando el error es cero. Ese polo en el origen sube en uno el \clave{tipo} del sistema: el error ante escalón, en la referencia o en una perturbación de carga, se anula.

\begin{ideaclave}
\textbf{La integral ajusta el \emph{bias} sola.}\enspace En un PI, $u_0$ se reemplaza por $\frac{\Kc}{\Ti}\int_0^t e\,\dif\tau$, que crece o decrece hasta encontrar el valor que deja el error en cero, sin conocer el proceso. Así se construían los PI neumáticos, con un lazo de \emph{reset automático} (\cref{fig:pn-reset}): $U_c=\Kc E+\frac{1}{\Ti s+1}U_c$, es decir, $\frac{U_c}{E}=\Kc\Bigl(1+\frac{1}{\Ti s}\Bigr)$.
\end{ideaclave}

\begin{figure}[ht]
\centering
\includegraphics[width=0.55\linewidth]{cap07/pi_reset}
\caption{PI con reajuste automático (\emph{automatic reset}).}
\label{fig:pn-reset}
\end{figure}

La acción integral tiene memoria: cuando el error llega a cero, la integral no es cero, porque guarda todo lo anterior. La \cref{fig:pn-integral} (izquierda) lo muestra con $\Gc=1/s$ y $\Gp=1/(s+1)$: la salida del controlador sobrepasa su valor final y solo empieza a bajar cuando el error se vuelve negativo, es decir, después de que la salida ya pasó la referencia. De ahí viene el sobrepico. A la derecha, al reducir $\Ti$ (más acción integral) la respuesta es más rápida pero menos amortiguada: el lazo es $s^2+s+1/\Ti$, con $\zeta=\tfrac12\sqrt{\Ti}$.

\begin{figure}[ht]
\centering
\includegraphics[width=\linewidth]{cap07/f_integral}
\caption{Izquierda: control integral puro; la salida del controlador ($u$) sigue subiendo mientras el error es positivo. Derecha: menos $\Ti$, más rápido y más oscilatorio.}
\label{fig:pn-integral}
\end{figure}

El PI resuelve los dos problemas anteriores: la oscilación del on/off y el error del proporcional. Por eso es, de lejos, la estructura más usada en la industria; si no hay exigencias fuertes de velocidad, la acción derivativa sobra. Su precio es una pérdida de estabilidad (un polo más en el origen) y un fenómeno que veremos en la \cref{sec:pn-modificaciones}: el \emph{wind-up}. Algunos fabricantes usan, en lugar de $\Ti$, la \clave{tasa de reajuste} $1/\Ti$, en repeticiones por minuto: cuántas veces por minuto la integral repite la acción proporcional ante un error constante.

\subsection{Acción derivativa}

El PI mira el pasado y el presente, pero no se anticipa. En la \cref{fig:pn-derivada} los dos errores tienen el mismo valor en $t_1$ y la misma área hasta $t_1$, así que un PI daría exactamente la misma salida. Sin embargo, en (a) el error cae rápido y el controlador debería frenar para no pasarse, y en (b) el error está volviendo a subir y debería reaccionar ya.

\begin{figure}[ht]
\centering
\includegraphics[width=0.8\linewidth]{cap07/f_derivada}
\caption{Dos errores con el mismo valor y la misma área en $t_1$: un PI no los distingue; la derivada sí.}
\label{fig:pn-derivada}
\end{figure}

La acción derivativa es proporcional a la pendiente del error: $u_c=K_d\,\dot e$, con $K_d=\Kc\Td$. Casi nunca se usa sola, sino con la proporcional, y entonces tiene una interpretación muy clara. Por Taylor, $e(t+\Td)\approx e(t)+\Td\,\dot e(t)$, así que
\[
  \Kc\bigl(e(t)+\Td\,\dot e(t)\bigr)\approx\Kc\,e(t+\Td):
\]
\begin{ideaclave}
\textbf{El PD es un proporcional que mira $\Td$ segundos adelante.}\enspace La salida en $t$ se basa en el error que se \emph{predice} para $t+\Td$, extrapolando en línea recta. Por eso los fabricantes la llamaron \emph{pre-act} o \emph{rate action}.
\end{ideaclave}

La derivada sirve cuando el lazo se beneficia de anticiparse, como en procesos de temperatura de orden alto. Pero tiene dos problemas. $\Gc=K_d s$ es una función de transferencia \alerta{impropia} (más ceros que polos): no se puede realizar físicamente y amplifica sin límite las altas frecuencias, es decir, el ruido. Y reacciona con violencia a los saltos de la referencia. Las dos cosas se resuelven en la \cref{sec:pn-modificaciones}.

%% ---------------------------------------------------------------------
\section{Tres formas de escribir un PID}
\label{sec:pn-estructuras}
%% ---------------------------------------------------------------------

Los controladores comerciales no implementan todos el mismo algoritmo, y los parámetros que da un método de sintonización dependen de la forma. Las tres clásicas son:
\begin{align}
  \text{ideal (ISA, estándar):}\qquad & \Gc(s)=\Kc^I\left(1+\frac{1}{\Ti^I s}+\Td^I s\right), \label{eq:pn-forma-ideal}\\
  \text{paralelo:}\qquad & \Gc(s)=\Kc^P+\frac{K_i^P}{s}+K_d^P s, \label{eq:pn-forma-paralelo}\\
  \text{serie (interactiva):}\qquad & \Gc(s)=\Kc^S\left(1+\frac{1}{\Ti^S s}\right)\left(1+\Td^S s\right). \label{eq:pn-forma-serie}
\end{align}
Los métodos de sintonización del capítulo~8 dan los parámetros en forma ideal; si el controlador es de otra forma, hay que convertirlos.

\paragraph{Paralelo e ideal.} Las tres acciones son independientes en las dos, y la conversión es inmediata:
\[
  \Kc^I=\Kc^P,\qquad \Ti^I=\frac{\Kc^P}{K_i^P},\qquad \Td^I=\frac{K_d^P}{\Kc^P}.
\]
La forma paralelo es la más general: es la única que admite $\Kc^P=0$, es decir, un controlador integral o integral-derivativo puro. (Quitar la integral o la derivada se puede en cualquiera de las tres, con $\Ti\to\infty$ o $\Td=0$.) La usan, entre otros, Siemens, Modicon y Allen-Bradley.

\paragraph{Serie e ideal.} Multiplicando la forma serie,
\[
  \Kc^S\left(1+\frac{1}{\Ti^S s}\right)\left(1+\Td^S s\right)
  =\Kc^S\left(1+\frac{\Td^S}{\Ti^S}+\frac{1}{\Ti^S s}+\Td^S s\right),
\]
y agrupando como en la ideal,
\begin{equation}
  \Kc^I=\Kc^S\left(1+\frac{\Td^S}{\Ti^S}\right),\qquad
  \Ti^I=\Ti^S+\Td^S,\qquad
  \Td^I=\frac{\Ti^S\Td^S}{\Ti^S+\Td^S}.
  \label{eq:pn-serie-ideal}
\end{equation}
Las acciones interactúan: cambiar $\Td^S$ cambia la ganancia y el tiempo integral equivalentes. El camino inverso no siempre existe. Imponiendo \eqref{eq:pn-serie-ideal}, $\Ti^S+\Td^S=\Ti^I$ y $\Ti^S\Td^S=\Ti^I\Td^I$, así que $\Ti^S$ y $\Td^S$ son las raíces de $z^2-\Ti^Iz+\Ti^I\Td^I=0$:
\begin{equation}
  \Kc^S=\frac{\Kc^I}{2}\bigl(1+\rho\bigr),\qquad
  \Ti^S=\frac{\Ti^I}{2}\bigl(1+\rho\bigr),\qquad
  \Td^S=\frac{\Ti^I}{2}\bigl(1-\rho\bigr),\qquad
  \rho=\sqrt{1-4\,\frac{\Td^I}{\Ti^I}} ,
  \label{eq:pn-ideal-serie}
\end{equation}
que solo son reales si \alerta{$\Ti^I\ge4\Td^I$}.

\begin{ideaclave}
\textbf{Los ceros del PID.}\enspace Un PID tiene un polo en el origen y dos ceros, las raíces de $\Ti\Td s^2+\Ti s+1$. La forma serie solo puede tener ceros reales (en $-1/\Ti^S$ y $-1/\Td^S$). La ideal admite ceros complejos, cuando $\Ti<4\Td$: es más general. En $\Ti=4\Td$ los dos ceros coinciden.
\end{ideaclave}

La forma serie fue la primera que se construyó con tecnología neumática, y muchos fabricantes la conservaron por tradición y porque se dice que es más fácil de sintonizar a mano.

\begin{ejemplo}[label=ej:pn-conversion]{De serie a ideal, y de vuelta}
Un controlador en forma serie tiene $\Kc^S=2$, $\Ti^S=8$\,s y $\Td^S=2$\,s. Halle la forma ideal y regrese.
\solucion Con \eqref{eq:pn-serie-ideal}: $\Kc^I=2(1+2/8)=2.5$, $\Ti^I=10$\,s y $\Td^I=16/10=1.6$\,s. De vuelta, $\rho=\sqrt{1-4\cdot1.6/10}=0.6$, y \eqref{eq:pn-ideal-serie} da $\Kc^S=2.5\cdot0.8=2$, $\Ti^S=5\cdot1.6=8$\,s y $\Td^S=5\cdot0.4=2$\,s.
\end{ejemplo}

\codigorepo{cap07_pid_nociones/C7_Ej3_Estructuras.m}{conversion}{ejemplo~\ref{ej:pn-conversion} con la función \cmd{convertirPID}}

\begin{pausa}
¿Tiene forma serie el PID ideal con $\Kc=5$, $\Ti=4$\,s y $\Td=2$\,s?
\respuesta{No: $4\Td=8>\Ti=4$. Los ceros, raíces de $8s^2+4s+1$, son $-0.25\pm j0.25$.}
\end{pausa}

%% ---------------------------------------------------------------------
\section{Un ejemplo para verlas actuar}
\label{sec:pn-ejemplo}
%% ---------------------------------------------------------------------

\begin{ejemplo}[label=ej:pn-mra]{Masa, resorte y amortiguador}
El sistema $M\ddot x+b\dot x+kx=F$, con $M=1$\,kg, $b=10$\,N\,s/m y $k=20$\,N/m, tiene
\[
  \Gp(s)=\frac{X(s)}{F(s)}=\frac{1}{s^2+10s+20}.
\]
Se quiere, con retroalimentación unitaria y un PID en forma paralelo, un tiempo de subida corto, el menor sobrepico posible y error cero ante escalón.
\solucion \textbf{Lazo abierto.} Los polos son $-2.76$ y $-7.24$ ($\zeta=1.12$): sin sobrepico, con $\tr=0.88$\,s, $\ts=1.60$\,s y valor final $1/20$. Respecto a una referencia unitaria, el error es del $95\%$.

\textbf{P.} El lazo es $s^2+10s+(20+\Kc)$: $\wn=\sqrt{20+\Kc}$, $\zeta=5/\sqrt{20+\Kc}$ y $\ess=20/(20+\Kc)$. Subir $\Kc$ reduce el tiempo de subida y el error, pero baja $\zeta$ y sube el sobrepico. Con $\Kc=300$: $\ess=0.0625$, $\zeta=0.28$ y $\OS=40\%$ (\cref{fig:pn-mra}, arriba a la izquierda).

\textbf{PD.} Con $\Kc=300$, la derivada agrega amortiguamiento: el lazo es $s^2+(10+K_d)s+320$. Con $K_d=20$, $\zeta=0.84$ y el sobrepico cae al $5.8\%$; el error no cambia ($K_ds$ se anula en $s=0$). El PD introduce además un cero en $-\Kc/K_d=-15$; por eso el sobrepico real es mayor que el que predice el segundo orden sin cero ($0.8\%$).

\textbf{PI.} Con $\Kc=300$, la integral anula el error, pero el lazo es de tercer orden, $s^3+10s^2+320s+K_i$, estable solo si $K_i<3200$ (Routh). Con $K_i=200$: $\ess=0$ y $\OS=36\%$.

\textbf{PID.} Con $\Kc=350$, $K_i=300$ y $K_d=50$, la respuesta no tiene sobrepico, sube en $0.05$\,s, se establece en $0.83$\,s y no tiene error. Los ceros del controlador ($-1$ y $-6$) casi cancelan los polos más lentos del lazo cerrado.
\end{ejemplo}

\begin{figure}[ht]
\centering
\includegraphics[width=\linewidth]{cap07/f_mra}
\caption{Ejemplo~\ref{ej:pn-mra}: efecto de cada acción sobre la respuesta al escalón.}
\label{fig:pn-mra}
\end{figure}

\codigorepo{cap07_pid_nociones/C7_Ej2_Acciones.m}{pid}{ejemplo~\ref{ej:pn-mra}: P, PD, PI y PID sobre el sistema masa-resorte-amortiguador}

\begin{table}[ht]
\centering
\caption{Efecto típico de aumentar cada parámetro de un PID sobre una planta estable~\citep{ACL05}. Son tendencias, no leyes: el problema~\ref{prob:pn-mra} pide comprobarlas.}
\label{tab:pn-efectos}
\small
\begin{tabular}{@{}lccccc@{}}
\toprule
Aumentar & $\tr$ & $\OS$ & $\ts$ & $\ess$ & Estabilidad\\
\midrule
$\Kc$ & baja & sube & sube un poco & baja & empeora\\
$K_i=\Kc/\Ti$ & baja un poco & sube & sube & se anula & empeora\\
$K_d=\Kc\Td$ & baja un poco & baja & baja & casi no cambia & mejora (si no hay ruido)\\
\bottomrule
\end{tabular}
\end{table}

\begin{receta}[Ajuste manual de un PID]
\begin{enumerate}
  \item Mire la respuesta en lazo abierto y decida qué hay que mejorar.
  \item Agregue acción proporcional para acelerar la respuesta.
  \item Agregue acción derivativa si el sobrepico es excesivo.
  \item Agregue acción integral para eliminar el error en estado estacionario.
  \item Ajuste las tres hasta lograr la respuesta deseada, y verifique la estabilidad.
  \item No use las tres si no hace falta: la derivada trae problemas (\cref{sec:pn-modificaciones}). Mantenga el controlador tan simple como se pueda.
\end{enumerate}
\end{receta}

En resumen: el \textbf{P} acelera pero deja error (salvo en procesos integradores); el \textbf{PI} elimina el error pero oscila más y pierde estabilidad al subir $\Kc$; el \textbf{PD} anticipa y amortigua, pero deja el error del P; el \textbf{PID} junta las ventajas, a costa de amplificar el ruido. En el capítulo~10 (\cref{sec:rf-pid}) veremos lo mismo en frecuencia: un PI es un compensador de atraso, un PD uno de adelanto y un PID uno de atraso-adelanto.

%% ---------------------------------------------------------------------
\section{Del PID de libro al PID de planta}
\label{sec:pn-modificaciones}
%% ---------------------------------------------------------------------

El PID ideal tiene problemas que aparecen apenas se conecta a un proceso real: saltos en la salida cuando cambia la referencia, ruido amplificado y un comportamiento extraño cuando el actuador se satura. Casi todos los controladores comerciales traen opciones para resolverlos; veamos cuáles y por qué funcionan.

\subsection{El \emph{proportional kick} y la estructura I-P}

Si el lazo está en equilibrio ($e\approx0$) y la referencia salta, el error salta con ella, y la acción proporcional lleva ese salto completo a la salida del controlador: un golpe al actuador que se llama \clave{proportional kick}. La solución es que el término proporcional actúe sobre la salida del proceso y no sobre el error, mientras la integral sigue sobre el error (estructura \clave{I-P}, \cref{fig:pn-ip}):
\[
  U_c(s)=\frac{\Kc}{\Ti s}E(s)-\Kc Y(s).
\]
La salida del controlador ya no salta, y el rechazo de perturbaciones es el mismo (para $R=0$ las dos estructuras coinciden). A cambio, la respuesta a la referencia es más lenta.

\begin{figure}[ht]
\centering
\includegraphics[width=0.72\linewidth]{cap07/estructura_ip}
\caption{Estructura I-P.}
\label{fig:pn-ip}
\end{figure}

\subsection{La derivada y el ruido: el filtro derivativo}

Un ruido de medida $n(t)=A\sin\omega t$ que pase por la derivada pura produce $u_c=AK_d\,\omega\cos\omega t$: \alerta{la amplitud crece con la frecuencia}. Con $\Kc=3$, $\Td=2$ y $A=0.01$, un ruido de $100$\,rad/s aparece en la salida del controlador con amplitud $6$. La solución es limitar la ganancia de alta frecuencia con un \clave{filtro derivativo}:
\begin{equation}
  \Gc(s)=\Kc\left(1+\frac{1}{\Ti s}+\frac{\Td s}{\dfrac{\Td}{N}s+1}\right).
  \label{eq:pn-filtro}
\end{equation}
Para $\omega\ll N/\Td$ la derivada funciona como antes; por encima, su ganancia se aplana en $\Kc N$, y la del controlador en $\Kc(1+N)$ (\cref{fig:pn-ruido}). El controlador queda propio (grado relativo cero) y realizable. $N$ suele estar entre $8$ y $16$~\citep{ACL05}; otros textos dan $5\le N\le20$~\citep{DoBi08}. Con $N=10$ el ruido del ejemplo baja de $6$ a $0.30$.

\begin{figure}[ht]
\centering
\includegraphics[width=0.8\linewidth]{cap07/f_ruido}
\caption{Magnitud del PID ideal ($\Kc=3$, $\Ti=8$, $\Td=2$) sin filtro y con filtro derivativo. Las líneas punteadas son las cotas $\Kc(1+N)$.}
\label{fig:pn-ruido}
\end{figure}

\begin{profundiza}[Otros filtros]
Hay variantes: un filtro de segundo orden en la derivada, $\Td s/\bigl(1+\frac{\Td s}{N}+\frac12(\frac{\Td s}{N})^2\bigr)$~\citep{gerry2004pid}; la forma serie con filtro, $\Kc\bigl(1+\frac{1}{\Ti s}\bigr)\frac{\Td s+1}{\Td s/N+1}$; o un filtro sobre toda la salida, $\Kc\bigl(1+\frac{1}{\Ti s}+\Td s\bigr)\frac{1}{(\tau_fs+1)^k}$ con $k=1$ o $2$. El diseño del filtro es un compromiso entre atenuar el ruido y no perder la acción derivativa útil.
\end{profundiza}

\subsection{El \emph{derivative kick} y la estructura I-PD}

Aun con filtro, si la derivada actúa sobre el error, un escalón en la referencia produce un pico enorme: la derivada (filtrada) de un escalón es un pulso de altura $\Kc\Td/(\Td/N)=\Kc N$. Es el \clave{derivative kick}. Con el PID $350/300/50$ del ejemplo~\ref{ej:pn-mra} y un filtro de constante $0.01$\,s, el pico es de unas $5300$ unidades (\cref{fig:pn-kicks}). Como la referencia es constante a trozos, su derivada no aporta nada útil entre saltos; por eso se deriva solo la salida. Con la proporcional también sobre la salida se obtiene la estructura \clave{I-PD} (\cref{fig:pn-ipd}).

\begin{figure}[ht]
\centering
\includegraphics[width=0.72\linewidth]{cap07/estructura_ipd}
\caption{Estructura I-PD: solo la integral ve el error.}
\label{fig:pn-ipd}
\end{figure}

\begin{figure}[ht]
\centering
\includegraphics[width=\linewidth]{cap07/f_kicks}
\caption{PID $350/300/50$ sobre $1/(s^2+10s+20)$. Estándar: pico de unas $5300$ unidades en la salida del controlador. PI-D (derivada sobre la salida): queda el salto proporcional de $350$. I-PD: sin saltos, pero más lento.}
\label{fig:pn-kicks}
\end{figure}

\codigorepo{cap07_pid_nociones/C7_Ej5_Kicks.m}{kicks}{las tres estructuras ante un escalón de referencia}

\subsection{\emph{Wind-up} y anti-windup}

Todo actuador tiene límites: una válvula no abre más del $100\%$. Si un cambio grande de referencia lleva la salida del controlador al límite, el lazo queda, en la práctica, abierto: el actuador ya no responde a lo que pida el controlador. Pero el error sigue siendo grande, y la integral sigue acumulándolo (\emph{winds up}). Cuando la salida del proceso por fin alcanza la referencia, la integral está inflada y mantiene el actuador saturado hasta que un error negativo igual de grande la descarga: sobrepico grande y establecimiento lento. Es el \clave{wind-up} del integrador.

\begin{ejemplo}[label=ej:pn-windup]{\emph{Wind-up} en un proceso con tiempo muerto}
Un PID ideal con $\Kc=3$, $\Ti=8$\,s, $\Td=2$\,s (y $N=10$) controla $\Gp(s)=e^{-4s}/(10s+1)$, con el actuador limitado a $[0,\,1.5]$~\citep{Vis06}.
\solucion Ante un escalón unitario, el actuador queda saturado $25$\,s en total; el término integral llega a más de $3$ cuando el actuador solo puede dar $1.5$. Resultado: $\OS=38\%$ y $\ts=45$\,s (\cref{fig:pn-windup}, en rojo).
\end{ejemplo}

Las correcciones \clave{anti-windup} más usadas son dos.
\begin{itemize}
  \item \textbf{Integración condicional} (\emph{clamping}): se detiene la integral mientras el actuador esté saturado y el error empuje hacia el mismo lado ($u_c$ saturado y $(v-u)\,e>0$). No hay que sintonizar nada. En el ejemplo: $\OS=3.7\%$, $\ts=29$\,s.
  \item \textbf{Recálculo} (\emph{back-calculation}, \cref{fig:pn-antiwindup}): se retroalimenta a la integral la diferencia entre la salida saturada $u$ y la que pedía el controlador $v$,
  \[
    \frac{\dif I}{\dif t}=\frac{\Kc}{\Ti}\,e+\frac{1}{T_t}\,(u-v).
  \]
  Mientras no hay saturación, $u=v$ y no pasa nada; con saturación, la integral se descarga con constante de tiempo $T_t$. Funciona como un observador: estima el estado correcto del controlador. Se recomienda $T_t=\sqrt{\Ti\Td}$~\citep{astrom1995pid}, o $T_t=\Ti$ para un PI~\citep{atherton1995pid}. En el ejemplo, con $T_t=4$: $\OS=3.3\%$, $\ts=30$\,s; con $T_t=8$: $\OS=11\%$.
\end{itemize}

\begin{figure}[ht]
\centering
\includegraphics[width=0.85\linewidth]{cap07/antiwindup}
\caption{PID con anti-windup por recálculo (\emph{back-calculation}).}
\label{fig:pn-antiwindup}
\end{figure}

\begin{figure}[ht]
\centering
\includegraphics[width=\linewidth]{cap07/f_windup}
\caption{Ejemplo~\ref{ej:pn-windup}. Sin corrección, la integral (punteada) sube mucho más de lo que el actuador puede dar; con anti-windup se mantiene acotada.}
\label{fig:pn-windup}
\end{figure}

\codigorepo{cap07_pid_nociones/C7_Ej6_Windup.m}{windup}{ejemplo~\ref{ej:pn-windup}: sin corrección, con recálculo y con integración condicional}

\begin{modelo}[AntiWindup.slx]
El script \cmd{crearModelo\_AntiWindup.m} arma el ejemplo con el bloque \emph{PID Controller} de \simulink, que trae la saturación y el anti-windup incorporados (parámetros \emph{Limit output} y \emph{Anti-windup method}). Compare las dos copias del lazo y cambie $T_t$.
\end{modelo}

\begin{cuidado}[Dónde aparece el \emph{wind-up}]
El \emph{wind-up} es más grave ante cambios grandes de la referencia que ante perturbaciones de carga~\citep{VPPH97}, y en procesos de orden bajo con tiempo muerto pequeño respecto a la constante de tiempo, donde el controlador puede ser agresivo. Un controlador que nunca satura no lo necesita, pero en planta casi todos saturan alguna vez (arranques, cambios de producto).
\end{cuidado}

\subsection{Ponderación de la referencia: dos grados de libertad}

Un lazo sintonizado para rechazar perturbaciones rápido suele responder con mucho sobrepico a los cambios de referencia. Con un solo controlador sobre el error no se pueden tener las dos cosas. La salida es separar cómo actúa el controlador sobre la referencia y sobre la salida (\cref{fig:pn-2gdl}):
\[
  u_c=\Kc\Bigl(\beta\,r-y\Bigr)+\frac{\Kc}{\Ti}\int(r-y)\,\dif\tau+\Kc\Td\,\deriv{}{t}\bigl(\gamma\,r-y\bigr),
  \qquad 0\le\beta,\gamma\le1 .
\]
La integral siempre actúa sobre el error completo, para no perder el error cero. Con $\beta=\gamma=1$ es el PID estándar; con $\beta=1$, $\gamma=0$, el PI-D; con $\beta=\gamma=0$, el I-PD. Esta es la estructura del \clave{PID ISA} (o controlador $\beta$-$\gamma$). Equivale a un filtro sobre la referencia,
\[
  G_f(s)=\frac{1+\bigl(\beta\Ti+\frac{\Td}{N}\bigr)s+\Ti\Td\bigl(\gamma+\frac{\beta}{N}\bigr)s^2}{1+\bigl(\Ti+\frac{\Td}{N}\bigr)s+\Ti\Td\bigl(1+\frac{1}{N}\bigr)s^2},
\]
que \alerta{no aparece en la función de transferencia de la perturbación a la salida}: el rechazo de perturbaciones se sintoniza con $\Kc$, $\Ti$, $\Td$, y la respuesta a la referencia, después, con $\beta$ y $\gamma$.

\begin{figure}[ht]
\centering
\includegraphics[width=0.62\linewidth]{cap07/dos_gdl}
\caption{Controlador de dos grados de libertad: $G_{sp}=\Kc\bigl(\beta+\frac{1}{\Ti s}+\gamma\Td s\bigr)$ sobre la referencia y $\Gc$ sobre la salida.}
\label{fig:pn-2gdl}
\end{figure}

\begin{ejemplo}[label=ej:pn-beta]{Un PI agresivo con ponderación}
Un PI con $\Kc=20$ y $\Ti=0.3$\,s controla $\Gp=1/\bigl((s+1)(s+5)\bigr)$; está sintonizado para rechazar perturbaciones (pico de $0.032$ ante un escalón unitario a la entrada de la planta, frente a $0.11$ de un PI conservador con $\Kc=2$, $\Ti=1$).
\solucion Ante la referencia, con $\beta=1$ el sobrepico es del $51\%$; con $\beta=0.5$, del $33\%$, y con $\beta=0$, del $27\%$ (\cref{fig:pn-beta}). La respuesta a la perturbación es la misma para todos. Note que ningún $\beta$ baja el sobrepico de $27\%$: con $\beta=0$ queda solo la integral sobre la referencia, y la oscilación la imponen los polos del lazo cerrado, que $\beta$ no mueve.
\end{ejemplo}

\begin{figure}[ht]
\centering
\includegraphics[width=\linewidth]{cap07/f_beta}
\caption{Ejemplo~\ref{ej:pn-beta}: $\beta$ cambia la respuesta a la referencia y no la respuesta a la perturbación.}
\label{fig:pn-beta}
\end{figure}

\codigorepo{cap07_pid_nociones/C7_Ej7_Ponderacion.m}{beta}{ejemplo~\ref{ej:pn-beta}: ponderación de la referencia}

%% ---------------------------------------------------------------------
\section{Implementación digital}
\label{sec:pn-digital}
%% ---------------------------------------------------------------------

Hoy casi todos los PID corren en un computador que muestrea el error cada $T_s$ segundos. Hay varias formas de discretizar~\citep{AsWi97}; la más sencilla aproxima la integral por una suma y la derivada por una diferencia hacia atrás:
\[
  \int_0^t e\,\dif\tau\approx T_s\sum_{i=1}^{k}e[i],\qquad \deriv{e}{t}\approx\frac{e[k]-e[k-1]}{T_s}.
\]
El resultado es el \clave{algoritmo posicional}, que calcula la salida completa en cada muestra:
\begin{equation}
  u_c[k]=\Kc\left(e[k]+\frac{T_s}{\Ti}\sum_{i=1}^{k}e[i]+\frac{\Td}{T_s}\bigl(e[k]-e[k-1]\bigr)\right).
  \label{eq:pn-posicional}
\end{equation}
Restando \eqref{eq:pn-posicional} en $k$ y en $k-1$ se obtiene el \clave{algoritmo incremental} (o de velocidad), que calcula solo el cambio:
\begin{equation}
  u_c[k]=u_c[k-1]+\Kc\left[\Bigl(1+\frac{T_s}{\Ti}+\frac{\Td}{T_s}\Bigr)e[k]-\Bigl(1+\frac{2\Td}{T_s}\Bigr)e[k-1]+\frac{\Td}{T_s}\,e[k-2]\right].
  \label{eq:pn-incremental}
\end{equation}

\begin{ideaclave}
\textbf{Por qué gusta el incremental.}\enspace No guarda la suma de todos los errores, se lleva bien con actuadores que reciben incrementos (motores de paso) y, si la salida se satura y se recorta el valor que se guarda como $u_c[k-1]$, no acumula integral: es naturalmente resistente al \emph{wind-up}. Note el factor $\Td/T_s$: con muestreo rápido amplifica mucho cualquier cambio entre dos muestras, incluido el ruido y la cuantización del conversor.
\end{ideaclave}

\begin{ejemplo}[label=ej:pn-digital]{Coeficientes del algoritmo incremental}
Para $\Kc=3$, $\Ti=8$\,s, $\Td=2$\,s y $T_s=0.1$\,s,
\[
  u_c[k]=u_c[k-1]+63.04\,e[k]-123\,e[k-1]+60\,e[k-2].
\]
Simulado con $\Gp=1/(10s+1)$, da la misma señal de control que el posicional (diferencias del orden de $10^{-13}$).
\end{ejemplo}

\codigorepo{cap07_pid_nociones/C7_Ej8_Digital.m}{digital}{ejemplo~\ref{ej:pn-digital}: posicional e incremental}

%% ---------------------------------------------------------------------
\section{Para llevar}
\label{sec:pn-cierre}
%% ---------------------------------------------------------------------

\begin{parallevar}
\begin{enumerate}
  \item PID ideal: $\Kc\bigl(1+\frac{1}{\Ti s}+\Td s\bigr)$. P mira el presente, I el pasado (área), D el futuro (tendencia): un PD es un proporcional que mira $\Td$ segundos adelante.
  \item El on/off oscila; el P deja error en procesos de tipo~0; la I lo elimina ajustando el \emph{bias} sola, a costa de estabilidad; la D amortigua, pero amplifica el ruido.
  \item Formas ideal, paralelo y serie. Serie $\to$ ideal siempre; ideal $\to$ serie solo si $\Ti\ge4\Td$. Los métodos de sintonización dan la forma ideal.
  \item Filtro derivativo con $N\approx8$--$16$: ganancia acotada en $\Kc(1+N)$ y controlador realizable.
  \item Los \emph{kicks} se evitan sacando la P y la D del camino de la referencia (I-P, PI-D, I-PD). La ponderación $\beta$, $\gamma$ separa la respuesta a la referencia del rechazo de perturbaciones.
  \item Con actuador saturado, la integral se infla (\emph{wind-up}). Anti-windup: integración condicional o recálculo con $T_t\approx\sqrt{\Ti\Td}$.
  \item Digital: algoritmo posicional o incremental; el incremental resiste mejor el \emph{wind-up}.
\end{enumerate}
\end{parallevar}

\begin{profundiza}[Correcciones respecto a la versión anterior del capítulo]
(i) La conversión ideal $\to$ serie daba $\Td^S=\frac{\Td^I}{2}(1-\rho)$; es $\frac{\Ti^I}{2}(1-\rho)$. (ii) La expansión de la forma serie tenía los términos mal agrupados. (iii) La forma paralelo no es ``la única que permite desactivar la integral'': es la única que admite $\Kc=0$. (iv) Historia: la Taylor introdujo una derivada fija hacia 1936; en 1939 salieron el Fulscope y el Stabilog con las tres acciones ajustables. Se retiró que la ganancia del Stabilog ``podía ajustarse entre 0.7 y 100'' y la mención al \emph{Trans Scan Log}, que no se encontraron en las fuentes. (v) Se completaron las figuras y simulaciones pendientes (on/off, kicks, \emph{wind-up}, ponderación) y se calcularon todos los números de los ejemplos. (vi) Notación: $\tau_i,\tau_d\to\Ti,\Td$; MP $\to\OS$; ``realimentar'' $\to$ retroalimentar; ``Allan Braedley'' $\to$ Allen-Bradley.
\end{profundiza}

%% ---------------------------------------------------------------------
\section{Problemas}
\label{sec:pn-problemas}
%% ---------------------------------------------------------------------

Los problemas marcados con \etiquetaM{} requieren \matlab/\simulink; los marcados con \etiquetaEstrella{} son de profundización.

\begin{problema}[formas del PID]\label{prob:pn-formas}
(a) Expanda la forma serie y deduzca \eqref{eq:pn-serie-ideal}. (b) Un controlador en forma serie tiene $\Kc^S=2$, $\Ti^S=8$\,s y $\Td^S=2$\,s; halle la forma ideal. (c) Muestre que $\Ti^S$ y $\Td^S$ son las raíces de $z^2-\Ti^Iz+\Ti^I\Td^I=0$, obtenga \eqref{eq:pn-ideal-serie} y verifique con (b). (d) Muestre, con el discriminante, que la conversión exige $\Ti^I\ge4\Td^I$. (e) Verifique que el PID ideal con $\Kc=5$, $\Ti=4$\,s y $\Td=2$\,s no tiene forma serie y calcule sus ceros. (f) Pase el PID paralelo $350/300/50$ del ejemplo~\ref{ej:pn-mra} a forma ideal y diga si tiene forma serie.
\end{problema}

\begin{problema}[las tendencias, con números]\label{prob:pn-mra}
Con el sistema del ejemplo~\ref{ej:pn-mra}: (a) en lazo abierto, calcule el valor final, $\zeta$, $\wn$, $\tr$ y $\ts$, y verifique con \cmd{stepinfo}; (b) con un P, halle $\ess$, $\zeta$, $\wn$ y $\OS$ en función de $\Kc$ y tabule para $\Kc=10$, $100$ y $300$; (c) con $\Kc=300$ y $K_d=10$, $20$, ubique el cero, explique por qué $\ess$ no cambia y compare el $\OS$ con el que predice el segundo orden sin el cero; (d) con $\Kc=300$, halle por Routh el rango de $K_i$ y evalúe $K_i=70$ y $200$; (e) \etiquetaM{} arme una tabla P, PD, PI, PID con $\tr$, $\OS$, $\ts$, $\ess$ y $|u_c|_{\max}$, y verifique una por una las tendencias de la \cref{tab:pn-efectos}. ¿Alguna no se cumple aquí?
\end{problema}

\begin{problema}[banda proporcional y \emph{bias}]\label{prob:pn-banda}
Un lazo de temperatura tiene un transmisor de $0$--$200\,\grados$C, una válvula de $0$--$100\%$ y un proceso de ganancia estática $\Kproc=2.5\,\grados$C/\%, con control proporcional. (a) Con $PB=25\%$, calcule $\Kc$ con las dos convenciones y explique por qué difieren. (b) Con la segunda y $u_0=50\%$, halle el error para $r=120\,\grados$C. (c) Halle el $u_0$ que lo anula. (d) Con ese $u_0$, ¿cuál es el error si la ganancia cae a $2.2$? ¿Y si sube a $2.8$? (e) Explique, con esos números, por qué el \emph{bias} manual no es robusto y qué hace la integral. (f) Con $u_0=50\%$ y $r=150\,\grados$C, ¿qué banda proporcional deja el error por debajo de $2\,\grados$C? ¿Es razonable?
\end{problema}

\begin{problema}[ruido y filtro derivativo]\label{prob:pn-ruido}
Con el PID ideal $\Kc=3$, $\Ti=8$\,s, $\Td=2$\,s: (a) para un ruido $0.01\sin\omega t$, calcule su amplitud en la salida del controlador con la derivada pura para $\omega=0.1$, $10$ y $100$\,rad/s; (b) muestre que la derivada filtrada tiene ganancia acotada en alta frecuencia y halle la cota; (c) halle la frecuencia a partir de la cual actúa el filtro y la amplitud del ruido en $\omega=100$ para $N=5$, $10$ y $20$; (d) \etiquetaM{} grafique $|\Gc(j\omega)|$ con y sin filtro; (e) verifique que el controlador filtrado es propio; (f) \etiquetaM{} simule el lazo con $\Gp=1/(10s+1)$ y ruido blanco en la medida, y mida la desviación estándar de $u_c$ para cada $N$.
\end{problema}

\begin{problema}[algoritmo digital]\label{prob:pn-digital}
Con el mismo PID: (a) escriba el algoritmo posicional y el incremental con $T_s=0.1$\,s; (b) deduzca el incremental del posicional; (c) \etiquetaM{} compruebe que dan la misma salida con $\Gp=1/(10s+1)$; (d) \etiquetaM{} repita con $T_s=0.01$, $0.5$ y $1$\,s y halle el mayor $T_s$ con el que $\OS$ y $\ts$ empeoran menos del $5\%$ frente al continuo; (e) \etiquetaEstrella{} cuantice el error a 12 bits en $\pm10$ y mida el ruido en $u_c$ para cada $T_s$; (f) agregue una saturación y explique por qué el incremental sufre menos \emph{wind-up}.
\end{problema}

\begin{problema}[\emph{kicks}]\label{prob:pn-kicks}
Con $\Gp=1/(s^2+10s+20)$ y el PID paralelo $350/300/50$, compare las estructuras estándar, PI-D e I-PD. (a) Escriba $U_c(s)$ en función de $R$ y $Y$ para cada una. (b) Con el teorema del valor inicial, halle $u_c(0^+)$ ante un escalón de referencia; muestre que en la estándar no está acotado. (c) Repita con la derivada filtrada de constante $\tau_f$ y evalúe para $\tau_f=0.01$, $0.05$ y $0.1$\,s. (d) \etiquetaM{} Simule y tabule $|u_c|_{\max}$, $\OS$ y $\ts$. (e) \etiquetaM{} Agregue una saturación en $[-100,\,100]$ y diga cuáles saturan y cuánto tiempo. (f) Muestre que PI-D e I-PD son casos de la ponderación con $\beta$ y $\gamma$.
\end{problema}

\begin{problema}[\emph{wind-up}]\label{prob:pn-windup}
Con el ejemplo~\ref{ej:pn-windup}: (a) \etiquetaM{} simule sin saturación y reporte $\OS$, $\ts$ y el pico de $u_c$; (b) agregue la saturación y reporte cuánto tiempo satura y cuánto empeoran $\OS$ y $\ts$; (c) grafique el término integral y marque el instante en que el error cambia de signo y el instante en que el actuador sale de saturación; explique la diferencia; (d) implemente el recálculo con $T_t=\sqrt{\Ti\Td}$ y $T_t=\Ti$; (e) recorra $T_t$ entre $0.5$ y $20$\,s y grafique $\OS$, $\ts$ e IAE; (f) implemente la integración condicional y compárela con el mejor recálculo.
\end{problema}

\begin{problema}[dos grados de libertad]\label{prob:pn-beta}
Con el ejemplo~\ref{ej:pn-beta}: (a) verifique por Routh que el PI agresivo es estable; (b) \etiquetaM{} compare el rechazo de un escalón de perturbación a la entrada con el PI conservador ($\Kc=2$, $\Ti=1$); (c) simule la respuesta a la referencia y constate el conflicto; (d) obtenga $G_f(s)$ para un PI con ponderación $\beta$; (e) muestre que $G_f$ no aparece en la función de transferencia de la perturbación; (f) tabule $\OS$ y $\ts$ para $\beta=0$, $0.3$, $0.5$ y $1$; ¿qué valores de $\beta$ dan $\OS\le30\%$? Explique por qué ningún $\beta$ baja de unos $27\%$.
\end{problema}

%% ---------------------------------------------------------------------
\section{Soluciones seleccionadas}
%% ---------------------------------------------------------------------

\paragraph{Problema~\ref{prob:pn-formas}}
(b) $\Kc^I=2.5$, $\Ti^I=10$\,s, $\Td^I=1.6$\,s. (e) $\Ti=4<8$; ceros $-0.25\pm j0.25$. (f) $\Kc=350$, $\Ti=1.167$\,s, $\Td=0.143$\,s; como $\Ti\ge4\Td$, sí: $\Kc^S=300$, $\Ti^S=1$\,s, $\Td^S=0.167$\,s (ceros en $-1$ y $-6$).

\paragraph{Problema~\ref{prob:pn-mra}}
(a) Valor final $0.05$, $\zeta=1.12$, $\wn=4.47$\,rad/s, $\tr=0.88$\,s, $\ts=1.60$\,s. (b) $\ess=\frac{20}{20+\Kc}$, $\wn=\sqrt{20+\Kc}$, $\zeta=\frac{5}{\sqrt{20+\Kc}}$; para $\Kc=10$, $100$, $300$: $\ess=0.667$, $0.167$, $0.0625$; $\OS=0.1\%$, $20\%$, $40\%$. (c) Cero en $-30$ y $-15$; $\OS=15.3\%$ y $5.8\%$ frente a $12.0\%$ y $0.8\%$ sin el cero. (d) $0<K_i<3200$; $\OS=33\%$ ($\ts=4.7$\,s) y $36\%$ ($\ts=1.2$\,s). (e) Aquí, subir $\Kc$ de $10$ a $100$ \emph{baja} $\ts$ (de $0.88$ a $0.76$\,s) en lugar de subirlo; de $100$ a $300$ casi no lo cambia. Las tendencias de la tabla no son leyes.

\paragraph{Problema~\ref{prob:pn-banda}}
(a) $\Kc=4$ (adimensional) y $\Kc=2$\,\%/$\grados$C. (b) $e=-0.83\,\grados$C (el \emph{bias} de $50\%$ llevaría la PV a $125\,\grados$C). (c) $u_0=48\%$. (d) $+2.67$ y $-2.18\,\grados$C. (f) $\Kc\ge4.6$\,\%/$\grados$C, es decir, $PB\le10.9\%$.

\paragraph{Problema~\ref{prob:pn-ruido}}
(a) $0.006$, $0.6$ y $6$. (b) $\Kc N$ en la derivada, $\Kc(1+N)$ en el controlador. (c) A partir de $N/\Td$: $2.5$, $5$ y $10$\,rad/s; amplitudes $0.15$, $0.30$ y $0.60$.

\paragraph{Problema~\ref{prob:pn-digital}}
(a) $u_c[k]=3\bigl(e[k]+0.0125\sum e[i]+20(e[k]-e[k-1])\bigr)$ y $u_c[k]=u_c[k-1]+63.04\,e[k]-123\,e[k-1]+60\,e[k-2]$.

\paragraph{Problema~\ref{prob:pn-kicks}}
(b) $u_c(0^+)=\Kc+K_d\,\delta$: un impulso. (c) $u_c(0^+)=\Kc+K_d/\tau_f=5350$, $1350$ y $850$. (d) $|u_c|_{\max}\approx5300$, $350$ y $20$; $\OS=0$, $5.3$ y $0\%$; $\ts=0.84$, $2.06$ y $4.29$\,s. (f) PI-D: $\beta=1$, $\gamma=0$; I-PD: $\beta=\gamma=0$.

\paragraph{Problema~\ref{prob:pn-windup}}
(a) Sin saturación, $\OS=67\%$ y el pico de $u_c$ es de $32$ (el \emph{derivative kick} con $N=10$). (b) $25$\,s saturado, $\OS=38\%$, $\ts=45$\,s. (d) $T_t=4$: $\OS=3.3\%$, $\ts=30$\,s; $T_t=8$: $\OS=11\%$, $\ts=32$\,s. (f) $\OS=3.7\%$, $\ts=29$\,s, sin parámetro que sintonizar.

\paragraph{Problema~\ref{prob:pn-beta}}
(a) $s^3+6s^2+25s+66.7$: $6\cdot25=150>66.7$. (b) Pico $0.032$ frente a $0.114$. (f) $\OS=27\%$, $29\%$, $33\%$ y $51\%$; $\OS\le30\%$ con $\beta\lesssim0.35$.

\respuestaspausas

                   % migrado (Handout6a_Ctrl.tex)
%% =====================================================================
%%  Capítulo 8 -- Controladores PID: sintonización
%%  Migrado al formato nuevo (noviembre de 2026) desde Handout6b_Ctrl.tex.
%%  Correcciones de fondo (ver MIGRACION.md y la caja "Correcciones" de la
%%  sección 8.8): pendiente máxima de un primer orden; PI de Ziegler-Nichols
%%  de malla cerrada (0.45 Kcr, Pcr/1.2); Tyreus-Luyben no es de "un cuarto";
%%  efecto de lambda en síntesis directa; IMC inestable F(1/tau_u) = 1; fila
%%  FOPDT de la tabla IMC; notación tau_i, tau_d -> T_i, T_d; theta -> L.
%% =====================================================================

\capitulo[cap:ps]{Controladores PID: sintonización}
  {La autoridad y el dinero se rodean de calma, y se puede medir el alcance de la influencia de una persona por la intensidad del silencio que la rodea.}
  {Hernán Díaz}

\begin{objetivos}[partes/08-pid-sintonizacion.html]
  \item Usar la terminología del control de procesos y reconocer los elementos de un lazo industrial.
  \item Clasificar un proceso por su respuesta al escalón e identificar un modelo de primer orden con tiempo muerto, con la ganancia normalizada.
  \item Sintonizar un PID con reglas de malla abierta (Ziegler-Nichols, Cohen-Coon, coeficiente de ajustabilidad) y de malla cerrada (Ziegler-Nichols, Tyreus-Luyben), y juzgar el resultado por desempeño \emph{y} por robustez.
  \item Diseñar un controlador por síntesis directa y por modelo interno (IMC), y obtener de ellos un PID con un solo parámetro de ajuste, $\lambda$.
\end{objetivos}

%% ---------------------------------------------------------------------
\section{El lazo de un proceso industrial}
\label{sec:ps-terminos}
%% ---------------------------------------------------------------------

El capítulo~\ref{cap:pn} explicó qué hace cada acción del PID; falta decidir cuánto vale cada una. A eso se le llama \textbf{sintonizar}. El problema no está resuelto en la práctica: un estudio muy citado de Van Overschee y De Moor reportó que cerca del $80$\,\% de los PID industriales estaban mal sintonizados, que un $30$\,\% operaba en manual y que un $25$\,\% conservaba los parámetros de fábrica~\citep{VaDe00}. Y no es por falta de reglas: el compendio de O'Dwyer reúne más de cuatrocientas~\citep{ODw06}. Este capítulo presenta unas pocas, las que más se usan y las que mejor enseñan, empezando por el vocabulario de los procesos, que es donde vive el PID.

\paragraph{Terminología.} La \textbf{variable del proceso} (PV) es la que se quiere controlar: nivel, presión, flujo, temperatura, pH. La \textbf{referencia} o \emph{set point} (SP) es el valor deseado. La \textbf{variable manipulada} (MV) es la que se puede mover para corregir, y la \textbf{salida del controlador} (CO, la $u_c$ del capítulo anterior) es la orden que llega al elemento final. A veces se mide algo distinto de la PV: en un control de nivel que manipula el caudal de entrada, la PV es el nivel y la MV es el caudal. El \textbf{offset} es una desviación permanente entre la PV y la SP, y una \textbf{perturbación de carga} es un cambio indeseado en alguna entrada que afecta la PV. Cuando una persona hace los ajustes, el control es \textbf{manual}; cuando los hace el controlador, \textbf{automático}.

\begin{figure}[ht]
\centering
\includegraphics[width=0.9\linewidth]{cap08/lazo_terminos}
\caption{Lazo de un proceso: controlador, elemento final $\Gv$, proceso $\Gp$, perturbación a través de $G_d$ y medición $G_m$ (sensor y transmisor). El controlador solo ve el producto $\Gv\Gp G_m$.}
\label{fig:ps-lazo}
\end{figure}

\paragraph{Componentes.} La \cref{fig:ps-lazo} muestra los bloques. El \textbf{sensor} o elemento primario es el primero que mide la PV; el \textbf{transductor} convierte una magnitud en otra (una presión en una capacitancia, por ejemplo); un \textbf{convertidor} cambia el tipo de señal (corriente a presión, analógica a digital); el \textbf{transmisor} lleva la medida a una señal estándar. El \textbf{elemento final de control} (válvula, bomba, variador) es el que mueve físicamente la MV, y su \textbf{actuador} lo acciona con energía neumática, hidráulica o eléctrica. Completan el lazo los \textbf{indicadores}, que muestran la medida a un operador, y los \textbf{registradores}, que la guardan.

\paragraph{Señales estándar.} La señal neumática va de $3$ a $15$\,psig; la analógica, de $4$ a $20$\,mA (o de $1$ a $5$\,V), con el cero ``vivo'' en $4$\,mA para distinguir un valor mínimo de un cable roto; y la digital usa protocolos como HART, Fieldbus o Modbus. El extremo inferior del rango se llama LRV y el superior URV.

\begin{ejemplo}[label=ej:ps-4a20]{Un transmisor de 4 a 20 mA}
Un transmisor de temperatura está calibrado entre $95$ y $105\,\grados$C. ¿Qué corriente corresponde a $100\,\grados$C? ¿Y a $98\,\grados$C?
\solucion La relación es lineal: $i=4+16\,\dfrac{T-95}{105-95}$\,mA. Para $100\,\grados$C, $i=12$\,mA, la mitad del rango; para $98\,\grados$C, $i=8.8$\,mA.
\end{ejemplo}

\begin{ideaclave}
\textbf{Lo que se sintoniza es todo el lazo.}\enspace El controlador no distingue la válvula, el proceso y el transmisor: ve un solo sistema, de la CO a la medida. Por eso las pruebas de este capítulo se hacen moviendo la CO y mirando la PV medida, y el ``proceso'' identificado incluye la válvula y el transmisor.
\end{ideaclave}

%% ---------------------------------------------------------------------
\section{Tipos de procesos e identificación}
\label{sec:ps-identificacion}
%% ---------------------------------------------------------------------

La forma más común de conocer la dinámica de un proceso industrial es la \textbf{prueba escalón en lazo abierto} (\cref{fig:ps-prueba}): con el controlador en \emph{manual}, se espera a que la PV esté en estado estacionario, se cambia la CO de golpe y se registra la PV. La respuesta dice dónde estarán las dificultades del control: un tiempo muerto grande, una respuesta inversa, una no linealidad. La misma prueba puede hacerse en lazo cerrado, con el controlador en automático y un escalón de referencia; la usaremos en la \cref{sec:ps-malla-cerrada}.

\begin{figure}[ht]
\centering
\includegraphics[width=0.75\linewidth]{cap08/prueba_abierta}
\caption{Prueba escalón en lazo abierto: el controlador en manual entrega un escalón $\Delta u_c$ y se registra la PV.}
\label{fig:ps-prueba}
\end{figure}

\subsection{Seis tipos de respuesta}

La \cref{fig:ps-tipos} reúne las respuestas típicas.
\begin{description}
  \item[(a) Primer orden.] Sube de inmediato y se acerca exponencialmente al valor final, como un termómetro de vidrio delgado sumergido en agua hirviendo. Se comporta como un filtro pasabajos.
  \item[(b) Orden superior.] La más común: varios almacenamientos en serie (procesos \emph{multicapacidad}) dan una curva en forma de S que arranca con pendiente nula.
  \item[(c) Integrador.] No tiene valor final: la PV crece como una rampa mientras la CO esté desplazada. Aparece donde hay almacenamiento sin autorregulación: nivel, presión en tanques cerrados, concentración en lotes, temperatura en recipientes bien aislados.
  \item[(d) Oscilatorio.] Raro en procesos sin controlar, común en sistemas mecánicos y en lazos ya controlados.
  \item[(e) Tiempo muerto.] Nada pasa durante $L$ segundos: transporte por tuberías o bandas, la ducha del capítulo~\ref{cap:in}, retardos de comunicación. Casi siempre aparece combinado con otro tipo.
  \item[(f) Respuesta inversa.] Arranca en la dirección contraria. Es la marca de un cero en el semiplano derecho (\emph{fase no mínima}). El caso clásico es el nivel de una caldera: más agua fría hace colapsar burbujas de vapor y el nivel primero baja.
\end{description}
Los procesos de primer orden e integradores son los más fáciles de controlar, porque la PV responde de inmediato a la CO; los de tiempo muerto y respuesta inversa, los más difíciles. Los oscilatorios son manejables mientras no se les pida rapidez.

\begin{figure}[ht]
\centering
\includegraphics[width=\linewidth]{cap08/f_tipos}
\caption{Respuesta escalón de los seis tipos de proceso: $1/(4s+1)$, $1/(2s+1)^4$, $1/(s(s+1))$, $1/(s^2+0.4s+1)$, $e^{-6s}/(4s+1)$ y $(1-3s)/((2s+1)(4s+1))$.}
\label{fig:ps-tipos}
\end{figure}

\subsection{Ganancia normalizada}

En los controladores industriales la PV y la CO suelen expresarse en \emph{porcentaje de su rango}, y la ganancia del controlador es adimensional. La ganancia del proceso debe estar en las mismas unidades:
\begin{equation}
  \Kproc=\frac{\Delta y/\text{rango de la PV}}{\Delta u_c/\text{rango de la CO}} .
  \label{eq:ps-ganancia}
\end{equation}

\begin{ejemplo}[label=ej:ps-ganancia]{Ganancia de un proceso de temperatura}
La CO va de $0$ a $100$\,\% y la PV se mide entre $10$ y $70\,\grados$C. La CO pasa de $20$ a $40$\,\% y la PV, de $30$ a $60\,\grados$C. Halle $\Kproc$.
\solucion $\Kproc=\dfrac{30/60}{20/100}=2.5$. En unidades físicas serían $30/20=1.5\,\grados$C/\%; la diferencia es el factor $100/60$ de los rangos.
\end{ejemplo}

\codigorepo{cap08_pid_sintonizacion/C8_Ej1_Identificacion.m}{ganancia}{ejemplo~\ref{ej:ps-ganancia}}

\begin{cuidado}
Si no se dice nada, en este capítulo las señales están normalizadas ($0$--$100$\,\%) y $\Kproc=\Delta y/\Delta u_c$. Mezclar una ganancia física con una regla de sintonización pensada para ganancias normalizadas da un $\Kc$ errado por el cociente de los rangos.
\end{cuidado}

\subsection{Modelo de primer orden con tiempo muerto}

Casi todas las reglas de sintonización resumen el proceso con tres números, los del modelo FOPDT (\emph{first order plus dead time}):
\begin{equation}
  \Gp(s)\approx\frac{\Kproc\,e^{-Ls}}{\tau s+1}.
  \label{eq:ps-fopdt}
\end{equation}
Con la curva de reacción (\cref{fig:ps-identificacion}) se obtienen así:
\begin{receta}[Método de la tangente]
\begin{enumerate}
  \item Ubique el punto de máxima pendiente de la PV y trace la tangente por él.
  \item Donde la tangente corta el valor inicial de la PV, lea $L$ (desde el instante del escalón).
  \item \textbf{Método del 63\,\%:} $\tau$ es el tiempo entre $L$ y el instante en que la PV alcanza el $63.2$\,\% de su cambio total. \textbf{Variante de Ziegler y Nichols:} $\tau$ es el tiempo entre $L$ y el corte de la tangente con el valor final.
  \item $\Kproc=\Delta y/\Delta u_c$ con \eqref{eq:ps-ganancia}.
\end{enumerate}
\end{receta}

\begin{figure}[ht]
\centering
\includegraphics[width=0.85\linewidth]{cap08/f_identificacion}
\caption{Ejemplo~\ref{ej:ps-ident}: tangente de máxima pendiente sobre la respuesta de $e^{-2s}/((8s+1)(2s+1))$. El tiempo muerto aparente ($3.0$\,s) es mayor que el retardo puro ($2$\,s), y $\tau_{ZN}>\tau_{63}$.}
\label{fig:ps-identificacion}
\end{figure}

\begin{ejemplo}[label=ej:ps-ident]{Un segundo orden visto como primer orden}
Identifique un FOPDT para $\Gp(s)=e^{-2s}/((8s+1)(2s+1))$ con los dos métodos.
\solucion La pendiente máxima es $0.079$ por segundo, en $t=5.7$\,s, con la PV en $0.213$; la tangente corta el cero en $L\approx3.0$\,s. El $63.2$\,\% se alcanza en $t=12.3$\,s, así que $\tau_{63}\approx9.3$\,s, mientras que la tangente llega al valor final $\tau_{ZN}\approx12.7$\,s después de $L$. El retardo aparente absorbe el polo rápido ($2$\,s de retardo puro más parte de la constante de $2$\,s), y la variante de Ziegler y Nichols da una constante de tiempo mayor.
\end{ejemplo}

\codigorepo{cap08_pid_sintonizacion/C8_Ej1_Identificacion.m}{tangente}{identificación con la función \cmd{identificarFOPDT}}

Dos observaciones. Primera, en un proceso de orden superior el tiempo muerto identificado suele ser mayor que el retardo puro: incluye el efecto de los polos rápidos. Segunda, la variante de Ziegler y Nichols tiende a sobrestimar $\tau$; las reglas posteriores a la suya usan el método del $63$\,\%.

\begin{cuidado}[La pendiente máxima de un primer orden]
En un proceso de primer orden, $y(t)=\Kproc\bigl(1-e^{-(t-L)/\tau}\bigr)$ tiene pendiente $\frac{\Kproc}{\tau}e^{-(t-L)/\tau}$, que es \emph{máxima al comienzo}, en $t=L$, y no en el $63$\,\% (allí vale solo el $37$\,\% del máximo). Por eso en un primer orden puro la tangente arranca justo en $L$ y los dos métodos dan el mismo $\tau$.
\end{cuidado}

\begin{pausa}
Con un primer orden puro sin retardo, $1/(5s+1)$, ¿qué $L$ y qué $\tau$ dan la tangente de Ziegler y Nichols y el método del 63\,\%?
\respuesta{$L=0$ en los dos. La tangente en $t=0$ tiene pendiente $1/5$ y llega a $1$ en $t=5$: $\tau_{ZN}=5$\,s. El $63.2$\,\% se alcanza en $t=5$\,s: $\tau_{63}=5$\,s.}
\end{pausa}

\subsection{Procesos integradores}

En un integrador lo difícil es empezar: hay que encontrar una CO que deje la PV quieta (el caudal de salida igual al de entrada, por ejemplo). Desde ahí, un escalón $\Delta u_c$ produce, tras un tiempo muerto, una rampa de pendiente $v$ (\cref{fig:ps-integrador}). Como la ganancia estática es infinita, se usa la \textbf{ganancia de velocidad}
\begin{equation}
  K_v=\frac{v}{\Delta u_c},\qquad \Gp(s)\approx\frac{K_v\,e^{-Ls}}{s},
  \label{eq:ps-kv}
\end{equation}
con $L$ medido donde la asíntota de la rampa corta el nivel inicial.

\begin{figure}[ht]
\centering
\includegraphics[width=0.6\linewidth]{cap08/f_integrador}
\caption{Respuesta de un proceso integrador con tiempo muerto: rampa de pendiente $v$ tras un tiempo muerto $L$.}
\label{fig:ps-integrador}
\end{figure}

%% ---------------------------------------------------------------------
\section{Reglas de malla abierta}
\label{sec:ps-malla-abierta}
%% ---------------------------------------------------------------------

Un buen PID debe dar un lazo estable y robusto, rechazar las perturbaciones, seguir la referencia con rapidez y suavidad, y no dejar error en estado estacionario. Las reglas de malla abierta buscan eso a partir del modelo \eqref{eq:ps-fopdt}: se hace la prueba escalón, se identifican $\Kproc$, $L$ y $\tau$, y una tabla da $\Kc$, $\Ti$ y $\Td$ del PID ideal $\Kc\bigl(1+\frac{1}{\Ti s}+\Td s\bigr)$. Todas dependen de la \textbf{razón de controlabilidad} $L/\tau$: cuanto mayor, más difícil el proceso.

\subsection{Ziegler y Nichols}

\begin{historia}[Ziegler y Nichols, 1942]
John G. Ziegler y Nathaniel B. Nichols, ingenieros de la Taylor Instruments, publicaron en 1942 dos métodos para ajustar los controladores que vendía su empresa: uno a partir de la curva de reacción en lazo abierto y otro a partir de la oscilación sostenida en lazo cerrado~\citep{ZiNi42}. Antes hubo reglas sueltas (Callender y colaboradores, 1936), pero fueron las suyas las que se volvieron estándar. Su criterio fue empírico: que la respuesta a una perturbación decayera con una razón de un cuarto entre picos sucesivos.
\end{historia}

La regla supone una curva de reacción en forma de S, es decir, un proceso sin integradores ni polos dominantes complejos; con otra forma de respuesta no debe usarse. Mide $\tau$ con la tangente hasta el $100$\,\% (\cref{sec:ps-identificacion}).

\begin{table}[ht]
\centering
\caption{Ziegler-Nichols de malla abierta para $\Kproc e^{-Ls}/(\tau s+1)$ y para el integrador $K_v e^{-Ls}/s$. La fila serie es la misma sintonía escrita en forma serie (capítulo~\ref{cap:pn}).}
\label{tab:ps-zn-abierta}
\renewcommand{\arraystretch}{1.35}
\begin{tabular}{@{}lcccccc@{}}
\toprule
 & \multicolumn{3}{c}{FOPDT} & \multicolumn{3}{c}{Integrador}\\
\cmidrule(lr){2-4}\cmidrule(l){5-7}
 & $\Kc$ & $\Ti$ & $\Td$ & $\Kc$ & $\Ti$ & $\Td$\\
\midrule
P & $\dfrac{\tau}{\Kproc L}$ & --- & --- & $\dfrac{1}{K_vL}$ & --- & ---\\
PI & $\dfrac{0.9\tau}{\Kproc L}$ & $3.33L$ & --- & $\dfrac{0.9}{K_vL}$ & $3L$ & ---\\
PID ideal & $\dfrac{1.2\tau}{\Kproc L}$ & $2L$ & $0.5L$ & $\dfrac{1.2}{K_vL}$ & $2L$ & $0.5L$\\
PID serie & $\dfrac{0.6\tau}{\Kproc L}$ & $L$ & $L$ & $\dfrac{0.6}{K_vL}$ & $L$ & $L$\\
\bottomrule
\end{tabular}
\end{table}

Note que $\Kc$ crece como $\tau/L$: con un tiempo muerto pequeño frente a la constante de tiempo, la regla pide ganancias enormes y tiempos integrales cortos, y el lazo oscila. Por eso no se recomienda para $L/\tau$ pequeño.

\subsection{Cohen y Coon}

Cohen y Coon partieron del mismo modelo, pero ajustaron las fórmulas para ubicar los polos dominantes del lazo con razón de decaimiento de un cuarto en un rango amplio de $L/\tau$~\citep{CoCo53}. Geraldine Coon verificó el método con simulaciones en un computador analógico. Una ventaja es que incluye un PD. Con $r=L/\tau$:

\begin{table}[ht]
\centering
\caption{Cohen-Coon, con $r=L/\tau$.}
\label{tab:ps-cc}
\renewcommand{\arraystretch}{1.5}
\begin{tabular}{@{}lccc@{}}
\toprule
 & $\Kc$ & $\Ti$ & $\Td$\\
\midrule
P & $\dfrac{\tau}{\Kproc L}\Bigl(1+\dfrac{r}{3}\Bigr)$ & --- & ---\\
PI & $\dfrac{\tau}{\Kproc L}\Bigl(0.9+\dfrac{r}{12}\Bigr)$ & $L\,\dfrac{30+3r}{9+20r}$ & ---\\
PD & $\dfrac{\tau}{\Kproc L}\Bigl(\dfrac54+\dfrac{r}{6}\Bigr)$ & --- & $L\,\dfrac{6-2r}{22+3r}$\\
PID & $\dfrac{\tau}{\Kproc L}\Bigl(\dfrac43+\dfrac{r}{4}\Bigr)$ & $L\,\dfrac{32+6r}{13+8r}$ & $\dfrac{4L}{11+2r}$\\
\bottomrule
\end{tabular}
\end{table}

El precio de ser tan agresivo es la fragilidad: un cambio pequeño en el proceso puede llevar el lazo a la inestabilidad. La \cref{sec:ps-robustez} lo cuantifica.

\subsection{Coeficiente de ajustabilidad}

Como Ziegler-Nichols y Cohen-Coon dan respuestas muy oscilatorias, en la práctica se usan versiones más suaves. Una regla empírica de instrumentistas usa el \textbf{coeficiente de ajustabilidad} $r_c=L/\tau$ (\cref{tab:ps-ajust}); se recomienda para $0.1<r_c<1$, y para $r_c$ grandes un PID solo no basta. La tabla no tiene una derivación publicada que la respalde: tómela como una regla de taller que da resultados conservadores, no como un resultado.

\begin{table}[H]
\centering
\caption{Regla del coeficiente de ajustabilidad, $r_c=L/\tau$.}
\label{tab:ps-ajust}
\renewcommand{\arraystretch}{1.5}
\begin{tabular}{@{}lccc@{}}
\toprule
$r_c$ & $\Kc$ & $\Ti$ & $\Td$\\
\midrule
$0$ a $0.1$ & $5/\Kproc$ & $\tau$ & $0$\\
$0.1$ a $0.2$ & $\dfrac{0.5}{\Kproc r_c}$ & $\tau$ & $0$\\
$0.2$ a $0.5$ & $\dfrac{0.5(1+0.5r_c)}{\Kproc r_c}$ & $\tau(1+0.5r_c)$ & $\tau\,\dfrac{0.5r_c}{0.5r_c+1}$\\
\bottomrule
\end{tabular}
\end{table}

\begin{ejemplo}[label=ej:ps-abierta]{Tres reglas sobre el mismo proceso}
Un proceso se identificó como $2e^{-2s}/(10s+1)$. Sintonice un PID ideal con las tres reglas y compare la respuesta a un escalón de referencia y a un escalón de perturbación a la entrada del proceso. Use la estructura PI-D (derivada solo sobre la medida) y filtro derivativo con $N=10$.
\solucion Con $L/\tau=0.2$:
\begin{center}
\renewcommand{\arraystretch}{1.2}
\begin{tabular}{@{}lccccccc@{}}
\toprule
Regla & $\Kc$ & $\Ti$ (s) & $\Td$ (s) & $\OS$ & $\ts$ (s) & pico pert. & IAE pert.\\
\midrule
Ziegler-Nichols & $3$ & $4$ & $1$ & $64$\,\% & $19.7$ & $0.38$ & $1.6$\\
Cohen-Coon & $3.46$ & $4.55$ & $0.70$ & $84$\,\% & $24.2$ & $0.39$ & $1.6$\\
Ajustabilidad & $1.38$ & $11$ & $0.91$ & $2$\,\% & $20.4$ & $0.48$ & $8.0$\\
\bottomrule
\end{tabular}
\end{center}
Ziegler-Nichols y Cohen-Coon responden rápido pero con sobrepicos enormes; la regla de ajustabilidad casi no sobrepasa, pero deja crecer más la perturbación y tarda mucho más en eliminarla: su IAE ante la perturbación es cinco veces mayor, porque $\Ti$ es casi tres veces más largo (\cref{fig:ps-abierta}). Ante la perturbación, Ziegler-Nichols decae con una razón cercana a $0.19$ entre picos sucesivos: cerca del cuarto que buscaba.
\end{ejemplo}

\begin{figure}[ht]
\centering
\includegraphics[width=\linewidth]{cap08/f_malla_abierta}
\caption{Ejemplo~\ref{ej:ps-abierta}: tres sintonías de malla abierta sobre $2e^{-2s}/(10s+1)$.}
\label{fig:ps-abierta}
\end{figure}

\codigorepo{cap08_pid_sintonizacion/C8_Ej2_MallaAbierta.m}{reglas}{las tres reglas con la función \cmd{reglasPID}}
\codigorepo{cap08_pid_sintonizacion/C8_Ej2_MallaAbierta.m}{simulacion}{referencia y perturbación con \cmd{simularPID}}

\begin{ideaclave}
\textbf{Seguir la referencia o rechazar perturbaciones.}\enspace Ziegler-Nichols y Cohen-Coon se diseñaron para rechazar perturbaciones, y en eso son buenos; ante cambios de referencia sobrepasan demasiado. Una regla conservadora hace lo contrario. No hay una sintonía mejor en todo: hay que decidir para qué se sintoniza.
\end{ideaclave}

\begin{pausa}
En el ejemplo~\ref{ej:ps-abierta}, ¿qué pasa con el $\Kc$ de Ziegler-Nichols si el tiempo muerto fuera de $0.5$\,s en lugar de $2$\,s?
\respuesta{Se multiplica por cuatro: $\Kc=1.2\cdot10/(2\cdot0.5)=12$, y $\Ti=1$\,s. Un $L/\tau$ pequeño lleva a ganancias muy altas.}
\end{pausa}

%% ---------------------------------------------------------------------
\section{Reglas de malla cerrada}
\label{sec:ps-malla-cerrada}
%% ---------------------------------------------------------------------

Las reglas de malla abierta exigen sacar el lazo a manual, y eso no siempre es posible: un proceso inestable en lazo abierto (un reactor exotérmico, ciertos niveles) se escapa durante la prueba. Las de malla cerrada dejan el controlador en \emph{automático} y caracterizan el proceso con dos números medidos en el propio lazo.

\subsection{Ajuste manual}

Antes de las reglas, el método que se usa en campo: ensayo y error con una secuencia ordenada. Es lento y la acción derivativa es difícil de ajustar así, pero no necesita modelo y se inspira en la misma idea que la regla de Ziegler y Nichols.
\begin{receta}[Ajuste manual en lazo cerrado]
\begin{enumerate}
  \item Arranque con $\Kc$ y $\Td$ pequeños y $\Ti$ grande.
  \item Duplique $\Kc$ y observe la respuesta, hasta que aparezcan oscilaciones sostenidas con $\Kc=K_{cu}$. Deje $\Kc=K_{cu}/2$.
  \item Reduzca $\Ti$ a la mitad sucesivamente hasta que vuelva a oscilar, con $\Ti=\Ti^*$. Deje $\Ti=2\Ti^*$.
  \item Aumente $\Td$ hasta que la respuesta oscile, con $\Td=\Td^*$. Deje $\Td=\Td^*/3$.
\end{enumerate}
\end{receta}

\subsection{Ziegler y Nichols de malla cerrada}

Con el controlador en solo proporcional ($\Ti\to\infty$, $\Td=0$), se sube $\Kc$ desde cero hasta que la salida oscila sin crecer ni decaer. Esa es la \textbf{ganancia crítica} $\Kcr$, y el periodo de la oscilación es el \textbf{periodo crítico} $\Pcr$. En ese punto dos polos del lazo cerrado están sobre el eje imaginario, en $\pm j\wcr$ con $\wcr=2\pi/\Pcr$. Si el lazo nunca oscila al subir la ganancia (un primer o segundo orden sin retardo, por ejemplo), el método no se aplica. Con un modelo, $\Kcr$ y $\wcr$ salen de sustituir $s=\jw$ en la ecuación característica, o de la fila nula de Routh (capítulo~\ref{cap:es}).

\begin{table}[ht]
\centering
\caption{Reglas de malla cerrada: Ziegler-Nichols~\citep{ZiNi42} y Tyreus-Luyben~\citep{TyLu92}.}
\label{tab:ps-cerrada}
\renewcommand{\arraystretch}{1.3}
\begin{tabular}{@{}llccc@{}}
\toprule
Regla & & $\Kc$ & $\Ti$ & $\Td$\\
\midrule
Ziegler-Nichols & P & $0.5\Kcr$ & --- & ---\\
 & PI & $0.45\Kcr$ & $\Pcr/1.2$ & ---\\
 & PID ideal & $0.6\Kcr$ & $0.5\Pcr$ & $0.125\Pcr$\\
 & PID serie & $0.3\Kcr$ & $0.25\Pcr$ & $0.25\Pcr$\\
\midrule
Tyreus-Luyben & PI & $\Kcr/3.2$ & $2.2\Pcr$ & ---\\
 & PID ideal & $\Kcr/2.2$ & $2.2\Pcr$ & $\Pcr/6.3$\\
\bottomrule
\end{tabular}
\end{table}

\begin{ejemplo}[label=ej:ps-critico]{Ganancia y periodo críticos}
Halle $\Kcr$ y $\Pcr$ para $\Gp(s)=1/((s+1)(0.5s+1)(0.2s+1))$ con control proporcional.
\solucion La ecuación característica es $0.1s^3+0.8s^2+1.7s+1+\Kc=0$, y Routh da estabilidad para $-1<\Kc<12.6$. Con $s=\jw$:
\[
  \underbrace{-0.8\omega^2+1+\Kc}_{\text{real}}+j\,\underbrace{\bigl(-0.1\omega^3+1.7\omega\bigr)}_{\text{imaginaria}}=0 .
\]
La parte imaginaria da $\wcr=\sqrt{17}=4.12$\,rad/s, y la real, $\Kcr=0.8\cdot17-1=12.6$. Entonces $\Pcr=2\pi/\sqrt{17}=1.52$\,s.
\end{ejemplo}

La meta de Ziegler y Nichols era la \textbf{razón de decaimiento de un cuarto}: que cada pico de la oscilación fuera la cuarta parte del anterior. Es una respuesta bastante oscilatoria, con sobrepicos de $40$ a $60$\,\% ante la referencia.

\subsection{Tyreus y Luyben}

Tyreus y Luyben propusieron sus reglas para procesos con integrador y tiempo muerto, donde Ziegler-Nichols resulta demasiado agresivo~\citep{TyLu92}. Usan los mismos $\Kcr$ y $\Pcr$, pero con menos ganancia y tiempos integrales mucho más largos: dan respuestas más amortiguadas y menos sensibles a cambios en el proceso. \alerta{No buscan la razón de un cuarto}: son deliberadamente más conservadoras.

\begin{ejemplo}[label=ej:ps-zntl]{Ziegler-Nichols contra Tyreus-Luyben}
Para $\Gp(s)=1/((3s+1)(2s+1)(s+1))$, halle $\Kcr$ y $\Pcr$, sintonice un PID ideal con las dos reglas y compare.
\solucion La ecuación característica es $6s^3+11s^2+6s+1+\Kc=0$. La parte imaginaria, $-6\omega^3+6\omega=0$, da $\wcr=1$\,rad/s y $\Pcr=2\pi=6.28$\,s; la real, $\Kcr=11-1=10$. Las reglas dan
\begin{center}
\renewcommand{\arraystretch}{1.2}
\begin{tabular}{@{}lcccccc@{}}
\toprule
Regla & $\Kc$ & $\Ti$ (s) & $\Td$ (s) & $\OS$ & $\ts$ (s) & decaimiento\\
\midrule
Ziegler-Nichols & $6$ & $3.14$ & $0.785$ & $46.5$\,\% & $16.3$ & $0.17$\\
Tyreus-Luyben & $4.55$ & $13.8$ & $1.00$ & $9.0$\,\% & $30.0$ & sin segundo pico\\
\bottomrule
\end{tabular}
\end{center}
(simulación con filtro derivativo $N=10$). Ziegler-Nichols oscila con una razón de $0.17$, algo mejor que un cuarto; Tyreus-Luyben casi no sobrepasa, a cambio de un establecimiento más lento (\cref{fig:ps-cerrada}). Si la ganancia del proceso sube un $50$\,\% sin retocar el controlador, el sobrepico pasa a $55$\,\% con Ziegler-Nichols y a $22$\,\% con Tyreus-Luyben.
\end{ejemplo}

\begin{figure}[ht]
\centering
\includegraphics[width=\linewidth]{cap08/f_malla_cerrada}
\caption{Ejemplo~\ref{ej:ps-zntl}. Izquierda: lazo proporcional por debajo, en y por encima de $\Kcr$. Derecha: PID ideal sintonizado por las dos reglas.}
\label{fig:ps-cerrada}
\end{figure}

\codigorepo{cap08_pid_sintonizacion/C8_Ej3_MallaCerrada.m}{critico}{$\Kcr$ y $\Pcr$ por sustitución y con \cmd{tablaRouth}}
\codigorepo{cap08_pid_sintonizacion/C8_Ej3_MallaCerrada.m}{zntl}{las dos sintonías y la razón de decaimiento}

\begin{cuidado}[Llevar un lazo real al límite]
La prueba de malla cerrada lleva el proceso al borde de la inestabilidad. En una planta real eso puede ser inaceptable (o peligroso), y el actuador puede saturar antes de que la oscilación sea limpia. Por eso se prefiere la variante con relé de Åström y Hägglund, que produce una oscilación controlada de amplitud pequeña~\citep{AsHa84}.
\end{cuidado}

\begin{pausa}
Para el ejemplo~\ref{ej:ps-zntl}, ¿qué PI daría Ziegler-Nichols de malla cerrada?
\respuesta{$\Kc=0.45\cdot10=4.5$ y $\Ti=2\pi/1.2=5.24$\,s.}
\end{pausa}

%% ---------------------------------------------------------------------
\section{Síntesis directa}
\label{sec:ps-sintesis}
%% ---------------------------------------------------------------------

Las reglas anteriores fijan la estructura (un PID) y ajustan sus tres números. La síntesis directa invierte el razonamiento: se escoge la respuesta de lazo cerrado que se quiere, $T(s)$, y se despeja el controlador que la produce. De $T=\Gc\Gp/(1+\Gc\Gp)$,
\begin{equation}
  \Gc(s)=\frac{1}{\Gp(s)}\,\frac{T(s)}{1-T(s)} ,
  \label{eq:ps-sintesis}
\end{equation}
la \clave{ecuación de síntesis}. El controlador contiene el inverso del modelo: cancela la dinámica del proceso y la reemplaza por la deseada. Nada garantiza que el resultado sea implementable, y lo difícil es escoger $T(s)$. Debe dar error nulo en estado estacionario ($T(0)=1$), ser rápida sin sobrepasar mucho, y ser simple. Las elecciones usuales son
\begin{equation}
  T(s)=\frac{1}{\lambda s+1}
  \qquad\text{o}\qquad
  T(s)=\frac{1}{(\lambda_1 s+1)(\lambda_2 s+1)},
  \label{eq:ps-T-deseada}
\end{equation}
que llevadas a \eqref{eq:ps-sintesis} dan
\[
  \Gc(s)=\frac{1}{\Gp(s)}\,\frac{1}{\lambda s}
  \qquad\text{y}\qquad
  \Gc(s)=\frac{1}{\Gp(s)}\,\frac{1}{s\bigl(\lambda_1\lambda_2 s+\lambda_1+\lambda_2\bigr)} .
\]
En los dos aparece un integrador: por eso el error ante el escalón es nulo.

\begin{ideaclave}
\textbf{$\lambda$ es la constante de tiempo del lazo cerrado.}\enspace Con $T=1/(\lambda s+1)$, el lazo responde como un primer orden de constante $\lambda$: se establece en unos $4\lambda$. Un $\lambda$ \emph{pequeño} da una respuesta \emph{rápida}, a cambio de una ganancia del controlador grande y una acción de control mayor; un $\lambda$ grande da un lazo lento y suave. Es un único parámetro de ajuste con significado físico.
\end{ideaclave}

\paragraph{Proceso estático.} Si $\Gp=\Kproc$, $\Gc=1/(\lambda\Kproc s)$: un controlador integral puro.

\paragraph{Primer orden.} Si $\Gp=\Kproc/(\tau s+1)$,
\[
  \Gc(s)=\frac{\tau s+1}{\Kproc\lambda s}=\frac{\tau}{\Kproc\lambda}\left(1+\frac{1}{\tau s}\right):
  \qquad \Kc=\frac{\tau}{\Kproc\lambda},\quad \Ti=\tau ,
\]
un PI cuyo cero cancela el polo del proceso.

\paragraph{Segundo orden.} Si $\Gp=\Kproc/((\tau_1s+1)(\tau_2s+1))$,
\[
  \Gc(s)=\frac{(\tau_1s+1)(\tau_2s+1)}{\Kproc\lambda s}:
  \qquad \Kc=\frac{\tau_1+\tau_2}{\Kproc\lambda},\quad \Ti=\tau_1+\tau_2,\quad \Td=\frac{\tau_1\tau_2}{\tau_1+\tau_2},
\]
un PID ideal. En la práctica se le agrega el filtro derivativo.

\begin{ejemplo}[label=ej:ps-sintesis]{Un PID por síntesis directa}
Para $\Gp(s)=3/((2s+1)(10s+1))$ y $T(s)=1/(\lambda s+1)$, halle el PID y evalúe $\lambda=1$, $3$ y $5$\,s.
\solucion $\Kc=12/(3\lambda)=4/\lambda$, $\Ti=12$\,s y $\Td=20/12=1.67$\,s. Con filtro $N=10$, la simulación da
\begin{center}
\renewcommand{\arraystretch}{1.2}
\begin{tabular}{@{}cccc@{}}
\toprule
$\lambda$ (s) & $\Kc$ & $\ts$ (s) & $|u_c|_{\max}$\\
\midrule
$1$ & $4$ & $3.9$ & $43.8$\\
$3$ & $1.33$ & $12.1$ & $14.6$\\
$5$ & $0.8$ & $19.8$ & $8.8$\\
\bottomrule
\end{tabular}
\end{center}
El tiempo de establecimiento es casi exactamente $4\lambda$, y la acción de control crece como $1/\lambda$.
\end{ejemplo}

\codigorepo{cap08_pid_sintonizacion/C8_Ej4_SintesisDirecta.m}{pi}{síntesis directa de un PI}
\codigorepo{cap08_pid_sintonizacion/C8_Ej4_SintesisDirecta.m}{pid}{ejemplo~\ref{ej:ps-sintesis}}

\subsection{Procesos de fase no mínima}

Un cero en el semiplano derecho no puede invertirse: se convertiría en un polo inestable del controlador. Lo mismo vale para el tiempo muerto, que es de fase no mínima (su aproximación de Padé de la \cref{sec:es-retardo} tiene un cero en $+2/L$). La salida es dejar ese factor en la respuesta deseada: lo que el proceso no permite cambiar debe aparecer en $T(s)$.

\begin{ejemplo}[label=ej:ps-nomin]{Síntesis directa con un cero inestable}
Para $\Gp(s)=(1-9s)/((15s+1)(3s+1))$, diseñe por síntesis directa.
\solucion Con $T=1/(\lambda s+1)$,
\[
  \Gc(s)=\frac{(15s+1)(3s+1)}{\lambda s(1-9s)},
\]
que tiene un polo en $s=+1/9$: no es implementable. Si se conserva el cero en la respuesta deseada y se sube el orden para mantener el grado relativo, $T(s)=(1-9s)/(\lambda s+1)^2$, entonces
\[
  \Gc(s)=\frac{45s^2+18s+1}{s\bigl(\lambda^2s+2\lambda+9\bigr)}
  =\Kc\left(1+\frac{1}{\Ti s}+\Td s\right)\frac{1}{\tau_f s+1},
\]
un PID ideal en serie con un filtro de primer orden, con
\[
  \Kc=\frac{18}{2\lambda+9},\qquad \Ti=18\text{ s},\qquad \Td=2.5\text{ s},\qquad \tau_f=\frac{\lambda^2}{2\lambda+9}.
\]
Para $\lambda=5$: $\Kc=0.947$ y $\tau_f=1.32$\,s. La respuesta inversa sigue ahí (\cref{fig:ps-nomin}): ningún controlador lineal estable la elimina.
\end{ejemplo}

\begin{figure}[ht]
\centering
\includegraphics[width=\linewidth]{cap08/f_sintesis_nomin}
\caption{Ejemplo~\ref{ej:ps-nomin}: salida y acción de control para tres valores de $\lambda$. Un $\lambda$ menor acelera la respuesta pero profundiza el arranque en sentido contrario y exige más acción de control.}
\label{fig:ps-nomin}
\end{figure}

\codigorepo{cap08_pid_sintonizacion/C8_Ej4_SintesisDirecta.m}{nomin}{ejemplo~\ref{ej:ps-nomin}}

\begin{pausa}
¿Por qué, en el ejemplo~\ref{ej:ps-nomin}, la respuesta deseada tiene denominador de segundo orden y no de primero?
\respuesta{Para conservar el grado relativo. $\Gp$ tiene grado relativo $1$, y un lazo cerrado no puede tener menos. Con $T=(1-9s)/(\lambda s+1)$, de grado relativo $0$, saldría $\Gc=(15s+1)(3s+1)/((\lambda+9)s)$: impropio, un PID ideal sin filtro que pediría a la salida saltar instantáneamente.}
\end{pausa}

%% ---------------------------------------------------------------------
\section{Control por modelo interno}
\label{sec:ps-imc}
%% ---------------------------------------------------------------------

La síntesis directa usa un modelo del proceso para diseñar. El \clave{control por modelo interno} (IMC, \emph{internal model control}) va más lejos: mete el modelo \emph{dentro} del controlador, y con eso maneja de forma explícita los errores de modelado y las perturbaciones~\citep{GaMo82,RMS86}. Para ciertos modelos, el resultado es un PID con un solo parámetro de ajuste.

\subsection{Control basado en modelo, en lazo abierto}

Supongamos un proceso estable $\Gp$ y un controlador $Q(s)$ en lazo abierto: $Y=Q\Gp R$. Con $\Gp=\Kproc/(\tau s+1)$, el controlador estático $Q=1/\Kproc$ lleva la salida al valor de la referencia, pero con la dinámica del proceso. Para acelerar habría que invertir el proceso, $Q=1/\Gp$, y la salida seguiría a la referencia de manera perfecta. Pero
\[
  Q(s)=\frac{\tau s+1}{\Kproc}
  \quad\Longrightarrow\quad
  U(s)=\frac{1}{\Kproc}(\tau s+1)R(s),
  \qquad
  u(t)=\frac{\tau}{\Kproc}\frac{\dif r}{\dif t}+\frac{r(t)}{\Kproc},
\]
que es impropio: ante un escalón de referencia de tamaño $A$ pide un impulso $\frac{A\tau}{\Kproc}\impulso(t)$. La solución es un filtro que haga $Q$ propio,
\begin{equation}
  F(s)=\frac{1}{\lambda s+1},\qquad Q(s)=\frac{F(s)}{\Gp(s)}=\frac{1}{\Kproc}\,\frac{\tau s+1}{\lambda s+1},
  \label{eq:ps-q-filtro}
\end{equation}
un adelanto-atraso. Ahora $Y=F R$: un primer orden de constante $\lambda$, más rápido que el proceso si $\lambda<\tau$. Como en síntesis directa, $\lambda$ es la constante de tiempo deseada.

\subsection{Estabilidad interna}

¿Se puede invertir cualquier proceso? Con $\Gp=(1-9s)/((15s+1)(3s+1))$ y $\lambda=5$, compare
\[
  Q_1(s)=\frac{(15s+1)(3s+1)}{(1-9s)(5s+1)}
  \qquad\text{y}\qquad
  Q_2(s)=\frac{(15s+1)(3s+1)}{(5s+1)^2} .
\]
$Q_1$ invierte todo el proceso y $\Gp Q_1=1/(5s+1)$: una salida impecable. Pero $Q_1$ tiene un polo en $s=+1/9$, y la acción de control crece como $e^{t/9}$: supera $50$ unidades en $t=28$\,s y diverge (\cref{fig:ps-interna}). La salida no lo muestra porque el cero del proceso cancela exactamente ese polo. $Q_2$ deja el cero sin invertir; la salida conserva la respuesta inversa, $Y=(1-9s)/(5s+1)^2\,R$, pero todas las señales están acotadas.

\begin{figure}[ht]
\centering
\includegraphics[width=\linewidth]{cap08/f_imc_interna}
\caption{Estabilidad interna: si $Q$ invierte el cero en $+1/9$, la salida es perfecta mientras la acción de control diverge.}
\label{fig:ps-interna}
\end{figure}

\begin{definicion}[label=def:ps-interna]{Estabilidad interna}
Un sistema de control es \textbf{internamente estable} si toda señal del lazo (salida, acción de control, error, señales internas) permanece acotada ante cualquier entrada acotada aplicada en cualquier punto. Revisar solo $Y/R$ no basta: una cancelación entre un polo inestable y un cero inestable la esconde.
\end{definicion}

Por eso el modelo se factoriza en una parte no invertible y una invertible,
\begin{equation}
  \tilde G_p(s)=\tilde G_p^+(s)\,\tilde G_p^-(s),
  \label{eq:ps-factorizacion}
\end{equation}
donde $\tilde G_p^+$ contiene los ceros en el semiplano derecho y el tiempo muerto, con $\tilde G_p^+(0)=1$. Hay dos formas: la \textbf{simple}, que deja en $\tilde G_p^+$ solo el factor $(1-as)$; y la \textbf{pasa-todo} (\emph{all-pass}), que lo divide por su espejo, $\tilde G_p^+=(1-as)/(1+as)$, de modo que $|\tilde G_p^+(\jw)|=1$ y $\tilde G_p^-$ gana un polo estable. Solo $\tilde G_p^-$ se invierte:
\begin{equation}
  Q(s)=\bigl(\tilde G_p^-(s)\bigr)^{-1}F(s),
  \label{eq:ps-q}
\end{equation}
con el orden del filtro escogido para que $Q$ sea propio.

\subsection{La estructura IMC}

En lazo abierto, cualquier error del modelo pasa directo a la salida. Si el proceso real es $2/(8s+1)$ y se diseña con el modelo $1/(10s+1)$, $Q=(10s+1)/(\lambda s+1)$ y la salida se va a $2$ en lugar de $1$. Y una perturbación no se corrige nunca. La estructura IMC (\cref{fig:ps-imc}) corre el modelo $\tilde G_p$ en paralelo con el proceso y retroalimenta solo la diferencia, $\tilde D=Y-\tilde G_pU$: lo que el modelo no explica, sea perturbación o error de modelado. De la figura,
\begin{equation}
  Y=\frac{\Gp Q}{1+Q(\Gp-\tilde G_p)}\,R+\frac{1-\tilde G_pQ}{1+Q(\Gp-\tilde G_p)}\,D .
  \label{eq:ps-imc-lazo}
\end{equation}
Con un modelo perfecto, $\Gp=\tilde G_p$, el denominador es $1$: $Y=\Gp QR+(1-\Gp Q)D$. El lazo es estable si $Q$ y $\Gp$ lo son, y $Q(0)=1/\tilde G_p(0)$ asegura que la perturbación constante se rechaza por completo. Con modelo imperfecto, la estabilidad depende de $1+Q(\Gp-\tilde G_p)=0$, y $\lambda$ se vuelve una perilla de robustez.

\begin{figure}[ht]
\centering
\includegraphics[width=0.85\linewidth]{cap08/imc}
\caption{Estructura IMC. Solo vuelve a la entrada lo que el modelo no explica.}
\label{fig:ps-imc}
\end{figure}

\begin{receta}[Diseño IMC]
\begin{enumerate}
  \item Obtenga un modelo nominal $\tilde G_p$ y factorícelo como en \eqref{eq:ps-factorizacion}.
  \item Forme $Q=(\tilde G_p^-)^{-1}F$ con un filtro que lo haga propio:
  \begin{itemize}
    \item cambios escalón de referencia: $F=1/(\lambda s+1)^n$;
    \item referencias tipo rampa: $F=(n\lambda s+1)/(\lambda s+1)^n$;
    \item rechazo de perturbaciones, procesos integradores o inestables: $F=(\gamma s+1)/(\lambda s+1)^n$, con $\gamma$ escogido para cancelar el polo más lento o inestable.
  \end{itemize}
  \item Ajuste $\lambda$ revisando la estabilidad con el proceso real o con los peores modelos posibles.
\end{enumerate}
\end{receta}

\begin{ejemplo}[label=ej:ps-bequette]{El modelo no coincide con el proceso}
El proceso real es $\Gp=1.5(1-s)/((s+1)(4s+1))$, pero el modelo es $\tilde G_p=2/(5s+1)$ (ejemplo de Bequette~\citep{Beq03}). ¿Para qué $\lambda$ es estable el lazo IMC?
\solucion Con $Q=(5s+1)/(2(\lambda s+1))$, la ecuación $1+Q(\Gp-\tilde G_p)=0$ da, multiplicando por el denominador común,
\[
  8\lambda s^3+(10\lambda-7.5)s^2+(2\lambda+6)s+1.5=0 .
\]
Routh exige $\lambda>0.75$ por el segundo coeficiente y, por la fila $s^1$, $(10\lambda-7.5)(2\lambda+6)>12\lambda$, es decir, $20\lambda^2+33\lambda-45>0$: \alerta{$\lambda>0.887$}. La \cref{fig:ps-robusto} lo confirma: con $\lambda=0.85$ el lazo es inestable; con $1.5$ responde rápido con algo de oscilación; con $3$, suave.
\end{ejemplo}

\begin{figure}[ht]
\centering
\includegraphics[width=0.85\linewidth]{cap08/f_imc_robusto}
\caption{Ejemplo~\ref{ej:ps-bequette}: el $\lambda$ del filtro decide si el error de modelado desestabiliza el lazo.}
\label{fig:ps-robusto}
\end{figure}

\codigorepo{cap08_pid_sintonizacion/C8_Ej5_IMC.m}{interna}{estabilidad interna con $Q_1$ y $Q_2$}
\codigorepo{cap08_pid_sintonizacion/C8_Ej5_IMC.m}{robusto}{ejemplo~\ref{ej:ps-bequette}}

\begin{modelo}[Modelo de Simulink: \texttt{crearModelo\_IMC}]
\texttt{codigo/cap08\_pid\_sintonizacion/crearModelo\_IMC.m} arma la \cref{fig:ps-imc} con el proceso y el modelo del ejemplo~\ref{ej:ps-bequette}, un escalón de referencia y una perturbación a la salida. Cambie \texttt{lambda} en el espacio de trabajo y vuelva a simular: por debajo de $0.887$ el lazo se inestabiliza.
\end{modelo}

\subsection{El PID que hay dentro del IMC}

Reordenando la \cref{fig:ps-imc}, el lazo IMC equivale a un lazo clásico con
\begin{equation}
  \Gc(s)=\frac{Q(s)}{1-\tilde G_p(s)Q(s)} .
  \label{eq:ps-imc-pid}
\end{equation}
Para muchos modelos sencillos, $\Gc$ resulta un PID (a veces con un filtro de primer orden): es la relación \clave{IMC-PID}~\citep{RMS86}.

\begin{receta}[IMC-PID]
\begin{enumerate}
  \item Forme $Q$ con un filtro que lo deje \emph{semipropio}: se admite que el numerador tenga un grado más, porque eso es justamente la acción derivativa. Para procesos integradores o inestables use $F=(\gamma s+1)/(\lambda s+1)^n$; si el modelo tiene un polo inestable en $s=p_u$, imponga $F(p_u)=1$.
  \item Calcule $\Gc$ con \eqref{eq:ps-imc-pid} y escríbalo como cociente de polinomios.
  \item Identifique $\Kc$, $\Ti$, $\Td$ y, si aparece, $\tau_f$ en $\Kc\Bigl(\dfrac{\Ti\Td s^2+\Ti s+1}{\Ti s}\Bigr)\dfrac{1}{\tau_f s+1}$.
  \item Ajuste $\lambda$ con simulaciones, con el modelo nominal y con modelos perturbados. Un punto de partida usual es entre un tercio y la mitad de la constante de tiempo dominante.
\end{enumerate}
\end{receta}

\paragraph{Primer orden.} Con $\tilde G_p=\Kproc/(\tau s+1)$ y $F=1/(\lambda s+1)$, $Q$ es \eqref{eq:ps-q-filtro} y
\[
  \Gc(s)=\frac{\tau s+1}{\Kproc\lambda s}:\qquad \Kc=\frac{\tau}{\Kproc\lambda},\quad\Ti=\tau ,
\]
el mismo PI de la síntesis directa: con modelo perfecto, el IMC es una síntesis directa con $T=F$. Queda un único parámetro; un $\lambda$ pequeño da un $\Kc$ grande, y al revés.

\paragraph{Segundo orden.} Con $\tilde G_p=\Kproc/((\tau_1s+1)(\tau_2s+1))$, $\Gc=(\tau_1s+1)(\tau_2s+1)/(\Kproc\lambda s)$: el PID ideal con $\Kc=(\tau_1+\tau_2)/(\Kproc\lambda)$, $\Ti=\tau_1+\tau_2$ y $\Td=\tau_1\tau_2/(\tau_1+\tau_2)$.

\paragraph{Primer orden con tiempo muerto.} El retardo no se invierte. Con la aproximación de Padé de primer orden de la \cref{sec:es-retardo}, $e^{-Ls}\approx(1-\frac{L}{2}s)/(1+\frac{L}{2}s)$, y la factorización pasa-todo,
\[
  \tilde G_p^+=\frac{1-\frac L2s}{1+\frac L2s},\qquad
  \tilde G_p^-=\frac{\Kproc}{(\tau s+1)(\frac L2s+1)},\qquad
  Q=\frac{(\tau s+1)(\frac L2s+1)}{\Kproc(\lambda s+1)} .
\]
Entonces $\Gc=\dfrac{(\tau s+1)(Ls+2)}{\Kproc(2\lambda+L)s}$, un PID ideal con
\begin{equation}
  \Kc=\frac{\tau+L/2}{\Kproc(\lambda+L/2)},\qquad \Ti=\tau+\frac L2,\qquad \Td=\frac{\tau L}{2\tau+L} .
  \label{eq:ps-imc-fopdt}
\end{equation}
Como Padé solo es buena a bajas frecuencias, $\lambda$ no puede ser arbitrariamente pequeño: Rivera, Morari y Skogestad recomiendan $\lambda>0.8L$~\citep{RMS86}.

\begin{ejemplo}[label=ej:ps-imc-fopdt]{IMC-PID con el retardo exacto}
Con $\Kproc=1.2$, $\tau=20$\,s y $L=4$\,s, evalúe $\lambda=3.2$, $6.7$ y $10$\,s simulando con el retardo exacto (estructura estándar, $N=10$).
\solucion $\Ti=22$\,s y $\Td=1.82$\,s en los tres casos; $\Kc=3.53$, $2.11$ y $1.53$. El sobrepico es $16$\,\%, $0$ y $0$, con $\ts=16.5$, $28.3$ y $41.1$\,s (\cref{fig:ps-imc-fopdt}). Bajando $\lambda$, el lazo con retardo exacto se inestabiliza hacia $\lambda\approx0.75$\,s: la regla $\lambda>0.8L=3.2$\,s deja un margen holgado.
\end{ejemplo}

\begin{figure}[ht]
\centering
\includegraphics[width=0.85\linewidth]{cap08/f_imc_fopdt}
\caption{Ejemplo~\ref{ej:ps-imc-fopdt}: IMC-PID sobre $1.2e^{-4s}/(20s+1)$ con el retardo exacto.}
\label{fig:ps-imc-fopdt}
\end{figure}

\codigorepo{cap08_pid_sintonizacion/C8_Ej5_IMC.m}{fopdt}{ejemplo~\ref{ej:ps-imc-fopdt}}

\paragraph{Proceso inestable.} Con $\tilde G_p=\Kproc/(1-\tau_u s)$, $\tau_u>0$, el polo está en $p_u=+1/\tau_u$. Un filtro $F=(\gamma s+1)/(\lambda s+1)^2$ con $F(1/\tau_u)=1$ exige
\[
  \frac{\gamma}{\tau_u}+1=\Bigl(\frac{\lambda}{\tau_u}+1\Bigr)^2
  \quad\Longrightarrow\quad
  \gamma=\lambda\Bigl(\frac{\lambda}{\tau_u}+2\Bigr),
\]
y con $Q=(1-\tau_us)(\gamma s+1)/(\Kproc(\lambda s+1)^2)$ resulta un PI:
\[
  \Kc=-\frac{\lambda+2\tau_u}{\Kproc\lambda},\qquad \Ti=\gamma .
\]
La ganancia es negativa porque, a alta frecuencia, $\Kproc/(1-\tau_us)\approx-\Kproc/(\tau_us)$ invierte el signo. Con $\Kproc=\tau_u=1$ los polos del lazo quedan dobles en $-1/\lambda$; para $\lambda=0.5$, $1$ y $2$\,s el sobrepico es $28$, $45$ y $79$\,\%: el cero del PI, en $-1/\gamma$, se acerca al origen a medida que $\lambda$ crece. Con un proceso inestable, el IMC se implementa como el lazo clásico \eqref{eq:ps-imc-pid}: la estructura de la \cref{fig:ps-imc} correría el modelo inestable en lazo abierto.

\begin{table}[ht]
\centering
\caption{Sintonías IMC-PID para el PID ideal (algunas con un filtro $1/(\tau_fs+1)$). Las filas con tiempo muerto usan aproximaciones: Taylor de primer orden, $e^{-Ls}\approx1-Ls$, para el PI, y Padé de primer orden para el PID; la fila con $\tau_3$ trunca además la serie de $1/(\tau_3s+1)$. Fuentes:~\citep{RMS86,Beq03}.}
\label{tab:ps-imc}
\renewcommand{\arraystretch}{1.7}
\begin{tabular}{@{}lccc@{}}
\toprule
Modelo $\tilde G_p$ & $\Kc\Kproc$ & $\Ti$ & $\Td$\\
\midrule
$\dfrac{\Kproc}{\tau s+1}$ & $\dfrac{\tau}{\lambda}$ & $\tau$ & ---\\
$\dfrac{\Kproc e^{-Ls}}{\tau s+1}$ (PI) & $\dfrac{\tau}{\lambda+L}$ & $\tau$ & ---\\
$\dfrac{\Kproc e^{-Ls}}{\tau s+1}$ (PID) & $\dfrac{\tau+L/2}{\lambda+L/2}$ & $\tau+\dfrac L2$ & $\dfrac{\tau L}{2\tau+L}$\\
$\dfrac{\Kproc}{(\tau_1s+1)(\tau_2s+1)}$ & $\dfrac{\tau_1+\tau_2}{\lambda}$ & $\tau_1+\tau_2$ & $\dfrac{\tau_1\tau_2}{\tau_1+\tau_2}$\\
$\dfrac{\Kproc(\tau_3s+1)e^{-Ls}}{(\tau_1s+1)(\tau_2s+1)}$ & $\dfrac{\tau_1+\tau_2-\tau_3}{\lambda+L}$ & $\tau_1+\tau_2-\tau_3$ & $\dfrac{\tau_1\tau_2}{\tau_1+\tau_2-\tau_3}-\tau_3$\\
$\dfrac{\Kproc}{\tau^2s^2+2\zeta\tau s+1}$ & $\dfrac{2\zeta\tau}{\lambda}$ & $2\zeta\tau$ & $\dfrac{\tau}{2\zeta}$\\
$\dfrac{\Kproc}{s}$ & $\dfrac{1}{\lambda}$ & --- & ---\\
$\dfrac{\Kproc e^{-Ls}}{s}$ & $\dfrac{2\lambda+L}{(\lambda+L)^2}$ & $2\lambda+L$ & ---\\
\bottomrule
\end{tabular}
\end{table}

%% ---------------------------------------------------------------------
\section{Desempeño contra robustez}
\label{sec:ps-robustez}
%% ---------------------------------------------------------------------

Una sintonía se calcula con un modelo, pero se aplica a un proceso que no es exactamente ese modelo: la ganancia cambia con el punto de operación, el tiempo muerto con el caudal, la constante de tiempo con el nivel. Una buena sintonía debe aguantar esas diferencias. Comparemos las reglas con el proceso del ejercicio de identificación, $1.2e^{-4s}/(20s+1)$, que tiene $L/\tau=0.2$.

\begin{ejemplo}[label=ej:ps-robustez]{Cuatro sintonías, nominal y perturbado}
Compare ZN, CC, ajustabilidad e IMC-PID ($\lambda=6.7$\,s) en estructura estándar con $N=10$: respuesta nominal, márgenes con el retardo exacto y estabilidad ante una malla de procesos con $\Kproc$ en $\pm30$\,\%, $\tau$ en $\pm30$\,\% y $L$ en $\pm50$\,\% (cinco valores de cada uno: 125 procesos).
\solucion Los márgenes de ganancia y de fase son los del capítulo~\ref{cap:rf} (definiciones~\ref{def:rf-gm} y \ref{def:rf-pm}); $\Delta L=\PM/\omega_{PM}$ es el retardo adicional que el lazo tolera antes de volverse inestable.
\begin{center}
\renewcommand{\arraystretch}{1.2}
\begin{tabular}{@{}lccccccc@{}}
\toprule
Sintonía & $\OS$ & $\ts$ (s) & IAE & GM & PM & $\Delta L$ (s) & estables\\
\midrule
Ziegler-Nichols & $87$\,\% & $42$ & $9.7$ & $1.33$ & $39^\circ$ & $2.2$ & $98/125$\\
Cohen-Coon & $95$\,\% & $48$ & $11.5$ & $1.31$ & $27^\circ$ & $1.3$ & $90/125$\\
Ajustabilidad & $0$\,\% & $26$ & $8.0$ & $2.91$ & $75^\circ$ & $10.1$ & $125/125$\\
IMC, $\lambda=6.7$ & $0$\,\% & $28$ & $8.7$ & $3.16$ & $76^\circ$ & $11.2$ & $125/125$\\
\bottomrule
\end{tabular}
\end{center}
Cohen-Coon, la más agresiva, tolera solo $1.3$\,s más de retardo (un $33$\,\% de $L$); en la esquina con $\Kproc$ alto, $\tau$ corto y $L$ largo, Ziegler-Nichols y Cohen-Coon se vuelven inestables. Las dos sintonías conservadoras toleran toda la malla. Para que Cohen-Coon la tolere hay que multiplicar su $\Kc$ por $0.48$; y entonces, curiosamente, \emph{mejora} también el desempeño nominal ($\OS=24$\,\%, IAE $9.4$), porque la sintonía original oscilaba demasiado.
\end{ejemplo}

\codigorepo{cap08_pid_sintonizacion/C8_Ej6_Robustez.m}{malla}{márgenes y malla de procesos con \cmd{allmargin}}

\begin{ideaclave}
\textbf{La agresividad se paga en robustez.}\enspace Las reglas que buscan la razón de un cuarto dejan márgenes de ganancia cercanos a $1.3$ y márgenes de fase de $25$ a $40^\circ$: muy poco para un proceso que cambia. Ante la perturbación, en cambio, son las mejores (IAE de $1.9$ contra $9.6$ y $10.4$ de las conservadoras). Escoger una sintonía es escoger un punto en ese compromiso, y el IMC lo hace explícito con un solo parámetro: $\lambda$.
\end{ideaclave}

%% ---------------------------------------------------------------------
\section{Para llevar}
\label{sec:ps-cierre}
%% ---------------------------------------------------------------------

\begin{parallevar}
\begin{enumerate}
  \item El controlador ve todo el lazo: válvula, proceso y transmisor. Se identifica moviendo la CO y midiendo la PV, con ganancias normalizadas a los rangos.
  \item Modelo FOPDT $\Kproc e^{-Ls}/(\tau s+1)$ con la tangente de máxima pendiente: $L$ en el corte con el valor inicial, $\tau$ hasta el $63$\,\% (o hasta el $100$\,\% en la variante de Ziegler y Nichols, que sobrestima $\tau$). En un primer orden la pendiente máxima está en $t=L$.
  \item Malla abierta: Ziegler-Nichols, Cohen-Coon (agresivas, buscan el decaimiento de un cuarto) y la regla empírica del coeficiente de ajustabilidad (conservadora). Todas dependen de $L/\tau$.
  \item Malla cerrada: $\Kcr$ y $\Pcr$ de la oscilación sostenida. Ziegler-Nichols (PI: $0.45\Kcr$, $\Pcr/1.2$) es oscilatoria; Tyreus-Luyben es más conservadora y \emph{no} es de un cuarto.
  \item Síntesis directa: $\Gc=\frac{1}{\Gp}\frac{T}{1-T}$. Con $T=1/(\lambda s+1)$, $\lambda$ es la constante de tiempo del lazo: menor $\lambda$, más rápido. Los ceros inestables y el retardo no se invierten: van en $T$.
  \item IMC: $Q=(\tilde G_p^-)^{-1}F$, estructura que retroalimenta solo el error del modelo, y $\Gc=Q/(1-\tilde G_pQ)$, que para modelos sencillos es un PID con un solo parámetro, $\lambda$. Revise siempre la estabilidad interna.
  \item Desempeño y robustez se oponen: las reglas agresivas dejan márgenes de ganancia cercanos a $1.3$ y se desestabilizan con cambios moderados del proceso.
\end{enumerate}
\end{parallevar}

\begin{profundiza}[Correcciones respecto a la versión anterior del capítulo]
(i) Se decía que en un primer orden la pendiente máxima coincide con el $63$\,\%; está al comienzo de la respuesta, en $t=L$. (ii) El PI de Ziegler-Nichols de malla cerrada es $0.45\Kcr$ y $\Pcr/1.2$, no $0.4\Kcr$ y $0.8\Pcr$. (iii) Tyreus-Luyben no es ``un método de caída de un cuarto'': es más conservador que Ziegler-Nichols a propósito. (iv) En síntesis directa se decía que con $\lambda$ pequeño ``el sistema no responde rápido''; es al revés: $\ts\approx4\lambda$. (v) En el ejemplo de Tyreus-Luyben el proceso aparecía como $\Gc$ y $\Pcr=6.2$ en lugar de $6.28$\,s; en el de fase no mínima, $\tau_2=2.5$ era $\Td=2.5$; en el IMC faltaba la $s$ en $U(s)=\frac{1}{\Kproc}(\tau s+1)R(s)$. (vi) Proceso inestable: la condición es $F(1/\tau_u)=1$, no $F(\tau_u)=1$. (vii) En la tabla IMC, la fila del primer orden con retardo mezclaba la ganancia del PID con el $\Ti$ del PI; se separó en una fila PI y otra PID, y se retiró la fila del subamortiguado con cero y retardo, cuyas fórmulas no se pudieron verificar. (viii) La ganancia de velocidad $K_v$ se define con la pendiente de la rampa. (ix) La tabla del coeficiente de ajustabilidad se conserva, marcada como regla empírica. (x) Se completaron las simulaciones pendientes y se calcularon todos los números. (xi) Notación: $\tau_i,\tau_d\to\Ti,\Td$; $\theta\to L$; $T$ del FOPDT $\to\tau$, para no confundirla con $T(s)$.
\end{profundiza}

%% ---------------------------------------------------------------------
\section{Problemas}
\label{sec:ps-problemas}
%% ---------------------------------------------------------------------

Los problemas marcados con \etiquetaM{} requieren \matlab/\simulink; los marcados con \etiquetaEstrella{} son de profundización.

\begin{table}[ht]
\centering
\caption{Respuesta experimental del proceso de temperatura del problema~\ref{prob:ps-ident}.}
\label{tab:ps-datos}
\small
\begin{tabular}{@{}l*{7}{c}@{}}
\toprule
$t$ (s) & 0 & 2 & 4 & 6 & 8 & 10 & 15\\
PV ($\grados$C) & 80.0 & 80.0 & 80.1 & 81.7 & 83.3 & 84.6 & 87.7\\
\midrule
$t$ (s) & 20 & 24 & 30 & 40 & 50 & 60 & 100\\
PV ($\grados$C) & 89.9 & 91.4 & 93.1 & 95.0 & 96.2 & 96.9 & 97.9\\
\bottomrule
\end{tabular}
\end{table}

\begin{problema}[identificación]\label{prob:ps-ident}
En un proceso de temperatura, la CO ($0$--$100$\,\%) pasa de $30$ a $45$\,\% y el transmisor está calibrado entre $50$ y $150\,\grados$C; los datos están en la \cref{tab:ps-datos}. (a) Halle $\Kproc$ normalizada con \eqref{eq:ps-ganancia} y en unidades físicas; ¿por qué coinciden aquí? (b) Con la tangente y el método del $63$\,\%, halle $L$ y $\tau$. (c) Repita con la variante de Ziegler y Nichols y compare $\tau$. (d) Calcule $r_c=L/\tau$ y ubique el proceso en la \cref{tab:ps-ajust}. (e) \etiquetaM{} Superponga el modelo $1.2e^{-4s}/(20s+1)$ sobre los datos y calcule el error cuadrático medio. (f) Derive $y(t)=\Kproc(1-e^{-t/\tau})$ y diga dónde está la pendiente máxima.
\end{problema}

\begin{problema}[reglas de malla abierta]\label{prob:ps-abierta}
Con $\Kproc=1.2$, $L=4$\,s y $\tau=20$\,s: (a) halle las sintonías P, PI y PID de Ziegler-Nichols; (b) las P, PI, PD y PID de Cohen-Coon; (c) la del coeficiente de ajustabilidad. (d) \etiquetaM{} Simule los tres PID con el retardo exacto (estructura estándar, $N=10$) y tabule $\OS$, $\ts$, IAE y $|u_c|_{\max}$ ante un escalón de referencia. (e) Ordénelos de más a menos agresivo. (f) \etiquetaM{} Repita (d) con un escalón de perturbación a la entrada del proceso. ¿Se mantiene el orden? ¿Por qué una regla puede ser mejor para seguir la referencia y peor para rechazar perturbaciones?
\end{problema}

\begin{problema}[malla cerrada]\label{prob:ps-cerrada}
Con $\Gp=1/((s+1)(2s+1)(4s+1))$ y control proporcional: (a) halle $\Kcr$, $\wcr$ y $\Pcr$ por sustitución $s=\jw$; (b) verifique con Routh y \etiquetaM{} simulando con $\Kc$ un poco por debajo, igual y un poco por encima de $\Kcr$; (c) sintonice un PID ideal con Ziegler-Nichols y con Tyreus-Luyben; (d) \etiquetaM{} simule y mida la relación de decaimiento; ¿cuál da el cuarto? (e) \etiquetaM{} multiplique la ganancia del proceso por $1.5$ y repita; ¿cuál tolera mejor el cambio?
\end{problema}

\begin{problema}[síntesis directa]\label{prob:ps-sintesis}
(a) Para $\Gp=2/(5s+1)$ y $T=1/(\lambda s+1)$, muestre que resulta un PI y halle $\Kc$ y $\Ti$. (b) Para $\Gp=3/((2s+1)(10s+1))$, muestre que resulta un PID ideal. (c) \etiquetaM{} Simule (b) con $\lambda=1$, $3$ y $5$\,s; reporte $\ts$ y $|u_c|_{\max}$ y enuncie la relación entre $\lambda$ y la velocidad. (d) Con $\Gp=(1-9s)/((15s+1)(3s+1))$, muestre que $T=1/(\lambda s+1)$ da un controlador inestable y que $T=(1-9s)/(\lambda s+1)^2$ no. (e) Lleve el último a la forma PID en serie con filtro y verifique para $\lambda=5$. (f) Explique por qué el cero inestable debe ir en $T$ y qué pasa con el grado relativo si no.
\end{problema}

\begin{problema}[diseño IMC]\label{prob:ps-imc}
(a) Para $\tilde G_p=2/(5s+1)$, forme $Q$ con un filtro de primer orden y verifique que con modelo perfecto $Y=FR$. (b) Halle $\Gc$ con \eqref{eq:ps-imc-pid} y compare con el problema~\ref{prob:ps-sintesis}(a). (c) Repita para $\Kproc/((\tau_1s+1)(\tau_2s+1))$ y verifique la fila de la \cref{tab:ps-imc}. (d) Para $\tilde G_p=1.5(1-s)/((s+1)(4s+1))$, factorice de forma simple y pasa-todo, forme los dos $Q$ y compare sus órdenes. (e) \etiquetaM{} Simule los dos diseños con $\lambda=2$ ante referencia y ante una perturbación $1/(4s+1)$ a la salida; ¿cuál prefiere? (f) ¿Qué característica de la respuesta queda fijada por el cero inestable, haga lo que haga el controlador?
\end{problema}

\begin{problema}[robustez]\label{prob:ps-robustez}
Con las tres sintonías PID del problema~\ref{prob:ps-abierta}: (a) \etiquetaM{} calcule los márgenes de ganancia y de fase con el retardo exacto; (b) halle el retardo adicional $\Delta L$ que tolera cada una, en segundos y en porcentaje de $L$; (c) \etiquetaM{} varíe $\Kproc$ en $\pm30$\,\%, $\tau$ en $\pm30$\,\% y $L$ en $\pm50$\,\% (cinco valores de cada uno) y cuente los lazos estables; (d) identifique el peor caso y confírmelo por simulación; (e) ordene las sintonías por robustez y compare con el orden por agresividad; (f) \etiquetaEstrella{} halle el factor por el que hay que multiplicar el $\Kc$ de la menos robusta para que tolere toda la malla, y diga cuánto desempeño nominal cuesta.
\end{problema}

\begin{problema}[IMC-PID con tiempo muerto]\label{prob:ps-imc-fopdt}
Con $\Kproc=1.2$, $\tau=20$\,s y $L=4$\,s: (a) aproxime el retardo por Padé de primer orden, factorice y forme $Q$; (b) obtenga $\Gc$ y deduzca \eqref{eq:ps-imc-fopdt}; (c) compare con la fila FOPDT-PI de la \cref{tab:ps-imc}: ¿de qué aproximación sale cada una? (d) \etiquetaM{} Evalúe $\lambda=3.2$, $6.7$ y $10$\,s con el retardo exacto y tabule $\OS$, $\ts$, IAE y $|u_c|_{\max}$. (e) \etiquetaM{} Halle el menor $\lambda$ estable y compárelo con $0.8L$. (f) Compare la mejor sintonía IMC con las del problema~\ref{prob:ps-abierta} en IAE y margen de fase.
\end{problema}

\begin{problema}[estabilidad interna y errores de modelado]\label{prob:ps-interna}
(a) Para el ejemplo~\ref{ej:ps-bequette}, deduzca la ecuación característica y la cota de $\lambda$ por Routh. (b) \etiquetaM{} Simule por debajo, en y por encima de la cota. (c) \etiquetaM{} Con $\Gp=(1-9s)/((15s+1)(3s+1))$ y $\lambda=5$, simule $Q_1$ y $Q_2$ y grafique $y$ y $u$ durante $200$\,s. (d) Halle el polo responsable de la divergencia, su constante de tiempo y el instante en que $|u|$ supera $50$; explique por qué $y$ no lo revela. (e) Defina estabilidad interna. (f) Si el proceso es $K/(5s+1)$ y el modelo $2/(5s+1)$, halle la ecuación característica del lazo IMC y el rango de $K$ estable. ¿Cuánto error de ganancia tolera?
\end{problema}

%% ---------------------------------------------------------------------
\section{Soluciones seleccionadas}
%% ---------------------------------------------------------------------

\paragraph{Problema~\ref{prob:ps-ident}}
(a) $\Delta y\approx17.9\,\grados$C: $\Kproc=(17.9/100)/(15/100)=1.19\approx1.2$, y en unidades físicas $1.19\,\grados$C/\%; coinciden porque el rango de la PV es de $100\,\grados$C. (b) La pendiente máxima de los datos, $0.8\,\grados$C/s entre $4$ y $6$\,s, da $L=3.9$\,s; el $63.2$\,\% ($91.3\,\grados$C) se alcanza en $t\approx23.8$\,s, así que $\tau_{63}\approx19.9$\,s. (c) $\tau_{ZN}=17.9/0.8=22.4$\,s, mayor. (d) $r_c\approx0.2$: justo en el límite entre la fila PI y la PID; se usa la PID. (e) ECM $=0.0019\,(\grados\text{C})^2$. (f) La pendiente $\frac{\Kproc}{\tau}e^{-t/\tau}$ es máxima en $t=0$; en $t=\tau$ vale el $37$\,\% de ese máximo.

\paragraph{Problema~\ref{prob:ps-abierta}}
(a) P: $\Kc=4.17$; PI: $3.75$, $13.3$\,s; PID: $5$, $8$\,s, $2$\,s. (b) P: $4.44$; PI: $3.82$, $9.42$\,s; PD: $5.35$, $\Td=0.99$\,s; PID: $5.76$, $9.10$\,s, $1.40$\,s. (c) $2.29$, $22$\,s, $1.82$\,s. (d) ZN: $\OS=87$\,\%, $\ts=42$\,s, IAE $9.7$, $|u_c|_{\max}=54$; CC: $95$\,\%, $48$\,s, $11.5$, $61$; ajustabilidad: $0$\,\%, $26$\,s, $8.0$, $25$. (e) CC, ZN, ajustabilidad. (f) Picos $0.229$, $0.231$ y $0.288$; IAE $1.90$, $1.89$ y $9.6$: el orden se invierte.

\paragraph{Problema~\ref{prob:ps-cerrada}}
(a) $8s^3+14s^2+7s+1+\Kc=0$: $\wcr=\sqrt{7/8}=0.935$\,rad/s, $\Kcr=11.25$, $\Pcr=6.72$\,s. (c) ZN: $6.75$, $3.36$\,s, $0.84$\,s; TL: $5.11$, $14.8$\,s, $1.07$\,s. (d) ZN: $\OS=48$\,\%, decaimiento $0.16$; TL: $\OS=11$\,\%, sin segundo pico. (e) Con la ganancia $\times1.5$: $56$\,\% y $23$\,\%.

\paragraph{Problema~\ref{prob:ps-sintesis}}
(a) $\Kc=5/(2\lambda)$, $\Ti=5$\,s. (b) $\Kc=4/\lambda$, $\Ti=12$\,s, $\Td=5/3$\,s. (c) $\ts=3.9$, $12.1$ y $19.8$\,s; $|u_c|_{\max}=43.8$, $14.6$ y $8.8$: menor $\lambda$, más rápido. (d) Polo en $+1/9$. (e) $\Kc=18/(2\lambda+9)$, $\Ti=18$\,s, $\Td=2.5$\,s, $\tau_f=\lambda^2/(2\lambda+9)$; con $\lambda=5$, $0.947$ y $1.32$\,s.

\paragraph{Problema~\ref{prob:ps-imc}}
(d) Simple: $\tilde G_p^+=1-s$, $Q=(s+1)(4s+1)/(1.5(\lambda s+1)^2)$, de orden $2$. Pasa-todo: $\tilde G_p^+=(1-s)/(1+s)$, $\tilde G_p^-=1.5/(4s+1)$, $Q=(4s+1)/(1.5(\lambda s+1))$, de orden $1$. (e) Con $\lambda=2$: $\ts=12.5$ y $10.0$\,s, IAE $5$ y $4$ ($2\lambda+1$ y $\lambda+2$); ante la perturbación, picos $0.52$ y $0.47$. Se prefiere la pasa-todo. (f) La respuesta inversa.

\paragraph{Problema~\ref{prob:ps-robustez}}
(a) GM $1.33$, $1.31$, $2.91$; PM $39^\circ$, $27^\circ$, $75^\circ$. (b) $\Delta L=2.2$\,s ($54$\,\%), $1.3$\,s ($33$\,\%) y $10.1$\,s ($252$\,\%). (c) $98$, $90$ y $125$ de $125$. (d) $1.3\Kproc$, $0.7\tau$, $1.5L$. (f) CC con $\Kc\times0.48=2.77$: $\OS=24$\,\%, $\ts=31$\,s, IAE $9.4$, mejor que la original.

\paragraph{Problema~\ref{prob:ps-imc-fopdt}}
(c) La fila PI sale de $e^{-Ls}\approx1-Ls$; la PID, de Padé. La versión anterior de la tabla combinaba la ganancia de la segunda con el $\Ti$ de la primera. (d) $\Kc=3.53$, $2.11$, $1.53$; $\OS=16$\,\%, $0$, $0$; $\ts=16.5$, $28.3$, $41.1$\,s; IAE $5.8$, $8.7$, $12.0$; $|u_c|_{\max}=38$, $23$, $16$. (e) $\lambda_{\min}\approx0.75$\,s frente a $3.2$\,s. (f) IMC con $\lambda=6.7$: IAE $8.7$ y PM $76^\circ$, frente a $9.7$ y $39^\circ$ de Ziegler-Nichols.

\paragraph{Problema~\ref{prob:ps-interna}}
(a) $\lambda>0.887$. (d) Polo en $+1/9$, constante de tiempo $9$\,s; $|u|>50$ en $t=27.95$\,s. El cero del proceso cancela el polo en la función $Y/R$. (f) $1+Q(\Gp-\tilde G_p)=0$ se reduce a $2\lambda s+K=0$: estable para todo $K>0$. Con un error puro de ganancia, el IMC con modelo de primer orden tolera cualquier ganancia positiva.

\respuestaspausas

              % migrado (Handout6b_Ctrl.tex)
%% =====================================================================
%%  Capítulo 9 -- Retroalimentación de estado, observadores y acción integral
%%  Migrado al formato nuevo (noviembre de 2026) desde Handout7_Ctrl.tex.
%%  Cambios de fondo respecto a la versión anterior (ver MIGRACION.md):
%%    el ESTIMADO del observador se escribe \tilde{x} (macro \est), porque
%%    \hat{x} es la variable desviación en todo el curso; el error de
%%    estimación es \varepsilon (e queda para el error de seguimiento);
%%    "Ackerman" -> Ackermann; signo del polinomio del servo corregido;
%%    K del ejemplo diag(-1,-3) calculado; Ke = 18 314 (exacto).
%% =====================================================================

\capitulo[cap:re]{Retroalimentación de estado y observadores}
  {[L]a virginidad o, mejor, la inocencia se pierde de otro modo. Se pierde, como se sabe, cuando entran en juego las emociones.}
  {Alberto Fuguet}

\begin{objetivos}[partes/09-retroalimentacion-estado.html]
  \item Explicar con controlabilidad y observabilidad por qué una función de transferencia puede esconder modos, y decidir si un sistema es estabilizable y detectable.
  \item Diseñar una retroalimentación de estado $u=r-\Km x$ que ubique los polos donde piden las especificaciones, por comparación de coeficientes y con la fórmula de Ackermann.
  \item Diseñar un observador de Luenberger y unirlo con la retroalimentación de estado usando el principio de separación.
  \item Agregar acción integral para anular el error en estado estacionario.
  \item Verificar cada diseño en \matlab/\simulink, incluido el precio que se paga en esfuerzo de control.
\end{objetivos}

%% ---------------------------------------------------------------------
\section{Lo que la función de transferencia esconde}
\label{sec:re-kalman}
%% ---------------------------------------------------------------------

Hasta ahora hemos diseñado controladores en el dominio de Laplace: un PID, un compensador ajustado sobre la función de transferencia. Volvemos ahora a la representación de estado del capítulo~4,
\begin{equation}
  \dot x=\Am x+\Bm u,\qquad y=\Cm x+\Dm u,
  \label{eq:re-estado}
\end{equation}
con $x\in\mathbb{R}^n$ (estado), $u\in\mathbb{R}^m$ (entradas), $y\in\mathbb{R}^p$ (salidas), $\Am$ de $n\times n$, $\Bm$ de $n\times m$, $\Cm$ de $p\times n$ y $\Dm$ de $p\times m$. La idea del capítulo es sencilla de enunciar: si medimos \emph{todo} el estado, podemos escoger $u$ como una combinación de los estados y poner los polos del lazo cerrado donde queramos. Antes de hacerlo hay que contestar dos preguntas: ¿la entrada alcanza a mover todos los estados? y ¿la salida alcanza a mostrarlos todos? Son la \clave{controlabilidad} y la \clave{observabilidad}.

\begin{historia}[Kalman y la teoría de sistemas de estado]
Rudolf E.\ Kalman introdujo los conceptos de controlabilidad y observabilidad, y su dualidad, en el primer congreso de la IFAC, en Moscú, en 1960~\citep{Kal60}. Una de sus motivaciones era explicar por qué fracasaba un método popular para controlar plantas inestables: cancelar los polos inestables con ceros del compensador. La cancelación algebraica podía ser perfecta y aun así el sistema real resultaba inestable. En 1963 mostró que todo sistema lineal se puede descomponer en cuatro partes según sea o no controlable y observable, y que la función de transferencia solo ``ve'' la parte que es ambas cosas~\citep{Kal63}.
\end{historia}

\begin{ejemplo}[label=ej:re-friedland]{Cuatro estados, una función de primer orden}
Considere el sistema~\citep{Fri86}
\[
\begin{aligned}
  \dot x_1&=2x_1+3x_2+2x_3+x_4+u, & \dot x_2&=-2x_1-3x_2-2u,\\
  \dot x_3&=-2x_1-2x_2-4x_3+2u,   & \dot x_4&=-2x_1-2x_2-2x_3-5x_4-u,
\end{aligned}
\qquad y=7x_1+6x_2+4x_3+2x_4.
\]
Calcule su función de transferencia.
\solucion $H(s)=\Cm(s\Id-\Am)^{-1}\Bm+\Dm$ da
\[
  H(s)=\frac{s^3+9s^2+26s+24}{s^4+10s^3+35s^2+50s+24}=\frac{(s+2)(s+3)(s+4)}{(s+1)(s+2)(s+3)(s+4)}=\frac{1}{s+1}.
\]
¡Un sistema de cuarto orden se ve desde afuera como uno de primer orden! Con un cambio de variables (una \emph{transformación de similaridad} que lleva $\Am$ a su forma diagonal) el mismo sistema queda
\[
  \dot z_1=-z_1+u,\qquad \dot z_2=-2z_2,\qquad \dot z_3=-3z_3+u,\qquad \dot z_4=-4z_4,\qquad y=z_1+z_2,
\]
y la \cref{fig:re-friedland} lo dice todo: $z_1$ es afectado por $u$ y visible en $y$; $z_2$ es visible pero $u$ no lo mueve; $z_3$ lo mueve $u$ pero no se ve en $y$; $z_4$ ni lo uno ni lo otro. Solo $z_1$ llega a la función de transferencia.
\end{ejemplo}

\begin{figure}[ht]
\centering
\includegraphics[width=0.85\linewidth]{cap09/modal_friedland}
\caption{Ejemplo~\ref{ej:re-friedland} en coordenadas modales. Solo el modo $z_1$ conecta la entrada con la salida; los otros tres existen, pero la función de transferencia no los muestra.}
\label{fig:re-friedland}
\end{figure}

\codigorepo{cap09_retroalimentacion_estado/C9_Ej1_Kalman.m}{kalman}{ejemplo \ref{ej:re-friedland}}

Kalman mostró que esto es general: todo sistema \eqref{eq:re-estado} se puede descomponer en (1) una parte controlable y observable, (2) una observable pero no controlable, (3) una controlable pero no observable y (4) una que no es ni lo uno ni lo otro~\citep{Kal63}. La función de transferencia está determinada solo por la primera. Un sistema que contiene una parte no controlable se dice \emph{no controlable}; uno que contiene una parte no observable, \emph{no observable}.

\begin{definicion}[label=def:re-estabilizable]{Estabilizable y detectable}
Un sistema es \clave{estabilizable} si todos sus modos no controlables son estables, y \clave{detectable} si todos sus modos no observables son estables.
\end{definicion}

En el ejemplo, los modos ``escondidos'' están en $-2$, $-3$ y $-4$: son estables, el sistema es estabilizable y detectable, y podemos ignorar esas dinámicas internas sin riesgo. Pero si uno de ellos estuviera en el semiplano derecho, ninguna medición de $y$ lo delataría y ningún diseño sobre $u$ lo corregiría: el sistema se desestabilizaría ``por dentro''.

\begin{cuidado}[Cancelar un polo inestable con un cero]
Cancelar un polo en el semiplano derecho con un cero del controlador produce, en el papel, una función de transferencia estable. En la realidad, el modo inestable sigue ahí: se volvió no controlable o no observable, y cualquier perturbación o condición inicial lo excita. Es exactamente lo que Kalman explicó. Nunca se cancelan polos inestables.
\end{cuidado}

%% ---------------------------------------------------------------------
\section{Retroalimentación de estado}
\label{sec:re-retro}
%% ---------------------------------------------------------------------

Supongamos que medimos todos los estados. La ley de control más sencilla que los usa es lineal:
\begin{equation}
  u(t)=r(t)-\Km\,x(t),
  \label{eq:re-ley}
\end{equation}
donde $r$ es una nueva entrada de referencia y $\Km$ (de $m\times n$; un vector fila en sistemas de una entrada) es lo \emph{único} que diseñamos (\cref{fig:re-lazo}). Reemplazando en \eqref{eq:re-estado},
\begin{equation}
  \dot x=(\Am-\Bm\Km)\,x+\Bm r,\qquad y=\Cm x.
  \label{eq:re-lazocerrado}
\end{equation}
El lazo cerrado es otro sistema de estado con matriz $\Am-\Bm\Km$: sus polos son las raíces de
\begin{equation}
  \det\bigl(\lambda\Id-(\Am-\Bm\Km)\bigr)=0,
  \label{eq:re-ecaracteristica}
\end{equation}
y escoger $\Km$ es escoger esa ecuación característica. Con $r=0$ el lazo es un \clave{regulador}: lleva el estado al equilibrio en el origen, y basta que $\Am-\Bm\Km$ sea Hurwitz (todos sus valores propios con parte real negativa). Con $r\neq0$ acotada, el punto de equilibrio se mueve, pero sigue siendo asintóticamente estable.

\begin{figure}[ht]
\centering
\includegraphics[width=0.8\linewidth]{cap09/lazo_estado}
\caption{Retroalimentación de estado $u=r-\Km x$. Con $r=0$ es un regulador.}
\label{fig:re-lazo}
\end{figure}

\begin{ejemplo}[label=ej:re-motivador]{Estabilizar una planta inestable}
Sea $\dot x=\begin{bmatrix}-1&a\\3&-2\end{bmatrix}x+\begin{bmatrix}0\\1\end{bmatrix}u$, $y=[1\;\;0]\,x$. Diseñe $u=-\Km x$ para que los polos queden en $-2$ y $-3$.
\solucion Con $\Km=[K_1\;\;K_2]$,
\[
  \Am-\Bm\Km=\begin{bmatrix}-1&a\\3-K_1&-2-K_2\end{bmatrix},\qquad
  \det\bigl(\lambda\Id-(\Am-\Bm\Km)\bigr)=\lambda^2+(3+K_2)\lambda+2+K_2+a(K_1-3).
\]
El polinomio deseado es $(\lambda+2)(\lambda+3)=\lambda^2+5\lambda+6$. Igualando coeficientes, $K_2=2$ y $K_1=3+\dfrac2a$ (con $a\neq0$).

Con $a=2$ la planta tiene valores propios en $1$ y $-4$: el equilibrio es inestable. Con $\Km=[4\;\;2]$ los estados regresan al origen (\cref{fig:re-motivador}); el precio es un esfuerzo inicial $u(0)=-\Km x(0)=-6$.
\end{ejemplo}

\begin{figure}[ht]
\centering
\includegraphics[width=\linewidth]{cap09/f_motivador}
\caption{Ejemplo~\ref{ej:re-motivador} con $a=2$ y $x(0)=[1\;\;1]^\top$: sin control los estados divergen; con $u=-\Km x$ vuelven al origen.}
\label{fig:re-motivador}
\end{figure}

\codigorepo{cap09_retroalimentacion_estado/C9_Ej2_Motivador.m}{motivador}{ejemplo \ref{ej:re-motivador}}

\begin{pausa}
En el ejemplo~\ref{ej:re-motivador}, ¿qué pasa si $a=0$? ¿Qué polo no se puede mover?
\respuesta{Con $a=0$ el polinomio queda $\lambda^2+(3+K_2)\lambda+2+K_2=(\lambda+1)(\lambda+2+K_2)$: el factor $(\lambda+1)$ no depende de $\Km$. El polo en $-1$ no se mueve, porque con $a=0$ la ecuación de $x_1$ es $\dot x_1=-x_1$ y la entrada no llega a $x_1$. Es el tema de la próxima sección.}
\end{pausa}

%% ---------------------------------------------------------------------
\section{Controlabilidad}
\label{sec:re-controlabilidad}
%% ---------------------------------------------------------------------

La pausa anterior muestra que no siempre se pueden poner los polos donde uno quiera. ¿Cuándo sí?

\begin{definicion}[label=def:re-controlable]{Controlabilidad}
El sistema $\dot x=\Am x+\Bm u$ es \clave{controlable} si, para \emph{cualquier} estado inicial $x(0)=x_0$ y \emph{cualquier} estado final $x_T$, existe una entrada $u(t)$ que lleva el sistema de $x_0$ a $x_T$ en un tiempo \emph{finito} $T$.
\end{definicion}

Hay una caracterización equivalente que es la que nos interesa para diseñar: el sistema es controlable si y solo si, para cualquier polinomio mónico $\alpha_c(\lambda)$ de grado $n$ con coeficientes reales, existe $\Km$ tal que el polinomio característico de $\Am-\Bm\Km$ es $\alpha_c(\lambda)$ (\emph{ubicación arbitraria de polos}). En sistemas de una entrada ese $\Km$ es, además, único. Y hay una prueba algebraica que no exige calcular nada más que productos de matrices.

\begin{definicion}[label=def:re-ctrb]{Matriz de controlabilidad}
El sistema con $\Am\in\mathbb{R}^{n\times n}$ es controlable si y solo si la \clave{matriz de controlabilidad}
\begin{equation}
  \Ctrb=\bigl[\,\Bm\;\;\Am\Bm\;\;\Am^2\Bm\;\;\cdots\;\;\Am^{n-1}\Bm\,\bigr]
  \label{eq:re-ctrb}
\end{equation}
tiene rango $n$. Con una sola entrada $\Ctrb$ es cuadrada y basta revisar que $\det\Ctrb\neq0$.
\end{definicion}

La intuición: $\Bm u$ empuja el estado en la dirección de $\Bm$; un instante después la dinámica ha llevado ese empujón a la dirección $\Am\Bm$, luego a $\Am^2\Bm$, y así. Si esas $n$ direcciones generan todo $\mathbb{R}^n$, la entrada puede llevar el estado a cualquier parte. En el ejemplo~\ref{ej:re-motivador}, $\Ctrb=\begin{bmatrix}0&a\\1&-2\end{bmatrix}$ y $\det\Ctrb=-a$: controlable si y solo si $a\neq0$.

\begin{ejemplo}[label=ej:re-diag]{Dos modos y una entrada}
Para $\Am=\begin{bmatrix}-1&0\\0&-3\end{bmatrix}$, $\Bm=\begin{bmatrix}1\\1\end{bmatrix}$: ¿es controlable? Si lo es, halle $\Km$ para tener los dos polos en $-2$.
\solucion $\Ctrb=[\Bm\;\;\Am\Bm]=\begin{bmatrix}1&-1\\1&-3\end{bmatrix}$, con $\det\Ctrb=-2\neq0$: controlable (la entrada llega a los dos modos). Los polos de la planta son $-1$ y $-3$. Con $\Km=[K_1\;\;K_2]$,
\[
  \det\bigl(\lambda\Id-(\Am-\Bm\Km)\bigr)=\lambda^2+(4+K_1+K_2)\lambda+3+3K_1+K_2,
\]
e igualando con $(\lambda+2)^2=\lambda^2+4\lambda+4$: $K_1+K_2=0$ y $3K_1+K_2=1$, de donde $\Km=[0.5\;\;-0.5]$.
\end{ejemplo}

\codigorepo{cap09_retroalimentacion_estado/C9_Ej3_Controlabilidad.m}{rango}{ejemplo \ref{ej:re-diag}}

\begin{profundiza}[¿Cuál modo falla? La prueba PBH]
El rango de $\Ctrb$ dice \emph{si} el sistema es controlable, pero no \emph{qué} modo falla. La prueba de Popov-Belevitch-Hautus lo dice: $(\Am,\Bm)$ es controlable si y solo si
\[
  \operatorname{rango}\bigl[\,\Am-\lambda\Id\;\;\Bm\,\bigr]=n\quad\text{para cada valor propio }\lambda\text{ de }\Am.
\]
El valor propio para el cual el rango cae es el modo no controlable. En el ejemplo~\ref{ej:re-motivador} con $a=0$, el rango cae en $\lambda=-1$. Y el sistema es estabilizable si y solo si los $\lambda$ que fallan tienen parte real negativa.
\end{profundiza}
\codigorepo{cap09_retroalimentacion_estado/C9_Ej3_Controlabilidad.m}{pbh}{prueba PBH}

\subsection{Controlable no quiere decir barato}

La controlabilidad es una propiedad de sí o no, pero en la práctica importa \emph{cuánto}. En el ejemplo~\ref{ej:re-motivador}, $K_1=3+2/a$: cuando $a\to0$ el sistema sigue siendo controlable, pero la ganancia necesaria crece sin límite, y con ella el esfuerzo de control (\cref{fig:re-condicion}). El número de condición de $\Ctrb$ es un buen termómetro: si es muy grande, $\Ctrb$ está ``casi'' sin rango completo y mover los polos sale caro (actuadores que se saturan, ruido amplificado).

\begin{figure}[ht]
\centering
\includegraphics[width=0.85\linewidth]{cap09/f_condicion}
\caption{Ejemplo~\ref{ej:re-motivador} con los polos siempre en $-2$ y $-3$: a medida que $a\to0$, el número de condición de $\Ctrb$ y el pico de $|u|$ crecen como $1/a$.}
\label{fig:re-condicion}
\end{figure}

\codigorepo{cap09_retroalimentacion_estado/C9_Ej3_Controlabilidad.m}{condicion}{\cref{fig:re-condicion}}

%% ---------------------------------------------------------------------
\section{La fórmula de Ackermann}
\label{sec:re-ackermann}
%% ---------------------------------------------------------------------

Comparar coeficientes funciona bien con $n=2$ o $3$, pero con más estados el álgebra se vuelve pesada. Para sistemas de una entrada hay una fórmula cerrada.

\begin{historia}[Jürgen Ackermann]
La fórmula lleva el nombre del ingeniero alemán Jürgen Ackermann, que la publicó en 1972 en un artículo sobre el diseño de sistemas de control lineales en el espacio de estado~\citep{Ack72}. Ackermann trabajó luego en el Centro Aeroespacial Alemán (DLR) en control robusto, con aplicaciones a la dirección de automóviles.
\end{historia}

\begin{definicion}[label=def:re-ackermann]{Fórmula de Ackermann}
Si el sistema de una entrada es controlable y el polinomio característico deseado es $q(\lambda)=\lambda^n+\alpha_{n-1}\lambda^{n-1}+\cdots+\alpha_1\lambda+\alpha_0$, entonces
\begin{equation}
  \Km=\bigl[\,0\;\;0\;\;\cdots\;\;0\;\;1\,\bigr]\,\Ctrb^{-1}\,q(\Am),
  \qquad
  q(\Am)=\Am^n+\alpha_{n-1}\Am^{n-1}+\cdots+\alpha_1\Am+\alpha_0\Id.
  \label{eq:re-ackermann}
\end{equation}
\end{definicion}

Observe que $q(\Am)$ es el polinomio deseado evaluado en la \emph{matriz} $\Am$ (con $\Id$ en lugar del término independiente), y que la fórmula exige $\Ctrb^{-1}$: solo funciona si el sistema es controlable.

\begin{ejemplo}[label=ej:re-ackermann]{El doble integrador}
Para $Y(s)/U(s)=1/s^2$, halle $\Km$ para tener los polos en $-1\pm j$.
\solucion En variables de fase, $\Am=\begin{bmatrix}0&1\\0&0\end{bmatrix}$, $\Bm=\begin{bmatrix}0\\1\end{bmatrix}$. El polinomio deseado es $q(\lambda)=\lambda^2+2\lambda+2$. Entonces
\[
  \Ctrb=[\Bm\;\;\Am\Bm]=\begin{bmatrix}0&1\\1&0\end{bmatrix}=\Ctrb^{-1},\qquad
  q(\Am)=\Am^2+2\Am+2\Id=\begin{bmatrix}2&2\\0&2\end{bmatrix},
\]
\[
  \Km=[0\;\;1]\begin{bmatrix}0&1\\1&0\end{bmatrix}\begin{bmatrix}2&2\\0&2\end{bmatrix}=[1\;\;0]\begin{bmatrix}2&2\\0&2\end{bmatrix}=[2\;\;2].
\]
Comprobación: $\Am-\Bm\Km=\begin{bmatrix}0&1\\-2&-2\end{bmatrix}$, cuyo polinomio es $\lambda^2+2\lambda+2$.
\end{ejemplo}

\codigorepo{cap09_retroalimentacion_estado/C9_Ej4_Ackermann.m}{ackermann}{ejemplo \ref{ej:re-ackermann}: la fórmula paso a paso, \cmd{acker} y \cmd{place}}

\begin{cuidado}[\texttt{acker} o \texttt{place}]
\cmd{acker} implementa \eqref{eq:re-ackermann}: es exacta en papel pero numéricamente frágil para $n$ grande (invierte $\Ctrb$, que puede estar mal condicionada). \cmd{place} usa un algoritmo más robusto, sirve para varias entradas, pero no acepta polos repetidos más veces que el número de entradas. Para los ejemplos del curso dan lo mismo.
\end{cuidado}

\begin{receta}[Ubicación de polos por retroalimentación de estado]
\begin{enumerate}
  \item Escribir el modelo de estado y verificar que $\operatorname{rango}\Ctrb=n$.
  \item Traducir las especificaciones en polos deseados (\cref{sec:re-especificaciones}) y armar $q(\lambda)$.
  \item Calcular $\Km$: comparando los coeficientes de $\det(\lambda\Id-\Am+\Bm\Km)$ con $q(\lambda)$, o con \eqref{eq:re-ackermann}.
  \item Verificar \cmd{eig(A-B*K)}, simular, y mirar el esfuerzo de control $u(t)$.
\end{enumerate}
\end{receta}

\begin{pausa}
Use la fórmula de Ackermann para $1/(s(s+3))$ en variables de fase, con polos en $-4\pm j4$.
\respuesta{$\Am=\begin{bmatrix}0&1\\0&-3\end{bmatrix}$, $\Ctrb=\begin{bmatrix}0&1\\1&-3\end{bmatrix}$, $\Ctrb^{-1}=\begin{bmatrix}3&1\\1&0\end{bmatrix}$, $q(\Am)=\Am^2+8\Am+32\Id=\begin{bmatrix}32&5\\0&17\end{bmatrix}$. Entonces $\Km=[1\;\;0]\,q(\Am)=[32\;\;5]$. Comprobación: $\lambda^2+(3+K_2)\lambda+K_1=\lambda^2+8\lambda+32$.}
\end{pausa}

%% ---------------------------------------------------------------------
\section{De las especificaciones a los polos}
\label{sec:re-especificaciones}
%% ---------------------------------------------------------------------

Con $n$ estados hay que escoger $n$ polos. La estrategia es la del \cref{cap:so}: dos polos \emph{dominantes} fijados por las especificaciones del segundo orden ($\OS$ fija $\zeta$; $\tp$ o $\ts$ fijan $\wn$) y los demás bastante más a la izquierda, para que sus modos mueran rápido. Una regla usual es ponerlos entre cinco y diez veces más lejos que la parte real de los dominantes. Si la planta tiene un cero de fase mínima cerca de los dominantes, conviene usar uno de los polos sobrantes para \emph{cancelarlo} (sin modificar la ganancia).

\begin{ejemplo}[label=ej:re-trespolos]{Un tercer orden con especificaciones}
La planta $G_p(s)=\dfrac{10}{(s+10)(s^2+0.5s+1)}$ tiene $\OS=44\%$ y $\tp=3.35$\,s. Diseñe $u=r-\Km x$ para $\OS=10\%$ y $\tp=1$\,s.
\solucion $\OS=10\%$ da $\zeta=0.591$, y $\tp=1$ da $\wn=\pi/\sqrt{1-\zeta^2}=3.90$\,rad/s: polos dominantes en $-2.30\pm j3.14$. El tercero va diez veces más lejos, en $-23.0$. Con la representación que entrega \cmd{tf2ss},
\[
  \Km=[17.131\;\;115.21\;\;339.34].
\]
La respuesta (\cref{fig:re-trespolos}) tiene $\OS=9.8\%$ y $\tp=1.05$\,s: el tercer polo, aunque lejano, la desvía un poco de las fórmulas. Pero el valor final es $0.029$: la salida no sigue la referencia. Lo arreglamos en la \cref{sec:re-integral}.
\end{ejemplo}

\begin{figure}[ht]
\centering
\includegraphics[width=\linewidth]{cap09/f_tres_polos}
\caption{Ejemplo~\ref{ej:re-trespolos}. La forma de la respuesta cumple; el valor final (que a la derecha se normalizó para comparar la forma), no.}
\label{fig:re-trespolos}
\end{figure}

\codigorepo{cap09_retroalimentacion_estado/C9_Ej5_Especificaciones.m}{tercerorden}{ejemplo \ref{ej:re-trespolos}}

\begin{cuidado}[Las ganancias dependen de la representación; los polos, no]
El mismo sistema tiene infinitas representaciones de estado, y $\Km$ cambia con cada una porque los estados son otros. \cmd{tf2ss} entrega la forma con la primera fila $[-a_{n-1}\;\cdots\;-a_0]$; las \emph{variables de fase} del capítulo de espacio de estado tienen la última fila $[-a_0\;\cdots\;-a_{n-1}]$. En el ejemplo~\ref{ej:re-trespolos}, en variables de fase sale $\Km=[339.34\;\;115.21\;\;17.131]$: las mismas cifras al revés, porque los estados están numerados al revés. Los polos del lazo cerrado son idénticos. Diga siempre en qué representación trabaja; la función del curso \cmd{variablesFase} construye la de las notas.
\end{cuidado}

\subsection{Cuando la referencia entra por otro lado}

No siempre la referencia se suma a $u$. Si la salida es un estado (por ejemplo $y=x_1$), lo natural es compararla con $r$ como en un lazo clásico.

\begin{ejemplo}[label=ej:re-servo]{Un servo de tercer orden}
La planta es $\dfrac{Y(s)}{U(s)}=\dfrac{1}{s(s+1)(s+2)}$ con la configuración de la \cref{fig:re-servo}: $u=K_1(r-x_1)-K_2x_2-K_3x_3$. Ubique los polos en $-2\pm j3.464$ y $-10$.
\solucion En variables de fase, $\Am=\begin{bmatrix}0&1&0\\0&0&1\\0&-2&-3\end{bmatrix}$, $\Bm=\begin{bmatrix}0\\0\\1\end{bmatrix}$, $\Cm=[1\;\;0\;\;0]$. La matriz de controlabilidad,
\[
  \Ctrb=[\Bm\;\;\Am\Bm\;\;\Am^2\Bm]=\begin{bmatrix}0&0&1\\0&1&-3\\1&-3&7\end{bmatrix},
\]
tiene rango $3$. La ley es $u=-\Km x+K_1r$ con $\Km=[K_1\;\;K_2\;\;K_3]$, y el polinomio característico es
\[
  \lambda^3+(3+K_3)\lambda^2+(2+K_2)\lambda+K_1.
\]
El deseado es $(\lambda+10)(\lambda^2+4\lambda+16)=\lambda^3+14\lambda^2+56\lambda+160$, luego $\Km=[160\;\;54\;\;11]$ y
\[
  \dot x=\begin{bmatrix}0&1&0\\0&0&1\\-160&-56&-14\end{bmatrix}x+\begin{bmatrix}0\\0\\160\end{bmatrix}r,\qquad y=[1\;\;0\;\;0]\,x.
\]
Ahora el valor final \emph{sí} es $1$: la planta tiene un integrador y la referencia se compara con la salida, como en un lazo tipo 1 (\cref{cap:er}).
\end{ejemplo}

\begin{figure}[ht]
\centering
\includegraphics[width=0.9\linewidth]{cap09/servo_k1}
\caption{Ejemplo~\ref{ej:re-servo}: la referencia se compara con $x_1=y$; los otros estados se retroalimentan con $K_2$ y $K_3$.}
\label{fig:re-servo}
\end{figure}

\codigorepo{cap09_retroalimentacion_estado/C9_Ej6_Servo.m}{servo}{ejemplo \ref{ej:re-servo}}

\begin{pausa}
Para $T(s)=\dfrac{100(s+1)}{s(s+2)(s+3)}$ se pide $\OS=10\%$ y $\tp=1$\,s con $u=r-\Km x$, cancelando el cero con el tercer polo. ¿Dónde va el tercer polo y qué función de transferencia queda en lazo cerrado?
\respuesta{En $-1$, sobre el cero. Queda $\dfrac{100}{s^2+4.605s+15.17}$: la forma pedida ($\zeta=0.591$, $\wn=3.90$), pero con valor final $100/15.17=6.59$, no $1$. Otra vez hace falta la acción integral.}
\end{pausa}

%% ---------------------------------------------------------------------
\section{Observadores de estado}
\label{sec:re-observador}
%% ---------------------------------------------------------------------

Todo lo anterior supone que medimos todos los estados. Casi nunca es así: medir cuesta, y algunos estados (una velocidad interna, una concentración) no tienen sensor. La salida de este problema es \emph{estimarlos}: construir dentro del controlador una copia del modelo que, alimentada con la entrada y con la salida medida, reproduzca el estado.

\begin{historia}[David Luenberger]
David G.\ Luenberger acuñó el término \emph{observador} en su tesis doctoral en Stanford (1963) y lo publicó en 1964~\citep{Lue64}. Mostró que para cualquier sistema lineal observable se puede construir un observador cuyo error de estimación tienda a cero tan rápido como se quiera, y que su diseño es el ``dual'' de la ubicación de polos por retroalimentación de estado.
\end{historia}

\paragraph{La idea}
Una copia del modelo, $\dot{\est{x}}=\Am\est{x}+\Bm u$, reproduciría el estado si arrancara exactamente en $x(0)$ y el modelo fuera perfecto. Como no es así, la corregimos con lo que sí sabemos: la diferencia entre la salida medida $y$ y la que predice la copia, $\Cm\est{x}$ (\cref{fig:re-observador}):
\begin{equation}
  \dot{\est{x}}=\Am\est{x}+\Bm u+\Lm\,(y-\Cm\est{x})=(\Am-\Lm\Cm)\,\est{x}+\Bm u+\Lm y.
  \label{eq:re-observador}
\end{equation}
$\Lm$ (de $n\times p$) es la \clave{ganancia del observador}. Restando \eqref{eq:re-observador} de \eqref{eq:re-estado}, el \clave{error de estimación} $\errest=x-\est{x}$ cumple
\begin{equation}
  \dot{\errest}=(\Am-\Lm\Cm)\,\errest.
  \label{eq:re-errorobs}
\end{equation}
Si $\Am-\Lm\Cm$ es Hurwitz, $\est{x}\to x$ \emph{sin importar la entrada}, y la velocidad la fijan los valores propios de $\Am-\Lm\Cm$, que escogemos con $\Lm$.

\begin{figure}[ht]
\centering
\includegraphics[width=0.8\linewidth]{cap09/observador}
\caption{Observador de Luenberger: una copia del modelo corregida con el error de salida $y-\Cm\est{x}$.}
\label{fig:re-observador}
\end{figure}

\begin{cuidado}[$\est{x}$ y no $\hat{x}$]
En muchos libros el estimado se escribe $\hat{x}$. En este curso $\hat{x}$ es la \emph{variable desviación} (\cref{cap:so}), y los dos símbolos aparecen juntos cuando se diseña un observador para un modelo linealizado. Por eso el estimado es $\est{x}$ y el error de estimación es $\errest$ (la letra $e$ queda para el error de seguimiento $r-y$).
\end{cuidado}

\begin{definicion}[label=def:re-observable]{Observabilidad}
El sistema \eqref{eq:re-estado} es \clave{observable} si el estado inicial $x(0)$ (y con él todo $x(t)$) se puede determinar a partir de la entrada $u(\tau)$ y la salida $y(\tau)$ en un intervalo finito $0\leq\tau\leq t$. Equivalentemente, si y solo si la \clave{matriz de observabilidad}
\begin{equation}
  \Obsv=\begin{bmatrix}\Cm\\ \Cm\Am\\ \vdots\\ \Cm\Am^{n-1}\end{bmatrix}
  \label{eq:re-obsv}
\end{equation}
tiene rango $n$. Y los polos del observador (los valores propios de $\Am-\Lm\Cm$) se pueden ubicar arbitrariamente si y solo si $(\Am,\Cm)$ es observable.
\end{definicion}

\paragraph{Dualidad}
Los valores propios de $\Am-\Lm\Cm$ son los de su transpuesta, $\Am^\top-\Cm^\top\Lm^\top$, que tiene la forma $\Am-\Bm\Km$ con $\Am\to\Am^\top$, $\Bm\to\Cm^\top$ y $\Km\to\Lm^\top$. Diseñar un observador es ubicar polos en el sistema ``dual''. En \matlab: \cmd{L = acker(A',C',p)'} (o \cmd{place}). La fórmula de Ackermann para el observador es
\begin{equation}
  \Lm=\rho(\Am)\,\Obsv^{-1}\begin{bmatrix}0\\ \vdots\\ 0\\ 1\end{bmatrix},\qquad
  \rho(\Am)=\Am^n+\beta_{n-1}\Am^{n-1}+\cdots+\beta_0\Id,
  \label{eq:re-ackermannobs}
\end{equation}
donde $\rho(\lambda)=\lambda^n+\beta_{n-1}\lambda^{n-1}+\cdots+\beta_0$ es el polinomio deseado del observador.

\begin{ejemplo}[label=ej:re-observador]{Un observador con Ackermann}
Para $\Am=\begin{bmatrix}2&3\\-1&4\end{bmatrix}$, $\Cm=[1\;\;0]$, halle $\Lm$ para que el observador tenga $\rho(\lambda)=\lambda^2+2\zeta\wn\lambda+\wn^2$ con $\zeta=0.8$ y $\wn=10$.
\solucion $\rho(\lambda)=\lambda^2+16\lambda+100$. Entonces
\[
  \rho(\Am)=\Am^2+16\Am+100\Id=\begin{bmatrix}133&66\\-22&177\end{bmatrix},\qquad
  \Obsv=\begin{bmatrix}1&0\\2&3\end{bmatrix},\qquad
  \Obsv^{-1}=\begin{bmatrix}1&0\\-2/3&1/3\end{bmatrix},
\]
y $\Lm=\rho(\Am)\Obsv^{-1}[0\;\;1]^\top=\rho(\Am)\,[0\;\;\tfrac13]^\top=[22\;\;59]^\top$. Con la planta $\Am=\begin{bmatrix}0&1\\-10&-2\end{bmatrix}$ y el mismo $\rho$, $\Obsv=\Id$ y sale $\Lm=[14\;\;62]^\top$; la \cref{fig:re-obs} muestra cómo el estimado alcanza al estado real partiendo de condiciones iniciales distintas.
\end{ejemplo}

\begin{figure}[ht]
\centering
\includegraphics[width=\linewidth]{cap09/f_observador}
\caption{Ejemplo~\ref{ej:re-observador} con $\Am=\begin{bmatrix}0&1\\-10&-2\end{bmatrix}$: la planta arranca en $x(0)=[1\;\;0]^\top$ y el observador en cero. En medio segundo el error de estimación ya es menor que el $4\%$.}
\label{fig:re-obs}
\end{figure}

\codigorepo{cap09_retroalimentacion_estado/C9_Ej7_Observador.m}{observador}{ejemplo \ref{ej:re-observador}}

\begin{ideaclave}
¿Qué tan rápido debe ser el observador? Lo bastante para que el estimado esté listo antes de que el controlador lo necesite: es usual ponerlo de dos a cinco veces más rápido que los polos de control. Más rápido no es gratis: un $\Lm$ grande amplifica el ruido de medición que entra por $y$. Es el mismo compromiso entre velocidad y ruido de todo el curso.
\end{ideaclave}

%% ---------------------------------------------------------------------
\section{Retroalimentación de estado con observador}
\label{sec:re-separacion}
%% ---------------------------------------------------------------------

El diseño completo se hace en cuatro pasos: (1) diseñar $\Km$ como si todos los estados se midieran; (2) diseñar un observador para los que no se miden; (3) conectar los dos, usando el estimado en lugar del estado:
\begin{equation}
  u=-\Km\est{x}\;(+Nr),
  \label{eq:re-leyobs}
\end{equation}
y (4) agregar la entrada de referencia. Al conjunto ``observador más $\Km$'' se le llama \clave{compensador} (\cref{fig:re-compensador}).

\begin{figure}[ht]
\centering
\includegraphics[width=0.9\linewidth]{cap09/compensador}
\caption{Compensador por variables de estado: el observador produce $\est{x}$ y la ley de control lo usa en lugar de $x$.}
\label{fig:re-compensador}
\end{figure}

¿Sigue siendo estable el lazo si usamos $\est{x}$ en lugar de $x$? Con $\est{x}=x-\errest$, la ley \eqref{eq:re-leyobs} (con $r=0$) da $\dot x=(\Am-\Bm\Km)x+\Bm\Km\errest$, y el error sigue cumpliendo \eqref{eq:re-errorobs}. Juntando,
\begin{equation}
  \begin{bmatrix}\dot x\\ \dot{\errest}\end{bmatrix}=
  \begin{bmatrix}\Am-\Bm\Km & \Bm\Km\\ 0 & \Am-\Lm\Cm\end{bmatrix}
  \begin{bmatrix}x\\ \errest\end{bmatrix}.
  \label{eq:re-aumentado}
\end{equation}
La matriz es triangular por bloques, así que su ecuación característica es el producto
\begin{equation}
  \det\bigl(\lambda\Id-(\Am-\Bm\Km)\bigr)\cdot\det\bigl(\lambda\Id-(\Am-\Lm\Cm)\bigr)=0.
\end{equation}

\begin{ideaclave}
\textbf{Principio de separación.} Los polos del lazo con observador son la unión de los polos de control (los de $\Am-\Bm\Km$) y los del observador (los de $\Am-\Lm\Cm$). Por eso $\Km$ y $\Lm$ se diseñan \emph{por separado}: si cada parte es estable, el conjunto lo es.
\end{ideaclave}

Visto desde afuera, el compensador es simplemente una función de transferencia entre la salida medida y la entrada a la planta (con $r=0$):
\begin{equation}
  \frac{U(s)}{Y(s)}=-\Km\bigl(s\Id-(\Am-\Bm\Km-\Lm\Cm)\bigr)^{-1}\Lm.
  \label{eq:re-compensador}
\end{equation}

\begin{ejemplo}[label=ej:re-separacion]{Un compensador completo}
Para $\Am=\begin{bmatrix}0&1\\-2&-3\end{bmatrix}$, $\Bm=\begin{bmatrix}0\\1\end{bmatrix}$, $\Cm=[1\;\;0]$ (la planta $1/((s+1)(s+2))$, de la que solo se mide la posición), diseñe $\Km$ para polos en $-2\pm j2$ y un observador con polos dobles en $-10$.
\solucion Control: $\lambda^2+(3+K_2)\lambda+2+K_1=\lambda^2+4\lambda+8$, luego $\Km=[6\;\;1]$. Observador: $\Am-\Lm\Cm=\begin{bmatrix}-L_1&1\\-2-L_2&-3\end{bmatrix}$ tiene polinomio $\lambda^2+(3+L_1)\lambda+3L_1+2+L_2$; igualando con $(\lambda+10)^2$, $\Lm=[17\;\;47]^\top$. Los cuatro polos de \eqref{eq:re-aumentado} son $-2\pm j2$, $-10$ y $-10$. Y el compensador \eqref{eq:re-compensador} resulta
\[
  \frac{U(s)}{Y(s)}=-\frac{149s+554}{s^2+21s+123},
\]
un compensador de segundo orden con un cero en $-3.72$ y polos en $-10.5\pm j3.57$: un ``adelanto con filtro'' que salió del espacio de estado sin que lo buscáramos.
\end{ejemplo}

\codigorepo{cap09_retroalimentacion_estado/C9_Ej8_Compensador.m}{separacion}{ejemplo \ref{ej:re-separacion}}

\begin{figure}[ht]
\centering
\includegraphics[width=0.85\linewidth]{cap09/f_separacion}
\caption{Ejemplo~\ref{ej:re-separacion} con la referencia escalada por $N$ para que el valor final sea $1$. Con un error inicial de estimación, el observador rápido (polos en $-10$) casi no se nota; uno lento (en $-4$) deforma el transitorio.}
\label{fig:re-separacionfig}
\end{figure}

\begin{modelo}[Observador.slx]
El script \cmd{crearModelo\_Observador.m} construye la planta (que arranca en $x(0)=[1\;\;-1]^\top$), el observador con entradas $[u;\,y]$ y la ley $u=-\Km\est{x}+Nr$. Un interruptor manual cambia $\est{x}$ por el estado real para comparar. Mueva los polos del observador (\cmd{po}) y observe cuánto tarda $\est{x}$ en alcanzar a $x$.
\end{modelo}

\begin{cuidado}[La separación no garantiza robustez]
El principio de separación supone que el modelo es exacto. Si la planta real difiere del modelo con que se diseñó el observador, el error de estimación ya no obedece \eqref{eq:re-errorobs} y los márgenes de estabilidad del compensador completo pueden ser mucho menores que los de la retroalimentación de estado sola. Hay que verificar el diseño con la planta perturbada (problema~\ref{prob:re-robustez}).
\end{cuidado}

%% ---------------------------------------------------------------------
\section{Acción integral: error nulo en estado estacionario}
\label{sec:re-integral}
%% ---------------------------------------------------------------------

Los ejemplos~\ref{ej:re-trespolos} y la pausa de la \cref{sec:re-especificaciones} mostraron el problema: $u=r-\Km x$ pone los polos donde queremos, pero nada obliga a que $y$ llegue a $r$. En el \cref{cap:er} vimos que el remedio es un integrador en el lazo (\cref{sec:er-estado}). En espacio de estado se hace igual.

\begin{ejemplo}[label=ej:re-sinintegral]{El error de la retroalimentación de estado pura}
Para $\Am=\begin{bmatrix}0&1\\-3&-5\end{bmatrix}$, $\Bm=\begin{bmatrix}0\\1\end{bmatrix}$, $\Cm=[1\;\;0]$, diseñe $u=r-\Km x$ con $\OS=10\%$ y $\ts=0.5$\,s, y calcule el error en estado estacionario ante un escalón unitario.
\solucion $\zeta=0.591$ y $\zeta\wn=4/0.5=8$, así que $\wn=13.53$ y el polinomio deseado es $s^2+16s+183.1$ (polos en $-8\pm j10.91$). El lazo cerrado tiene $s^2+(5+K_2)s+(3+K_1)$, luego $\Km=[180.1\;\;11]$. Con $\bar{\Am}=\Am-\Bm\Km$, $Y(s)/R(s)=\Cm(s\Id-\bar{\Am})^{-1}\Bm=1/(s^2+16s+183.1)$ y
\[
  \ess=\lim_{s\to0}s\,\frac1s\bigl(1-\Cm(s\Id-\bar{\Am})^{-1}\Bm\bigr)=1-\frac{1}{183.1}=0.9945.
\]
¡La salida llega a menos del $1\%$ de la referencia!
\end{ejemplo}

\paragraph{El estado integrador}
Agregamos un estado nuevo, la integral del error de seguimiento (\cref{fig:re-integral}):
\begin{equation}
  z(t)=\int_0^t e(\tau)\,\dif\tau,\qquad \dot z=e=r-\Cm x,
\end{equation}
y la ley de control $u=-\Km x+K_e z=-\bigl[\Km\;\;-K_e\bigr]\begin{bmatrix}x\\z\end{bmatrix}$. El sistema aumentado es
\begin{equation}
  \begin{bmatrix}\dot x\\ \dot z\end{bmatrix}=
  \begin{bmatrix}\Am-\Bm\Km & \Bm K_e\\ -\Cm & 0\end{bmatrix}
  \begin{bmatrix}x\\ z\end{bmatrix}+\begin{bmatrix}0\\1\end{bmatrix}r,
  \qquad y=\bigl[\Cm\;\;0\bigr]\begin{bmatrix}x\\ z\end{bmatrix}.
  \label{eq:re-integral}
\end{equation}
Si el lazo aumentado es estable, en estado estacionario $\dot z=0$, es decir $y=r$: el error es nulo para \emph{cualquier} valor de los parámetros de la planta, mientras el lazo siga estable. Hay un polo más que ubicar: si no hay ceros cerca, de cinco a diez veces más a la izquierda que los dominantes; si los hay, sobre el cero.

\begin{figure}[ht]
\centering
\includegraphics[width=0.85\linewidth]{cap09/integral_estado}
\caption{Retroalimentación de estado con acción integral: $u=-\Km x+K_ez$, con $\dot z=r-y$.}
\label{fig:re-integral}
\end{figure}

\begin{ejemplo}[label=ej:re-integral]{Ahora con integrador}
Repita el ejemplo~\ref{ej:re-sinintegral} con acción integral, con el tercer polo en $-100$.
\solucion El polinomio característico de \eqref{eq:re-integral} es $s^3+(5+K_2)s^2+(3+K_1)s+K_e$, y el deseado, $(s+100)(s^2+16s+183.1)=s^3+116s^2+1783.1s+18\,310$. Luego
\[
  K_1=1780.1,\qquad K_2=111,\qquad K_e=18\,314
\]
(con $\wn^2$ sin redondear, $183.137$; con $183.1$ saldría $18\,310$). La \cref{fig:re-integralfig} muestra que el error se anula. El precio es el esfuerzo: un tercer polo tan rápido pide picos de $u$ de más de cien veces el escalón.
\end{ejemplo}

\begin{figure}[ht]
\centering
\includegraphics[width=\linewidth]{cap09/f_integral}
\caption{Ejemplos~\ref{ej:re-sinintegral} y~\ref{ej:re-integral}: sin integrador la salida no sigue la referencia; con él la sigue, a cambio de un esfuerzo de control mucho mayor.}
\label{fig:re-integralfig}
\end{figure}

\codigorepo{cap09_retroalimentacion_estado/C9_Ej9_Integral.m}{integrador}{ejemplo \ref{ej:re-integral}}

\begin{cuidado}[Cuando la salida combina estados]
Si $\Cm$ tiene más de un elemento no nulo (por ejemplo $\Cm=[4\;\;1]$), $K_e$ aparece en \emph{varios} coeficientes del polinomio aumentado, no solo en el independiente. Calcule siempre $\det(\lambda\Id-\Am_{aum})$ completo; no copie la forma del ejemplo anterior (problema~\ref{prob:re-integralcero}).
\end{cuidado}

\begin{modelo}[RetroEstado.slx]
El script \cmd{crearModelo\_RetroEstado.m} diseña las dos leyes del ejemplo y construye dos copias de la planta: una con $u=r-\Km x$ y otra con $u=-\Km x+K_ez$. Los \emph{Scope} comparan las salidas y los esfuerzos. Pruebe a cambiar el tercer polo y a perturbar $\Am$ (la planta real): la versión con integrador sigue llegando a $1$.
\end{modelo}

\begin{profundiza}[Retroalimentación de estado con integrador = un PID]
Para una planta de segundo orden en variables de fase, $x_1=y$ y $x_2=\dot y$. La ley $u=-K_1y-K_2\dot y+K_e\int e$ es exactamente un controlador \textbf{I-PD}: la acción integral actúa sobre el error y las acciones proporcional y derivativa, solo sobre la salida (capítulo del PID). La formulación de estado agrega dos cosas: una manera sistemática de ubicar \emph{todos} los polos, y la posibilidad de usar un observador cuando $\dot y$ no se mide o la planta es de orden mayor (problema~\ref{prob:re-ipd}).
\end{profundiza}

%% ---------------------------------------------------------------------
\section{Para llevar}
\label{sec:re-cierre}
%% ---------------------------------------------------------------------

\begin{parallevar}
\begin{enumerate}
  \item La función de transferencia solo muestra la parte controlable y observable. Estabilizable: los modos no controlables son estables; detectable: los no observables lo son. Nunca se cancelan polos inestables.
  \item Retroalimentación de estado: $u=r-\Km x$ y lazo cerrado $\dot x=(\Am-\Bm\Km)x+\Bm r$. Se pueden poner los $n$ polos donde se quiera si y solo si $\operatorname{rango}\Ctrb=n$, con $\Ctrb=[\Bm\;\Am\Bm\;\cdots\;\Am^{n-1}\Bm]$.
  \item $\Km$ por comparación de coeficientes o por Ackermann: $\Km=[0\cdots0\;1]\,\Ctrb^{-1}q(\Am)$. Las ganancias dependen de la representación; los polos, no.
  \item Polos: dos dominantes por las especificaciones y los demás de cinco a diez veces más lejos (o sobre un cero). Más lejos cuesta más esfuerzo.
  \item Observador: $\dot{\est{x}}=(\Am-\Lm\Cm)\est{x}+\Bm u+\Lm y$; el error cumple $\dot\errest=(\Am-\Lm\Cm)\errest$. Requiere $\operatorname{rango}\Obsv=n$. $\Lm$ se diseña como $\Km$ en el sistema dual.
  \item Separación: los polos del lazo con observador son los de control más los del observador. El compensador resultante es una función de transferencia $U/Y$.
  \item Acción integral: $\dot z=r-y$, $u=-\Km x+K_ez$. Un polo más y error nulo ante escalones.
\end{enumerate}
\end{parallevar}

\begin{profundiza}[Correcciones respecto a la versión anterior del capítulo]
(i) El estimado del observador se escribe $\est{x}$ y el error de estimación $\errest$ (antes $\hat{x}$ y $e$, que chocaban con la variable desviación y con el error de seguimiento). (ii) ``Ackerman'' $\to$ Ackermann. (iii) En el servo de tercer orden, el polinomio es $\lambda^3+(3+K_3)\lambda^2+(2+K_2)\lambda+K_1$ (decía $-\lambda(2-K_2)$). (iv) La ubicación de polos da un $\Km$ \emph{único} solo con una entrada. (v) El ejemplo $\operatorname{diag}(-1,-3)$ ahora tiene solución: $\Km=[0.5\;\;-0.5]$. (vi) $K_e=18\,314$ con $\wn^2$ exacto. (vii) Se agregó la prueba PBH y la noción cuantitativa de controlabilidad. Todos los números se verificaron con \texttt{herramientas/cap09\_numeros.py} y \cmd{verificarCap09.m}.
\end{profundiza}

%% ---------------------------------------------------------------------
\section{Problemas}
\label{sec:re-problemas}
%% ---------------------------------------------------------------------

Los problemas marcados con \etiquetaM{} requieren \matlab/\simulink; los marcados con \etiquetaEstrella{} son de profundización.

\subsection*{Controlabilidad y observabilidad}

\begin{problema}[estructura]\label{prob:re-estructura}
Sean (i) $\Am=\operatorname{diag}(-1,-3)$, $\Bm=[1\;\;1]^\top$, $\Cm=[1\;\;0]$ y (ii) $\Am=\operatorname{diag}(-1,-2,-3,-4)$, $\Bm=[1\;\;0\;\;1\;\;0]^\top$, $\Cm=[1\;\;1\;\;0\;\;0]$. (a) Para (i), construya $\Ctrb$ y $\Obsv$ y concluya. (b) Para (i), halle $\Km$ para polos en $\{-2,-2\}$ y simule desde $x(0)=[1\;\;1]^\top$. (c) El sistema (ii) es la forma modal del ejemplo~\ref{ej:re-friedland}: verifique con los rangos que exactamente dos modos son controlables y dos observables, y diga cuáles. (d) Calcule su función de transferencia y relacione su orden con (c). (e) ¿Es estabilizable? ¿Detectable? (f) Reemplace el valor propio $-2$ por $+2$ y repita (e). ¿Qué consecuencia práctica tiene?
\end{problema}

\begin{problema}[controlabilidad cuantitativa]\label{prob:re-cuantitativa}
\etiquetaEstrella{} Para el sistema del ejemplo~\ref{ej:re-motivador}: (a) calcule $\det\Ctrb$ en función de $a$. (b) Con $a=0$, identifique el modo no controlable escribiendo las ecuaciones y concluya si es estabilizable. (c) Con $a=0$ y $a_{11}=+1$, muestre que \emph{ningún} $\Km$ estabiliza el sistema. (d) Aplique la prueba PBH a (b) y (c). (e) Con polos deseados en $-2$ y $-3$, evalúe $\|\Km\|$ y el número de condición de $\Ctrb$ para $a=2$, $1$, $0.1$ y $0.01$. (f) \etiquetaM{} Simule desde $x(0)=[1\;\;1]^\top$ y reporte el pico de $|u|$. ¿Por qué ``controlable sí o no'' no basta para juzgar un diseño?
\end{problema}

\subsection*{Ubicación de polos}

\begin{problema}[Ackermann y especificaciones]\label{prob:re-ackermann}
(a) Para $1/(s(s+3))$ en variables de fase, verifique controlabilidad y calcule $\Km$ a mano con Ackermann para polos en $-4\pm j4$ (reporte $\Ctrb^{-1}$ y $q(\Am)$); verifique con \cmd{acker} y \cmd{place}. (b) Para $1/(s(s+1)(s+2))$ con la configuración de la \cref{fig:re-servo}, diseñe $\Km$ para $\OS\leq5\%$ y $\ts=2$\,s, con el tercer polo diez veces más a la izquierda. (c) \etiquetaM{} Simule y reporte el pico de $u$. (d) Repita con el tercer polo a $5$ y $50$ veces y haga una tabla con $\OS$, $\ts$ y $|u|_{\max}$. ¿Qué precio se paga por alejarlo?
\end{problema}

\begin{problema}[dos representaciones]\label{prob:re-representaciones}
Para $G_p(s)=\dfrac{10}{(s+10)(s^2+0.5s+1)}$ (ejemplo~\ref{ej:re-trespolos}): (a) \etiquetaM{} reproduzca $\Km$ con \cmd{tf2ss} y diga en qué forma está; (b) repita en variables de fase y explique por qué $\Km$ cambia y los polos no; (c) calcule el error ante escalón y agregue el integrador con un cuarto polo en $-50$.
\end{problema}

\subsection*{Observadores y separación}

\begin{problema}[observador]\label{prob:re-observador}
Para $\Am=\begin{bmatrix}0&1\\-2&-3\end{bmatrix}$, $\Cm=[1\;\;0]$: (a) verifique observabilidad; (b) halle $\Lm$ para polos dobles en $-10$, por coeficientes y por Ackermann; (c) \etiquetaM{} simule planta y observador desde $x(0)=[1\;\;0]^\top$, $\est{x}(0)=0$, sin entrada, y mida cuándo el error cae por debajo del $2\%$; (d) repita con polos en $-4$, $-25$ y $-50$; (e) \etiquetaM{} agregue a $y$ ruido de desviación $0.01$ y mida la desviación del error de estimación en régimen permanente para cada observador. Grafíquela contra la ubicación de los polos y discuta.
\end{problema}

\begin{problema}[compensador completo]\label{prob:re-compensador}
Con la planta del problema anterior, diseñe $\Km$ para polos en $-3\pm j3$ y un observador con polos dobles en $-15$. (a) Halle $\Km$ y $\Lm$. (b) Arme la matriz \eqref{eq:re-aumentado} y verifique que sus valores propios son la unión de los dos conjuntos. (c) Obtenga $U(s)/Y(s)$ con \eqref{eq:re-compensador}: ¿a qué compensador clásico se parece? (d) \etiquetaM{} Simule con $u=-\Km x$ y con el observador partiendo de $\errest(0)=[1\;\;-1]^\top$, y cuantifique el deterioro. (e) Verifique que $\errest(t)$ no depende de la referencia.
\end{problema}

\begin{problema}[robustez del compensador]\label{prob:re-robustez}
\etiquetaEstrella{}\etiquetaM{} Para la planta del problema~\ref{prob:re-representaciones}, con el integrador y un observador cinco veces más rápido que los polos de control: (a) verifique la separación; (b) suponga que el polo real verdadero está en $-8$ y no en $-10$, pero $\Km$ y $\Lm$ se calcularon con el modelo nominal: ¿sigue estable? Encuentre el rango de ubicaciones de ese polo con el que lo es. (c) Perturbe ahora la ganancia de la planta y halle el margen de ganancia efectivo, con y sin observador. ¿Qué dice esto sobre el principio de separación?
\end{problema}

\subsection*{Acción integral}

\begin{problema}[integrador con un cero]\label{prob:re-integralcero}
Para $\Am=\begin{bmatrix}0&1\\-7&-9\end{bmatrix}$, $\Bm=[0\;\;1]^\top$, $\Cm=[4\;\;1]$: (a) verifique controlabilidad y observabilidad y halle $Y/U$ y su cero. (b) Diseñe $u=r-\Km x$ para $\OS=10\%$ y $\tp=2$\,s y calcule el error ante escalón. (c) Agregue el integrador y escriba el polinomio aumentado: ¿por qué $K_e$ aparece en dos coeficientes? (d) Ubique el tercer polo sobre el cero y halle $\Km$ y $K_e$. (e) \etiquetaM{} Verifique que el error es nulo y compare con (b). (f) Repita (d) con el tercer polo en $-8$ y en $-20$ y explique por qué cancelar el cero da el mejor transitorio.
\end{problema}

\begin{problema}[la conexión con el PID]\label{prob:re-ipd}
Para $G_p=\dfrac{1}{(s+1)(s+2)}$ en variables de fase ($x_1=y$, $x_2=\dot y$): (a) escriba $\Am$, $\Bm$, $\Cm$; (b) arme el sistema aumentado con el integrador y su polinomio en función de $K_1$, $K_2$, $K_e$; (c) diseñe para $\OS=10\%$, $\ts=2$\,s y un tercer polo cinco veces más lejos; (d) escriba la ley en el tiempo con $x_1=y$, $x_2=\dot y$, $z=\int e$ y muestre que es un I-PD, identificando $K_c$, $K_i$ y $K_d$; (e) \etiquetaM{} compruébelo numéricamente; (f) ¿qué aporta la formulación de estado si la planta es de orden mayor o si $\dot y$ no se mide?
\end{problema}

%% ---------------------------------------------------------------------
\section{Soluciones seleccionadas}
%% ---------------------------------------------------------------------

\paragraph{Problema~\ref{prob:re-estructura}}
(a) $\Ctrb=\begin{bmatrix}1&-1\\1&-3\end{bmatrix}$ (rango 2: controlable); $\Obsv=\begin{bmatrix}1&0\\-1&0\end{bmatrix}$ (rango 1: el modo $-3$ no se ve). (b) $\Km=[0.5\;\;-0.5]$. (c) Controlables: $-1$ y $-3$; observables: $-1$ y $-2$. (d) $1/(s+1)$: solo el modo controlable y observable. (e) Estabilizable y detectable (los modos escondidos son estables). (f) El modo $+2$ es observable pero no controlable: el sistema es detectable pero \emph{no estabilizable}; ningún controlador lo estabiliza, y la salida crecerá aunque $u$ haga lo que haga.

\paragraph{Problema~\ref{prob:re-cuantitativa}}
(a) $\det\Ctrb=-a$. (b) Con $a=0$, $\dot x_1=-x_1$: el modo $-1$ no es controlable, pero es estable: estabilizable. (c) Con $a_{11}=+1$ el polinomio es $(\lambda-1)(\lambda+2+K_2)$ para todo $\Km$: no estabilizable. (e) $\|\Km\|=4.47$, $5.39$, $23.1$ y $203$; número de condición $4.3$, $5.8$, $50$ y $500$. (f) Pico de $|u|$: $6$, $7$, $25$ y $205$.

\paragraph{Problema~\ref{prob:re-ackermann}}
(a) $\Km=[32\;\;5]$ (pausa de la \cref{sec:re-ackermann}). (b) $\zeta=0.690$, $\wn=2.90$, polos $-2\pm j2.10$ y $-20$: $\Km=[168.0\;\;86.4\;\;21]$. (d) Con el tercer polo a $5$, $10$ y $50$ veces: $\OS=4.73$, $4.94$ y $5.00\%$; $\ts=2.17$, $2.12$ y $2.08$\,s; $|u|_{\max}=84$, $168$ y $840$. Alejar el polo casi no mejora la respuesta y multiplica el esfuerzo.

\paragraph{Problema~\ref{prob:re-observador}}
(b) $\Lm=[17\;\;47]^\top$. (d) Con polos dobles en $-4$, $-25$ y $-50$: $\Lm=[5\;\;-1]^\top$, $[47\;\;482]^\top$ y $[97\;\;2207]^\top$: la ganancia crece con el cuadrado de la velocidad, y con ella la amplificación del ruido.

\paragraph{Problema~\ref{prob:re-compensador}}
(a) $\Km=[16\;\;3]$ y $\Lm=[27\;\;142]^\top$. (b) $\{-3\pm j3,\,-15,\,-15\}$.

\paragraph{Problema~\ref{prob:re-integralcero}}
(a) $\dfrac{s+4}{s^2+9s+7}$, cero en $-4$. (b) $\zeta=0.591$, $\wn=1.95$: $\Km=[-3.21\;\;-6.70]$ y valor final $1.055$ (error $-0.055$). (c) $\lambda^3+(9+K_2)\lambda^2+(7+K_1+K_e)\lambda+4K_e$: $K_e$ entra por los dos elementos de $\Cm$. (d) $\Km=[2.21\;\;-2.70]$, $K_e=3.79$; $\OS=10\%$ y $\tp=2$\,s exactos. (f) Con $-8$: $\OS=11.4\%$, $\tp=1.82$\,s; con $-20$: $11.7\%$ y $1.73$\,s. Cancelando el cero, el lazo es exactamente el segundo orden especificado.

\paragraph{Problema~\ref{prob:re-ipd}}
(b) $\lambda^3+(3+K_2)\lambda^2+(2+K_1)\lambda+K_e$. (c) $\zeta\wn=2$, $\wn=3.38$; con el polo en $-10$: $K_1=49.45$, $K_2=11$, $K_e=114.46$. (d) $u=K_e\int_0^t e\,\dif\tau-K_1y-K_2\dot y$: I-PD con $K_i=K_e$, $K_c=K_1$ y $K_d=K_2$ (proporcional y derivativa sobre la salida).

\respuestaspausas

       % migrado (Handout7_Ctrl.tex)
%% =====================================================================
%%  Capítulo 10 -- Respuesta frecuencial y diagramas de Bode
%%  PILOTO del formato nuevo (noviembre de 2026).
%%  Integra: Handout8 (versión anterior), notas de la Clase 1, las cuatro
%%  presentaciones Bode1-Bode4 (v3) y los códigos de Bode_4Clases/matlab.
%%  Correcciones respecto al handout anterior (ver \S\ref{sec:rf-cierre}):
%%    adelanto w_M = 8.85 (no 14); atraso con holgura w = 1.68, PM = 44.8°;
%%    rango de Kc: GM = +40 dB (no -40 dB); w_GM asintótico 31.6 (no 50);
%%    Kc exacto 33 056; se usa M_r (no M_P) para el pico de resonancia.
%%  Sin retardos (decisión del curso): el PM se explica como presupuesto de
%%  fase frente a dinámicas no modeladas (un sensor de primer orden).
%% =====================================================================

\capitulo[cap:rf]{Respuesta frecuencial y diagramas de Bode}
  {Compramos libros porque creemos que compramos el tiempo para leerlos.}
  {Warren Zevon}

\begin{objetivos}[partes/10-respuesta-frecuencial.html]
  \item Explicar por qué un sistema lineal y estable responde a una senoide con otra senoide, y calcular su amplitud y su fase.
  \item Dibujar a mano el diagrama de Bode asintótico de una malla abierta con ganancia, integradores y polos y ceros reales, y estimar cuánto se equivoca.
  \item Leer en un diagrama de Bode los márgenes de ganancia y de fase, decidir si el lazo cerrado es estable y cuánta ganancia admite.
  \item Predecir el sobrepico, la velocidad y el error en estado estacionario del lazo cerrado mirando solamente la malla abierta.
  \item Diseñar un controlador proporcional, un compensador en adelanto, en atraso o atraso-adelanto para cumplir especificaciones de error y de margen de fase, y verificar el diseño en \matlab/\simulink.
\end{objetivos}

%% ---------------------------------------------------------------------
\section{¿Por qué la frecuencia?}
\label{sec:rf-porque}
%% ---------------------------------------------------------------------

Hasta ahora hemos analizado los sistemas de control de dos maneras: con la función de transferencia, mirando dónde quedan los polos en el plano-$s$, y con la representación de estado, mirando los valores propios de $\Am$. Las dos exigen \emph{tener el modelo}. La respuesta frecuencial cambia la pregunta. En lugar de preguntar dónde están los polos, pregunta
\begin{center}
\sffamily\large\color{ctrlazul}¿qué le hace el sistema a una senoide?
\end{center}
La respuesta, como veremos en la \cref{sec:rf-senoide}, son dos números por cada frecuencia: cuánto cambia la amplitud y cuánto se corre la fase. Lo valioso es que esos dos números se pueden \emph{medir} en un laboratorio con un generador de señales y un osciloscopio, aun cuando nadie haya escrito nunca la función de transferencia. Y, sobre todo, con esos dos números dibujados en una sola gráfica, el \clave{diagrama de Bode}, vamos a poder responder las tres preguntas del curso (estabilidad, respuesta transitoria y error) y además \emph{diseñar}: la \cref{fig:rf-mapa} es el mapa del capítulo.

\begin{figure}[ht]
\centering
\begin{tikzpicture}[font=\sffamily\small,
  caja/.style={draw=#1,fill=#1!8,thick,rounded corners=2pt,align=center,minimum height=10mm,text width=27mm},
  flecha/.style={-{Stealth[length=2.2mm]},thick,draw=ctrlgris}]
  \node[caja=ctrlnoche] (sen) {una senoide entra\\ y sale otra};
  \node[caja=ctrlazul,right=6mm of sen] (bode) {diagrama de Bode\\ de $\Gma(\jw)$};
  \node[caja=ctrlverde,right=9mm of bode,yshift=13mm] (est) {estabilidad\\ \footnotesize $\GM$, $\PM$ (\S\ref{sec:rf-estabilidad})};
  \node[caja=ctrlvioleta,right=9mm of bode] (tra) {transitorio\\ \footnotesize $\PM\to\zeta$, $\wbw$ (\S\ref{sec:rf-lazocerrado})};
  \node[caja=ctrlambar,right=9mm of bode,yshift=-13mm] (err) {error\\ \footnotesize tipo, $\kp,\kv,\ka$ (\S\ref{sec:rf-tipo})};
  \node[caja=ctrlcoral,below=7mm of bode] (dis) {diseño: $\Kc$, adelanto,\\ atraso (\S\ref{sec:rf-ajuste}--\ref{sec:rf-pid})};
  \draw[flecha] (sen) -- (bode);
  \draw[flecha] (bode.east) -- (est.west);
  \draw[flecha] (bode.east) -- (tra.west);
  \draw[flecha] (bode.east) -- (err.west);
  \draw[flecha] (bode) -- (dis);
\end{tikzpicture}
\caption{Mapa del capítulo. Todo sale de un mismo dibujo: la malla abierta $\Gma(\jw)$ en magnitud y fase.}
\label{fig:rf-mapa}
\end{figure}

%% ---------------------------------------------------------------------
\subsection{Un poco de historia: el control nace en un laboratorio de teléfonos}

En los años veinte, las líneas telefónicas transcontinentales de AT\&T necesitaban cientos de amplificadores de tubos de vacío en cascada. La distorsión de cada uno se acumulaba y la ganancia derivaba con la temperatura y con el envejecimiento de los tubos. Harold S.\ Black propuso en 1927 una solución que hoy parece obvia: \emph{retroalimentar} negativamente el amplificador, sacrificando ganancia a cambio de precisión y linealidad~\citep{black34}.\nota{Black cuenta que la idea se le ocurrió el 2 de agosto de 1927 en el ferry de Lackawanna, camino al trabajo, y que la anotó en una página del \emph{New York Times}~\citep{black77}. La patente, presentada en 1928, se concedió en 1937~\citep{blackpat}.} El precio fue un problema nuevo: con mucha ganancia de lazo, los amplificadores \emph{oscilaban}.

\notamargen{\textbf{Pregunta para la clase.} ¿Por qué una compañía telefónica necesitaba una teoría de estabilidad?}
Harry Nyquist dio en 1932 el criterio para predecir la estabilidad del lazo cerrado a partir de mediciones de la malla abierta~\citep{Nyq32}, y Hendrik W.\ Bode, entre 1940 y 1945, convirtió esas ideas en un método de diseño de regla y lápiz: las gráficas de magnitud en decibeles y de fase contra el logaritmo de la frecuencia que hoy llevan su nombre, las aproximaciones asintóticas, los márgenes de ganancia y de fase, y la relación entre la pendiente de la magnitud y la fase en los sistemas de fase mínima~\citep{Bod40,bode45}. Durante la Segunda Guerra Mundial estos métodos salieron de los teléfonos y llegaron a los servomecanismos de radares y cañones~\citep{JaNiPh47}; el lugar de las raíces de Evans llegaría después~\citep{Eva50}. La \cref{fig:rf-linea} resume la cronología.

\begin{figure}[ht]
\centering
\includegraphics[width=\linewidth]{cap10/f_linea_tiempo}
\caption{De Black a Evans: casi todo el control clásico en frecuencia salió de Bell Labs (en ámbar, los aportes de Bode).}
\label{fig:rf-linea}
\end{figure}

\begin{historia}[Hendrik Wade Bode (1905--1982)]
Ingeniero de redes y amplificadores en Bell Labs. Su problema era diseñar retroalimentación con mucha ganancia y sin oscilaciones. En 1940 mostró que en los sistemas de fase mínima la magnitud determina la fase: basta mirar la pendiente. En la guerra trabajó en directores de tiro antiaéreo y en 1967 pasó a Harvard como profesor de ingeniería de sistemas. Su gran aporte práctico fueron las aproximaciones con líneas rectas, que permiten dibujar a mano, en minutos, el comportamiento de un sistema en \emph{todas} las frecuencias.
\end{historia}

%% ---------------------------------------------------------------------
\subsection{El lenguaje: malla abierta, lazo abierto y \emph{open loop}}
\label{sec:rf-lenguaje}

En el curso distinguimos dos funciones que en los libros en inglés suelen aparecer con el mismo nombre. Considere el lazo de la \cref{fig:rf-mallalazo}(a), con camino directo $G(s)$ y retroalimentación $\Hs(s)$.

\begin{figure}[ht]
\centering
\includegraphics[width=\linewidth]{cap10/malla_vs_lazo}
\caption{(a) La malla abierta es lo que recorre una señal al dar una vuelta completa al lazo. (b) El lazo abierto aparece al reescribir el mismo lazo con retroalimentación unitaria.}
\label{fig:rf-mallalazo}
\end{figure}

\begin{definicion}[label=def:rf-malla]{Malla abierta y lazo abierto}
La \clave{malla abierta} es el producto de todo lo que recorre una señal al dar una vuelta completa al lazo, \emph{sin} el signo del comparador:
\begin{equation}
  \Gma(s)=G(s)\,\Hs(s).
  \label{eq:rf-mallaabierta}
\end{equation}
Es la función con la que se estudia la estabilidad del lazo cerrado y la que se dibuja en el diagrama de Bode. Si el lazo tiene válvula, planta, sensor y transmisor, $\Gma=\Gv\,\Gp\,\Gs\,\Gt$.

El \clave{lazo abierto} es la función que relaciona el error con la salida cuando el lazo se reescribe con retroalimentación unitaria:
\begin{equation}
  \Gla(s)=\frac{G(s)}{1+G(s)\Hs(s)-G(s)}.
  \label{eq:rf-lazoabierto}
\end{equation}
Es la que se usa para el tipo del sistema y los errores en estado estacionario (capítulo de error).
\end{definicion}

Las dos llevan a la misma ecuación característica,
\begin{equation}
  1+G(s)\Hs(s)=0\quad\Longleftrightarrow\quad 1+\Gla(s)=0,
  \label{eq:rf-ecuacioncaracteristica}
\end{equation}
y cuando $\Hs(s)=1$ (el caso de casi todos nuestros ejemplos) coinciden: $\Gma(s)=\Gla(s)=G(s)$. Por eso los textos las mezclan. Ogata, por ejemplo, llama \emph{open-loop transfer function} a $B(s)/E(s)=G(s)H(s)$, que es nuestra malla abierta, y \emph{feedforward transfer function} a $C(s)/E(s)=G(s)$~\citep{ogata10}; Åström y Murray la llaman \emph{loop transfer function}, $L=PC$~\citep{AsMu10}; en electrónica se dice \emph{loop gain}, $A\beta$. La regla para no confundirse es preguntar \emph{qué se corta y qué se mide}.

\begin{cuidado}
Con $\Gma(s)$ estudiamos estabilidad y diseñamos en frecuencia; con $\Gla(s)$ calculamos el tipo y el error. Cuando $\Hs\neq1$ son funciones \emph{distintas}: un sensor con dinámica cambia la malla abierta y, con ella, los márgenes, aunque el error se calcule con otra función.
\end{cuidado}

\paragraph{Cortar el lazo}
Imagine que corta el lazo justo después del comparador (\cref{fig:rf-corte}), inyecta una señal $x$ en el corte y mide lo que regresa, $z$. La señal recorre $G$, luego $\Hs$, y llega al comparador, que le cambia el signo:
\begin{equation}
  \frac{Z(s)}{X(s)}=-\,G(s)\,\Hs(s)=-\,\Gma(s).
\end{equation}
Si para alguna frecuencia $\omega$ lo que vuelve es exactamente $-x$ (misma amplitud, invertido), es decir, si $\Gma(\jw)=-1$, el comparador lo convierte en $+x$: el lazo se alimenta solo y oscila sin necesidad de entrada. Por eso el punto crítico del análisis en frecuencia es $-1$ (magnitud $1=0\dB$ y fase $-180\grados$), y por eso el signo del comparador \emph{no} se incluye en $\Gma(s)$. Bode llamó \clave{diferencia de retorno} a $1+\Gma(s)$: la diferencia entre la señal inyectada y la que regresa. Sus ceros son los polos del lazo cerrado,
\begin{equation}
  \Tlc(s)=\frac{Y(s)}{R(s)}=\frac{G(s)}{1+\Gma(s)}.
\end{equation}

\begin{figure}[ht]
\centering
\includegraphics[width=0.62\linewidth]{cap10/cortar_lazo}
\caption{Cortar el lazo: se inyecta $x$ y se mide lo que regresa, $z=-\Gma x$.}
\label{fig:rf-corte}
\end{figure}

\paragraph{El mismo lenguaje en electrónica}
\notamargen{En las hojas de datos de los amplificadores operacionales, la \emph{open-loop gain} es $A$, del orden de $10^5$.}
En el amplificador de Black (\cref{fig:rf-black}), $A$ es la ganancia del amplificador sin retroalimentar y $\beta$ la fracción de la salida que se resta a la entrada:
\begin{equation}
  A_f=\frac{A}{1+A\beta}\ \xrightarrow{\ A\beta\gg 1\ }\ \frac{1}{\beta}.
\end{equation}
Con $A=10^5$ y $\beta=0.01$ resulta $A_f=99.90$; si los tubos envejecen y $A$ cae a la mitad, $A_f=99.80$. La ganancia del amplificador cayó un $50\%$ y la del amplificador retroalimentado apenas un $0.1\%$. Es la misma razón por la que retroalimentamos plantas: la ganancia de lazo $A\beta\leftrightarrow\Gma(s)$ compra insensibilidad a las variaciones del sistema.

\begin{figure}[ht]
\centering
\includegraphics[width=0.5\linewidth]{cap10/black}
\caption{El amplificador retroalimentado de Black: $A\beta$ es la malla abierta.}
\label{fig:rf-black}
\end{figure}

%% ---------------------------------------------------------------------
\section{Una senoide entra, una senoide sale}
\label{sec:rf-senoide}
%% ---------------------------------------------------------------------

\begin{definicion}[label=def:rf-respuesta]{Respuesta en frecuencia}
La \clave{respuesta en frecuencia} de un sistema lineal, invariante en el tiempo y \emph{estable} es su respuesta en estado estacionario a una entrada senoidal. La salida es una senoide de la \emph{misma} frecuencia; solo cambian la amplitud y la fase.
\end{definicion}

El supuesto de estabilidad (en malla abierta, cuando lo que se excita es $\Gma(s)$) es esencial: si hay polos en el semiplano derecho, los transitorios no se apagan y no hay un estado estacionario que medir. Lo mantendremos durante todo el capítulo.

\subsection{¿Por qué sale una senoide? Deducción}
\label{sec:rf-deduccion}

Sea $x(t)=x_0\sin(\omega t)$, cuya transformada es
\begin{equation}
  X(s)=\frac{x_0\,\omega}{s^2+\omega^2}=\frac{x_0\,\omega}{(s-\jw)(s+\jw)}.
\end{equation}
Si $\Gma(s)$ tiene polos $-p_1,\dots,-p_n$ (distintos, para simplificar), la salida se expande en fracciones parciales como
\begin{equation}
  Z(s)=\Gma(s)X(s)=\underbrace{\frac{K_1}{s+p_1}+\cdots+\frac{K_n}{s+p_n}}_{\text{modos del sistema}}
  +\underbrace{\frac{a}{s-\jw}+\frac{\bar a}{s+\jw}}_{\text{modo de la entrada}}.
\end{equation}
En el tiempo, los términos $K_ie^{-p_it}$ se extinguen si todos los polos tienen parte real negativa. Lo que queda es el modo de la entrada. El residuo en $s=\jw$ es
\begin{equation}
  a=\Big[(s-\jw)\,\Gma(s)\,\frac{x_0\,\omega}{(s-\jw)(s+\jw)}\Big]_{s=\jw}
  =\Gma(\jw)\,\frac{x_0}{2j},
\end{equation}
y el de $s=-\jw$ es su conjugado. Escribiendo $\Gma(\jw)=|\Gma(\jw)|\,e^{j\phi}$, con $\phi=\angle\Gma(\jw)$, y usando que $\Gma(-\jw)$ es el conjugado de $\Gma(\jw)$ (coeficientes reales):
\begin{equation}
  z_{ss}(t)=a\,e^{\jw t}+\bar a\,e^{-\jw t}
  =x_0\,|\Gma(\jw)|\,\frac{e^{j(\omega t+\phi)}-e^{-j(\omega t+\phi)}}{2j}
  =x_0\,|\Gma(\jw)|\,\sin(\omega t+\phi).
  \label{eq:rf-zss}
\end{equation}

\begin{ideaclave}
Para cada frecuencia, el sistema \emph{multiplica} la amplitud por $|\Gma(\jw)|$ y \emph{desplaza} la fase en $\phi=\angle\Gma(\jw)$. El diagrama de Bode es la gráfica de esos dos números en función de $\omega$.
\end{ideaclave}

\begin{ejemplo}[label=ej:rf-seno]{Un primer orden excitado con $\sin t$}
Calcule la salida en estado estacionario de $\Gma(s)=1/(s+10)$ cuando $x(t)=\sin t$.
\solucion Con $\omega=1$\,rad/s,
\[
  \Gma(j1)=\frac{1}{10+j}\ \Rightarrow\ |\Gma(j1)|=\frac{1}{\sqrt{101}}=0.0995,\qquad
  \phi=-\tan^{-1}\frac{1}{10}=-5.71\grados,
\]
y por \eqref{eq:rf-zss}, $z_{ss}(t)=0.0995\,\sin(t-5.71\grados)$. La salida es diez veces más pequeña y llega $5.7\grados$ tarde. Eso es \emph{todo} lo que el diagrama de Bode guarda para $\omega=1$. La \cref{fig:rf-seno} muestra que la simulación completa coincide con esta predicción en cuanto se apaga el transitorio $e^{-10t}$ (unas cuatro constantes de tiempo, $0.4$\,s).
\end{ejemplo}

\begin{figure}[ht]
\centering
\includegraphics[width=\linewidth]{cap10/f_seno_primer_orden}
\caption{$1/(s+10)$ excitado con $\sin t$: la simulación completa (azul) y la predicción de estado estacionario \eqref{eq:rf-zss} (coral) coinciden después del transitorio.}
\label{fig:rf-seno}
\end{figure}

\script[lsim y la fórmula dan lo mismo]{cap10_respuesta_frecuencial/Clase1_RespuestaFrecuencial/C1_Ej1_RespuestaSenoidal.m}
\codigorepo{cap10_respuesta_frecuencial/Clase1_RespuestaFrecuencial/C1_Ej1_RespuestaSenoidal.m}{senoide}{ejemplo \ref{ej:rf-seno}}

\video{V1_maquina_bode}{La máquina de Bode: se barre la frecuencia, se miden los dos números y el Bode se dibuja punto a punto.}
\begin{modelo}[RespuestaSenoidal.slx]
El script \cmd{crearModelo\_RespuestaSenoidal.m} construye por código el diagrama \emph{Sine Wave} $\to$ $G(s)$ $\to$ \emph{Scope}, lo simula y mide automáticamente la amplitud y el desfase en el último periodo, comparándolos con \cmd{bode(G,w)}. Cambie \cmd{w}, \cmd{num} y \cmd{den} en el \emph{workspace} y repita: está midiendo el diagrama de Bode como se hacía en Bell Labs, un punto a la vez.
\end{modelo}

\begin{pausa}
Sin calcular nada: si la entrada de $1/(s+10)$ fuera $\sin(100t)$, ¿la salida sería más grande o más pequeña que la del \cref{ej:rf-seno}? ¿Llegaría más o menos atrasada?
\respuesta{Más pequeña y más atrasada: $|\Gma(j100)|=1/\sqrt{10\,100}=0.00995$ (diez veces menor que en $\omega=1$) y $\phi=-\tan^{-1}10=-84.3\grados$. Un polo real atenúa y atrasa más a medida que la frecuencia sube.}
\end{pausa}

\subsection{¿Qué ganamos trabajando en frecuencia?}
\begin{enumerate}
  \item \textbf{Se mide.} Hay generadores de senoides en rangos enormes de frecuencia y amplitud: el modelo se puede obtener experimentalmente.
  \item \textbf{Se multiplica sumando.} En escala logarítmica, los bloques en cascada se suman (\cref{sec:rf-decibel}).
  \item \textbf{Se diseña a ojo.} La estabilidad relativa (márgenes), el error en estado estacionario y la velocidad de respuesta se leen en el mismo diagrama.
  \item \textbf{Se filtra.} El ruido de medición y las perturbaciones suelen vivir en bandas de frecuencia conocidas; el diseñador ve de inmediato qué deja pasar el lazo.
\end{enumerate}

%% ---------------------------------------------------------------------
\section{El decibel}
\label{sec:rf-decibel}
%% ---------------------------------------------------------------------

\subsection{Un invento de ingenieros telefónicos}

Antes de 1923, la pérdida de una línea telefónica se expresaba en \emph{millas de cable estándar} (MSC): cuántas millas de un cable de referencia atenuaban lo mismo a $5000$\,rad/s ($\approx796$\,Hz, una frecuencia representativa de la voz). La unidad dependía de la frecuencia y del cable. En 1924, W.~H.~Martin, de Bell Labs, propuso reemplazarla por la \emph{transmission unit} (TU): diez veces el logaritmo decimal de la razón de potencias~\citep{martin24}. La TU casi coincidía con la milla de cable,
\begin{equation}
  1\ \text{MSC}\approx\frac{1}{1.056}\ \text{TU}\approx 0.95\ \text{dB},
\end{equation}
de modo que los ingenieros no tuvieron que cambiar su intuición. En 1928 el Bell System adoptó el nombre \emph{decibel}, en honor a Alexander Graham Bell, siguiendo la propuesta internacional de llamar \emph{bel} a la unidad decimal y \emph{neper} (por John Napier) a la natural~\citep{martin29}:
\begin{gather}
  \text{bel}=\log_{10}\frac{P_2}{P_1},\qquad
  \text{dB}=10\log_{10}\frac{P_2}{P_1},\qquad
  1\ \text{Np}=20\log_{10}e=8.686\ \text{dB}.
\end{gather}
El bel resultó demasiado grande para el trabajo diario; sobrevivió su décima parte.

\subsection{¿Por qué 10 log y por qué 20 log?}

El decibel se define para \emph{potencias}. Si la potencia se disipa sobre la misma resistencia, $P=V^2/R$ y
\begin{equation}
  10\log_{10}\frac{V_2^2/R}{V_1^2/R}=20\log_{10}\frac{V_2}{V_1}.
\end{equation}
En el diagrama de Bode $|\Gma(\jw)|$ es una razón de \emph{amplitudes} (salida sobre entrada); por eso siempre se usa
\begin{equation}
  |\Gma|_{\mathrm{dB}}=20\log_{10}|\Gma(\jw)|.
\end{equation}

\begin{table}[ht]
\centering\small
\begin{tabular}{@{}cc l@{}}
\toprule
$|G|$ & dB & para recordar\\
\midrule
$10$ & $+20$ & una década de amplitud\\
$2$ & $+6.02$ & el doble\\
$\sqrt2$ & $+3.01$ & la esquina de un cero real\\
$1$ & $0$ & cruce de ganancia\\
$1/\sqrt2$ & $-3.01$ & ancho de banda; esquina de un polo real\\
$0.1$ & $-20$ & una década menos\\
\bottomrule
\end{tabular}
\caption{Valores que conviene memorizar. La mitad de la \emph{potencia} es $-3$\,dB: por eso la frecuencia de corte de un filtro se define donde la magnitud cae $3$\,dB.}
\label{tab:rf-db}
\end{table}

\paragraph{Tres razones para usar logaritmos}
\begin{enumerate}
  \item \textbf{Cascadas.} Las ganancias se multiplican; en dB se suman, y las fases también se suman:
  \[
    20\log|G_1G_2G_3|=20\log|G_1|+20\log|G_2|+20\log|G_3|,\qquad
    \angle(G_1G_2G_3)=\angle G_1+\angle G_2+\angle G_3 .
  \]
  Por ejemplo, $\times10\cdot\times0.5\cdot\times4=\times20$ en amplitud equivale a $20-6.02+12.04=26.02$\,dB.
  \item \textbf{Rangos enormes.} Una razón de amplitudes de $10^6$ son ``solo'' $120$\,dB.
  \item \textbf{Percepción.} La sonoridad que percibimos crece aproximadamente con el logaritmo de la intensidad; un cambio de $1$\,dB apenas se nota.
\end{enumerate}

\notamargen{Cuando se agrega una referencia, el decibel deja de ser una razón y pasa a ser un \emph{nivel}: dBm ($1$\,mW), dBV ($1$\,V), dB SPL ($20\,\mu$Pa). Un número en dB sin referencia es siempre una razón.}

\begin{pausa}
Un amplificador de $+26$\,dB se conecta en cascada con un filtro que, a cierta frecuencia, atenúa $-32$\,dB. ¿Por cuánto queda multiplicada la amplitud de una senoide de esa frecuencia?
\respuesta{$26-32=-6$\,dB, es decir, $10^{-6/20}=0.50$: la amplitud queda a la mitad. En dB las cascadas se suman.}
\end{pausa}

%% ---------------------------------------------------------------------
\section{La forma de Bode y los cuatro factores}
\label{sec:rf-factores}
%% ---------------------------------------------------------------------

Toda función de transferencia racional se puede escribir en la \clave{forma de Bode}:
\begin{equation}
  \Gma(s)=\frac{\Kb\prod_{i=1}^{m}(1+\tau_i s)}
  {s^N\prod_{l=1}^{M}(1+\tau_l s)\prod_{k=1}^{n-M-N}\Big(\frac{s^2}{\omega_{n_k}^2}+\frac{2\zeta_k}{\omega_{n_k}}s+1\Big)},
  \label{eq:rf-formabode}
\end{equation}
con $m$ ceros reales, $N$ polos en el origen, $M$ polos reales y el resto en pares complejos conjugados; el grado relativo es $n-m$. (Si hay ceros complejos se tratan igual que los polos complejos, con el signo cambiado.) Lo importante es que cada factor está escrito como $(1+\tau s)$ y \emph{no} como $(s+a)$: así cada factor vale $1$ ($0\dB$) en bajas frecuencias y la ganancia $\Kb$ aparece sola. Tomando $20\log$ y el ángulo,
\begin{align}
  20\log|\Gma|&=20\log \Kb+\sum_i 20\log|1+\jw\tau_i|-20N\log\omega
  -\sum_l 20\log|1+\jw\tau_l|-\sum_k 20\log\Big|1-\tfrac{\omega^2}{\omega_{n_k}^2}+j\tfrac{2\zeta_k\omega}{\omega_{n_k}}\Big|,\notag\\
  \phi(\omega)&=\sum_i\tan^{-1}(\omega\tau_i)-N\cdot 90\grados
  -\sum_l\tan^{-1}(\omega\tau_l)
  -\sum_k\tan^{-1}\frac{2\zeta_k\omega_{n_k}\omega}{\omega_{n_k}^2-\omega^2}.
  \label{eq:rf-sumas}
\end{align}
El diagrama completo es la \emph{suma} de los diagramas de cada factor. Basta entonces aprender cuatro recetas.

%% ---------------------------------------------------------------------
\subsection{Factor 1: ganancia constante}
\[
  \Gma(s)=\Kb:\qquad 20\log|\Kb|\ \text{dB},\qquad
  \phi=\begin{cases}0\grados & \Kb>0,\\ \pm180\grados & \Kb<0.\end{cases}
\]
Una ganancia sube o baja \emph{toda} la curva de magnitud y no toca la fase. Esta idea, que parece trivial, es la base del diseño por ajuste de ganancia (\cref{sec:rf-ajuste}).

\subsection{Factor 2: polos y ceros en el origen}
\[
  \Gma(\jw)=\frac{1}{\jw}\ \Rightarrow\ 20\log\Big|\frac{1}{\jw}\Big|=-20\log\omega,\qquad\phi=-90\grados.
\]
En escala logarítmica es una recta de $-20\dBdec$ que pasa por $0\dB$ en $\omega=1$. En general $1/s^N$ da $-20N\dBdec$ y $-90N\grados$, y $s^N$ da $+20N\dBdec$ y $+90N\grados$. Todas estas rectas pasan por $0\dB$ en $\omega=1$ (\cref{fig:rf-simples}). El número de integradores, $N$, es el \emph{tipo} del sistema, y se lee en la pendiente de baja frecuencia (\cref{sec:rf-tipo}).

\begin{figure}[ht]
\centering
\includegraphics[width=\linewidth]{cap10/f_factores_simples}
\caption{Factores 1 y 2: una ganancia es una recta horizontal que no toca la fase; un integrador es una recta de $-20\dBdec$ que pasa por $0\dB$ en $\omega=1$ con fase constante de $-90\grados$.}
\label{fig:rf-simples}
\end{figure}

\subsection{Factor 3: polos y ceros reales}
Para el cero $\Gma(s)=s/\omega_1+1$,
\[
  20\log|\Gma(\jw)|=20\log\sqrt{\Big(\frac{\omega}{\omega_1}\Big)^2+1}\ \approx\
  \begin{cases}0\dB, & \omega\ll\omega_1,\\[2pt]
  20\log\dfrac{\omega}{\omega_1}\ \ (+20\dBdec), & \omega\gg\omega_1,\end{cases}
\]
y en la esquina $\omega=\omega_1$ el valor real es $20\log\sqrt2=3\dB$: el error máximo de la asíntota. La fase es $\phi=\tan^{-1}(\omega/\omega_1)$; su aproximación asintótica vale $0\grados$ hasta $\omega_1/10$, sube $45\grados$ por década pasando por $45\grados$ en $\omega_1$ y se queda en $90\grados$ desde $10\,\omega_1$ (el error en los extremos es $\tan^{-1}0.1=5.7\grados$). El polo es el espejo: $-20\dBdec$ desde $\omega_1$, $-3\dB$ en la esquina y fase de $0\grados$ a $-90\grados$ (\cref{fig:rf-polo}).

\begin{figure}[ht]
\centering
\includegraphics[width=0.88\linewidth]{cap10/f_polo_real}
\caption{Polo real: asíntotas (negro) y curva real (verde). El error máximo en magnitud es de $3\dB$ en la esquina; en fase, de $5.7\grados$ en $\omega_1/10$ y en $10\,\omega_1$.}
\label{fig:rf-polo}
\end{figure}

\begin{cuidado}[La trampa clásica: la ganancia escondida]
La regla solo vale si el factor está escrito como $s/\omega_1+1$. Como $s+10=10\,(s/10+1)$, el factor $1/(s+10)$ trae una ganancia $1/10=-20\dB$ escondida. Si se olvida sacarla, toda la curva queda corrida $20\dB$. Es el error más frecuente en las evaluaciones.
\end{cuidado}

\subsection{Factor 4: polos complejos (lo esencial)}
\label{sec:rf-complejos}
Para el factor normalizado de segundo orden,
\begin{equation}
  \Gma(s)=\frac{\wn^2}{s^2+2\zeta\wn s+\wn^2}\ \Rightarrow\
  |\Gma(\jw)|=\frac{\wn^2}{\sqrt{(\wn^2-\omega^2)^2+4\zeta^2\wn^2\omega^2}}.
\end{equation}
Las asíntotas son sencillas: $0\dB$ hasta $\wn$ y $-40\dBdec$ después; la fase va de $0\grados$ a $-180\grados$ y pasa por $-90\grados$ en $\wn$. Lo que la asíntota \emph{no} ve es el pico de resonancia. En $\omega=\wn$ el término $\wn^2-\omega^2$ se anula y
\begin{equation}
  \Gma(\jw_n)=\frac{1}{j2\zeta}\ \Rightarrow\ |\Gma(\jw_n)|=\frac{1}{2\zeta},\qquad
  \angle \Gma(\jw_n)=-90\grados,
  \label{eq:rf-pico}
\end{equation}
exactamente y para cualquier $\zeta$. Con $\zeta=0.1$ la curva real está $14\dB$ por encima de la asíntota (\cref{fig:rf-complejos}).

\begin{figure}[ht]
\centering
\includegraphics[width=0.88\linewidth]{cap10/f_complejos}
\caption{Polos complejos: magnitud y fase reales para varios $\zeta$, con las asíntotas en negro. Todas las fases pasan por $-90\grados$ en $\wn$.}
\label{fig:rf-complejos}
\end{figure}

\begin{ideaclave}
$\zeta$ pequeño $\Rightarrow$ pico grande en $\wn$ $\Rightarrow$ sobrepico grande en el tiempo. En el curso \alerta{no dibujamos asintóticos de segundo orden}: se dibujan las asíntotas, se marca el punto $20\log\frac{1}{2\zeta}$ en $\wn$ para recordar que la curva real se separa, y para la curva real se usa \cmd{bode}.
\end{ideaclave}

\video{V2_resonancia_zeta}{Tres vistas del mismo $\zeta$: los polos se acercan al eje imaginario, crece el pico del Bode y crece el sobrepico.}

\subsection{¿Por qué no basta la asíntota? Un ejemplo}
\label{sec:rf-sorpresa}

\begin{ejemplo}[label=ej:rf-sorpresa]{La asíntota dice ``estable''; la realidad, no}
¿Es estable el lazo cerrado con retroalimentación unitaria si
\begin{equation*}
  \Gma(s)=\frac{1000}{(s+1)(s^2+4s+100)}=\frac{10}{(s+1)\big(\frac{s^2}{100}+0.04s+1\big)}\ ?
\end{equation*}
\solucion Aquí $\wn=10$ y $\zeta=0.2$. Las asíntotas dicen: $20\dB$ hasta $\omega=1$, $-20\dBdec$ hasta $\omega=10$ (donde llegan exactamente a $0\dB$) y $-60\dBdec$ después. La fase real en $\omega=10$ es $-84.3\grados-90\grados=-174.3\grados$. Solo con asíntotas, el lazo cerrado parece estable por poco.

La curva real cuenta otra historia (\cref{fig:rf-sorpresa}): la resonancia suma $20\log\frac{1}{2(0.2)}=+7.96\dB$ en $\wn=10$, de modo que $|\Gma(j10)|=7.92\dB$. El cruce por $0\dB$ se corre a $12.7$\,rad/s, donde la fase ya pasó $-180\grados$: el margen de fase es $-45.5\grados$ y el lazo cerrado es \alerta{inestable}. Se comprueba con Routh-Hurwitz sobre $s^3+5s^2+104s+1100$: como $5\cdot104=520<1100$, hay dos cambios de signo.
\end{ejemplo}

\begin{figure}[ht]
\centering
\includegraphics[width=0.88\linewidth]{cap10/f_sorpresa}
\caption{Asíntotas contra curva real para $1000/[(s+1)(s^2+4s+100)]$: la resonancia empuja el cruce de ganancia hasta donde la fase ya pasó $-180\grados$.}
\label{fig:rf-sorpresa}
\end{figure}

\codigorepo{cap10_respuesta_frecuencial/Clase1_RespuestaFrecuencial/C1_Ej3_SegundoOrden.m}{sorpresa}{ejemplo \ref{ej:rf-sorpresa}}

\begin{profundiza}[El factor de segundo orden en detalle]
\textbf{Pico de resonancia.} El máximo de la magnitud y su frecuencia, para $\zeta<1/\sqrt2$, son
\begin{equation}
  \Mr=\frac{1}{2\zeta\sqrt{1-\zeta^2}},\qquad \wres=\wn\sqrt{1-2\zeta^2}.
  \label{eq:rf-Mr}
\end{equation}
Para $\zeta\geq1/\sqrt2\approx0.707$ no hay pico. La \cref{tab:rf-zeta} reúne los valores útiles, junto con el sobrepico de la respuesta al escalón, $\OS=100\,e^{-\pi\zeta/\sqrt{1-\zeta^2}}$.

\begin{center}\small
\begin{tabular}{@{}cccccc@{}}
\toprule
$\zeta$ & $20\log\frac{1}{2\zeta}$ (dB) & $\Mr$ & $\Mr$ (dB) & $\wres/\wn$ & $\OS$\\
\midrule
0.05 & 20.00 & 10.01 & 20.01 & 0.997 & 85.4\\
0.1  & 13.98 & 5.03 & 14.02 & 0.990 & 72.9\\
0.2  & 7.96  & 2.55 & 8.14  & 0.959 & 52.7\\
0.3  & 4.44  & 1.75 & 4.85  & 0.906 & 37.2\\
0.5  & 0.00  & 1.15 & 1.25  & 0.707 & 16.3\\
0.707& $-3.01$ & 1 & 0 & -- & 4.3\\
1    & $-6.02$ & 1 & 0 & -- & 0\\
\bottomrule
\end{tabular}
\captionof{table}{Corrección del factor de segundo orden en $\wn$, pico de resonancia y sobrepico al escalón.}
\label{tab:rf-zeta}
\end{center}

\textbf{La pendiente de la fase.} Cerca de $\wn$ la fase cae con pendiente
$\left.\frac{\dif\phi}{\dif\log_{10}\omega}\right|_{\wn}=-\frac{180\grados}{\pi}\,\frac{\ln10}{\zeta}\approx-\frac{132\grados}{\zeta}$ por década. Con $\zeta=0.2$ son $-660\grados$ por década: la fase cae casi como un escalón de $-180\grados$. Por eso una asíntota de fase de segundo orden dependería de $\zeta$, y por eso no la usamos.

\textbf{El factor de calidad.} En física y electrónica se usa $Q=1/(2\zeta)$, que es precisamente $|\Gma(\jw_n)|$. En un RLC serie $\zeta=\frac{R}{2}\sqrt{C/L}$, y en un sistema masa-resorte-amortiguador $\zeta=b/(2\sqrt{km})$.

\textbf{Leer $\zeta$ en un Bode medido.} La fase cruza $-90\grados$ en $\wn$; el pico da $\Mr$ y de ahí $\zeta$. Si la fase vale $-90\grados$ en $\omega=10$ y el pico es de $8.14\dB$, entonces $\wn=10$, $\Mr=10^{8.14/20}=2.55$ y $\zeta=0.20$ (para $\zeta$ pequeño, $\zeta\approx\frac{1}{2\Mr}=0.196$). La planta era $100/(s^2+4s+100)$.

\textbf{Ceros complejos: el filtro \emph{notch}.} Ceros sobre el eje imaginario anulan la ganancia en una frecuencia: $N(s)=\dfrac{s^2+\omega_0^2}{s^2+\frac{\omega_0}{Q}s+\omega_0^2}$ con $\omega_0=2\pi(60)=377$\,rad/s es el filtro que quita el zumbido de $60$\,Hz de la red en señales de ECG y de audio; $Q$ controla qué tan angosta es la muesca.
\end{profundiza}

\codigorepo{cap10_respuesta_frecuencial/Clase1_RespuestaFrecuencial/C1_Ej3_SegundoOrden.m}{leerzeta}{leer $\wn$ y $\zeta$ de un Bode medido}

%% ---------------------------------------------------------------------
\section{Ejemplo resuelto: $\Gma(s)=s/(s+10)$}
\label{sec:rf-ejemplo}
%% ---------------------------------------------------------------------

Este ejemplo recoge todo lo anterior y es el que más conviene dominar. Lo resolvemos en cinco pasos que sirven para cualquier malla abierta.

\begin{ejemplo}[label=ej:rf-ss10]{Diagrama asintótico de $s/(s+10)$}
Dibuje el diagrama de Bode asintótico de magnitud y fase de $\Gma(s)=s/(s+10)$ y compárelo con el real.
\solucion
\textbf{Paso 1: forma de Bode.} Se saca el $10$ del polo:
\[
  \Gma(s)=\frac{s}{10\big(\frac{s}{10}+1\big)}
  =\underbrace{\textcolor{ctrlazul}{\frac{1}{10}}}_{\text{ganancia}}\cdot\underbrace{\textcolor{ctrlcoral}{s}}_{\text{cero en el origen}}\cdot
  \underbrace{\textcolor{ctrlvioleta}{\frac{1}{\frac{s}{10}+1}}}_{\text{polo en }-10}.
\]
\textbf{Paso 2: identificar los factores.}
\begin{itemize}
  \item \textcolor{ctrlazul}{$1/10$}: $20\log\frac{1}{10}=-20\dB$ y $0\grados$ (recta horizontal).
  \item \textcolor{ctrlcoral}{$s$}: $+20\dBdec$ pasando por $0\dB$ en $\omega=1$, y $+90\grados$.
  \item \textcolor{ctrlvioleta}{$1/(s/10+1)$}: $0\dB$ hasta $\omega=10$ y $-20\dBdec$ después; fase de $0\grados$ a $-90\grados$ entre $\omega=1$ y $\omega=100$.
\end{itemize}
\textbf{Paso 3: dibujar cada factor por separado} (líneas delgadas de la \cref{fig:rf-ss10}).

\textbf{Paso 4: sumar.} En magnitud,
\[
  \omega<10:\ -20+20\log\omega,\qquad
  \omega>10:\ -20+20\log\omega-20\log\frac{\omega}{10}=0\dB .
\]
Desde $\omega=10$ las pendientes $+20$ y $-20$ se cancelan y la curva queda plana en $0\dB$. En $\omega=1$ la asíntota vale $-20\dB$ y en $\omega=0.1$, $-40\dB$. En fase, $\phi(\omega)=90\grados-\tan^{-1}(\omega/10)$: $90\grados$ para $\omega\le1$, $45\grados$ en $\omega=10$ y $0\grados$ para $\omega\ge100$.

\textbf{Paso 5: comparar con la curva real.} En la esquina,
\[
  \Big|\frac{j10}{10+j10}\Big|=\frac{1}{\sqrt2}\to-3\dB,\qquad\phi=90\grados-45\grados=45\grados.
\]
Otros valores reales: en $\omega=1$, $-20.04\dB$ y $84.3\grados$; en $\omega=20$, $-0.97\dB$ y $26.6\grados$; en $\omega=0.1$, $-40.0\dB$ y $89.4\grados$.
\end{ejemplo}

\begin{figure}[ht]
\centering
\includegraphics[width=\linewidth]{cap10/f_ej_s_s10}
\caption{Los tres factores de $s/(s+10)$ (delgados), su suma asintótica (negro) y la curva real (verde). Los puntos marcan la esquina: $-3\dB$ y $45\grados$.}
\label{fig:rf-ss10}
\end{figure}

\codigorepo{cap10_respuesta_frecuencial/Clase1_RespuestaFrecuencial/C1_Ej2_Asintotico_s_s10.m}{factores}{ejemplo \ref{ej:rf-ss10}}

\video{V10_asintotico_s_s10}{El diagrama se arma término a término con el mismo código de colores; al final, la curva real y los $-3\dB$ de la esquina.}

\paragraph{Interpretación: un filtro pasa-altos}
$s/(s+10)$ deja pasar con ganancia $1$ las frecuencias por encima de $10$\,rad/s y atenúa las de abajo ($-20\dB$ por cada década hacia la izquierda). En bajas frecuencias la salida, además de pequeña, va adelantada casi $90\grados$: el cero en el origen es un \emph{derivador}. Por ejemplo, para $x(t)=\sin(0.1t)$ la salida es $0.01\sin(0.1t+89.4\grados)$. Con una entrada compuesta $x(t)=\sin t+0.3\sin 20t$, la componente lenta sale multiplicada por $0.0995$ ($-20\dB$) y la rápida por $0.894$: en la salida domina la componente de $20$\,rad/s (\cref{fig:rf-pasaaltos}).

\begin{figure}[ht]
\centering
\includegraphics[width=\linewidth]{cap10/f_pasaaltos}
\caption{$s/(s+10)$ como filtro pasa-altos: de $\sin t+0.3\sin20t$ sobrevive casi solo la componente rápida.}
\label{fig:rf-pasaaltos}
\end{figure}

\video{V11_pasaaltos_s_s10}{Barrido de $0.1$ a $1000$\,rad/s mostrando entrada y salida junto al punto del Bode; al final, la entrada compuesta.}

\subsection{La receta y la hoja de trucos}

\begin{receta}[Diagrama de Bode asintótico, paso a paso]
\begin{enumerate}
  \item Llevar $\Gma(s)$ a la forma de Bode \eqref{eq:rf-formabode}: cada factor como $(s/\omega_i+1)$, y la ganancia $\Kb$ sola.
  \item Ordenar las esquinas de menor a mayor.
  \item Magnitud: arrancar con la recta de baja frecuencia, $20\log \Kb-20N\log\omega$, y recorrer de izquierda a derecha sumando $+20\dBdec$ por cada cero y $-20$ por cada polo ($-40$ por cada par complejo).
  \item Fase: factor por factor ($\pm45\grados$/dec entre $\omega_i/10$ y $10\,\omega_i$ para los reales; $-90N\grados$ de los integradores) y sumar.
  \item Verificar con \cmd{bode}. Con polos complejos y $\zeta$ pequeño, la curva real puede estar muy por encima de la asíntota.
\end{enumerate}
\end{receta}

\begin{anchocompleto}
\centering\small
\begin{tabular}{@{}llll@{}}
\toprule
Factor & Magnitud & Fase & Cuidado\\
\midrule
$\Kb$ & $20\log|\Kb|$, plana & $0\grados$ ($\pm180\grados$ si $\Kb<0$) & calcular $\Kb$ en forma de Bode\\
$1/s^N$ & $-20N\dBdec$, $0\dB$ en $\omega=1$ & $-90N\grados$ & define el tipo del sistema\\
$(s/\omega_1+1)^{\pm1}$ & $\pm20\dBdec$ desde $\omega_1$ & $0\to\pm90\grados$ ($\pm45\grados$ en $\omega_1$) & real: $\pm3\dB$ en $\omega_1$\\
polos complejos & $-40\dBdec$ desde $\wn$; pico $\approx\frac{1}{2\zeta}$ & $0\to-180\grados$ ($-90\grados$ en $\wn$) & curva real: \cmd{bode}\\
\bottomrule
\end{tabular}
\captionof{table}{Hoja de trucos de Bode.}
\label{tab:rf-trucos}
\end{anchocompleto}

\begin{pausa}
Sin dibujar: ¿cuánto vale la asíntota de magnitud de $\Gma(s)=\dfrac{100(s+1)}{s(s+10)}$ en $\omega=1$? ¿Y la curva real?
\respuesta{Forma de Bode: $\Gma=\dfrac{10(s+1)}{s(s/10+1)}$, $\Kb=10$. En $\omega=1$ la asíntota vale $20\log10-20\log1=20\dB$. La curva real está $3\dB$ por encima por el cero en $1$ (y un poco menos por el polo en $10$): $22.97\dB$.}
\end{pausa}

%% ---------------------------------------------------------------------
\section{Leer un Bode al revés: de la curva a la función de transferencia}
\label{sec:rf-identificacion}
%% ---------------------------------------------------------------------

La receta de la sección anterior también funciona en reversa. Si alguien mide el diagrama de magnitud de un sistema (con la máquina de Bode, por ejemplo), podemos reconstruir $\Gma(s)$ leyendo las pendientes y las esquinas de izquierda a derecha; la ganancia sale al final de un valor conocido.

\begin{ejemplo}[label=ej:rf-ident]{¿Qué función de transferencia es esta?}
El sistema es de fase mínima y de ganancia positiva. Su diagrama asintótico de magnitud (negro en la \cref{fig:rf-ident}) baja con $-20\dBdec$ en bajas frecuencias, se vuelve plano en $36\dB$ desde $\omega=1$ y cae con $-40\dBdec$ desde $\omega=10$. Determine $\Gma(s)$ y dibuje la fase.
\solucion Leemos las pendientes de izquierda a derecha:
\begin{itemize}
  \item $-20\dBdec$ desde el principio: un \textbf{polo en el origen}.
  \item En $\omega=1$ la pendiente pasa a $0$ ($-20+20$): un \textbf{cero en $-1$}.
  \item En $\omega=10$ la pendiente pasa a $-40$ ($0-40$): un \textbf{polo doble en $-10$}.
\end{itemize}
En forma de Bode, $\Gma(s)=\Kb\,\dfrac{s+1}{s\,(s/10+1)^2}$. Entre $\omega=1$ y $\omega=10$ la asíntota vale $20\log\Kb$ (el integrador y el cero se cancelan), de modo que
\[
  20\log \Kb=36\dB\ \Rightarrow\ \Kb=10^{36/20}=63.1,\qquad
  \Gma(s)=\frac{63.1\,(s+1)}{s\,(s/10+1)^2}=\frac{6309.6\,(s+1)}{s\,(s+10)^2}.
\]
La fase asintótica arranca en $-90\grados$ (integrador), sube hasta $-45\grados$ en $\omega=1$ por el cero, vuelve a $-90\grados$ en $\omega=10$ y termina en $-180\grados$ desde $\omega=100$.
\end{ejemplo}

\begin{figure}[ht]
\centering
\includegraphics[width=0.88\linewidth]{cap10/f_identificacion}
\caption{Identificación: el asintótico de magnitud es el dato; la fase sale ``gratis'' porque el sistema es de fase mínima.}
\label{fig:rf-ident}
\end{figure}

\codigorepo{cap10_respuesta_frecuencial/Clase2_Margenes/C2_Ej1_Identificacion.m}{identificacion}{ejemplo \ref{ej:rf-ident}}

\subsection{La fase sale gratis\ldots{} si el sistema es de fase mínima}
\label{sec:rf-faseminima}

En el ejemplo usamos una hipótesis fuerte: \emph{fase mínima}, es decir, sin ceros ni polos en el semiplano derecho. Bode demostró que, bajo esa hipótesis, la magnitud determina la fase~\citep{Bod40}. Sin ella, la magnitud no basta: $(1+s)/(1+2s)$ y $(1-s)/(1+2s)$ tienen \emph{exactamente} el mismo diagrama de magnitud, porque $|1+\jw|=|1-\jw|$, pero sus fases son muy distintas (\cref{fig:rf-nominima}). El cero en $+1$ agrega atraso de fase sin cambiar la magnitud.

\begin{figure}[ht]
\centering
\includegraphics[width=0.82\linewidth]{cap10/f_fase_no_minima}
\caption{Misma magnitud, distinta fase: el cero en el semiplano derecho (coral) atrasa la fase hasta $-180\grados$.}
\label{fig:rf-nominima}
\end{figure}

\begin{cuidado}
Antes de reconstruir la fase a partir de la magnitud, hay que saber (por la física del problema o por una medición de fase) que el sistema es de fase mínima. Los sistemas de fase no mínima aparecen en la práctica: la respuesta inversa de una caldera, un avión que cabecea, o el control de nivel con ciertas configuraciones. Para ellos el Bode de magnitud solo no basta y se requiere un estudio más detallado (criterio de Nyquist).
\end{cuidado}

\begin{pausa}
Una malla abierta de fase mínima tiene pendiente de $-40\dBdec$ en bajas frecuencias. ¿Cuántos integradores tiene y cuánto vale su fase a muy baja frecuencia?
\respuesta{Dos integradores ($-20\dBdec$ cada uno) y fase $-180\grados$. Es un sistema tipo 2 (\cref{sec:rf-tipo}).}
\end{pausa}

%% ---------------------------------------------------------------------
\section{Estabilidad desde la malla abierta: márgenes de ganancia y de fase}
\label{sec:rf-estabilidad}
%% ---------------------------------------------------------------------

\subsection{La frontera: $|\Gma|=1$ y fase $-180\grados$}

La ecuación característica del lazo cerrado es $1+\Gma(s)=0$, es decir, $\Gma(s)=-1$. Evaluada sobre el eje imaginario, la frontera entre estabilidad e inestabilidad es una frecuencia en la que, a la vez,
\begin{equation}
  |\Gma(\jw)|=1\ \ (0\dB)\qquad\text{y}\qquad \angle\Gma(\jw)=-180\grados .
  \label{eq:rf-frontera}
\end{equation}
La intuición es la de la \cref{sec:rf-lenguaje}: una senoide que da la vuelta al lazo regresa del mismo tamaño ($|\Gma|=1$) e invertida ($-180\grados$); el signo menos del comparador la vuelve a invertir, se suma a sí misma y el lazo oscila solo. Si en la frecuencia donde la fase vale $-180\grados$ la senoide regresa \emph{más pequeña} ($|\Gma|<1$), cada vuelta la achica y la oscilación muere.

\begin{ideaclave}
Si la malla abierta es estable y de fase mínima, el lazo cerrado es asintóticamente estable cuando la magnitud de $\Gma(\jw)$ es menor que $0\dB$ en la frecuencia en que su fase vale $-180\grados$.
\end{ideaclave}

\subsection{Márgenes de ganancia y de fase}

Saber que el lazo es estable no basta: queremos saber \emph{qué tan lejos} está de la frontera. Los dos márgenes miden esa distancia, cada uno en una dirección (\cref{fig:rf-margenes}).

\begin{definicion}[label=def:rf-gm]{Margen de ganancia}
El \clave{margen de ganancia} ($\GM$) es cuánto se puede aumentar la ganancia de la malla abierta, en decibeles, en la frecuencia donde la fase vale $-180\grados$ ($\wgm$, cruce de fase), antes de que el lazo cerrado se vuelva inestable:
\begin{equation}
  \GM=-20\log|\Gma(\jw_{GM})|\quad\text{dB},\qquad \angle\Gma(\jw_{GM})=-180\grados.
\end{equation}
\end{definicion}

\begin{definicion}[label=def:rf-pm]{Margen de fase}
El \clave{margen de fase} ($\PM$) es cuánto atraso de fase adicional se tolera en la frecuencia donde $|\Gma|=1$ ($\wpm$, cruce de ganancia) antes de llegar a $-180\grados$:
\begin{equation}
  \PM=180\grados+\angle\Gma(\jw_{PM}),\qquad |\Gma(\jw_{PM})|=1\ (0\dB).
\end{equation}
\end{definicion}

\begin{figure}[ht]
\centering
\includegraphics[width=0.88\linewidth]{cap10/f_margenes_def}
\caption{Lectura de los márgenes. El $\GM$ se mide en magnitud, en la frecuencia donde la fase cruza $-180\grados$; el $\PM$ se mide en fase, en la frecuencia donde la magnitud cruza $0\dB$.}
\label{fig:rf-margenes}
\end{figure}

\begin{cuidado}[Los nombres de MATLAB invitan a confundirse]
En \cmd{[Gm,Pm,Wcg,Wcp] = margin(G)}, \cmd{Wcg} (\emph{gain crossover} en otros textos) es en realidad la frecuencia donde se mide el $\GM$, es decir, nuestro $\wgm$, y \cmd{Wcp} es la frecuencia del $\PM$, nuestro $\wpm$. Además \cmd{Gm} sale en \emph{veces}: el $\GM$ en dB es \cmd{20*log10(Gm)}. La función del curso \cmd{margenesTabla(G)} imprime todo ya convertido.
\end{cuidado}

\begin{ejemplo}[label=ej:rf-margenes]{Márgenes asintóticos contra exactos}
Determine el $\GM$ y el $\PM$ de $\Gma(s)=\dfrac{10\,000}{(s+1)(s+10)(s+100)}$ con el diagrama asintótico y compare con \cmd{margin}.
\solucion \textbf{Forma de Bode:} $\Gma(s)=\dfrac{10}{(s+1)(s/10+1)(s/100+1)}$, es decir, $20\dB$ en bajas frecuencias y tres esquinas.

\textbf{Cruce por $0\dB$ (PM).} La asíntota baja $-20\dBdec$ desde $\omega=1$ y llega a $0\dB$ justo en $\omega=10$. La fase asintótica allí: el polo en $1$ ya aportó $-90\grados$, el de $10$ está en su esquina ($-45\grados$) y el de $100$ todavía no empieza. Total $-135\grados$: $\PM\approx45\grados$.

\textbf{Cruce por $-180\grados$ (GM).} Entre $\omega=10$ y $\omega=100$ la fase asintótica cae $-45\grados$ por década por el polo en $10$ y otros $-45\grados$ por década por el de $100$, desde $-135\grados$: llega a $-180\grados$ media década después de $10$, en $\wgm=10^{1.5}=31.6$\,rad/s. La magnitud asintótica allí es $0-40\times0.5=-20\dB$: $\GM\approx20\dB$.

\textbf{Comparación.}
\begin{center}\small
\begin{tabular}{@{}lcc@{}}
\toprule
 & asintótico & exacto (\cmd{margin})\\
\midrule
$\PM$ & $45\grados$ ($\omega=10$) & $54.9\grados$ ($\omega=7.80$)\\
$\GM$ & $20\dB$ ($\omega=31.6$) & $21.7\dB$ ($\omega=33.3$)\\
\bottomrule
\end{tabular}
\end{center}
El asintótico es conservador aquí: sirve para estimar en el tablero; el número final siempre con \cmd{margin}.
\end{ejemplo}

\begin{figure}[ht]
\centering
\includegraphics[width=0.9\linewidth]{cap10/f_margenes_asint}
\caption{Ejemplo~\ref{ej:rf-margenes}: los cuatro factores, su suma asintótica y la curva real, con los márgenes asintóticos marcados.}
\label{fig:rf-margenes-asint}
\end{figure}

\codigorepo{cap10_respuesta_frecuencial/Clase2_Margenes/C2_Ej2_Margenes.m}{margenes}{ejemplo \ref{ej:rf-margenes}}
\video{V3_asintotico_margenes}{El Bode se arma término a término; al final, PM y GM sobre la asintótica.}

%% ---------------------------------------------------------------------
\subsection{¿Qué significan los márgenes en la práctica?}
\label{sec:rf-significado}

\paragraph{El GM: el micrófono y el parlante}
\notamargen{El ``aullido'' de un sistema de sonido es un lazo cerrado que perdió su margen de ganancia.}
Un micrófono frente a un parlante forma un lazo: lo que sale del parlante vuelve a entrar al micrófono. Subir el volumen es subir la ganancia de la malla abierta \emph{sin cambiar la fase}. Mientras haya margen de ganancia no pasa nada; cuando el volumen sube más que el $\GM$, en la frecuencia donde la fase acústica vale $-180\grados$ aparece el aullido: una oscilación sostenida. El $\GM$ es, literalmente, cuánto se puede subir el volumen.

\paragraph{El PM: un presupuesto de fase para lo que no modelamos}
Ningún modelo es perfecto. Sensores, actuadores y filtros que dejamos por fuera agregan atraso de fase que la malla abierta modelada no ve. El $\PM$ es el \emph{presupuesto} con el que pagamos ese atraso. Para verlo, pongamos un sensor de primer orden $\Hs(s)=a/(s+a)$ en la retroalimentación del ejemplo~\ref{ej:rf-margenes} (\cref{fig:rf-lazosensor}). En el cruce, el sensor resta aproximadamente $\tan^{-1}(\wpm/a)$, de modo que
\begin{equation}
  \PM_{\text{nuevo}}\approx\PM-\tan^{-1}\frac{\wpm}{a}\qquad(\text{buena estimación si } a\gg\wpm).
\end{equation}

\begin{figure}[ht]
\centering
\includegraphics[width=0.55\linewidth]{cap10/lazo_sensor}
\caption{El lazo del ejemplo~\ref{ej:rf-margenes} con un sensor de primer orden. La malla abierta es ahora $\Gma=G\,\Hs$.}
\label{fig:rf-lazosensor}
\end{figure}

\begin{center}\small
\begin{tabular}{@{}ccccc@{}}
\toprule
$a$ (rad/s) & $\PM$ estimado & $\PM$ (\cmd{margin}) & $\GM$ & sobrepico\\
\midrule
sin sensor & -- & $54.9\grados$ & $21.7\dB$ & $18\%$\\
$50$ & $46.0\grados$ & $46.5\grados$ & $12.7\dB$ & $28\%$\\
$20$ & $33.6\grados$ & $36.4\grados$ & $8.6\dB$ & $44\%$\\
$10$ & $17.0\grados$ & $26.3\grados$ & $6.1\dB$ & $68\%$\\
$5$  & $-2.4\grados$ & $18.7\grados$ & $4.7\dB$ & $104\%$\\
\bottomrule
\end{tabular}
\captionof{table}{El sensor ``gasta'' fase. Con $a=50$, casi diez veces más rápido que el cruce, la estimación es buena; con sensores lentos la estimación es pesimista (el sensor también baja la magnitud y corre el cruce a la izquierda), pero la tendencia es la misma.}
\label{tab:rf-sensor}
\end{center}

\begin{figure}[ht]
\centering
\includegraphics[width=\linewidth]{cap10/f_sensor}
\caption{Un sensor más lento gasta más fase en el cruce (derecha) y el lazo cerrado se vuelve más oscilatorio (izquierda), aunque la planta sea la misma.}
\label{fig:rf-sensor}
\end{figure}

\codigorepo{cap10_respuesta_frecuencial/Clase2_Margenes/C2_Video_SensorMargenFase.m}{sensor}{tabla \ref{tab:rf-sensor}}
\video{V12_sensor_margen_fase}{El polo del sensor baja de $300$ a $5$\,rad/s: se ve cómo se consume el PM y crece el sobrepico.}

\begin{ideaclave}
El $\GM$ dice cuánta ganancia de más aguanta el lazo; el $\PM$, cuánto atraso de fase no modelado aguanta. Un buen diseño necesita los dos.
\end{ideaclave}

\begin{profundiza}[Los márgenes como distancia al punto $-1$]
Ambos márgenes miden qué tan cerca pasa $\Gma(\jw)$ del punto crítico $-1$, cada uno en una dirección: el $\GM$ sobre el eje real negativo y el $\PM$ sobre el círculo unitario. Una medida que los resume es la distancia mínima $\min_\omega|1+\Gma(\jw)|=1/M_s$, donde $M_s=\max_\omega|\Sens(\jw)|$ es el pico de la sensitividad $\Sens=1/(1+\Gma)$. Un $M_s$ entre $1.3$ y $2$ garantiza a la vez $\GM\geq20\log\frac{M_s}{M_s-1}$ y $\PM\geq2\sin^{-1}\frac{1}{2M_s}$; por ejemplo, $M_s=2$ asegura $\GM\geq6\dB$ y $\PM\geq29\grados$~\citep{AsMu10}.
\end{profundiza}

%% ---------------------------------------------------------------------
\subsection{¿Cuánto margen es suficiente?}

\begin{itemize}
  \item \textbf{Fase mínima:} $\GM$ y $\PM$ positivos $\Rightarrow$ estable. Márgenes negativos $\Rightarrow$ inestable.
  \item \textbf{Regla práctica:} $30\grados<\PM<60\grados$ y $\GM>6\dB$ dan una estabilidad que tolera variaciones razonables de la ganancia y de las constantes de tiempo.
  \item \textbf{Fase no mínima} (ceros o polos en el semiplano derecho): el Bode solo no basta; se requiere un estudio más detallado (criterio de Nyquist).
\end{itemize}

\begin{pausa}
Tres lazos de fase mínima: (A) $\PM=40\grados$, $\GM=12\dB$; (B) $\PM=-15\grados$, $\GM=-4\dB$; (C) $\PM=70\grados$, $\GM=2\dB$. ¿Cuáles son estables? ¿Cuál le preocupa más?
\respuesta{(A) estable, con márgenes cómodos. (B) inestable: hay que bajar la ganancia al menos $4\dB$. (C) estable pero frágil: $2\dB$ es un factor $10^{2/20}=1.26$, así que un $26\%$ más de ganancia lo inestabiliza. Un $\PM$ grande no basta: hay que mirar los dos márgenes.}
\end{pausa}

%% ---------------------------------------------------------------------
\subsection{¿Cuánta ganancia aguanta? Routh-Hurwitz contra Bode}
\label{sec:rf-rangokc}

\begin{figure}[ht]
\centering
\includegraphics[width=0.62\linewidth]{cap10/lazo_basico}
\caption{Lazo de control con retroalimentación. En los ejemplos de diseño, $\Hs(s)=1$.}
\label{fig:rf-lazo}
\end{figure}

\begin{ejemplo}[label=ej:rf-rangokc]{Rango de $\Kc$ por dos caminos}
En el lazo de la \cref{fig:rf-lazo}, con $\Hs=1$, $\Gc(s)=\Kc$ y $\Gp(s)=\dfrac{1}{(s+1)(s+10)(s+100)}$, determine el rango de $\Kc$ que garantiza estabilidad asintótica con Routh-Hurwitz y con el diagrama de Bode.
\solucion \textbf{Routh-Hurwitz.} El polinomio característico es $s^3+111s^2+1110s+(1000+\Kc)$. La fila $s^1$ exige $111\cdot1110-(1000+\Kc)>0$ y la fila $s^0$, $1000+\Kc>0$:
\[
  -1000<\Kc<122\,210 .
\]
En $\Kc=\Kcr=122\,210$ el polinomio auxiliar $111s^2+123\,210=0$ da $s=\pm j\sqrt{1110}$: el lazo oscila con $\wcr=33.32$\,rad/s y periodo $\Pcr=2\pi/\wcr=0.189$\,s.

\textbf{Bode.} Para dibujar hay que fijar algún $\Kc$. Conviene el que deja la ganancia DC en $0\dB$, $\Kc=1000$, con lo que $\Gma(s)=\dfrac{1}{(s+1)(s/10+1)(s/100+1)}$. El cruce de fase asintótico está otra vez en $\omega=31.6$, donde la magnitud asintótica vale $-40\dB$: hay $\GM\approx+40\dB$, es decir, la ganancia se puede multiplicar por $100$. Así $\Kcr\approx1000\times100=100\,000$ y el rango positivo es $0<\Kc<100\,000$. Con \cmd{margin} el $\GM$ exacto es $41.74\dB$ ($\times122.2$), que reproduce $\Kcr=122\,210$.

La lectura asintótica se equivoca hacia el lado seguro: subestima la ganancia crítica.
\end{ejemplo}

\begin{figure}[ht]
\centering
\includegraphics[width=0.9\linewidth]{cap10/f_rango_kc}
\caption{Respuesta al escalón del lazo cerrado al subir $\Kc$: en $\Kc=\Kcr$ la oscilación ya no se amortigua.}
\label{fig:rf-rangokc}
\end{figure}

\codigorepo{cap10_respuesta_frecuencial/Clase2_Margenes/C2_Ej3_RangoKc.m}{routh}{Routh-Hurwitz sin Symbolic Toolbox}
\video{V4_barrido_Kc}{Subir $\Kc$ solo sube la magnitud: la barra roja del GM se acorta hasta cero en $\Kc=122\,210$.}

\begin{modelo}[LazoKc.slx]
\cmd{crearModelo\_LazoKc.m} arma el lazo escalón $\to$ suma $\to$ $\Kc$ $\to$ $\Gp(s)$ con retroalimentación unitaria. Pruebe $\Kc=1000$, $50\,000$, $122\,210$ (marginal) y $150\,000$ (inestable). Para otra planta cambie el parámetro \emph{Poles} del bloque $\Gp$.
\end{modelo}

\paragraph{Conexión con lo que ya vimos: Ziegler-Nichols}
\notamargen{Ziegler y Nichols publicaron sus reglas en 1942~\citep{ZiNi42}.}
El $\GM$ entrega la ganancia crítica y $\wgm$ la frecuencia crítica: son justo los datos del método de oscilación sostenida de Ziegler y Nichols (capítulo de sintonización). Con nuestro ejemplo, $\Kcr=122\,210$ y $\Pcr=0.189$\,s, y la tabla de ZN para un PID da
\[
  \Kc=0.6\,\Kcr=73\,326,\qquad \Ti=\Pcr/2=0.0943\ \text{s},\qquad \Td=\Pcr/8=0.0236\ \text{s}.
\]
La ventaja del Bode es que \emph{no hay que llevar la planta real al borde de la inestabilidad} para medir $\Kcr$ y $\Pcr$: se leen en el diagrama. En la \cref{sec:rf-pid} veremos qué margen de fase deja esa sintonización.

%% ---------------------------------------------------------------------
\section{Del Bode de malla abierta al transitorio del lazo cerrado}
\label{sec:rf-lazocerrado}
%% ---------------------------------------------------------------------

Los márgenes dicen \emph{si} el lazo cerrado es estable. Ahora queremos saber \emph{cómo} se comporta: cuánto sobrepasa, qué tan rápido responde. Lo haremos con el caso de referencia que conocemos desde el capítulo de segundo orden y luego lo usaremos como aproximación de polos dominantes.

\subsection{El caso de referencia}
Con $\Gc=1$ y $\Gp(s)=\dfrac{\wn^2}{s(s+2\zeta\wn)}$, la malla abierta es tipo 1 y el lazo cerrado es el segundo orden estándar,
\begin{equation}
  \Tlc(s)=\frac{Y(s)}{R(s)}=\frac{\wn^2}{s^2+2\zeta\wn s+\wn^2}.
  \label{eq:rf-T2}
\end{equation}

\subsection{El lazo cerrado visto en frecuencia: pico de resonancia y ancho de banda}

El diagrama de magnitud del \emph{lazo cerrado} $|\Tlc(\jw)|$ tiene dos rasgos que resumen su comportamiento (\cref{fig:rf-Tzeta}).

\paragraph{Pico de resonancia}
Maximizar $|\Tlc(\jw)|=\wn^2/\sqrt{(\wn^2-\omega^2)^2+4\zeta^2\wn^2\omega^2}$ equivale a minimizar el radicando como función de $\omega^2$; derivando e igualando a cero se obtienen el pico $\Mr$ y su frecuencia $\wres$ de la ecuación~\eqref{eq:rf-Mr}:
\[
  \Mr=\frac{1}{2\zeta\sqrt{1-\zeta^2}},\qquad \wres=\wn\sqrt{1-2\zeta^2},\qquad 0<\zeta<\tfrac{1}{\sqrt2}.
\]
Un $\Mr$ grande indica polos dominantes poco amortiguados y una respuesta oscilatoria; para $\zeta>0.707$ no hay pico ($\Mr=1$).

\begin{cuidado}
No confundir el pico de resonancia $\Mr$ (frecuencia, lazo cerrado) con el sobrepico $\OS$ (tiempo, respuesta al escalón). Están relacionados porque ambos dependen solo de $\zeta$, pero son números distintos: con $\zeta=0.5$, $\Mr=1.15$ ($1.25\dB$) y $\OS=16.3\%$. En versiones anteriores de estas notas el pico de resonancia se llamaba $M_P$; para evitar la confusión con el sobrepico, ahora es $\Mr$.
\end{cuidado}

\paragraph{Ancho de banda}
El \clave{ancho de banda} $\wbw$ es la frecuencia donde $|\Tlc|$ cae $3\dB$ respecto a su valor de baja frecuencia, $|\Tlc(\jw_{BW})|=1/\sqrt2$. Mide la habilidad del lazo para reproducir fielmente la referencia:
\begin{equation}
  \wbw=\wn\sqrt{(1-2\zeta^2)+\sqrt{4\zeta^4-4\zeta^2+2}} .
  \label{eq:rf-BW}
\end{equation}
Como $\ts=4/(\zeta\wn)$ y $\tp=\pi/(\wn\sqrt{1-\zeta^2})$, para un $\zeta$ dado un ancho de banda mayor significa un sistema más rápido. Más ancho de banda también deja pasar más ruido de medición: nada es gratis.

\begin{figure}[ht]
\centering
\includegraphics[width=0.85\linewidth]{cap10/f_lazo_cerrado_zeta}
\caption{Magnitud del lazo cerrado $|\Tlc(\jw)|$ para varios $\zeta$: el pico $\Mr$ crece al bajar $\zeta$; la línea punteada marca los $-3\dB$ que definen $\wbw$.}
\label{fig:rf-Tzeta}
\end{figure}

\subsection{Del margen de fase a $\zeta$, paso a paso}

Ahora la idea central de la sección: obtener $\zeta$ \emph{sin} cerrar el lazo, mirando solo la malla abierta $\Gma(s)=\wn^2/[s(s+2\zeta\wn)]$.

\textbf{Paso 1: frecuencia de cruce.} Imponiendo $|\Gma(\jw_1)|=\wn^2/|{-\omega_1^2}+j2\zeta\wn\omega_1|=1$ y despejando,
\begin{equation}
  \omega_1=\wn\sqrt{-2\zeta^2+\sqrt{1+4\zeta^4}} .
  \label{eq:rf-omega1}
\end{equation}
\textbf{Paso 2: fase en el cruce y margen de fase.} $\angle\Gma(\jw_1)=-90\grados-\tan^{-1}\dfrac{\omega_1}{2\zeta\wn}$, de donde
\begin{equation}
  \PM=180\grados+\angle\Gma(\jw_1)=\tan^{-1}\frac{2\zeta}{\sqrt{-2\zeta^2+\sqrt{1+4\zeta^4}}}.
  \label{eq:rf-PMzeta}
\end{equation}
El margen de fase depende \emph{solo} de $\zeta$. La \cref{fig:rf-pmzeta} es la curva que usaremos todo el capítulo; para $\zeta<0.6$ vale la aproximación $\PM\approx100\,\zeta$.

\begin{figure}[ht]
\centering
\includegraphics[width=0.85\linewidth]{cap10/f_pm_zeta}
\caption{Margen de fase y sobrepico en función de $\zeta$.}
\label{fig:rf-pmzeta}
\end{figure}

\begin{center}\small
\begin{tabular}{@{}cccccc@{}}
\toprule
$\zeta$ & 0.25 & 0.40 & 0.50 & 0.60 & 0.707\\
\midrule
$\PM$ & $28.0\grados$ & $43.1\grados$ & $51.8\grados$ & $59.2\grados$ & $65.5\grados$\\
$\OS$ & $44.4\%$ & $25.4\%$ & $16.3\%$ & $9.5\%$ & $4.3\%$\\
\bottomrule
\end{tabular}
\captionof{table}{Valores de referencia de la relación $\PM\leftrightarrow\zeta\leftrightarrow\OS$. Un $\PM\geq65.5\grados$ asegura que el lazo cerrado no tiene pico de resonancia.}
\label{tab:rf-pmzeta}
\end{center}

\video{V5_PM_zeta}{Al subir $K$ en $K/[s(s+10)]$ el cruce se va a la derecha, el PM baja y el punto rojo recorre la curva $\PM(\zeta)$ mientras el escalón se vuelve más oscilatorio.}

\begin{receta}[Del Bode de malla abierta al transitorio]
\begin{enumerate}
  \item \textbf{Leer el PM} y su frecuencia $\wpm=\omega_1$ en la malla abierta.
  \item \textbf{Obtener $\zeta$} de \eqref{eq:rf-PMzeta} (función \cmd{zetaDePM}, o $\zeta\approx\PM/100$).
  \item \textbf{Obtener $\wn$} de \eqref{eq:rf-omega1}.
  \item \textbf{Transitorio:} $\OS=100\,e^{-\pi\zeta/\sqrt{1-\zeta^2}}$, $\tp=\pi/(\wn\sqrt{1-\zeta^2})$, $\ts=4/(\zeta\wn)$; o $\wbw$ con \eqref{eq:rf-BW}.
\end{enumerate}
Es exacto para el segundo orden de referencia; para órdenes mayores es una aproximación de polos dominantes. Siempre verificar con \cmd{step}.
\end{receta}

\begin{ejemplo}[label=ej:rf-pmzeta]{$\OS$, $\tp$ y $\ts$ sin simular}
Para $\Gma(s)=\dfrac{100}{s(s+10)}$ con retroalimentación unitaria, determine $\OS$, $\tp$ y $\ts$ del lazo cerrado a partir de la malla abierta.
\solucion Con \cmd{margin}: $\PM=51.8\grados$ en $\wpm=7.86$\,rad/s. De \eqref{eq:rf-PMzeta}, $\zeta=0.50$; de \eqref{eq:rf-omega1}, $\wn=7.86/0.786=10$\,rad/s. Entonces
\[
  \OS=16.3\%,\qquad \tp=\frac{\pi}{10\sqrt{1-0.25}}=0.363\ \text{s},\qquad \ts=\frac{4}{0.5\cdot10}=0.8\ \text{s},
\]
y además $\Mr=1.25\dB$ en $\wres=7.07$\,rad/s y $\wbw=12.7$\,rad/s. Como este sistema \emph{es} el segundo orden de referencia ($\wn^2=100$, $2\zeta\wn=10$), los valores son exactos (\cref{fig:rf-ejpmzeta}).

\textbf{¿Y si solo tengo el asintótico?} La asíntota $10/[s(s/10+1)]$ cruza $0\dB$ en $\omega=10$, donde la fase asintótica es $-135\grados$: $\PM\approx45\grados$, $\zeta\approx0.42$ y $\OS\approx23\%$. Son unos siete puntos de error en el sobrepico. La razón: cerca de la esquina $\omega=10$ la magnitud real está $3\dB$ por debajo de la asíntota, así que el cruce real ocurre antes ($7.86$) y con más fase.
\end{ejemplo}

\begin{figure}[ht]
\centering
\includegraphics[width=\linewidth]{cap10/f_ej_pm_zeta}
\caption{Ejemplo~\ref{ej:rf-pmzeta}: el PM de la malla abierta predice el sobrepico del lazo cerrado.}
\label{fig:rf-ejpmzeta}
\end{figure}

\codigorepo{cap10_respuesta_frecuencial/Clase3_LazoCerrado/C3_Ej1_PM_zeta.m}{receta}{ejemplo \ref{ej:rf-pmzeta}}

\begin{pausa}
Repita el ejemplo para $\Gma(s)=100/[s(s+5)]$. ¿Qué cambió y por qué?
\respuesta{Ahora $\wn=10$ y $2\zeta\wn=5$, así que $\zeta=0.25$: $\PM=28.0\grados$, $\OS=44.4\%$, $\ts=4/(0.25\cdot10)=1.6$\,s y $\tp=0.32$\,s. El polo de la planta está más cerca del origen: en el cruce aporta más atraso de fase, el PM baja y el lazo cerrado oscila más.}
\end{pausa}

\begin{ideaclave}
\begin{itemize}
  \item El $\OS$ baja al aumentar el $\PM$; la velocidad sube al aumentar $\wbw$ (en la práctica, al correr el cruce $\wpm$ a la derecha).
  \item En diseño es más común especificar $\PM$ y $\GM$ que $\Mr$.
  \item El error en estado estacionario mejora al subir la magnitud en \emph{bajas} frecuencias; el transitorio depende de lo que pasa \emph{en el cruce}. Esa separación es la que explotan los compensadores (\cref{sec:rf-compensadores}).
\end{itemize}
\end{ideaclave}

%% ---------------------------------------------------------------------
\section{Tipo y error: todo está en las bajas frecuencias}
\label{sec:rf-tipo}
%% ---------------------------------------------------------------------

Con retroalimentación unitaria, el error de seguimiento $E(s)=R(s)-Y(s)$ en estado estacionario depende del número de integradores de la malla abierta (el \emph{tipo}) y de una constante estática (\cref{cap:er}). Las dos cosas se leen en el extremo izquierdo del diagrama de magnitud (\cref{fig:rf-tipos}):

\begin{itemize}
  \item \textbf{Tipo 0:} la magnitud arranca plana. El nivel de baja frecuencia es $20\log\kp$.
  \item \textbf{Tipo 1:} arranca con $-20\dBdec$. La asíntota de baja frecuencia, $\kv/\omega$, cruza $0\dB$ en $\omega=\kv$.
  \item \textbf{Tipo 2:} arranca con $-40\dBdec$. La asíntota, $\ka/\omega^2$, cruza $0\dB$ en $\omega=\sqrt{\ka}$.
\end{itemize}

\begin{figure}[ht]
\centering
\includegraphics[width=\linewidth]{cap10/f_tipos}
\caption{El tipo se ve en la pendiente inicial y la constante estática, en un nivel o en un cruce de la asíntota de baja frecuencia.}
\label{fig:rf-tipos}
\end{figure}

\begin{ejemplo}[label=ej:rf-tipos]{Lea las constantes y el error}
Los tres diagramas de la \cref{fig:rf-tipos} corresponden a $\dfrac{0.01}{s/10+1}$, $\dfrac{1}{s(s/10+1)}$ y $\dfrac{100}{s^2(s/100+1)}$. Determine el tipo, la constante estática y el error ante escalón, rampa y parábola unitarios (los tres lazos cerrados son estables).
\solucion
\begin{center}\small
\begin{tabular}{@{}lccccc@{}}
\toprule
Sistema & lectura en el Bode & constante & escalón & rampa & parábola\\
\midrule
Tipo 0 & nivel $-40\dB$ & $\kp=10^{-40/20}=0.01$ & $\frac{1}{1+\kp}=0.99$ & $\infty$ & $\infty$\\
Tipo 1 & cruce en $\omega=1$ & $\kv=1$ & $0$ & $\frac{1}{\kv}=1$ & $\infty$\\
Tipo 2 & cruce en $\omega=10$ & $\ka=10^2=100$ & $0$ & $0$ & $\frac{1}{\ka}=0.01$\\
\bottomrule
\end{tabular}
\end{center}
\end{ejemplo}

\codigorepo{cap10_respuesta_frecuencial/Clase3_LazoCerrado/C3_Ej2_TipoSistema.m}{tipos}{ejemplo \ref{ej:rf-tipos}}

\begin{ideaclave}
Para mejorar el error hay que \emph{subir la curva en bajas frecuencias}; para el transitorio importa la \emph{zona del cruce}. Un compensador en atraso explota exactamente esta separación (\cref{sec:rf-atraso}).
\end{ideaclave}

\begin{pausa}
¿Por qué en un sistema tipo 2 la asíntota de baja frecuencia cruza $0\dB$ en $\sqrt{\ka}$ y no en $\ka$?
\respuesta{En bajas frecuencias $\Gma(\jw)\approx\ka/(\jw)^2$, cuya magnitud es $\ka/\omega^2$. Igualándola a $1$ resulta $\omega^2=\ka$, es decir, $\omega=\sqrt{\ka}$.}
\end{pausa}

%% ---------------------------------------------------------------------
\section{Diseño ajustando la ganancia}
\label{sec:rf-ajuste}
%% ---------------------------------------------------------------------

Primer diseño, con una sola perilla. Como una ganancia sube o baja la curva de magnitud \emph{sin tocar la fase} (factor 1), podemos escoger la frecuencia donde la fase ya tiene el margen que queremos y luego subir o bajar la magnitud hasta que esa frecuencia sea el cruce.

\begin{receta}[Ajuste de ganancia para un $\OS$ dado]
\begin{enumerate}
  \item De la especificación: $\OS\Rightarrow\zeta=\dfrac{-\ln(\OS/100)}{\sqrt{\pi^2+\ln^2(\OS/100)}}\Rightarrow\PMdes$ con \eqref{eq:rf-PMzeta}.
  \item Dibujar el Bode con una ganancia cómoda y buscar la frecuencia $\omega_d$ donde $\angle\Gma(\jw_d)=-180\grados+\PMdes$.
  \item Subir o bajar la magnitud hasta que valga $0\dB$ en $\omega_d$: $\Kc=1/|\Gp(\jw_d)|$.
  \item Verificar con \cmd{margin} y \cmd{step}.
\end{enumerate}
\end{receta}

\begin{ejemplo}[label=ej:rf-ajuste]{$\Kc$ para un $20\%$ de sobrepico}
En el lazo de la \cref{fig:rf-lazo} con $\Gc=\Kc$ y $\Gp(s)=\dfrac{1}{s(s+20)(s+100)}$, determine $\Kc$ para que el lazo cerrado tenga $\OS=20\%$.
\solucion Es un sistema de tercer orden y tipo 1. En lugar de despreciar un polo para parecerse a un segundo orden, usamos la receta directamente.

\textbf{Paso 1.} $\zeta=0.456$ y $\PMdes=48.15\grados$: la fase objetivo es $-180\grados+48.15\grados=-131.85\grados$.

\textbf{Paso 2.} Con $\Kc=2000$ (ganancia DC de los polos reales normalizada) se busca la frecuencia con fase $-131.85\grados$: $\omega_d=13.56$\,rad/s. Allí la magnitud vale $-24.36\dB$.

\textbf{Paso 3.} Hay que subir $24.36\dB$: $\Kc=2000\times10^{24.364/20}\approx33\,056$.

\textbf{Paso 4.} Con este $\Kc$ el lazo cerrado tiene $\OS=19.9\%$ y $\tp=0.21$\,s (\cref{fig:rf-ajuste}). El sistema es de tercer orden, pero el polo en $-100$ está lejos: domina el par complejo.
\end{ejemplo}

\begin{figure}[ht]
\centering
\includegraphics[width=\linewidth]{cap10/f_ajuste_ganancia}
\caption{Ejemplo~\ref{ej:rf-ajuste}: se busca la frecuencia con la fase objetivo y se sube la curva $24.36\dB$ para que allí esté el cruce.}
\label{fig:rf-ajuste}
\end{figure}

\codigorepo{cap10_respuesta_frecuencial/Clase3_LazoCerrado/C3_Ej3_AjusteGanancia.m}{ajuste}{ejemplo \ref{ej:rf-ajuste}: sin buscar a ojo}
\video{V6_ajuste_ganancia}{La curva de magnitud sube hasta cruzar $0\dB$ en $13.56$\,rad/s mientras la fase queda quieta.}

\begin{cuidado}[Leer a ojo cuesta cifras]
Leyendo el diagrama a mano se obtiene algo como $-24.2\dB$ y $\Kc=2000\times10^{24.2/20}\approx32\,400$ (el handout anterior reportaba $33\,192$ para esa lectura: era un error aritmético); el valor exacto, buscando la frecuencia con \cmd{fzero}, es $\Kc=33\,056$. La diferencia no importa para el diseño, pero sí para reportar: diga siempre de dónde sale cada cifra.
\end{cuidado}

\subsection{El límite del ajuste de ganancia}

Con una sola perilla solo se cumple \emph{una} especificación. Con $\Kc=33\,056$ la constante de velocidad es $\kv=\Kc/2000=16.5$. Si además se exigiera $\kv=40$, haría falta $\Kc=80\,000$, y entonces $\PM=25.4\grados$ y $\OS\approx49\%$: subir $\Kc$ para el error destruye el margen de fase. Necesitamos cambiar la \emph{forma} de la curva, no solo subirla o bajarla. Eso es un compensador.

%% ---------------------------------------------------------------------
\section{Compensadores en adelanto, en atraso y atraso-adelanto}
\label{sec:rf-compensadores}
%% ---------------------------------------------------------------------

\subsection{El dilema y la filosofía de diseño}

En el lazo de la \cref{fig:rf-lazo} con $\Gp(s)=K/[s(s+2)]$ queremos dos cosas: un error del $5\%$ ante una rampa ($\kv=20$, es decir, $K=40$) y un margen de fase de $45\grados$. Con $K=40$ el margen de fase es apenas $18\grados$ (sobrepico del $60\%$); bajando $K$ se recupera el margen pero se pierde el error. Con una sola perilla no alcanza: necesitamos cambiar la forma de la curva.

\begin{historia}[De las redes RC a los servomecanismos]
En Bell Labs las redes de adelanto y atraso eran simples combinaciones de resistencias y condensadores dentro del lazo de los amplificadores. Durante la Segunda Guerra Mundial los servomecanismos de radares y cañones se diseñaron en frecuencia en el MIT Radiation Laboratory, de donde salió el libro de James, Nichols y Phillips~\citep{JaNiPh47}. Hoy las mismas redes se implementan con amplificadores operacionales o en software: un adelanto es ``un PD con filtro'' y un atraso es ``casi un PI''. Lo que no ha cambiado es que se diseña mirando el Bode de $\Gc\Gp$: \emph{baja frecuencia = error; cruce = transitorio y estabilidad; alta frecuencia = ruido}.
\end{historia}

\begin{receta}[Cómo se reparte el trabajo del compensador]
\begin{enumerate}
  \item \textbf{$\Kc$ para el error.} Se escoge la ganancia DC más pequeña que cumpla $\ess$ (con $\kp$, $\kv$ o $\ka$).
  \item \textbf{$\Gcp(s)$ para la forma.} Se agregan términos con ganancia DC unitaria, $\Gcp(0)=1$: no cambian el error, pero sí el $\PM$ y la velocidad. El compensador es $\Gc(s)=\Kc\,\Gcp(s)$.
  \item \textbf{Verificar} en el Bode de $\Gc\Gp$ ($\PM$, $\GM$) y en el tiempo (escalón, rampa, esfuerzo de control).
\end{enumerate}
Supuestos: $\Gc\Gp$ sin polos en el semiplano derecho ni cancelaciones en el eje imaginario; así basta el Bode para decidir la estabilidad.
\end{receta}

%% ---------------------------------------------------------------------
\subsection{El compensador en adelanto: agregar fase justo donde hace falta}
\label{sec:rf-adelanto}

\begin{equation}
  \Gcp(s)=\frac{1+s/\wL}{1+s/\wH}=\alpha\,\frac{s+\wL}{s+\wH},\qquad
  \alpha=\frac{\wH}{\wL}>1 .
  \label{eq:rf-adelanto}
\end{equation}
Un cero antes que el polo ($\wL<\wH$): entre ellos la magnitud sube $20\dBdec$ (en total $20\log\alpha$) y la fase se vuelve positiva (\cref{fig:rf-comp}, izquierda). Efectos: más $\PM$ (menos sobrepico) y cruce más a la derecha (más ancho de banda, más rápido). Costo: amplifica el ruido de alta frecuencia en $20\log\alpha$.

\begin{figure}[ht]
\centering
\includegraphics[width=\linewidth]{cap10/f_compensadores_bode}
\caption{Bode del adelanto y del atraso con $\alpha=10$. El adelanto se usa por su fase positiva; el atraso, por su atenuación en alta frecuencia.}
\label{fig:rf-comp}
\end{figure}

\paragraph{¿Dónde está el máximo de fase? Paso a paso}
\begin{enumerate}
  \item La fase del adelanto es $\phi(\omega)=\tan^{-1}\dfrac{\omega}{\wL}-\tan^{-1}\dfrac{\omega}{\wH}$.
  \item Derivando respecto a $\omega$ e igualando a cero, $\dfrac{1/\wL}{1+\omega^2/\wL^2}=\dfrac{1/\wH}{1+\omega^2/\wH^2}$, de donde
  \begin{equation}
    \wM=\sqrt{\wL\,\wH}=\wL\sqrt\alpha .
  \end{equation}
  En escala logarítmica, $\wM$ queda justo en el centro entre el cero y el polo.
  \item Reemplazando, $\tan\phiM=\dfrac{\alpha-1}{2\sqrt\alpha}$, que es lo mismo que
  \begin{equation}
    \sin\phiM=\frac{\alpha-1}{\alpha+1}\qquad\Longleftrightarrow\qquad \alpha=\frac{1+\sin\phiM}{1-\sin\phiM}.
    \label{eq:rf-alpha}
  \end{equation}
  \item En $\wM$ la magnitud está a mitad de camino de su subida: $|\Gcp(\jw_M)|=\sqrt\alpha$, es decir, $+10\log\alpha\dB$.
\end{enumerate}

\paragraph{Rendimientos decrecientes}
La \cref{fig:rf-phimax} muestra $\phiM(\alpha)$. Pasar de $55\grados$ a $65\grados$ exige duplicar $\alpha$ (de $10$ a $20$) y con él el ruido de alta frecuencia. Si se necesitan más de unos $60\grados$, se usan dos etapas de adelanto en cascada o se cambia de estrategia.

\begin{figure}[ht]
\centering
\includegraphics[width=0.78\linewidth]{cap10/f_phimax}
\caption{Fase máxima del adelanto en función de $\alpha$: nunca llega a $90\grados$.}
\label{fig:rf-phimax}
\end{figure}

\paragraph{¿Por qué buscar $|\Gp(\jw_M)|=-10\log\alpha$?}
En $\wM$ el compensador suma $+10\log\alpha\dB$. Si en esa frecuencia la planta vale $-10\log\alpha\dB$, el producto vale $0\dB$: $\wM$ se vuelve el nuevo cruce y el máximo de fase cae exactamente donde se mide el $\PM$. \textbf{¿Y el ángulo de seguridad?} Al correr el cruce de $\wpm$ a $\wM$ (a la derecha) la fase de la planta baja un poco más; $\phiseg$ compensa esa pérdida. Si la pérdida supera a $\phiseg$, se aumenta $\phiseg$ de $5\grados$ en $5\grados$ y se repite.\notamargen{\emph{Caminante, no hay camino, se hace camino al andar.} El diseño en frecuencia es iterativo.}

\begin{receta}[Diseño en adelanto, en siete pasos]
\begin{enumerate}
  \item Escoger $\Kc$ para el error en estado estacionario; anotar $\PMdes$ (o $\zeta$, $\OS$).
  \item Bode de $\Kc\Gp$: medir $\PM$ y $\wpm$.
  \item $\phiM=\PMdes-\PM+\phiseg$, con $\phiseg$ entre $5\grados$ y $10\grados$.
  \item $\alpha=(1+\sin\phiM)/(1-\sin\phiM)$.
  \item Encontrar $\wM$ tal que $20\log|\Kc\Gp(\jw_M)|=-10\log\alpha$.
  \item $\wL=\wM/\sqrt\alpha$ y $\wH=\alpha\,\wL$.
  \item Bode de $\Gc\Gp$: verificar el $\PM$ (y simular). \emph{Verificación estricta:} comprobar que $\phiseg\geq\angle\Gp(\jw_{PM})-\angle\Gp(\jw_M)$; si no, volver al paso 3 con $\phiseg+5\grados$. Si $\phiM\geq90\grados$, el $\PM$ pedido no se alcanza con una sola etapa.
\end{enumerate}
\end{receta}

\begin{ejemplo}[label=ej:rf-adelanto]{Diseño en adelanto}
Diseñe un compensador en adelanto para $\Gp(s)=K/[s(s+2)]$ con error del $5\%$ ante rampa y $\PM=45\grados$.
\solucion
\textbf{Paso 1.} $\ess=1/\kv=0.05\Rightarrow\kv=20$; como $\kv=\lim_{s\to0}s\Gp(s)=K/2$, $K=40$.

\textbf{Paso 2.} Con $K=40$: $\PM=18.0\grados$ en $\wpm=6.17$\,rad/s. Sin compensar, $\OS=60\%$ y $\ts=3.7$\,s.

\textbf{Paso 3.} Con $\phiseg=10\grados$: $\phiM=45-18.0+10=37.0\grados$ (faltan unos $27\grados$; se piden $37\grados$ por la holgura).

\textbf{Paso 4.} $\alpha=(1+\sin37.0\grados)/(1-\sin37.0\grados)=4.03$.

\textbf{Paso 5.} $-10\log\alpha=-6.05\dB$. La magnitud de la planta vale eso en $\wM=8.85$\,rad/s.

\textbf{Paso 6.} $\wL=8.85/\sqrt{4.03}=4.41$ y $\wH=4.03\times4.41=17.76$:
\[
  \Gc(s)=4.03\,\frac{s+4.41}{s+17.76}\qquad(\text{con } K=40 \text{ en la planta}).
\]
\textbf{Paso 7.} $\PM=49.8\grados$ en $8.85$\,rad/s: se cumple con holgura. La verificación estricta también: la planta pierde $5.2\grados$ al correr el cruce de $6.17$ a $8.85$, menos que los $10\grados$ de seguridad. El sobrepico baja de $60\%$ a $22\%$ y $\ts$ de $3.7$ a unos $0.6$\,s (\cref{fig:rf-adelanto-ej}).
\end{ejemplo}

\begin{cuidado}[Una errata del handout anterior]
En la versión anterior de estas notas el texto decía $\wM=14$; era un error tipográfico. Su propio código usaba $8.84$, y los $\wL$ y $\wH$ que reportaba corresponden a $\wM=8.85$.
\end{cuidado}

\begin{figure}[ht]
\centering
\includegraphics[width=\linewidth]{cap10/f_adelanto_ej}
\caption{Ejemplo~\ref{ej:rf-adelanto}: la ``joroba'' de fase del adelanto se ubica en el nuevo cruce y el PM sube de $18.0\grados$ a $49.8\grados$, sin cambiar la ganancia en bajas frecuencias.}
\label{fig:rf-adelanto-ej}
\end{figure}

\codigorepo{cap10_respuesta_frecuencial/Clase4_Compensadores/C4_Ej1_Adelanto.m}{adelanto}{ejemplo \ref{ej:rf-adelanto}}
\codigorepo{cap10_respuesta_frecuencial/Clase4_Compensadores/disenoAdelanto.m}{}{los siete pasos como función}
\video{V7_adelanto_alpha}{$\alpha$ sube de $1$ a $4.03$ con $\wM$ fijo: la joroba crece, el PM sube y el sobrepico baja, con $\kv=20$ todo el tiempo.}

%% ---------------------------------------------------------------------
\subsection{El compensador en atraso: ganancia en bajas, atenuación en el cruce}
\label{sec:rf-atraso}

\begin{equation}
  \Gcp(s)=\frac{1+s/\wH}{1+s/\wL}=\frac{1}{\alpha}\,\frac{s+\wH}{s+\wL},\qquad
  \alpha=\frac{\wH}{\wL}>1 .
  \label{eq:rf-atraso}
\end{equation}
Un polo antes que el cero ($\wL<\wH$). No se usa por su fase (que es negativa) sino por su \emph{atenuación} de $-20\log\alpha$ en alta frecuencia (\cref{fig:rf-comp}, derecha). El truco es poner las esquinas una década por debajo del cruce: allí el atraso ya solo atenúa y ``roba'' apenas unos $5\grados$ de fase.

\begin{receta}[Diseño en atraso]
\begin{enumerate}
  \item Buscar la frecuencia $\omega_{des}$ donde la planta (con el $\Kc$ del error) ya tiene la fase que daría el margen deseado, más una holgura de $5\grados$: $\angle\Gp(\jw_{des})=-180\grados+\PMdes+5\grados$.
  \item Poner el cero una década abajo: $\wH=\omega_{des}/10$.
  \item Escoger $\alpha$ para que $\omega_{des}$ sea el nuevo cruce: $20\log\alpha=20\log|\Gp(\jw_{des})|$; el polo queda en $\wL=\wH/\alpha$.
\end{enumerate}
Es más sencillo que el adelanto (no hay que iterar sobre $\phiM$). El costo: el cruce se mueve a la izquierda y el sistema se vuelve más lento.
\end{receta}

\begin{ejemplo}[label=ej:rf-atraso]{El mismo problema, con atraso}
Repita el diseño del \cref{ej:rf-adelanto} con un compensador en atraso.
\solucion Con $K=40$, la fase vale $-180\grados+45\grados+5\grados=-130\grados$ en $\omega_{des}=1.68$\,rad/s. Allí $|\Gp(j1.68)|=9.13$ ($19.2\dB$), así que $\alpha=9.13$; $\wH=0.168$ y $\wL=0.168/9.13=0.0184$:
\[
  \Gc(s)=\frac{1}{9.13}\,\frac{s+0.168}{s+0.0184}\qquad(\text{con } K=40 \text{ en la planta}),
\]
con $\PM=44.8\grados$. La ganancia en bajas frecuencias no cambió ($\Gcp(0)=1$, $\kv=20$), pero el cruce pasó de $6.17$ a $1.68$\,rad/s (\cref{fig:rf-atraso-ej}).
\end{ejemplo}

\begin{cuidado}[La holgura de $5\grados$ no es opcional]
El handout anterior tomaba $\omega_{des}=2$, es decir, sin la holgura; con ese diseño el $\PM$ queda en unos $40\grados$, por debajo de lo pedido. Los $5\grados$ pagan la fase que el atraso roba en el cruce.
\end{cuidado}

\begin{figure}[ht]
\centering
\includegraphics[width=0.9\linewidth]{cap10/f_atraso_ej}
\caption{Ejemplo~\ref{ej:rf-atraso}: el atraso baja la curva $19.2\dB$ a partir de $\omega\approx0.2$ y conserva la ganancia de bajas frecuencias.}
\label{fig:rf-atraso-ej}
\end{figure}

\codigorepo{cap10_respuesta_frecuencial/Clase4_Compensadores/disenoAtraso.m}{}{los tres pasos como función}

%% ---------------------------------------------------------------------
\subsection{Adelanto contra atraso: el mismo $\kv$, dos personalidades}

\begin{center}\small
\begin{tabular}{@{}lcccccc@{}}
\toprule
 & cruce (rad/s) & $\PM$ & $\OS$ & $\ts$ (s) & $\wbw$ (rad/s) & $u(0^+)$\\
\midrule
sin compensar & 6.17 & $18.0\grados$ & $60\%$ & 3.7 & 9.8 & 1\\
adelanto & 8.85 & $49.8\grados$ & $22\%$ & 0.6 & 14.5 & 4.03\\
atraso & 1.68 & $44.8\grados$ & $27\%$ & 7.1 & 2.8 & 0.11\\
\bottomrule
\end{tabular}
\captionof{table}{Comparación de los dos diseños. $u(0^+)$ es el valor inicial de la señal de control ante un escalón unitario, que es $\Gcp(\infty)$: $\alpha$ para el adelanto y $1/\alpha$ para el atraso.}
\label{tab:rf-comparacion}
\end{center}

\begin{figure}[ht]
\centering
\includegraphics[width=\linewidth]{cap10/f_comparacion}
\caption{Escalón, error ante rampa y esfuerzo de control. Mismo error final ($5\%$): el adelanto llega rápido; el atraso sobrepasa y luego se arrastra hacia $0.05$ (la cola lenta del par polo-cero), pero con mucho menos esfuerzo.}
\label{fig:rf-comparacion}
\end{figure}

\codigorepo{cap10_respuesta_frecuencial/Clase4_Compensadores/C4_Ej3_Comparacion.m}{comparacion}{tabla \ref{tab:rf-comparacion}}
\video{V8_adelanto_vs_atraso}{Siguiendo una rampa: los dos llegan al mismo error, por caminos distintos.}

\begin{anchocompleto}
\centering\small
\begin{tabular}{@{}lll@{}}
\toprule
 & \textbf{Adelanto} & \textbf{Atraso}\\
\midrule
Estructura & $\alpha(s+\wL)/(s+\wH)$, $\alpha>1$ & $(1/\alpha)(s+\wH)/(s+\wL)$, $\alpha>1$\\
Se usa por su\ldots & fase positiva en el cruce & atenuación en el cruce\\
Cruce & se mueve a la derecha & se mueve a la izquierda\\
Velocidad, ancho de banda & aumentan & disminuyen\\
Ruido de alta frecuencia & lo amplifica ($+20\log\alpha$) & lo atenúa\\
Esfuerzo inicial & mayor ($\alpha$) & menor ($1/\alpha$)\\
Límite & unos $60\grados$ por etapa & cola lenta del par polo-cero\\
\bottomrule
\end{tabular}
\captionof{table}{Resumen: adelanto para velocidad; atraso para precisión con poco esfuerzo.}
\label{tab:rf-resumencomp}
\end{anchocompleto}

%% ---------------------------------------------------------------------
\subsection{Atraso-adelanto: cuando hacen falta las dos cosas}

\begin{ejemplo}[label=ej:rf-atrasoadelanto]{Una planta que con el $\kv$ pedido es inestable}
Para $\Gp(s)=K/[s(s+1)(s+2)]$ se pide $\kv=10$, $\PM\geq50\grados$ y $\GM\geq10\dB$.
\solucion $\kv=K/2=10\Rightarrow K=20$, y con esa ganancia el lazo cerrado es inestable: $\PM=-28.1\grados$, $\GM=-10.5\dB$. Un adelanto solo necesitaría unos $85\grados$ de fase; un atraso solo, un cruce muy lento. Se combinan:
\begin{enumerate}
  \item \textbf{Escoger el nuevo cruce}, $\omega_c=1.2$\,rad/s, donde la fase de la planta es $-171.2\grados$.
  \item \textbf{Adelanto centrado en $\omega_c$} con $\alpha=10$ (aporta $\phiM=54.9\grados$): $\wL=1.2/\sqrt{10}=0.38$ y $\wH=1.2\sqrt{10}=3.79$.
  \item \textbf{Atraso que baja la magnitud} en $\omega_c$ lo que sobra para que $\omega_c$ sea el cruce: con el adelanto, $|\Gc\Gp(j1.2)|=23.2\dB$, así que $\beta=10^{23.2/20}=14.5$; el cero va una década abajo ($0.12$) y el polo en $0.12/14.5=0.0083$.
\end{enumerate}
Resultado: $\PM=58.2\grados$ y $\GM=13.2\dB$, con $\kv=10$ intacto (\cref{fig:rf-atrasoadelanto}).
\end{ejemplo}

\begin{figure}[ht]
\centering
\includegraphics[width=\linewidth]{cap10/f_atraso_adelanto}
\caption{Ejemplo~\ref{ej:rf-atrasoadelanto}: el atraso baja la curva en el cruce y el adelanto pone la fase que falta.}
\label{fig:rf-atrasoadelanto}
\end{figure}

\codigorepo{cap10_respuesta_frecuencial/Clase4_Compensadores/C4_Ej4_AtrasoAdelanto.m}{atrasoadelanto}{ejemplo \ref{ej:rf-atrasoadelanto}}

%% ---------------------------------------------------------------------
\subsection{Implementación: de la función de transferencia al circuito}

\begin{figure}[ht]
\centering
\includegraphics[width=0.85\linewidth]{cap10/redes_rc}
\caption{Redes RC pasivas de adelanto y de atraso, con $\alpha=(R_1+R_2)/R_2$. En el adelanto, $\tau=\frac{R_1R_2}{R_1+R_2}C$ y $\alpha\tau=R_1C$; en el atraso, $\tau=R_2C$.}
\label{fig:rf-rc}
\end{figure}

La red de adelanto (\cref{fig:rf-rc}) atenúa en $1/\alpha$: hace falta una etapa de ganancia $\alpha$ (un amplificador operacional) para recuperar el $\kv$. Para el adelanto del \cref{ej:rf-adelanto}, el cero en $\wL=4.41$ da $\alpha\tau=R_1C=0.227$\,s y el polo en $\wH=17.76$ da $\tau=0.0563$\,s; con $C=1\,\mu$F resultan $R_1=227$\,k$\Omega$ y $R_2=R_1/(\alpha-1)=74.9$\,k$\Omega$.

\begin{modelo}[Compensadores.slx]
\cmd{crearModelo\_Compensadores.m} diseña los dos compensadores con \cmd{disenoAdelanto} y \cmd{disenoAtraso} y arma tres lazos en paralelo (sin compensar, adelanto, atraso) con la misma referencia, un selector escalón/rampa y dos osciloscopios: salidas y esfuerzos. Ideal para ver en vivo el error y el sobrepico. La herramienta interactiva \cmd{controlSystemDesigner(Gp)} permite además arrastrar polos y ceros del compensador sobre el Bode y ver el PM en tiempo real.
\end{modelo}

%% ---------------------------------------------------------------------
\section{El PID visto desde la frecuencia}
\label{sec:rf-pid}
%% ---------------------------------------------------------------------

Los compensadores de este capítulo y el PID de los capítulos anteriores son parientes cercanos:
\begin{itemize}
  \item un \textbf{PI}, $\Kc\big(1+\frac{1}{\Ti s}\big)=\Kc\,\frac{\Ti s+1}{\Ti s}$, es un atraso con el polo llevado al origen: ganancia infinita en DC (sube el tipo) a cambio de algo de fase en el cruce;
  \item un \textbf{PD con filtro}, $\Kc\,\frac{1+\Td s}{1+\Td s/N}$, es exactamente un adelanto con $\alpha=N$;
  \item un \textbf{PID} serie con filtro es un atraso-adelanto.
\end{itemize}
Esto permite sintonizar un PID \emph{mirando el Bode}, y entender qué hacen las reglas empíricas.

\paragraph{Ziegler-Nichols leído en frecuencia}
Con la planta del \cref{ej:rf-rangokc} y la sintonización de ZN ($\Kc=73\,326$, $\Ti=0.0943$\,s, $\Td=0.0236$\,s), la malla abierta en la frecuencia crítica vale $\Gc\Gp(\jw_{cr})=-0.6-0.28j$: ZN ``mueve'' el punto que estaba en $-1$ a un lugar fijo, sin importar la planta. El resultado es un lazo agresivo, con $\PM\approx20\grados$ y mucho sobrepico (\cref{fig:rf-pid}).

\paragraph{Moldeo del lazo}
Alternativa: escoger el cruce y el margen y construir el PID para cumplirlos. Con $\omega_c=20$\,rad/s y $\PMdes=50\grados$:
\begin{enumerate}
  \item cero del PI una década abajo del cruce: $\Ti=10/\omega_c=0.5$\,s;
  \item fase que debe aportar el PD filtrado ($N=10$): $-180\grados+50\grados-\big(\angle\Gp+\angle\mathrm{PI}\big)(j20)=37.6\grados$, de donde $\Td=0.0465$\,s;
  \item ganancia para que $|\Gc\Gp(j20)|=1$: $\Kc=33\,427$.
\end{enumerate}
El lazo queda con $\PM=50\grados$ y $\GM=24.4\dB$. La función \cmd{pidtune(Gp,'PIDF',wc,pidtuneOptions('PhaseMargin',50))} hace esto automáticamente.

\begin{figure}[ht]
\centering
\includegraphics[width=0.85\linewidth]{cap10/f_pid_frecuencia}
\caption{Dos PID para la misma planta: Ziegler-Nichols (agresivo) y moldeo del lazo con $\PM=50\grados$.}
\label{fig:rf-pid}
\end{figure}

\codigorepo{cap10_respuesta_frecuencial/Clase4_Compensadores/C4_Ej5_PID_Frecuencia.m}{moldeo}{moldeo del lazo}

\begin{profundiza}[El relé de Åström y Hägglund]
Para medir $\Kcr$ y $\Pcr$ sin llevar la planta al borde de la inestabilidad, Åström y Hägglund propusieron reemplazar el controlador por un relé de amplitud $d$~\citep{AsHa84}. El lazo oscila en la frecuencia donde la fase vale $-180\grados$, con amplitud $a\approx\frac{4d}{\pi\Kcr}$ (aproximación de la función descriptiva), de donde se despeja $\Kcr$. Es el principio de muchos autosintonizadores industriales.
\end{profundiza}

%% ---------------------------------------------------------------------
\section{Para llevar}
\label{sec:rf-cierre}
%% ---------------------------------------------------------------------

\begin{parallevar}
\begin{enumerate}
  \item Un sistema LTI estable responde a $x_0\sin\omega t$ con $x_0|\Gma(\jw)|\sin(\omega t+\angle\Gma(\jw))$. El diagrama de Bode guarda esos dos números para cada $\omega$.
  \item Antes de dibujar, \textbf{forma de Bode}: cada factor como $(s/\omega_i+1)$ y la ganancia $\Kb$ sola. La ganancia escondida es el error más frecuente.
  \item Magnitud: rectas de $\pm20\dBdec$ por factor; error de $3\dB$ en cada esquina real. Fase: $\pm45\grados$/dec entre $\omega_i/10$ y $10\omega_i$. Con polos complejos, \cmd{bode}: la resonancia puede cambiar la conclusión.
  \item La frontera de estabilidad es $\Gma(\jw)=-1$. $\GM$: cuánta ganancia de más se aguanta; $\PM$: cuánto atraso de fase no modelado se aguanta. Regla: $30\grados<\PM<60\grados$, $\GM>6\dB$.
  \item $\PM\leftrightarrow\zeta\leftrightarrow\OS$ (\cref{tab:rf-pmzeta}); el cruce y el ancho de banda dan la velocidad. Las bajas frecuencias dan el tipo y el error.
  \item Diseño: $\Kc$ para el error; luego la forma con $\Gcp(0)=1$. Adelanto para velocidad y margen; atraso para precisión con poco esfuerzo; atraso-adelanto cuando hacen falta las dos cosas. Siempre verificar con \cmd{margin} y en el tiempo.
\end{enumerate}
\end{parallevar}

\begin{profundiza}[Correcciones respecto a la versión anterior del capítulo]
Para quien tenga el handout anterior: (i) adelanto, $\wM=8.85$ (no $14$); (ii) atraso, con la holgura de $5\grados$, $\omega_{des}=1.68$ y $\PM=44.8\grados$; (iii) rango de $\Kc$, $\GM=+40\dB$ (no $-40\dB$); (iv) $\wgm$ asintótico del ejemplo $10\,000/[\cdots]$, $31.6$ (no $50$); (v) ajuste de ganancia, $\Kc$ exacto $33\,056$; (vi) el pico de resonancia se llama ahora $\Mr$ (antes $M_P$) para no confundirlo con el sobrepico. Todos los números del capítulo se verificaron de forma independiente (\texttt{herramientas/cap10\_numeros.py} en el repositorio y \cmd{verificarCap10.m} en \matlab).
\end{profundiza}

%% ---------------------------------------------------------------------
\section{Problemas}
\label{sec:rf-problemas}
%% ---------------------------------------------------------------------

Los problemas marcados con \etiquetaM{} requieren \matlab/\simulink; los marcados con \etiquetaEstrella{} son de profundización. Al final hay soluciones de los primeros.

\subsection*{Diagramas asintóticos y respuesta senoidal}

\begin{problema}[asintótico]\label{prob:rf-1}
Dibuje el diagrama asintótico de magnitud y fase de $\Gma(s)=\dfrac{100(s+1)}{s(s+10)}$. Marque el valor en $\omega=1$ y compare con \cmd{bode} en las esquinas.
\end{problema}

\begin{problema}[asintótico]\label{prob:rf-2}
Ídem para $\Gma(s)=\dfrac{50}{(s+2)(s+5)}$.
\end{problema}

\begin{problema}[respuesta senoidal]\label{prob:rf-3}
Para $\Gma(s)=\dfrac{10}{s+10}$ y $x(t)=2\sin(10t)$, calcule $z_{ss}(t)$ a mano y verifique con \cmd{lsim}.
\end{problema}

\begin{problema}[Simulink]\label{prob:rf-4}
\etiquetaM{} Ejecute \cmd{crearModelo\_RespuestaSenoidal.m} con $G(s)=100/(s^2+4s+100)$. Mida amplitud y desfase para $\omega=1$, $10$ y $50$ y ubique los tres puntos sobre el Bode de \matlab. ¿En cuál de los tres la salida es mayor que la entrada? ¿Por qué?
\end{problema}

\begin{problema}[repaso]\label{prob:rf-5}
Sea $\Gma(s)=\dfrac{20(s+2)}{s(s+20)}$. (a) Escríbala en forma de Bode. (b) Dibuje el diagrama asintótico de magnitud y fase. (c) Compare con los valores reales en $\omega=2$ y $\omega=20$. (d) ¿Cuál es el tipo del sistema y cuánto vale $\kv$?
\end{problema}

\begin{problema}[repaso]\label{prob:rf-6}
Sea $G(s)=\dfrac{50}{(s+5)(s+10)}$ con $x(t)=3\sin(5t)$. (a) Calcule $z_{ss}(t)$ exacta. (b) Estímela solo con el diagrama asintótico y compare. (c) ¿A qué frecuencia la salida es diez veces más pequeña que la entrada (asintótico y exacto)?
\end{problema}

\begin{problema}[el error de la asíntota]\label{prob:rf-7}
Para las mallas abiertas (i) $\dfrac{200}{(s+1)(s+20)}$, (ii) $\dfrac{50(s+2)}{s(s+10)(s+50)}$ y (iii) $\dfrac{5(s+1)}{s^2(s+5)}$:
(a) escríbalas en forma de Bode e identifique $\Kb$, el número de integradores y las esquinas;
(b) justifique a partir de la expresión de la magnitud las reglas asintóticas de un polo real simple;
(c) trace los tres diagramas indicando la pendiente de cada tramo;
(d) calcule el valor exacto de magnitud y fase en cada esquina y verifique que el error de magnitud en una esquina simple es siempre el mismo;
(e) para (ii), determine analíticamente dónde es máximo el error entre la asíntota y la curva real en el tramo $2<\omega<10$ y cuantifíquelo;
(f) superponga los diagramas sobre los de \cmd{bode} y comente dónde la aproximación es aceptable.
\end{problema}

\subsection*{Estabilidad y márgenes}

\begin{problema}[identificación]\label{prob:rf-8}
Un sistema de fase mínima y ganancia positiva tiene un diagrama asintótico de magnitud con pendiente de $-20\dBdec$ en bajas frecuencias, pasando por $40\dB$ en $\omega=0.1$; quiebre a $0\dBdec$ en $\omega=2$; quiebre a $-40\dBdec$ en $\omega=20$, y quiebre a $-20\dBdec$ en $\omega=200$. (a) Reconstruya $\Gma(s)$, justificando cada factor. (b) Trace la fase y determine $\PM$ y $\GM$. (c) Al medir la fase se encuentra que en alta frecuencia tiende a $-270\grados$ en lugar del valor predicho. ¿Qué modificación de $\Gma(s)$ es compatible con la misma magnitud y explica esa fase? (d) Demuestre que ambas funciones tienen diagramas de magnitud idénticos. (e) Compare sus fases en $\omega=20$, $200$ y $2000$\,rad/s y concluya cuál es estable en lazo cerrado con retroalimentación unitaria. (f) Enuncie qué información \emph{no} contiene el diagrama de magnitud y bajo qué condición la magnitud determina la fase.
\end{problema}

\begin{problema}[márgenes y rango de ganancia]\label{prob:rf-9}
Para $\Gma(s)=\dfrac{10\,000}{(s+1)(s+10)(s+100)}$:
(a) trace el asintótico de magnitud y halle $\wpm$;
(b) trace el asintótico de fase con sus puntos de quiebre y halle $\wgm$ y el $\GM$;
(c) compare los cuatro valores con \cmd{margin} y calcule el error porcentual de cada uno;
(d) con $\Gp=\dfrac{1}{(s+1)(s+10)(s+100)}$ y $\Gc=\Kc$, halle el rango de $\Kc$ con Routh-Hurwitz;
(e) obtenga $\Kcr$ del $\GM$ leído en el Bode, compare y explique por qué el error de la lectura asintótica va en la dirección conveniente para el diseño;
(f) repita (a)--(e) para $\Gp=\dfrac{1}{s(s+2)(s+10)}$, obteniendo además la frecuencia crítica por sustitución $s=\jw$; con $\Kcr$ y $\Pcr$ sintonice un PID por Ziegler-Nichols y determine el margen de fase resultante. ¿Es aceptable?
\end{problema}

\begin{problema}[presupuesto de fase]\label{prob:rf-10}
Al lazo del problema~\ref{prob:rf-9}(a) se le agrega el sensor $\Hs(s)=20/(s+20)$. (a) Estime el nuevo $\PM$ con $\PM-\tan^{-1}(\wpm/20)$. (b) Calcúlelo con \cmd{margin}. (c) ¿Por qué la estimación es pesimista? (d) \etiquetaM{} Simule el escalón con y sin sensor y compare los sobrepicos.
\end{problema}

\subsection*{Lazo cerrado, tipo y ajuste de ganancia}

\begin{problema}[PM y transitorio]\label{prob:rf-11}
Para $\Gma(s)=\dfrac{K}{s(s+5)}$ con retroalimentación unitaria:
(a) con $K=100$, identifique $\wn$ y $\zeta$ del lazo cerrado por comparación directa;
(b) obtenga lo mismo por la vía frecuencial ($\omega_1$, $\PM$, ecuación~\eqref{eq:rf-PMzeta});
(c) calcule $\OS$, $\ts$ y $\tp$ y valide simulando;
(d) calcule $\wbw$ con \eqref{eq:rf-BW} y verifíquelo en el Bode del lazo cerrado;
(e) ¿hay pico de resonancia? Si lo hay, calcule $\Mr$ y $\wres$;
(f) halle el $K$ que da $\PM=60\grados$ y compare con $K=100$ en términos de $\wbw$: ¿cuál es más rápido y cuál más amortiguado?
\end{problema}

\begin{problema}[tipo y error]\label{prob:rf-12}
Para (i) $\dfrac{20}{(s+2)(s+5)}$, (ii) $\dfrac{30}{s(s+3)(s+10)}$ y (iii) $\dfrac{400(s+2)}{s^2(s+8)}$ con retroalimentación unitaria:
(a) determine el tipo mirando solo la pendiente de baja frecuencia;
(b) lea gráficamente la constante de error;
(c) verifíquela analíticamente;
(d) calcule el error ante escalón, rampa y parábola, verificando primero la estabilidad;
(e) uno de los tres no admite el cálculo de (d): ¿cuál, por qué y qué rasgo de su Bode lo anticipa?
\end{problema}

\begin{problema}[ajuste de ganancia]\label{prob:rf-13}
Con $\Gc=\Kc$ y $\Gp=\dfrac{1}{s(s+20)(s+100)}$:
(a) halle $\zeta$ y el $\PM$ para $\OS=10\%$;
(b) en el Bode de $2000\,\Gp$ localice la frecuencia con la fase objetivo y lea la magnitud;
(c) determine $\Kc$ y $\kv$;
(d) \etiquetaM{} adapte \cmd{C3\_Ej3\_AjusteGanancia.m}, simule y mida el $\OS$ real;
(e) explique por qué no coincide exactamente con el especificado;
(f) calcule el $\GM$ del sistema ajustado y compárelo con el criterio $\GM>6\dB$.
\end{problema}

\subsection*{Compensadores}

\begin{problema}[adelanto]\label{prob:rf-14}
Diseñe un compensador en adelanto para $\Gp(s)=\dfrac{100K}{s(s+36)(s+100)}$ con $\kv=40$ y $\OS=20\%$. Use $\phiseg=10\grados$ y aplique la verificación estricta; si el $\PM$ no se cumple, itere. \etiquetaM{} Verifique con \cmd{disenoAdelanto(Gp,48.1,10,true)}.
\end{problema}

\begin{problema}[atraso y comparación]\label{prob:rf-15}
Para la planta del problema anterior diseñe un compensador en atraso con las mismas especificaciones (\cmd{disenoAtraso(Gp,48.1)}) y construya una tabla como la \cref{tab:rf-comparacion}. ¿En qué aplicación preferiría cada uno?
\end{problema}

\begin{problema}[implementación]\label{prob:rf-16}
Calcule $R_1$ y $R_2$ de la red de adelanto del problema~\ref{prob:rf-14} con $C=1\,\mu$F. ¿Qué ganancia debe tener el amplificador que la acompaña?
\end{problema}

\begin{problema}[límites del adelanto]\label{prob:rf-17}
\etiquetaEstrella{} Con $\Gp=K/[s(s+2)]$ exija $\kv\geq200$ y $\PM\geq45\grados$. (a) Halle $K$ y el $\PM$ sin compensar. (b) ¿Qué $\alpha$ necesitaría un adelanto de una sola etapa? (c) Repita con $\PM\geq65\grados$ y evalúe si es realizable (un adelanto de primer orden no aporta $90\grados$ y un $\alpha$ muy grande amplifica el ruido). (d) Si no lo es, diseñe una alternativa (dos etapas o atraso-adelanto) y verifique que cumple.
\end{problema}

%% ---------------------------------------------------------------------
\section{Soluciones seleccionadas}
%% ---------------------------------------------------------------------

\paragraph{Problema~\ref{prob:rf-1}}
$\Gma(s)=\dfrac{10(s+1)}{s(s/10+1)}$: $\Kb=10$, un integrador, un cero en $1$ y un polo en $10$. La magnitud baja a $-20\dBdec$ y pasa por $20\dB$ en $\omega=1$; entre $1$ y $10$ queda plana en $20\dB$; después cae a $-20\dBdec$. Fase asintótica: $-90\grados$ hasta $0.1$, sube hasta $-45\grados$ en $\omega=1$, plana en $-45\grados$ hasta $10$, y vuelve a $-90\grados$ en $100$. Valores reales: $22.97\dB$ en $\omega=1$ y $17.03\dB$ en $\omega=10$ (la asíntota dice $20\dB$ en ambos: error de $\pm3\dB$), y $-50.7\grados$ en ambas esquinas.

\paragraph{Problema~\ref{prob:rf-2}}
$\Gma(s)=\dfrac{5}{(s/2+1)(s/5+1)}$: $13.98\dB$ hasta $\omega=2$, $-20\dBdec$ entre $2$ y $5$, $-40\dBdec$ después; la fase va de $0\grados$ a $-180\grados$. Valor real en $\omega=2$: $10.33\dB$ y $-66.8\grados$.

\paragraph{Problema~\ref{prob:rf-3}}
$|\Gma(j10)|=10/|10+j10|=0.707$ y $\phi=-45\grados$: $z_{ss}(t)=1.414\sin(10t-45\grados)$.

\paragraph{Problema~\ref{prob:rf-4}}
$\omega=1$: $1.01$ ($+0.08\dB$), $-2.3\grados$; $\omega=10$: $2.5$ ($+7.96\dB$), $-90\grados$; $\omega=50$: $0.042$ ($-27.6\dB$), $-175.2\grados$. La salida es mayor que la entrada en $\omega=10$, la frecuencia natural ($\zeta=0.2$).

\paragraph{Problema~\ref{prob:rf-5}}
(a) $20(s+2)=40(s/2+1)$ y $s(s+20)=20\,s\,(s/20+1)$, luego $\Gma(s)=\dfrac{2\,(s/2+1)}{s\,(s/20+1)}$, $\Kb=2$ ($6.02\dB$).
(b) Magnitud asintótica: $6.02-20\log\omega$ para $\omega<2$; $0\dB$ entre $2$ y $20$ (el integrador y el cero se cancelan); $-20\log(\omega/20)$ después. Fase: $-90\grados$ hasta $0.2$; entre $2$ y $20$ el cero y el polo se compensan y la fase queda plana en $-45\grados$; vuelve a $-90\grados$ en $200$.
(c) $|\Gma(j2)|=+2.97\dB$ y $|\Gma(j20)|=-2.97\dB$ (la asíntota dice $0$: el error de cada esquina); fase real $-50.7\grados$ en ambas (asintótica: $-45\grados$).
(d) Tipo 1, $\kv=\lim_{s\to0}s\Gma(s)=2$; error ante rampa unitaria $1/\kv=0.5$.

\paragraph{Problema~\ref{prob:rf-6}}
(a) $|G(j5)|=\dfrac{50}{7.071\cdot11.18}=0.632$ ($-3.98\dB$), $\phi=-45\grados-26.6\grados=-71.6\grados$, luego $z_{ss}(t)=1.90\sin(5t-71.6\grados)$.
(b) En forma de Bode $G=1/[(s/5+1)(s/10+1)]$: en $\omega=5$ la asíntota vale $0\dB$ y la fase asintótica $-45\grados-45\grados\log_{10}5=-76.5\grados$, así que $z_{ss}\approx3\sin(5t-76.5\grados)$. La asíntota se equivoca en $3.98\dB$ (un factor $1.58$): sirve para estimar; para calcular se usa $G(\jw)$.
(c) Asintótico: $-6.02-40\log(\omega/10)=-20\Rightarrow\omega\approx22.4$\,rad/s. Exacto: $\omega^4+125\omega^2-247\,500=0\Rightarrow\omega=20.95$\,rad/s.

\codigorepo{cap10_respuesta_frecuencial/Clase1_RespuestaFrecuencial/C1_Ej4_Repaso.m}{}{verificación de los problemas \ref{prob:rf-5} y \ref{prob:rf-6}}

\respuestaspausas

          % piloto del formato nuevo (Handout8_Ctrl.tex)


\bibliographystyle{plainnat}
\bibliography{referencias}
\end{document}
